\section{The common marked slice and its dynamics}
\label{sec:q444-reconciliation}

Two Pro-mode packets dated 4 September 2026 describe a thirteen-dimensional
continuation of the marked nine-dimensional Hermitian carrier. One uses
the condition $u^*b=0$ in a complex-plus-triplet decomposition; the other
requires the observed quaternionic block $PHP^*$ to be complex. We prove
that these are the same embedded space, not merely isomorphic quadratic
models. We retain the original nonunitary marking throughout. We then
compute the omitted-coordinate force in the packet's scalar mechanics,
give a constrained alternative, and correct two action conventions.
The results below concern these specified matrices and dynamics; they do
not certify every theorem in either supplied packet.

The subsequent two-cross continuation computes the entire scalar--Gauss
domain in Section~\ref{q444:sgd:section}. Its classified parameters form
$\mathcal D\times\{1,-1\}^2\times U(1)^2$; retaining the frame phase
$\zeta$ gives the literal domain after the sole identification
$(\zeta,p_+,p_-)\sim(-\zeta,-p_+,-p_-)$.
Corollary~\ref{sgd:every-coupling} constructs tuples at every original
$t=g^2>0$. The two spatial residuals cannot vanish simultaneously on this class by
Theorem~\ref{q444:twocross:thm:exclusion}; existence on the remaining
unrestricted exceptional set is open.

Section~\ref{sgd:smooth-addon:section} proves smooth embeddedness and
the exact product
$\mathcal S_{\rm sg}\cong\{1,-1\}^2\times(0,\infty)
\times\mathbb R^2\times\mathbb T^3$, with $t=g^2$ the displayed
positive coordinate. The four literal components have dimension six;
every fixed positive coupling has four embedded five-dimensional
components $\mathbb R^2\times\mathbb T^3$. The inverse charts retain
all phases and both signs, without a physical gauge quotient or an
additional spatial equation.

\subsection{Conventions, determinant, and the retained marking}

Write $q=a+b\mathbf j$, with $a,b\in\mathbb C$,
$\mathbf j^2=-1$, $\mathbf ja=\bar a\mathbf j$. Thus
\[
 (a+b\mathbf j)(c+d\mathbf j)
 =ac-b\bar d+(ad+b\bar c)\mathbf j,
 \qquad \overline{a+b\mathbf j}=\bar a-b\mathbf j.
\]
For complex matrices of the appropriate sizes this gives
\[
 (A+B\mathbf j)(C+D\mathbf j)
 =AC-B\bar D+(AD+B\bar C)\mathbf j,
 \qquad (A+B\mathbf j)^*=A^*-B^{\mathsf T}\mathbf j.
\]
In particular a quaternionic Hermitian matrix is uniquely $A+B\mathbf j$
with $A=A^*$ and $B=-B^{\mathsf T}$. For $b\in\mathbb C^3$ define
\[
 B(b)=\begin{pmatrix}
 0&\bar b_3&-\bar b_2\\
 -\bar b_3&0&\bar b_1\\
 \bar b_2&-\bar b_1&0
 \end{pmatrix},\qquad \Theta(A,b)=A+B(b)\mathbf j.
\]
The map $\Theta:\operatorname{Herm}_3(\mathbb C)\oplus\mathbb C^3
\to\operatorname{Herm}_3(\mathbb H)$ is consequently a real-linear bijection.

We use the Moore determinant with value one on $I_3$. Its Hermitian
rank-three formula is
\begin{equation}\label{q444:moore}
 N\begin{pmatrix}
 a&q_3&\bar q_2\\ \bar q_3&d&q_1\\q_2&\bar q_1&f
 \end{pmatrix}
 =adf-a|q_1|^2-d|q_2|^2-f|q_3|^2
       +2\operatorname{Re}(q_1q_2q_3).
\end{equation}
The quaternionic spectral theorem and determinant congruence law used below
are recorded in S. Alesker, \emph{Non-commutative linear algebra and
plurisubharmonic functions of quaternionic variables}, Theorem 1.1.9,
\href{https://arxiv.org/abs/math/0104209}{arXiv:math/0104209}; see also
W. Wang, \emph{The linear algebra in the quaternionic pluripotential theory},
Propositions 2.1--2.3,
\href{https://arxiv.org/abs/1805.07759}{arXiv:1805.07759}.
In our signs the faithful complex representation is
\[
 \mathcal R(A+B\mathbf j)=
 \begin{pmatrix}A&B\\-\bar B&\bar A\end{pmatrix}.
\]
Wang's representation has the opposite signs on both off-diagonal blocks;
the exact conversion is conjugation by $\operatorname{diag}(I,-I)$.

\begin{lemma}[The full cubic]\label{q444:pair-cubic}
For every $A=A^*\in\operatorname{Mat}_3(\mathbb C)$ and $b\in\mathbb C^3$,
\begin{equation}\label{q444:pair}
 N(\Theta(A,b))=\det A-b^*Ab.
\end{equation}
\end{lemma}
\begin{proof}
Write the cyclic off-diagonal entries of $A$ as
$A_{23}=z_1$, $A_{31}=z_2$, $A_{12}=z_3$. Then
$q_i=z_i+\bar b_i\mathbf j$ in \eqref{q444:moore}, and
$|q_i|^2=|z_i|^2+|b_i|^2$. The complex part of their product is
\[
 z_1z_2z_3-\bar b_1b_2z_3-z_1\bar b_2b_3-\bar b_1\bar z_2b_3.
\]
The three additional cross terms in twice its real part are the negatives
of the off-diagonal terms of $b^*Ab$. The additional diagonal terms in
\eqref{q444:moore} are
$-a|b_1|^2-d|b_2|^2-f|b_3|^2$. All remaining terms are precisely
the ordinary determinant of $A$. This proves the identity with all signs.
\end{proof}

Retain the marked matrices
\begin{equation}\label{q444:marking}
 \begin{gathered}
 P=\begin{pmatrix}1&-1&0\\0&1&-1\end{pmatrix},\quad
 u=\begin{pmatrix}1\\1\\1\end{pmatrix},\quad
 C=PP^*=\begin{pmatrix}2&-1\\-1&2\end{pmatrix},\\
 R=P^*C^{-1}=\frac13\begin{pmatrix}2&1\\-1&1\\-1&-2\end{pmatrix},
 \quad X=(R\;u),\quad
 X^{-1}=\begin{pmatrix}P\\u^*/3\end{pmatrix},\quad \det X=1.
 \end{gathered}
\end{equation}
Direct multiplication gives $PR=I_2$, $u^*R=0$,
$RP=I_3-uu^*/3$. In particular the metric retained by this marking is
\begin{equation}\label{q444:D}
 D=X^*X=\operatorname{diag}(C^{-1},3),\qquad
 C^{-1}=\frac13\begin{pmatrix}2&1\\1&2\end{pmatrix}.
\end{equation}
It is not the identity matrix.

\subsection{The two constructions are the same embedded space}

Define, in the original quaternionic matrix space,
\[
 \mathcal J_P=\{H\in\operatorname{Herm}_3(\mathbb H):
                       PHP^*\in\operatorname{Herm}_2(\mathbb C)\},
 \qquad
 \mathcal L_u=\{\Theta(A,b):u^*b=0\}.
\]
For any $(A,b)$ put
\begin{equation}\label{q444:full-coordinates}
 X^{-1}AX^{-*}=\begin{pmatrix}G&z\\z^*&\lambda\end{pmatrix},
 \qquad X^*b=\binom{\chi}{\rho},\qquad
 \varepsilon=\begin{pmatrix}0&1\\-1&0\end{pmatrix}.
\end{equation}
Here $G$ is complex Hermitian, $z,\chi\in\mathbb C^2$,
$\lambda\in\mathbb R$, $\rho\in\mathbb C$, and $\rho=u^*b$.

\begin{theorem}[Embedded reconciliation]\label{q444:identification}
The two subspaces satisfy $\mathcal J_P=\mathcal L_u$. Their complete
coordinate relation, with no change of the ambient matrix $H$, is
\begin{equation}\label{q444:dictionary}
 H=X\begin{pmatrix}G&z+w\mathbf j\\(z+w\mathbf j)^*&\lambda\end{pmatrix}X^*,
 \qquad w=-\varepsilon\bar\chi,\qquad \chi=\varepsilon\bar w.
\end{equation}
The inverse in the original complex-plus-triplet coordinates is
\[
 A=X\begin{pmatrix}G&z\\z^*&\lambda\end{pmatrix}X^*,
 \qquad b=X^{-*}\binom{\varepsilon\bar w}{0}.
\]
Thus both thirteen-dimensional constructions have the same positive cone,
the same restricted determinant, and the same embedded old complex space,
which is $w=0$, equivalently $\chi=0$.
\end{theorem}
\begin{proof}
The skew matrices obey the following two identities in the fixed basis:
\begin{equation}\label{q444:skew}
 PB(b)P^{\mathsf T}=\overline{u^*b}\,\varepsilon,
 \qquad X^{-1}B(b)X^{-\mathsf T}=B(X^*b).
\end{equation}
For the first identity the upper-right entry is
$\bar b_3+\bar b_2+\bar b_1$; skew symmetry gives the other entries.
For the second, write $B(b)_{ij}=\sum_k\epsilon_{ijk}\bar b_k$.
For any real invertible $T$, expansion of the alternating determinant gives
\[
 \sum_{a,b}T_{ia}T_{jb}\epsilon_{abk}
 =\det(T)\sum_l\epsilon_{ijl}(T^{-1})_{kl}.
\]
Multiplication by $\bar b_k$ therefore gives
$TB(b)T^{\mathsf T}=\det(T)B(T^{-\mathsf T}b)$.
Set $T=X^{-1}$ and use $\det X=1$. This proves \eqref{q444:skew}.
Its first identity proves the asserted equality of subspaces.
At $\rho=0$ the second identity is the explicit matrix
\[
 B\binom{\chi}{0}
 =\begin{pmatrix}0&0&-\bar\chi_2\\
                  0&0&\bar\chi_1\\
                  \bar\chi_2&-\bar\chi_1&0\end{pmatrix}.
\]
Its upper-right column is $-\varepsilon\bar\chi$. This is exactly
\eqref{q444:dictionary}; $\varepsilon^2=-I_2$ gives the inverse. Every
claim about the embedded matrix and its positive cone follows from this
identity, rather than from matching dimensions or determinant levels.
\end{proof}

\begin{proposition}[The restricted and unrestricted cubics]\label{q444:cubics}
In \eqref{q444:full-coordinates} the unrestricted cubic is
\begin{align}\label{q444:N15}
 N_{15}={}&\lambda\det G-z^*\operatorname{adj}(G)z-\chi^*G\chi\notag\\
          &-2\operatorname{Re}(\rho\chi^*z)-\lambda|\rho|^2.
\end{align}
On the common thirteen-space this becomes
\begin{equation}\label{q444:N13}
 N_{13}=\lambda\det G-z^*\operatorname{adj}(G)z
                     -w^*\operatorname{adj}(G)w.
\end{equation}
\end{proposition}
\begin{proof}
Since $\det X=1$, the ordinary block determinant gives
$\det A=\lambda\det G-z^*\operatorname{adj}(G)z$ first for invertible
$G$ and then, as a polynomial identity, for all $G$. Also
\[
 b^*Ab=(\chi^*,\bar\rho)
       \begin{pmatrix}G&z\\z^*&\lambda\end{pmatrix}\binom{\chi}{\rho}.
\]
Expansion and Lemma~\ref{q444:pair-cubic} prove \eqref{q444:N15}.
For $G=\left(\begin{smallmatrix}a&c\\\bar c&d\end{smallmatrix}\right)$,
direct multiplication gives
$\varepsilon^*\operatorname{adj}(G)\varepsilon=G^{\mathsf T}$.
Substituting $w=-\varepsilon\bar\chi$ shows
$w^*\operatorname{adj}(G)w=\chi^*G\chi$, proving \eqref{q444:N13}.
In the second packet's notation $w=\varepsilon\bar\psi_R$, the exact
dictionary is consequently $\chi=-\psi_R$.
\end{proof}

\subsection{Positive fibres and the embedded sphere pair}

\begin{proposition}[All positive fibres]\label{q444:positive}
The matrix in \eqref{q444:dictionary} is positive semidefinite exactly when
\[
 G\succeq0,\qquad z,w\in\operatorname{Ran}_{\mathbb C}G,\qquad
 s:=\lambda-z^*G^\dagger z-w^*G^\dagger w\geq0.
\]
Its quaternionic rank is $\operatorname{rank}_{\mathbb C}G+\mathbf{1}_{s>0}$.
Here $G^\dagger$ is the inverse on $\operatorname{Ran}G$, extended by zero
on $\ker G$. In particular the fibre over $G=0$ is exactly
$\{\lambda uu^*: \lambda\geq0\}$.
\end{proposition}
\begin{proof}
Congruence by the invertible real matrix $X$ preserves positivity and rank.
Write $\psi=z+w\mathbf j$. For any complex Hermitian $T$,
quaternionic multiplication gives
\[
 \psi^*T\psi=z^*Tz+w^*Tw;
\]
the $\mathbf j$ coefficient is
$z^*Tw-w^{\mathsf T}\bar T\bar z=0$, while the scalar $w$ term is
the conjugate of the real number $w^*Tw$.
Testing the full Hermitian form on $(v,0)$ gives $G\succeq0$.
If $v\in\ker G$ and $v^*\psi\neq0$, testing on $(tv,q)$ with
$q$ chosen to make $\operatorname{Re}(v^*\psi q)\neq0$ makes the
form negative for one sufficiently large signed real $t$.
Thus $\psi$ is in the quaternionic range of $G$, exactly the two complex
range conditions above. With $a=G^\dagger\psi$, direct multiplication gives
\[
 \begin{pmatrix}I&0\\-a^*&1\end{pmatrix}
 \begin{pmatrix}G&\psi\\\psi^*&\lambda\end{pmatrix}
 \begin{pmatrix}I&-a\\0&1\end{pmatrix}
 =\begin{pmatrix}G&0\\0&s\end{pmatrix}.
\]
Both triangular factors are invertible, proving positivity and rank in
both directions. At $G=0$ the two spinors vanish, and return through $X$
gives the asserted ray.
\end{proof}

\begin{proposition}[The old boundary inside the new boundary]\label{q444:pair-topology}
Let $\mathcal C=\mathcal J_P\cap\{H\succeq0\}$. Its determinant-zero
locus is its boundary, and its one-point compactification satisfies
\[
 (\mathcal C^+,\partial\mathcal C^+)\cong(B^{13},S^{12}),\qquad
 ((\partial\mathcal C)^+,
   (\partial\operatorname{Herm}_3(\mathbb C)_{\succeq0})^+)
 \cong(S^{12},S^8).
\]
The sphere pair is the standard linear pair and
$S^{12}\setminus S^8\cong S^3\times\mathbb R^9$.
\end{proposition}
\begin{proof}
The identity belongs to $\mathcal J_P$, since $PI_3P^*=C$ is complex.
Set $\mathcal J_P^0=\{K\in\mathcal J_P:\operatorname{tr}K=0\}$.
If $\lambda_{\min}(K)$ is the least quaternionic Hermitian eigenvalue,
the mutually inverse maps are
\[
 K\longmapsto K-\lambda_{\min}(K)I_3,
 \qquad H\longmapsto H-\frac{\operatorname{tr}H}{3}I_3.
\]
They identify $\mathcal J_P^0$ with the positive singular locus. Indeed,
the first produces least eigenvalue zero; for singular positive $H$, the
least eigenvalue after the second map is $-\operatorname{tr}(H)/3$.
Eigenvalue continuity proves both maps continuous. More generally
\[
 (K,t)\longmapsto K+(t-\lambda_{\min}(K))I_3,
 \quad (K,t)\in\mathcal J_P^0\times[0,\infty)
\]
is a homeomorphism onto $\mathcal C$, taking $t=0$ to its boundary.
This also proves that no other boundary points occur. Its inverse sends
$H$ to $(H-(\operatorname{tr}H/3)I_3,\lambda_{\min}(H))$.
The one-point compactification of a closed half-space is a closed ball,
with the boundary compactifying to a sphere; for example inverse
stereographic projection from a point on the equator gives this pair.
On complex matrices the quaternionic eigenvalues are the same ordinary
Hermitian eigenvalues. The trace-free pair is therefore the linear pair
$(\mathbb R^{12},\mathbb R^8)$ under the displayed homeomorphisms.
Both compactifications have their infinity point in the smaller sphere,
so its complement is
$\mathbb R^8\times(\mathbb R^4\setminus\{0\})
\cong\mathbb R^8\times S^3\times\mathbb R$.
This proves the full pair assertion, not just the two dimensions.
\end{proof}

\subsection{The nonunitary Jordan-product correction}

\begin{theorem}[Completion with the retained metric]\label{q444:jordan}
The Jordan subalgebra generated by $\mathcal J_P$ under
$H\circ K=(HK+KH)/2$ is all of $\operatorname{Herm}_3(\mathbb H)$.
In the marked block coordinates the transported product is
\begin{equation}\label{q444:jordan-product}
 M\circ_D N=\frac12(MDN+NDM),\qquad D=\operatorname{diag}(C^{-1},3),
\end{equation}
not the ordinary block product.
\end{theorem}
\begin{proof}
For $F_X(M)=XMX^*$, multiplication gives
$F_X(M)\circ F_X(N)=F_X((MDN+NDM)/2)$ exactly.
Take
\[
 M=\begin{pmatrix}0&\psi\\\psi^*&0\end{pmatrix},\quad
 N=\begin{pmatrix}0&\phi\\\phi^*&0\end{pmatrix},\quad
 \psi=\binom10,\quad \phi=\binom0{\mathbf j}.
\]
The upper-left block of $M\circ_DN$ is
\[
 \frac32\begin{pmatrix}0&-\mathbf j\\\mathbf j&0\end{pmatrix}.
\]
The bottom scalar is zero: its two terms are
$\mathbf j/3$ and $-\mathbf j/3$. The other blocks vanish. Replacing
$\mathbf j$ by $\mathbf k=\mathbf i\mathbf j$ gives the independent
second missing direction. The common slice already contains every other
coordinate in $\operatorname{Herm}_2(\mathbb H)\oplus\mathbb H^2\oplus
\mathbb R$. Hence the linear span of the slice and these two products
is the whole fifteen-dimensional algebra. Any Jordan subalgebra containing
the slice contains that span. This proves the closure.
The $1/2$ product printed in the quaternionic packet is correct for the
ordinary product in the block frame, but the coefficient for the stated
embedded slice under the original product is $3/2$.
\end{proof}

\subsection{The corrected group-action convention}

Put $Z=(z\;w)$. The ordinary direct product has the following left action:
\begin{equation}\label{q444:action}
 (g,k)\cdot(G,Z,\lambda)=(gGg^*,gZk^{-1},\lambda),
 \quad g\in SL_2(\mathbb C),\quad k\in SU(2).
\end{equation}
Composition gives $g_1g_2Zk_2^{-1}k_1^{-1}
=(g_1g_2)Z(k_1k_2)^{-1}$, as required. The packet's $gZk$ describes
commuting left and right actions, or a left action of
$SL_2(\mathbb C)\times SU(2)^{\mathrm{op}}$; it is not a left action
of the usual direct product with those parameters. For instance
$k_1=\varepsilon$, $k_2=\operatorname{diag}(i,-i)$ do not commute.

Both conventions give the same transformations after $k\mapsto k^{-1}$.
To prove invariance, use
\[
 \operatorname{adj}(gGg^*)=(g^*)^{-1}\operatorname{adj}(G)g^{-1},
 \qquad \det g=1,
\]
first on invertible $G$ and then polynomially on all $G$.
The two spinor terms in \eqref{q444:N13} are
$\operatorname{tr}(Z^*\operatorname{adj}(G)Z)$; conjugation by $k$
does not change this trace. Positivity under the left factor follows by
quaternionic congruence with $\operatorname{diag}(g,1)$ in the block frame.
For the right factor, a unit quaternion $a+b\mathbf j$ acts on $Z$ by
\[
 Z\begin{pmatrix}a&b\\-\bar b&\bar a\end{pmatrix}.
\]
This follows by multiplying $(z+w\mathbf j)(a+b\mathbf j)$ explicitly;
it preserves $\psi^*G^\dagger\psi$, a real scalar, and the range conditions
in Proposition~\ref{q444:positive}.
Finally, if an element fixes every $G$, then $g$ is unitary and commutes
with every complex Hermitian matrix; hence $g=\pm I_2$.
Fixing every $Z$ then forces $k=g$. The effective group is therefore
$(SL_2(\mathbb C)\times SU(2))/\{(I,I),(-I,-I)\}$.

\subsection{The full carrier still has an exact four-dimensional observation}

Requiring the observed quaternionic block to be complex defines the
thirteen-slice above. A different operation retains every coordinate of
the full fifteen-dimensional carrier and applies a positive observation
map. These two operations have different domains, with an exact relation.

\begin{proposition}[Full positive observation and the hidden-coordinate bound]
\label{q444:full-observation}
Fix the same complex subfield $\mathbb C=\mathbb R+\mathbb R\mathbf i$
and real matrix $P$. For a quaternionic Hermitian $2\times2$ matrix $A$
put
\[
 \pi_{\mathbb C}(A)=\tfrac12\bigl(A+(\mathbf iI_2)^*A(\mathbf iI_2)\bigr).
\]
Then the real-linear map
\[
 \mathscr O:\operatorname{Herm}_3(\mathbb H)\longrightarrow
       \operatorname{Herm}_2(\mathbb C),\qquad
 \mathscr O(H)=\pi_{\mathbb C}(PHP^*)
\]
is positive, onto, and has positive section $G\mapsto RGR^*$.
Its real kernel has dimension eleven. On writing $H=\Theta(A_0,b)$,
its observed value is $G=PA_0P^*$ and $\rho=u^*b$ obeys
\begin{equation}\label{q444:rho-bound}
 H\succeq0\quad\Longrightarrow\quad
 G\succeq0,\qquad |\rho|^2\leq\det_{\mathbb C}G.
\end{equation}
Every pair $(G,\rho)$ satisfying these inequalities occurs. Thus the
possible hidden $\rho$ at a fixed $G$ are exactly the closed complex disc
of radius $\sqrt{\det G}$. At a null observed block, $\det G=0$ forces
$\rho=0$.
\end{proposition}
\begin{proof}
For $q=a+b\mathbf j$, $a,b\in\mathbb C$, direct multiplication gives
$(q-\mathbf i q\mathbf i)/2=a$.
Thus $\pi_{\mathbb C}$ extracts the complex matrix part. It is positive
as the average of the identity congruence and unitary congruence by
$\mathbf iI_2$. Compression by $P$ is positive as well. Since $PR=I_2$
and $G$ is complex, $\mathscr O(RGR^*)=G$; positivity of this section is
again congruence. Rank--nullity gives $15-4=11$.

The original skew identity \eqref{q444:skew} gives
\[
 PHP^*=G+\bar\rho\varepsilon\mathbf j
 =\begin{pmatrix}
 g_{11}&g_{12}+\bar\rho\mathbf j\\
 \bar g_{12}-\bar\rho\mathbf j&g_{22}
 \end{pmatrix}.
\]
Its Moore determinant is exactly
$g_{11}g_{22}-|g_{12}|^2-|\rho|^2=\det_{\mathbb C}G-|\rho|^2$.
A Hermitian quaternionic $2\times2$ matrix with real diagonal $a,d$
and upper entry $q$ is positive semidefinite precisely when
$a,d\geq0$ and $ad\geq|q|^2$: for $a>0$ its quadratic form is
$a|x+a^{-1}qy|^2+(d-|q|^2/a)|y|^2$; for $a=0$, positivity forces
$q=0$ by varying $x$, and then requires $d\geq0$.
Applying this criterion proves \eqref{q444:rho-bound} and its converse
for the displayed compressed block.

To realize every permitted pair in the original coordinates, take any
$\lambda\geq0$ and set
\[
 A_0=RGR^*+\lambda uu^*,\qquad b=(\rho/3)u,
 \qquad H=\Theta(A_0,b).
\]
The identities for $X=(R\;u)$ give, with $\chi=R^*b=0$,
\[
 X^{-1}HX^{-*}
   =\begin{pmatrix}G+\bar\rho\varepsilon\mathbf j&0\\0&\lambda\end{pmatrix}.
\]
The preceding criterion and congruence prove positivity of $H$, while
its observation and contraction are exactly $G$ and $\rho$.
This proves the claimed complete image and every disc fibre for $(G,\rho)$.
It does not discard the further $z,w,\lambda$ coordinates in a general $H$.
\end{proof}

In explicit Lorentz coordinates, retain
\[
 G=\begin{pmatrix}p_0+p_3&p_1-ip_2\\p_1+ip_2&p_0-p_3\end{pmatrix},
 \qquad \det G=p_0^2-p_1^2-p_2^2-p_3^2.
\]
The bound is therefore an exact inequality between this Lorentz quadratic
form and the omitted complex coordinate. It does not supply a positive
uniform gap: taking $G=\epsilon I_2$, $\rho=\epsilon/2$,
$\lambda=\epsilon$, $\epsilon>0$, in the realization above gives positive
definite $H$, nonzero $\rho$, and $\det G=\epsilon^2\to0$.
An interpretation as physical mass or as a sector with a nonzero threshold
requires a state-space, units, dynamics, and sector dictionary. The formula
identifies the exact quantity that such a dictionary must retain.
It also avoids imposing $\rho=0$ on unrestricted motion merely to obtain
a four-dimensional observation. Positivity of a chosen dynamical trajectory
itself must still be proved. In fact a positive semidefinite trace-free
Hermitian matrix has all three eigenvalues nonnegative with sum zero, so
it is zero. Thus the nonzero scalar configurations below cannot themselves
be positive matrices.

There is an exact bridge that retains this domain distinction. Put
$\mathcal J_{14}=\{Q\in\operatorname{Herm}_3(\mathbb H):\operatorname{tr}Q=0\}$.
The restriction $\mathscr O|_{\mathcal J_{14}}$ is onto all of
$\operatorname{Herm}_2(\mathbb C)$, not only its positive cone: its section is
\[
 j_0(G)=RGR^*-\frac{\operatorname{tr}(C^{-1}G)}3uu^*.
\]
Indeed $R^*R=C^{-1}$, $\operatorname{tr}(uu^*)=3$, and $Pu=0$,
so $\operatorname{tr}j_0(G)=0$ and $\mathscr O(j_0(G))=G$.
Its kernel therefore has dimension ten.

For a positive boundary observation of every scalar configuration retain
the separate, nonlinear map
\begin{equation}\label{q444:spectral-floor}
 \mathcal F(Q)=Q-\lambda_{\min}(Q)I_3,\qquad
 \widehat G(Q)=\mathscr O(\mathcal F(Q))
             =P\Phi P^*-\lambda_{\min}(Q)C,
 \quad Q=\Theta(\Phi,b).
\end{equation}
It maps $\mathcal J_{14}$ bijectively and continuously to the positive
singular boundary of $\operatorname{Herm}_3(\mathbb H)_{\succeq0}$.
The inverse is $H\mapsto H-(\operatorname{tr}H/3)I_3$: the least
eigenvalue of this trace-free matrix is $-\operatorname{tr}H/3$.
Eigenvalue continuity proves continuity in both directions. This is the
full-space version of the earlier sphere-pair construction. The map leaves
$b$, and hence $\rho$, unchanged, and therefore
\[
 \widehat G(Q)\succeq0,\qquad
 \det\widehat G(Q)\geq|u^*b|^2.
\]
The retained matrix $C=\left(\begin{smallmatrix}2&-1\\-1&2\end{smallmatrix}\right)$
is essential in \eqref{q444:spectral-floor}. This exact observation of a
scalar trajectory is not an assertion that the trajectory preserves a
positive linear slice; at repeated least eigenvalues the observation need
not even be differentiable. No Hamiltonian intertwining is asserted.

The fixed marked observation is not automatically a gauge-invariant
four-momentum. There is an exact covariance rule: under
$H'=UHU^*$, $U\in SU(3)$, transform the observer as $P'=PU^*$;
then $P'H'P'^*=PHP^*$ and hence its complex observation agrees.
Keeping $P$ fixed gives a chosen observer; specifying how that observer
and the integral angular lattice participate in a physical model is
additional data, not supplied by an equality of dimensions.

\subsection{The actual period-to-observation bridge}
The common marking also occurs in the rank-three toroidal extension of
the native rank-two cusp. To specify an actual map, retain its data
\[
 B_0=\begin{pmatrix}0&1\\-1&0\end{pmatrix},\quad
 Q_{\rm latt}=B_0^{-1}P=\begin{pmatrix}0&-1&1\\1&-1&0\end{pmatrix},
 \quad S=\begin{pmatrix}1&1\\0&1\\0&0\end{pmatrix}.
\]
Here $PS=I_2$. Let $C_2(t)$ be the retained bounded holomorphic period
correction at $t=0$. On the logarithmic cover $t=e^{2\pi i s}$ put
\[
 Z_2=sB_0+C_2(t),\quad C_3(t)=SC_2(t)Q_{\rm latt},\quad
 Z_3=sI_3+C_3(t),\quad \Phi_{\rm latt}=P\oplus Q_{\rm latt}.
\]
These are the local corrected periods of the toroidal elliptic bridge;
the underlying toric and compact-fibre morphisms are recorded in the
workbench's \emph{Higher-rung toroidal elliptic bridge}, equations
(TEB.16)--(TEB.35). The new calculation below only needs the displayed
matrices and retains the entire correction $C_2(t)$.

\begin{proposition}[Metric-dependent period observation]\label{q444:period-gram}
Set $\Omega_3=(I_3\ Z_3)$ and $\Omega_2=(I_2\ Z_2)$.
For every real symmetric positive-definite $6\times6$ matrix $W_3$,
\[
 H_3=\Omega_3W_3\Omega_3^*,\quad
 W_2=\Phi_{\rm latt}W_3\Phi_{\rm latt}^{\mathsf T}
 \quad\Longrightarrow\quad
 PH_3P^*=\Omega_2W_2\Omega_2^*.
\]
Both sides are positive definite. No fixed positive $W_3$ makes $H_3$
periodic under $s\mapsto s+1$. Nevertheless every retained positive
matrix $W_*$ gives an explicit descending choice
\[
 M(x)=\begin{pmatrix}I_3&xI_3\\0&I_3\end{pmatrix},\quad
 x=\operatorname{Re}s,\quad
 W_3(s)=M(-x)W_*M(-x)^{\mathsf T}.
\]
For this choice, with $y=\operatorname{Im}s$, the descended matrix is
\[
 H_3(t)=(I_3\;iyI_3+C_3(t))W_*(I_3\;iyI_3+C_3(t))^*.
\]
If $W_*=\left(\begin{smallmatrix}E&F\\F^{\mathsf T}&K\end{smallmatrix}\right)$,
then $G(t)=PH_3(t)P^*$ has leading term $y^2PKP^*$, with $PKP^*\succ0$,
and $\det G(t)=y^4\det(PKP^*)+O(y^3)$ at the cusp.
\end{proposition}
\begin{proof}
Direct multiplication gives $PZ_3=Z_2Q_{\rm latt}$, hence
$P\Omega_3=\Omega_2\Phi_{\rm latt}$ and the claimed Gram identity.
Both $P$ and $Q_{\rm latt}$ have rank two, so $\Phi_{\rm latt}$ is
surjective and $W_2$ is positive definite. Each period matrix has an
identity block, proving positivity of both Gram matrices.

Write a fixed $W_3$ in $3+3$ blocks. Since $Z_3(s+n)=Z_3(s)+nI_3$,
the coefficient of $n^2$ in $H_3(s+n)$ is its positive-definite lower
diagonal block. Therefore $H_3(s+n)$ cannot be constant in $n\in\mathbb Z$.
In particular the period-derived observation cannot be made single-valued
by silently retaining a constant lattice metric.

Put $M_3=M(1)$. The proposed varying metric is positive by congruence
and satisfies $W_3(s+1)=M_3^{-1}W_3(s)M_3^{-\mathsf T}$.
Since $\Omega_3(s+1)=\Omega_3(s)M_3$, its Gram matrix is invariant.
Also $\Omega_3(s)M(-x)=(I_3\;iyI_3+C_3(t))$, which proves the exact
descended formula. The complete expansion, without removing its lower
terms, is
\begin{align*}
 H_3={}&y^2K+iy\bigl(F^{\mathsf T}-F+KC_3^*-C_3K\bigr)\\
       &+E+FC_3^*+C_3F^{\mathsf T}+C_3KC_3^*.
\end{align*}
Boundedness of $C_3$ gives the asserted leading term and determinant
estimate. The principal block $K$ is positive definite and $P^*$ is
injective, so $PKP^*$ is positive definite. This proves every assertion.
\end{proof}

For the explicitly specified example $W_*=I_6$ the exact expression is
$H_3=y^2I_3+iy(C_3^*-C_3)+I_3+C_3C_3^*$ and
$\det G=3y^4+O(y^3)$, retaining $C=PP^*$ and $\det C=3$.
Thus this period-derived observation grows in a timelike, full-rank
direction, not toward a null determinant. It is an actual positive
matrix-valued map on the punctured cusp, not an inverse recovery of
periods or of individual torus points. Its extension through the
successor's finite-order overlaps and special fillings, and a dynamical
interpretation of the chosen lattice metric, are further constructions.

% Insertion-ready fragment for the existing quaternionic module.
% It uses the already defined original P,R,u,C,X,B_0,Q_latt,S,C_2,C_3,
% Z_3,Omega_3 and the monodromy M_3 from q444:period-gram.
% No canonical file was edited by the independent checking agent.

\subsection*{All positive period-metric fibres and the cost of rank loss}

The following calculation keeps the periods in
Proposition~\ref{q444:period-gram} fixed.  In particular
$C_3(t)=SC_2(t)Q_{\rm latt}$ still uses the displayed $S$, not the
different right inverse $R=P^*C^{-1}$.  Choose a punctured cusp on which
the retained correction is bounded, and fix a finite constant $K\geq0$
such that $\|C_3(t)\|\leq K$ there.  Such a constant exists because
$\|C_3(t)\|\leq\|S\|\|C_2(t)\|\|Q_{\rm latt}\|$.
All norms below are the Euclidean operator norms.  Write
\[
 t=e^{2\pi i s},\qquad s=x+iy,\qquad
 y=-\frac{\log|t|}{2\pi},
\]
and restrict to the nonempty smaller cusp $y>\max\{K,1\}$.
This changes no period or coefficient.  It identifies a domain where
the inverse used below is proved to exist.

\begin{theorem}[The complete positive metric fibre]
\label{q444:metric-fibre}
Retain
\[
 A_{\rm per}(s)=xI_3+\operatorname{Re}C_3(t),\qquad
 Y_{\rm per}(t)=yI_3+\operatorname{Im}C_3(t),
\]
and define the real period map and its inverse by
\[
 \mathcal N(s)=
 \begin{pmatrix}I_3&A_{\rm per}(s)\\0&Y_{\rm per}(t)\end{pmatrix},
 \qquad
 \mathcal N(s)^{-1}=
 \begin{pmatrix}
 I_3&-A_{\rm per}(s)Y_{\rm per}(t)^{-1}\\
 0&Y_{\rm per}(t)^{-1}
 \end{pmatrix}.                                             \tag{MF.1}
\]
For every prescribed Hermitian matrix $H=E+iF\succ0$, where
$E=E^{\mathsf T}$ and $F=-F^{\mathsf T}$ are real, all real symmetric
positive-definite lattice matrices giving the unchanged value
$\Omega_3W_3\Omega_3^*=H$ are exactly
\[
 W_3=\mathcal N^{-1}
 \left[
  \frac12\begin{pmatrix}E&-F\\F&E\end{pmatrix}
  +\begin{pmatrix}U&V\\V&-U\end{pmatrix}
 \right]\mathcal N^{-\mathsf T},                             \tag{MF.2}
\]
where $U,V$ are real symmetric $3\times3$ matrices and the bracket is
positive definite.  This is a nonempty open convex twelve-real-dimensional
fibre.  An explicit description of its entire parameter domain is
\[
 E_+=E/2+U\succ0,\qquad
 E/2-U-(-F/2+V)^{\mathsf T}E_+^{-1}(-F/2+V)\succ0.           \tag{MF.3}
\]
In particular $U=V=0$ always gives a positive lift.

Every smooth map $H_3(t)\succ0$ on the smaller punctured cusp has a
smooth positive lift (MF.2) with
\[
 W_3(s+1)=M_3^{-1}W_3(s)M_3^{-\mathsf T}.
                                                               \tag{MF.4}
\]
The full set of positive observations $G(t)=PH_3(t)P^*$ obtainable
from such metrics is exactly the set of all smooth
$\operatorname{Herm}_2(\mathbb C)_{\succ0}$-valued maps on that cusp.
For fixed $G\succ0$, all the intermediate positive $H$ are
\[
 H=X\begin{pmatrix}
       G&z\\ z^*&z^*G^{-1}z+\tau
     \end{pmatrix}X^*,\qquad
 z\in\mathbb C^2,\quad \tau>0.                                \tag{MF.5}
\]
The positive real lattice-metric fibre of $G$ has real dimension
seventeen.
\end{theorem}

\begin{proof}
For every real vector $v$,
\[
 \|Y_{\rm per}v\|
 \geq y\|v\|-\|\operatorname{Im}C_3\,v\|
 \geq(y-K)\|v\|.
\]
The inequality $\|\operatorname{Im}C_3\|\leq\|C_3\|$ follows from
$\operatorname{Im}C_3=(C_3-\bar C_3)/(2i)$ and
$\|\bar C_3\|=\|C_3\|$.  Since $y>K$, $Y_{\rm per}$ has zero kernel
and is an invertible square real matrix.  Block multiplication now proves
both inverse identities in (MF.1).  This is the original real period map:
if $a,b\in\mathbb R^3$, then
\[
 a+Z_3b=(a+A_{\rm per}b)+iY_{\rm per}b.
\]
Thus the inverse takes $v+iw$ to
$(v-A_{\rm per}Y_{\rm per}^{-1}w,\;Y_{\rm per}^{-1}w)$.
No correction in $C_3$ has been removed.

Put $\mathcal J=(I_3\;\;iI_3)$.  The exact factorization is
$\Omega_3=\mathcal J\mathcal N$.  Congruence gives a bijection between
real positive $W_3$ and real positive
$T=\mathcal NW_3\mathcal N^{\mathsf T}$.  For
\[
 T=\begin{pmatrix}E_0&F_0\\F_0^{\mathsf T}&K_0\end{pmatrix}
\]
one obtains by multiplication
\[
 \mathcal JT\mathcal J^*
   =E_0+K_0+i(F_0^{\mathsf T}-F_0).                            \tag{MF.6}
\]
Consequently this equals $E+iF$ precisely when
$E_0=E/2+U$, $K_0=E/2-U$, and $F_0=-F/2+V$ with $U,V$ symmetric.
Indeed, the first equality in (MF.6) fixes $K_0=E-E_0$,
and $F_0^{\mathsf T}-F_0=F$ is equivalent to symmetry of $F_0+F/2$.
These recover $U,V$ uniquely, proving the full fibre formula.

For real $v,w$, with $q=v+iw$,
\[
 \begin{pmatrix}v\\w\end{pmatrix}^{\mathsf T}
 \begin{pmatrix}E&-F\\F&E\end{pmatrix}
 \begin{pmatrix}v\\w\end{pmatrix}=q^*Hq.
\]
The right side is positive for every nonzero $(v,w)$, so the first
bracket term in (MF.2) is positive definite.  This proves nonemptiness.
The kernel parameters have dimension $6+6=12$.  Positive definiteness
is an open convex condition on real symmetric matrices, so its intersection
with this affine twelve-dimensional fibre is open and convex and has
dimension twelve.  Finally, for $B_+=-F/2+V$,
\[
 \begin{pmatrix}E_+&B_+\\B_+^{\mathsf T}&E/2-U\end{pmatrix}
 =
 \begin{pmatrix}I&0\\B_+^{\mathsf T}E_+^{-1}&I\end{pmatrix}
 \begin{pmatrix}E_+&0\\0&E/2-U-B_+^{\mathsf T}E_+^{-1}B_+\end{pmatrix}
 \begin{pmatrix}I&E_+^{-1}B_+\\0&I\end{pmatrix}.
\]
The two triangular factors are transposes of each other and are invertible;
the identity proves the necessity and sufficiency of
(MF.3), including every positivity restriction on the kernel parameters.

The original correction is single-valued in $t$, so
\[
 \mathcal N(s+1)=\mathcal N(s)M_3,\qquad
 \mathcal N(s+1)^{-1}=M_3^{-1}\mathcal N(s)^{-1}.               \tag{MF.7}
\]
For any smooth prescribed $H_3(t)$, its real and imaginary parts are smooth
and periodic on the cover.  The inverse of $Y_{\rm per}$ is smooth,
since its entries are its adjugate divided by its nonzero determinant.
Formula (MF.2) with $U=V=0$ is therefore smooth and satisfies (MF.4).
More generally, smooth periodic $U,V$ satisfying (MF.3) give all such
smooth lifts: conversely (MF.7) and (MF.4) make
$\mathcal NW_3\mathcal N^{\mathsf T}$ periodic, and its blocks recover
periodic $U,V$ by (MF.6).  No holomorphic condition is being imposed on
the real metric or on the prescribed Hermitian map.

If $W_3\succ0$ is real, it is also positive on nonzero complex columns:
for $p+iq$, its quadratic form is $p^{\mathsf T}W_3p+
q^{\mathsf T}W_3q$.  Since $\Omega_3^*$ is injective by the identity
block of $\Omega_3$, one has $H_3\succ0$.  Since $P^*$ is injective,
also $G=PH_3P^*\succ0$.  This proves necessity of strict positivity;
no singular $G$ occurs at an interior point for a positive metric.

For the converse, the marked inverse $X^{-1}=(P;\,u^*/3)$ makes the
upper-left block of $X^{-1}HX^{-*}$ equal to $PH P^*$.  With that block
fixed at $G\succ0$, the complete positive block matrices are
\[
 \begin{pmatrix}G&z\\z^*&\ell\end{pmatrix},\qquad
 \ell-z^*G^{-1}z=\tau>0.
\]
To prove this exact positivity statement, multiply
\[
 \begin{pmatrix}I_2&0\\z^*G^{-1}&1\end{pmatrix}
 \begin{pmatrix}G&0\\0&\tau\end{pmatrix}
 \begin{pmatrix}I_2&G^{-1}z\\0&1\end{pmatrix}.
\]
The result is the displayed block with $\ell=z^*G^{-1}z+\tau$.
The invertible triangular factors prove both directions, and congruence
by $X$ gives (MF.5).  Taking $z=0$ and any fixed $\tau=\lambda>0$
already gives the smooth lift $H_3(t)=RG(t)R^*+\lambda uu^*$.
Applying (MF.2) proves every asserted positive observation is realized.
The intermediate $H$ fibre has four real coordinates in $z$ and one in
$\tau$.  Equivalently, the real-linear metric-to-observation map has rank
four: the constructions above give a right inverse on all Hermitian
matrices before imposing positivity.  Its real symmetric domain has
dimension $6\cdot7/2=21$, so its affine kernel has dimension $21-4=17$.
The positive part is nonempty and relatively open, proving the stated
seventeen-dimensional fibre without a dimension-only surjectivity argument.
\end{proof}

The freedom above does not replace the retained periods by an identity
period matrix.  It computes the metric which compensates their full real
period map.  The correction itself has not lost information in passing
from rank two to rank three: the exact inverse on its prescribed image is
\[
 C_2=P C_3 S B_0,
\]
because $PS=I_2$ and $Q_{\rm latt}S=B_0^{-1}$.
To identify exactly what an unrestricted metric observation cannot recover,
take two retained corrections with real period maps $\mathcal N$ and
$\widetilde{\mathcal N}$ on a common sufficiently small cusp.  The
invertible real transformation
\[
 T_{\rm per}=\widetilde{\mathcal N}^{-1}\mathcal N,\qquad
 \widetilde W_3=T_{\rm per}W_3T_{\rm per}^{\mathsf T}
                                                               \tag{MF.8}
\]
preserves positivity and gives
$\widetilde\Omega_3\widetilde W_3\widetilde\Omega_3^*
=\Omega_3W_3\Omega_3^*$ by the two exact factorizations
$\Omega_3=\mathcal J\mathcal N$.
Furthermore $T_{\rm per}(s+1)=M_3^{-1}T_{\rm per}(s)M_3$, so (MF.8)
preserves the prescribed metric monodromy.  This is an explicit
correspondence between period--metric pairs with the same $H_3$ and $G$.
Thus these observations alone cannot recover the period correction when
the metric is allowed this freedom.  The constants remain in the metric
formula; they have not been set to zero.

\begin{proposition}[Relative-metric bounds obstruct projective rank loss]
\label{q444:metric-condition}
For an equivariant positive metric put
\[
 W_*(t)=M(x)W_3(s)M(x)^{\mathsf T},\qquad
 \widetilde\Omega_3(t)=(I_3\;\;iyI_3+C_3(t)).
\]
These are single-valued on the punctured cusp, and
$H_3=\widetilde\Omega_3W_*\widetilde\Omega_3^*$.
If fixed constants $0<a\leq b<\infty$ obey
$aI_6\preceq W_*(t)\preceq bI_6$, then
\[
 a\bigl(1+(y-K)^2\bigr)C
 \preceq G(t)\preceq
 b\bigl(1+(y+K)^2\bigr)C.                                    \tag{MC.1}
\]
In particular the unchanged eigenvalues satisfy
\[
 \lambda_{\min}(G)\geq a\bigl(1+(y-K)^2\bigr)\lambda_{\min}(C),
 \qquad
 \lambda_{\max}(G)\leq b\bigl(1+(y+K)^2\bigr)\lambda_{\max}(C).
                                                               \tag{MC.2}
\]
With the original $C$, $\lambda_{\min}(C)=1$ and
$\lambda_{\max}(C)=3$.  More generally, writing
$\kappa_*(t)=\lambda_{\max}(W_*)/\lambda_{\min}(W_*)$, one always has
\[
 \frac{\lambda_{\min}(G)}{\lambda_{\max}(G)}
 \geq\frac{1+(y-K)^2}{3\kappa_*(t)(1+(y+K)^2)}.               \tag{MC.3}
\]
Thus if the observed eigenvalue ratio tends to zero as $t\to0$, the
relative metric condition number tends to infinity.  Bounded relative
condition number excludes this projective rank loss even when the common
scale of the metric is not bounded above or below.
\end{proposition}

\begin{proof}
The identity $M(x+1)=M(x)M_3$, together with (MF.4), proves that $W_*$
is periodic.  Multiplication gives
$\Omega_3M(-x)=\widetilde\Omega_3$, which proves the stated Gram formula.
For every complex vector $v$,
\[
 (y-K)\|v\|\leq\|(iyI_3+C_3)^*v\|\leq(y+K)\|v\|.
\]
Both inequalities are the triangle inequality, its reverse, and
$\|C_3^*\|\leq K$.  Since
$\widetilde\Omega_3\widetilde\Omega_3^*
=I_3+(iyI_3+C_3)(iyI_3+C_3)^*$, they give
\[
 \bigl(1+(y-K)^2\bigr)I_3
 \preceq\widetilde\Omega_3\widetilde\Omega_3^*
 \preceq\bigl(1+(y+K)^2\bigr)I_3.
\]
Congruencing $aI_6\preceq W_*\preceq bI_6$ by
$\widetilde\Omega_3$, and then by $P$, proves (MC.1).
The minimum and maximum Rayleigh quotients prove (MC.2).
The characteristic polynomial of the original $C$ is
$(z-1)(z-3)$, retaining the factors $1$ and $3$.

At each point take
$a(t)=\lambda_{\min}(W_*)$ and $b(t)=\lambda_{\max}(W_*)$ in the
same inequalities.  Dividing the resulting positive scalar bounds gives
(MC.3); this division is a comparison of eigenvalues, not a rescaling or
replacement of $G$ or $W_*$.  Its period-dependent factor tends to $1$.
For fixed $a,b$, the lower bound tends to $a/(3b)>0$; for example it is
at least $a/(27b)$ once $y\geq2K$.
If the ratio on the left tends to zero, (MC.3) implies
\[
 \kappa_*(t)\geq
 \frac{1+(y-K)^2}
 {3(1+(y+K)^2)}
 \frac{\lambda_{\max}(G)}{\lambda_{\min}(G)}
 \longrightarrow+\infty .
\]
This proves the full asserted obstruction, including metrics whose
overall scale varies.
\end{proof}

\begin{proposition}[A retained-period rank-loss family and its exact cost]
\label{q444:rank-loss-metric}
Fix, without rescaling any matrix, four real constants
\[
 \alpha_{\rm obs}>0,\qquad \beta_{\rm obs}>0,\qquad
 \lambda>0,\qquad p>0,
\]
and set
\[
 \delta(t)=\beta_{\rm obs}y^{-p},\qquad
 G(t)=\begin{pmatrix}\alpha_{\rm obs}&0\\0&\delta(t)\end{pmatrix},
 \qquad H_3(t)=RG(t)R^*+\lambda uu^*.                         \tag{RL.1}
\]
Take the lift (MF.2) with $U=V=0$.  It is a smooth equivariant positive
real lattice metric at every point of the smaller punctured cusp, with
the original periods and correction unchanged.  Its observation has
\[
 \det G(t)=\alpha_{\rm obs}\beta_{\rm obs}y^{-p}
 =\alpha_{\rm obs}\beta_{\rm obs}
       \left(-\frac{\log|t|}{2\pi}\right)^{-p}\longrightarrow0,
 \qquad
 G(t)\longrightarrow\operatorname{diag}(\alpha_{\rm obs},0).
                                                               \tag{RL.2}
\]
Put $D_2(t)=\operatorname{Im}C_2(t)$ and
$Y=Y_{\rm per}(t)=yI_3+\operatorname{Im}C_3(t)$.  The relative metric
has the exact determinant
\[
 \det W_*(t)=
 \frac{\alpha_{\rm obs}^2\lambda^2\beta_{\rm obs}^2\,y^{-2p}}
      {64\,\det(Y)^2},\qquad
 \det(Y)=y\left[
 y^2+y\,\operatorname{tr}(D_2B_0^{-1})+\det D_2\right].        \tag{RL.3}
\]
In particular
\[
 \det W_*(t)\sim
 \frac{\alpha_{\rm obs}^2\lambda^2\beta_{\rm obs}^2}{64}
 y^{-2p-6}.                                                   \tag{RL.4}
\]
If $L_{\max}=\max\{2\alpha_{\rm obs}/3,\,3\lambda\}$, then
\[
 \lambda_{\min}(W_*)\sim\frac{\beta_{\rm obs}}4\,y^{-p-2},
 \qquad
 \lambda_{\max}(W_*)\longrightarrow\frac{L_{\max}}2,
 \qquad
 \kappa_*(t)\sim\frac{2L_{\max}}{\beta_{\rm obs}}\,y^{p+2}.
                                                               \tag{RL.5}
\]
Every constant and the complete correction in the exact formulas
(RL.1)--(RL.3) is retained.
\end{proposition}

\begin{proof}
Formula (MF.5) with $z=0$ and $\tau=\lambda$ gives $H_3\succ0$ and
$PH_3P^*=G$.  Formula (MF.2) proves the metric, smoothness, and monodromy
assertions.  Equation (RL.2) is the direct determinant of the displayed
matrix and uses the original logarithmic-cover factor $2\pi$.
The ratio of its least and greatest eigenvalues is eventually
$\delta/\alpha_{\rm obs}$, so this is projective rank loss in the
precise sense used in Proposition~\ref{q444:metric-condition}.

Write $A_c=\operatorname{Re}C_3$ and
\[
 \widetilde{\mathcal N}=
 \begin{pmatrix}I_3&A_c\\0&Y\end{pmatrix},\qquad
 \mathcal N=\widetilde{\mathcal N}M(x).
\]
Here $H_3$ is real symmetric.  The relative metric is exactly
\[
 W_*=\frac12\widetilde{\mathcal N}^{-1}
       \begin{pmatrix}H_3&0\\0&H_3\end{pmatrix}
       \widetilde{\mathcal N}^{-\mathsf T}
 =\frac12\begin{pmatrix}
 H_3+A_cY^{-1}H_3Y^{-\mathsf T}A_c^{\mathsf T}
       &-A_cY^{-1}H_3Y^{-\mathsf T}\\
 -Y^{-1}H_3Y^{-\mathsf T}A_c^{\mathsf T}
       &Y^{-1}H_3Y^{-\mathsf T}
 \end{pmatrix}.                                               \tag{RL.6}
\]
Since $H_3=X\operatorname{diag}(G,\lambda)X^*$ and $\det X=1$,
$\det H_3=\lambda\alpha_{\rm obs}\delta$.
Taking the real $6\times6$ determinant in (RL.6) gives
$\det W_*=(\det H_3)^2/(2^6\det Y^2)$, proving the first identity
in (RL.3).

The second identity retains the correction rather than replacing $Y$
by $yI_3$.  Indeed
$Y=yI_3+SD_2Q_{\rm latt}$ and $Q_{\rm latt}S=B_0^{-1}$.
For rectangular $U,V$, the identity
$\det(I+UV)=\det(I+VU)$ follows by eliminating either off-diagonal
block of $\left(\begin{smallmatrix}I&U\\-V&I\end{smallmatrix}\right)$
with triangular determinant-one row operations.  Applying it here gives
\[
 \det Y
 =y^3\det\left(I_2+y^{-1}D_2B_0^{-1}\right)
 =y\left[y^2+y\operatorname{tr}(D_2B_0^{-1})+\det D_2\right],
\]
because $\det B_0=1$.  If
$D_2=\left(\begin{smallmatrix}d_{11}&d_{12}\\d_{21}&d_{22}
\end{smallmatrix}\right)$, this is the complete polynomial
\[
 \det Y=y\left[y^2+y(d_{12}-d_{21})
                     +d_{11}d_{22}-d_{12}d_{21}\right].
\]
The entries of $D_2$ remain the prescribed bounded functions of $t$.
Thus $\det Y/y^3\to1$, proving (RL.4) with the displayed coefficient.

For the eigenvalue assertions, $H_3u=3\lambda u$ because $R^*u=0$.
On the two-dimensional invariant plane $u^\perp$, its trace is
\[
 \operatorname{tr}(C^{-1}G)
      =\frac23(\alpha_{\rm obs}+\delta)
\]
and the product of the two eigenvalues is
$\det H_3/(3\lambda)=\alpha_{\rm obs}\delta/3$.
The invariance of the plane follows also from Hermiticity and the
displayed $u$ eigenvector.  Its two eigenvalues are therefore exactly
\[
 h_\pm=
 \frac{\alpha_{\rm obs}+\delta
       \pm\sqrt{\alpha_{\rm obs}^2-\alpha_{\rm obs}\delta+\delta^2}}3.
                                                               \tag{RL.7}
\]
They are positive: their sum and product are positive.  Rationalizing
the numerator of $h_-$ gives
\[
 \frac{h_-}{\delta}
 =\frac{\alpha_{\rm obs}}
 {\alpha_{\rm obs}+\delta+
       \sqrt{\alpha_{\rm obs}^2-\alpha_{\rm obs}\delta+\delta^2}}
 \longrightarrow\frac12.
\]
Also $h_+\to2\alpha_{\rm obs}/3$.
Consequently $\lambda_{\min}(H_3)\sim\delta/2$ and
$\lambda_{\max}(H_3)\to L_{\max}$, with $H_3$ bounded.

The explicit blocks (RL.6), the bound
$\|Y^{-1}\|\leq(y-K)^{-1}$, and bounded $A_c,H_3$ imply
\[
 \left\|W_*-\tfrac12\operatorname{diag}(H_3,0)\right\|
 =O(y^{-2}).
\]
For real symmetric matrices the maximum eigenvalue changes by at most
the operator norm of a perturbation: apply the Rayleigh-quotient maximum
to each matrix and subtract.  Hence
$\lambda_{\max}(W_*)\to L_{\max}/2$.

To control the least eigenvalue without assuming a bounded condition
number of $H_3$, use the exact inverse
\[
 W_*^{-1}
 =2\widetilde{\mathcal N}^{\mathsf T}
      \operatorname{diag}(H_3^{-1},H_3^{-1})
      \widetilde{\mathcal N}.
\]
For $y>1$, the decomposition
$\widetilde{\mathcal N}=\operatorname{diag}(I_3,yI_3)
+\left(\begin{smallmatrix}0&A_c\\0&\operatorname{Im}C_3\end{smallmatrix}\right)$
gives $\|\widetilde{\mathcal N}\|\leq y+\sqrt2K$.
Thus
\[
 \lambda_{\max}(W_*^{-1})
 \leq\frac{2(y+\sqrt2K)^2}{\lambda_{\min}(H_3)}.
\]
Its lower-right principal block includes the positive term
$2Y^{\mathsf T}H_3^{-1}Y$.  If $v$ is a unit eigenvector of
$H_3^{-1}$ for its greatest eigenvalue, testing that block on
$Y^{-1}v/\|Y^{-1}v\|$ gives
\[
 \lambda_{\max}(W_*^{-1})
 \geq\frac{2(y-K)^2}{\lambda_{\min}(H_3)}.
\]
All denominators are positive.  Taking reciprocals yields
\[
 \frac{\lambda_{\min}(H_3)}{2(y+\sqrt2K)^2}
 \leq\lambda_{\min}(W_*)
 \leq\frac{\lambda_{\min}(H_3)}{2(y-K)^2}.                     \tag{RL.8}
\]
Both bounds are asymptotic to $\delta/(4y^2)$, proving the first
assertion of (RL.5).  Combining it with the maximum-eigenvalue limit
proves the condition-number asymptotic with its full coefficient.
\end{proof}

This is metric-choice freedom, not a selection of a physical momentum
or mass.  For the same unchanged periods the allowed choice $W_*=I_6$
instead gives exactly
\[
 H_3=I_3+(iyI_3+C_3)(iyI_3+C_3)^*,\qquad
 G=y^2C+O(y),\qquad \det G=3y^4+O(y^3),
\]
whereas (RL.1)--(RL.2) give a determinant tending to zero and retain the
same $C_2,C_3$.  These two explicit admissible choices, the complete
twelve-dimensional metric fibres, and the seventeen-dimensional observed
fibres prove that the specified period data and positivity alone do not
select an observed trajectory.  The eigenvalue ratios above only diagnose
rank loss; no field, period, metric, or parameter was replaced by a
normalized one.  Units, a metric-selection law, dynamics, and an
interpretation of the marked observation would be additional specified
data.  No extension across other finite-order overlaps or special
fillings, and no physical or quantum gap conclusion, is asserted here.

% Insertion-ready; uses the original P,R,u,C,X,B_0,S,Q_latt,C_2,C_3.
% Independent lane: no canonical file was edited.
\subsection*{Hermitian flat metrics, the original period lattice, and its dual spectrum}

Here a flat metric means a translation-invariant metric on each fixed
torus fibre. Its coefficients may depend smoothly on the cusp parameter
$t$, but not on the point of that fibre. The word \emph{physical} below
distinguishes the original real coordinates of $\mathbb C^3$ from integral
lattice coordinates; it does not assert a physical field theory.
Retain, without changing any basis or period,
\[
 B_0=\begin{pmatrix}0&1\\-1&0\end{pmatrix},\quad
 P=\begin{pmatrix}1&-1&0\\0&1&-1\end{pmatrix},\quad
 Q_{\rm latt}=B_0^{-1}P,\quad
 S=\begin{pmatrix}1&1\\0&1\\0&0\end{pmatrix},
\]
\[
 C_3=SC_2Q_{\rm latt},\qquad
 Z_3=sI_3+C_3,\qquad Z_2=sB_0+C_2,\qquad
 s=x+iy,\quad t=e^{2\pi i s},\quad
 y=-\frac{\log|t|}{2\pi}.
\]
Choose finite bounds $K_2,K\geq0$ on $\|C_2(t)\|,\|C_3(t)\|$,
respectively, and restrict to $y>\max\{K_2,K\}+1$. Norms are the
original Euclidean operator norms. Write
\[
 A=xI_3+\operatorname{Re}C_3,\quad
 Y=yI_3+\operatorname{Im}C_3,\quad
 \mathcal N=\begin{pmatrix}I_3&A\\0&Y\end{pmatrix},\quad
 \mathcal J=(I_3\;\;iI_3),\quad
 J_0=\begin{pmatrix}0&-I_3\\I_3&0\end{pmatrix}.
\]
Thus $\Omega_3=(I_3\;\;Z_3)=\mathcal J\mathcal N$ and
$\mathcal N^{-1}=\left(\begin{smallmatrix}
I_3&-AY^{-1}\\0&Y^{-1}\end{smallmatrix}\right)$.
Indeed $\|Yv\|\geq(y-K)\|v\|$ proves invertibility. The same argument,
using $\|B_0v\|=\|v\|$, proves invertibility of
$Y_2=\operatorname{Im}Z_2=yB_0+\operatorname{Im}C_2$.
For $r=2,3$ put
\[
 \Lambda_r=\Omega_r\mathbb Z^{2r},\qquad
 T_r=\mathbb C^r/\Lambda_r,\qquad
 \mathcal N_2=
 \begin{pmatrix}I_2&\operatorname{Re}Z_2\\0&Y_2\end{pmatrix}.
\]
The real coordinates of $z\in\mathbb C^3$ are
$q=(\operatorname{Re}z,\operatorname{Im}z)^{\mathsf T}
=\mathcal N\xi$, where $\xi\in\mathbb R^6/\mathbb Z^6$.
This is an exact change of coordinates, not a replacement of the lattice.
For any complex matrix $D$ define its realification by
\[
 \mathfrak r(D)=
 \begin{pmatrix}\operatorname{Re}D&-\operatorname{Im}D\\
                 \operatorname{Im}D&\operatorname{Re}D\end{pmatrix}.
\]
Block multiplication gives
$\mathfrak r(D_1D_2)=\mathfrak r(D_1)\mathfrak r(D_2)$ and
$\mathfrak r(D^*)=\mathfrak r(D)^{\mathsf T}$.

\begin{theorem}[The precise Hermitian metric in a period fibre]
\label{q444:hermitian-cometric}
Let $g_{\rm phys}\succ0$ be a real symmetric metric matrix and put
\[
 T=g_{\rm phys}^{-1},\qquad
 g_{\rm latt}=\mathcal N^{\mathsf T}g_{\rm phys}\mathcal N,\qquad
 W=g_{\rm latt}^{-1}=\mathcal N^{-1}T\mathcal N^{-\mathsf T}.
                                                               \tag{HS.1}
\]
Then the earlier period weight $W$ has the exact geometric meaning of a
\emph{lattice cometric}; the covariant lattice metric is its inverse.
The observation is $H=\Omega_3W\Omega_3^*=\mathcal JT\mathcal J^*$.
For a fixed positive Hermitian $H=E+iF$, its full real cometric fibre is
\[
 T=\frac12\mathfrak r(H)+
       \mathcal A(U,V),\qquad
 \mathcal A(U,V)=\begin{pmatrix}U&V\\V&-U\end{pmatrix},\qquad
 U=U^{\mathsf T},\quad V=V^{\mathsf T},\quad T\succ0.              \tag{HS.2}
\]
The real metric is compatible with the given complex structure, meaning
$J_0^{\mathsf T}g_{\rm phys}J_0=g_{\rm phys}$, if and only if
$U=V=0$. With the explicit Hermitian metric convention
\[
 g_{\rm phys}(v,w)=\operatorname{Re}(v^*h_3w),
       \qquad v,w\in\mathbb C^3,
\]
the unique compatible member over $H$ is
\[
 h_3=2H^{-1},\qquad h_3^{-1}=H/2,\qquad
 g_{\rm phys}=2\mathfrak r(H^{-1}),\qquad
 T=\tfrac12\mathfrak r(H),                                    \tag{HS.3}
\]
\[
 g_{\rm latt}=2\mathcal N^{\mathsf T}
                    \mathfrak r(H^{-1})\mathcal N,\qquad
 W=\tfrac12\mathcal N^{-1}\mathfrak r(H)\mathcal N^{-\mathsf T}.
                                                               \tag{HS.4}
\]
For an arbitrary positive real cometric, its observed Hermitian part is
precisely the positive $J_0$-invariant average
\[
 \frac{T+J_0TJ_0^{\mathsf T}}2=\frac12\mathfrak r(H).             \tag{HS.5}
\]
\end{theorem}
\begin{proof}
Substitution of $q=\mathcal N\xi$ into the covariant quadratic form
gives the middle formula in (HS.1); multiplication verifies its inverse.
The observation is positive: $\mathcal J^*$ is injective, and a real
positive matrix $T$ is also positive on complex vectors, since for
$a,b$ real, $(a+ib)^*T(a+ib)=a^{\mathsf T}Ta+b^{\mathsf T}Tb$.
For a real symmetric block matrix
$T=\left(\begin{smallmatrix}D&B\\B^{\mathsf T}&L\end{smallmatrix}\right)$,
\[
 \mathcal JT\mathcal J^*=D+L+i(B^{\mathsf T}-B).
\]
Thus equality to $E+iF$ is equivalent to
$D=E/2+U$, $L=E/2-U$, $B=-F/2+V$ with $U,V$ symmetric.
This proves both directions and uniqueness of the parameters in (HS.2).

Because $J_0^{\mathsf T}=-J_0$ and $J_0^2=-I_6$, metric compatibility
is equivalent to $g_{\rm phys}J_0=J_0g_{\rm phys}$, which is equivalent,
by multiplying by the inverse on both sides, to $TJ_0=J_0T$.
Direct multiplication shows that $\mathfrak r(H)$ commutes with $J_0$,
whereas $\mathcal A J_0=-J_0\mathcal A$. Hence compatibility is equivalent
to $J_0\mathcal A=0$, and the invertibility of $J_0$ makes this equivalent
to $\mathcal A=0$, or $U=V=0$.

For $z=v+iw$ one has
$(v,w)^{\mathsf T}\mathfrak r(H)(v,w)=z^*Hz>0$ unless $z=0$.
Thus $\mathfrak r(H)$ is positive. Realification respects products,
so inversion of $T=\mathfrak r(H)/2$ gives (HS.3), and (HS.1) gives
(HS.4). Finally
$J_0\mathfrak r(H)J_0^{\mathsf T}=\mathfrak r(H)$ and
$J_0\mathcal AJ_0^{\mathsf T}=-\mathcal A$, proving (HS.5).
This average is an average of cometrics, not an average of covariant
metrics. The metric associated to it is the inverse of the average.
\end{proof}

Thus the twelve real fibre coordinates in (HS.2) describe genuine
non-Hermitian translation-invariant real metrics, and compatibility
removes exactly those coordinates. For example, $H=2I_3$ and
$U=\varepsilon\operatorname{diag}(1,0,0)$, $V=0$ give
\[
 T_\varepsilon=\operatorname{diag}
 (1+\varepsilon,1,1,1-\varepsilon,1,1),\qquad -1<\varepsilon<1.
\]
All have the same $H$; only $\varepsilon=0$ is Hermitian-compatible.
They have different marked cometrics, and the next theorem proves
that some marked character eigenvalue therefore changes.

\begin{theorem}[The lattice cometric is the marked character spectrum]
\label{q444:torus-laplace}
For any $T\succ0$ as in (HS.1), define the nonnegative Laplacian by
$\Delta=-\operatorname{div}\operatorname{grad}$. Its complete character
eigenbasis and eigenvalues on $T_3$ are
\[
 \chi_k([\xi])=e^{2\pi i k^{\mathsf T}\xi},\qquad k\in\mathbb Z^6,
 \qquad
 \lambda_k=4\pi^2 k^{\mathsf T}Wk
          =4\pi^2 p_k^{\mathsf T}Tp_k,\quad
 p_k=\mathcal N^{-\mathsf T}k.                               \tag{HS.6}
\]
Here $p_k$ is in the integer-valued dual of
$\mathcal N\mathbb Z^6$: its pairing with $\mathcal Nn$ is
$k^{\mathsf T}n\in\mathbb Z$. Conversely every integer-valued dual
covector is of this form. The eigenspace multiplicity is exactly
the number of integer $k$ with the specified value of (HS.6).
The volume, in the original scales, is
\[
 \operatorname{Vol}(T_3,g_{\rm phys})
 =|\det Y|\sqrt{\det g_{\rm phys}}
 =|\det Y|\det h_3
 =\frac{8|\det Y|}{\det H}                                  \tag{HS.7}
\]
in the compatible case. Without compatibility the first equality
still holds.

The marked spectrum recovers every entry of $W$:
\[
 W_{aa}=\frac{\lambda_{e_a}}{4\pi^2},\qquad
 W_{ab}=
 \frac{\lambda_{e_a+e_b}-\lambda_{e_a}-\lambda_{e_b}}{8\pi^2}
 \quad(a\ne b).                                             \tag{HS.8}
\]
If $H(t)$ is single-valued and (HS.4) is used, all these metrics
descend with the original monodromy
\[
 \mathcal N(s+1)=\mathcal N(s)M_3,\quad
 M_3=\begin{pmatrix}I_3&I_3\\0&I_3\end{pmatrix},\quad
 W(s+1)=M_3^{-1}W(s)M_3^{-\mathsf T}.                         \tag{HS.9}
\]
The same character has labels $k$ and $k'=M_3^{\mathsf T}k$
in the two markings, and $\lambda'_{k'}=\lambda_k$.
\end{theorem}
\begin{proof}
The covariant metric and its determinant are constant in the fibre
coordinate $\xi$, so the defining divergence formula becomes
$\Delta=-\sum_{a,b=1}^6W_{ab}\partial_{\xi_a}\partial_{\xi_b}$.
Two derivatives of $\chi_k$ prove (HS.6). The dual-lattice assertion
follows from $p^{\mathsf T}\mathcal Nn\in\mathbb Z$ for all integer
$n$: taking the six unit vectors is equivalent to
$\mathcal N^{\mathsf T}p\in\mathbb Z^6$.

For completeness, no unlisted eigenfunctions are needed. Integration
over $[0,1]^6$ gives zero for $\chi_k\overline{\chi_l}$ when $k\ne l$;
the common squared norm in Riemannian measure is the volume.
To prove density, the one-dimensional Fejer kernels
\[
 F_M(\theta)=\frac1M\left|\sum_{r=0}^{M-1}e^{2\pi i r\theta}\right|^2
 =\sum_{|j|<M}\left(1-\frac{|j|}{M}\right)e^{2\pi i j\theta}
\]
are nonnegative and have integral one. For
$\operatorname{dist}(\theta,\mathbb Z)\geq\eta>0$ they are bounded by
$1/(M\sin^2(\pi\eta))$, using the finite geometric sum. The product
of six such kernels also has integral one; its mass outside a
coordinate cube of radius $\eta$ tends to zero by the union bound
over its six coordinates. Uniform continuity then proves that
convolution of any continuous periodic function with these product
kernels converges uniformly to that function. Each convolution is a
finite linear combination of characters. Continuous functions are dense
in $L^2([0,1]^6)$: step functions on finite unions of boxes are dense
by regular approximation in Lebesgue measure, and a box indicator is
approximated in $L^2$ by a continuous function that differs only in a
boundary strip of arbitrarily small measure. Multiplication by the
constant Riemannian density preserves this argument.

Consequently the characters, divided by the square root of the volume,
are a complete orthonormal basis. On Fourier coefficients the closure of
$\Delta$ is the real diagonal operator with domain
$\{(a_k):\sum_k\lambda_k^2|a_k|^2<\infty\}$; its adjoint has that same
domain, as testing against each coordinate vector forces the displayed
weighted square summability. It is therefore self-adjoint.
Finite Fourier sums are a core by truncation of both the ordinary
and weighted square sums. For a smooth function, integration by parts
identifies the Fourier coefficients of its smooth Laplacian with
$\lambda_k a_k$; their square sum is finite. Thus the operator on all
smooth functions lies in this diagonal operator and has the same
closure as its restriction to finite Fourier sums. Since
$\lambda_k\geq4\pi^2\lambda_{\min}(W)\|k\|^2$, bounded eigenvalue
intervals contain only finitely many integer indices. This proves
the completeness, discreteness and multiplicity assertions, including
a positive smallest eigenvalue away from $k=0$.

The parallelepiped $\mathcal N[0,1]^6$ has coordinate volume
$|\det\mathcal N|=|\det Y|$, proving the first equality in (HS.7).
For Hermitian $h_3$, the determinant identity
$\det\mathfrak r(h_3)=|\det h_3|^2=(\det h_3)^2$
follows by conjugating the realification over $\mathbb C$ to
$\operatorname{diag}(h_3,\overline{h_3})$ with the invertible change
$(v,w)\mapsto(v+iw,v-iw)$. Positive $h_3$ has positive determinant.
Now $\det(2H^{-1})=8/\det H$, proving (HS.7).
Expansion of the quadratic form at $e_a+e_b$ proves (HS.8).

The correction $C_3(t)$ is unchanged by $s\mapsto s+1$, giving the
first formula in (HS.9). The cometric formula gives the last.
Since $q=\mathcal N\xi=\mathcal N M_3\xi'$, one has
$\xi=M_3\xi'$. Hence $k'=M_3^{\mathsf T}k$ and substitution in
(HS.6) proves the final assertion. All matrix inversions depend
smoothly on the entries where their determinants are nonzero, so a
smooth $H(t)$ gives the asserted smooth descending family.
\end{proof}

\begin{theorem}[The original torus quotient is a metric and spectral map]
\label{q444:metric-submersion}
Let $H_3\succ0$, $G=PH_3P^*$, and equip $T_3$ with
$h_3=2H_3^{-1}$. Retain the original integral map
\[
 L=P\oplus Q_{\rm latt}:\mathbb Z^6\longrightarrow\mathbb Z^4,
 \qquad P_{\mathbb R}=\operatorname{diag}(P,P).
\]
Then $f_P:T_3\to T_2$, $[z]\mapsto[Pz]$, is a surjective holomorphic
homomorphism with connected kernel
\[
 \mathbb C u/(\mathbb Z u+s\mathbb Z u),\qquad
 u=(1,1,1)^{\mathsf T}.                                    \tag{HS.10}
\]
There is a unique target translation-invariant metric making this map
a Riemannian submersion. Its Hermitian metric, complex cometric, real
cometric and lattice cometric are, respectively,
\[
 h_2=2G^{-1},\qquad h_2^{-1}=G/2=Ph_3^{-1}P^*,\qquad
 \mathsf T_2=\tfrac12\mathfrak r(G)
       =P_{\mathbb R}\mathsf T_3P_{\mathbb R}^{\mathsf T},
                                                               \tag{HS.11}
\]
\[
 W_2=\mathcal N_2^{-1}\mathsf T_2\mathcal N_2^{-\mathsf T}
     =LW_3L^{\mathsf T}.                                    \tag{HS.12}
\]
Here $\mathsf T_3=\mathfrak r(H_3)/2$ is the real source cometric,
distinguished from the torus $T_3$. The horizontal lift is the exact complex map
\[
 \eta:\mathbb C^2\longrightarrow\mathbb C^3,\qquad
 \eta=H_3P^*G^{-1}.                                         \tag{HS.13}
\]
For every smooth function $f$ on the target,
\[
 \Delta_3(f\circ f_P)=(\Delta_2f)\circ f_P.                    \tag{HS.14}
\]
For a target character $\ell\in\mathbb Z^4$, its pullback is the source
character labelled by $L^{\mathsf T}\ell$, and its eigenvalue is unchanged.
These are exactly the source characters constant on the fibres. Thus
the target spectrum, with each character's multiplicity counted, is
a subcollection of the source character spectrum, not an identification
with all source modes.
\end{theorem}
\begin{proof}
Direct multiplication gives $PS=I_2$, $B_0Q_{\rm latt}=P$ and
$PZ_3=Z_2Q_{\rm latt}$. Thus
\[
 P\Omega_3=\Omega_2L,\qquad
 P_{\mathbb R}\mathcal N=\mathcal N_2L.                        \tag{HS.15}
\]
Moreover $\operatorname{diag}(S,SB_0)$ is an integral right inverse
of $L$, since $Q_{\rm latt}SB_0=I_2$. Therefore
$P\Lambda_3=\Lambda_2$, proving that the torus homomorphism is
well-defined and surjective.

If $Pz\in\Lambda_2$, surjectivity on lattices supplies
$\gamma\in\Lambda_3$ with $P\gamma=Pz$, so $z-\gamma\in\ker_{\mathbb C}P
=\mathbb C u$. To find its periods, if $\Omega_3(a,b)$ is in
$\mathbb C u$, then $\Omega_2(Pa,Q_{\rm latt}b)=0$.
The invertibility of $\mathcal N_2$ forces $Pa=Q_{\rm latt}b=0$.
Solving the two displayed difference equations gives
$a=ru$, $b=vu$ with $r,v\in\mathbb Z$. Since $C_3u=0$,
$\Omega_3(a,b)=(r+sv)u$. This proves (HS.10), including connectedness.
Its lattice has real rank two because $y>0$.

Put $K_3=h_3^{-1}=H_3/2$, $K_2=PK_3P^*=G/2$.
Surjectivity of $P$ makes $P^*$ injective, so $K_2\succ0$.
The lift in (HS.13) is $K_3P^*K_2^{-1}$ and satisfies
\[
 P\eta=I_2,\qquad
 \eta^*h_3 v=K_2^{-1}Pv=0\quad(v\in\ker P),\qquad
 \eta^*h_3\eta=K_2^{-1}=h_2.                                \tag{HS.16}
\]
Every $w\in\mathbb C^3$ has the orthogonal decomposition
$w=\eta(Pw)+(w-\eta(Pw))$, whose second term lies in $\ker P$.
Hence $P$ restricts to an isometry on the orthogonal complement of
its kernel; this is precisely the definition of a Riemannian
submersion. It also proves uniqueness of $h_2$, since each target
vector has a unique horizontal lift. Taking realifications proves
(HS.11), and multiplying (HS.15) by the inverses on both sides proves
(HS.12).

In lattice coordinates $f_P([\xi])=[L\xi]$. The chain rule gives
$\partial_{\xi_a}\partial_{\xi_b}f(L\xi)
=\sum_{i,j}L_{ia}L_{jb}(\partial_i\partial_jf)(L\xi)$.
Contraction with $-W_{3,ab}$ and (HS.12) proves (HS.14).
Also $e^{2\pi i\ell^{\mathsf T}L\xi}=\chi_{L^{\mathsf T}\ell}(\xi)$;
inserting this index in (HS.6) proves equality of the eigenvalues.

The lattice kernel is $\mathbb Z u\oplus\mathbb Z u$.
The characters constant on the connected kernel annihilate its real
tangent space, so their indices $(m,n)$ obey
$u^{\mathsf T}m=u^{\mathsf T}n=0$. Conversely these equalities make
the character constant there. Every integral zero-sum triple
$m=(m_1,m_2,m_3)$ is
$P^{\mathsf T}(m_1,-m_3)^{\mathsf T}$, and
$\operatorname{im}_{\mathbb Z}Q_{\rm latt}^{\mathsf T}
=\operatorname{im}_{\mathbb Z}P^{\mathsf T}$ because $B_0$ is
unimodular. Thus the stated indices are exactly
$L^{\mathsf T}\mathbb Z^4$. Injectivity of $L^{\mathsf T}$
preserves the number of target characters at each eigenvalue.
\end{proof}

The compatibility restriction still leaves genuine source-metric
freedom over a fixed target metric. Indeed, using the original
$R=P^*C^{-1}$, $C=PP^*$ and $X=(R\;\;u)$, all positive $H_3$
with $PH_3P^*=G$ are
\[
 H_3=X\begin{pmatrix}G&z\\z^*&z^*G^{-1}z+\tau\end{pmatrix}X^*,
 \qquad z\in\mathbb C^2,\quad \tau>0.                         \tag{HS.17}
\]
To prove the assertion, $PX=(I_2\;\;0)$ and
$X^{-1}=\left(\begin{smallmatrix}P\\u^*/3\end{smallmatrix}\right)$
show that the top-left block of $X^{-1}H_3X^{-*}$ is $G$.
For any remaining blocks $z,d$, its quadratic form is
\[
 (v+G^{-1}zw)^*G(v+G^{-1}zw)
       +(d-z^*G^{-1}z)|w|^2.
\]
It is positive exactly when $d=z^*G^{-1}z+\tau$ with $\tau>0$.
This proves (HS.17) in both directions and its five real parameters.
The injectivity of inversion makes distinct $H_3$ distinct compatible
source metrics; (HS.11)--(HS.14) prove that the quotient metric and
quotient spectrum stay fixed. This specifies exactly what the
remaining observation fibre means geometrically.

\begin{proposition}[Native Euclidean observation and expanding dual spectrum]
\label{q444:native-period-spectrum}
Equip the source with its native Euclidean metric in the retained
coordinates: $h_3=I_3$, $g_{\rm phys}=I_6$. Then
\[
 W_3=\mathcal N^{-1}\mathcal N^{-\mathsf T},\qquad
 H_3=2I_3,\qquad G=2C,\qquad
 C=\begin{pmatrix}2&-1\\-1&2\end{pmatrix}.                     \tag{HS.18}
\]
The induced target metric is $h_2=C^{-1}$, not $I_2$.
The source eigenvalues, with no correction suppressed, are
\[
 \lambda_{m,n}
 =4\pi^2\left(\|m\|^2+
       \|Y^{-\mathsf T}(n-A^{\mathsf T}m)\|^2\right),
       \qquad m,n\in\mathbb Z^3.                            \tag{HS.19}
\]
For $y>K+1$, its smallest positive eigenvalue is exactly
\[
 \lambda_1^{(3)}
 =4\pi^2\min_{0\ne n\in\mathbb Z^3}\|Y^{-\mathsf T}n\|^2,
 \qquad
 \frac{4\pi^2}{(y+K)^2}\leq\lambda_1^{(3)}
                  \leq\frac{4\pi^2}{(y-K)^2}.                \tag{HS.20}
\]
For the induced target metric, its smallest positive eigenvalue is
\[
 \lambda_1^{(2)}
 =4\pi^2\min_{\substack{0\ne n\in\mathbb Z^3\\u^{\mathsf T}n=0}}
                \|Y^{-\mathsf T}n\|^2,\qquad
 \frac{8\pi^2}{(y+K)^2}\leq\lambda_1^{(2)}
                  \leq\frac{8\pi^2}{(y-K)^2}.                \tag{HS.21}
\]
In particular the exact leading constants are
\[
 \lambda_1^{(3)}\sim\frac{4\pi^2}{y^2}
   =\frac{16\pi^4}{(\log|t|)^2},\qquad
 \lambda_1^{(2)}\sim\frac{8\pi^2}{y^2}
   =\frac{32\pi^4}{(\log|t|)^2}.                              \tag{HS.22}
\]
The source volume is the retained polynomial
\[
 \operatorname{Vol}(T_3,I)
 =\det Y
 =y\left[y^2+
 y\operatorname{tr}\bigl((\operatorname{Im}C_2)B_0^{-1}\bigr)
 +\det(\operatorname{Im}C_2)\right]\sim y^3.                  \tag{HS.23}
\]
\end{proposition}
\begin{proof}
In (HS.1) the physical cometric is $T=I_6$, so
$H_3=\mathcal J\mathcal J^*=2I_3$. Multiplication by $P,P^*$
and (HS.11) prove all of (HS.18). The exact inverse transpose is
\[
 \mathcal N^{-\mathsf T}
 =\begin{pmatrix}I_3&0\\-Y^{-\mathsf T}A^{\mathsf T}
                              &Y^{-\mathsf T}\end{pmatrix},
\]
which gives (HS.19). The singular values of $Y$ lie between
$y-K$ and $y+K$, since
$(y-K)\|v\|\leq\|Yv\|\leq(y+K)\|v\|$.
The same singular-value bounds for the inverse transpose give
\[
 \frac{\|n\|}{y+K}\leq\|Y^{-\mathsf T}n\|
                             \leq\frac{\|n\|}{y-K}.           \tag{HS.24}
\]
For source modes with $m\ne0$, (HS.19) is at least $4\pi^2$.
The mode $m=0,n=e_1$ has eigenvalue at most
$4\pi^2/(y-K)^2<4\pi^2$. Thus a minimizer has $m=0$.
The minimum exists because the inverse-transpose quadratic form is
positive definite and only finitely many integer vectors have bounded
value. Nonzero integer vectors have norm at least one, and a unit
vector exists, giving (HS.20).

Target modes have $m=P^{\mathsf T}a$, $n=Q_{\rm latt}^{\mathsf T}b$
for $a,b\in\mathbb Z^2$. A nonzero integer zero-sum triple has
squared norm at least two: it has at least one positive and at least
one negative integer coordinate. The value two is attained by
$(1,-1,0)^{\mathsf T}$. Therefore $a\ne0$ costs at least $8\pi^2$,
whereas a mode with $a=0$ and a shortest nonzero zero-sum $n$ has
eigenvalue at most $8\pi^2/(y-K)^2<8\pi^2$.
The complete index characterization in the preceding theorem now
proves (HS.21). Multiplying the bounds by $y^2$ and taking the limit
proves (HS.22), with the original definition of $y$.

Finally $D_2=\operatorname{Im}C_2$ is real and
$Y=yI_3+SD_2Q_{\rm latt}$. Since $Q_{\rm latt}S=B_0^{-1}$,
the identity
$\det(I+UV)=\det(I+VU)$ gives
\[
 \det Y=y^3\det(I_2+y^{-1}D_2B_0^{-1}).
\]
For completeness that determinant identity follows by taking the
determinant of
$\left(\begin{smallmatrix}I&U\\-V&I\end{smallmatrix}\right)$
and eliminating first the upper-right block or first the lower-left
block using block triangular matrices with determinant one.
Expansion of the two-by-two determinant and $\det B_0=1$ give
the displayed polynomial in (HS.23). Its sign is positive:
$yI_3+rSD_2Q_{\rm latt}$ is invertible for $0\leq r\leq1$ by
the same bound $y-rK>0$, so its determinant has the positive sign
of $y^3$. Formula (HS.7) and boundedness of $D_2$ finish the proof.
\end{proof}

Equations (HS.18)--(HS.23) exhibit the exact relationship rather than
a nonidentity alone: the constant observation is twice the native
physical cometric and twice its quotient cometric, while the unchanged
period map enters the dual lattice in (HS.19). The periods expand,
the volume grows, and both displayed positive spectral gaps tend to
zero. No period coefficient was replaced and no metric scale was
divided out. At fixed periods, $H_3$ uniquely specifies the
Hermitian-compatible source metric; at fixed $G$ the five-parameter
family (HS.17) specifies its precise residual freedom. When periods
are also allowed to vary, the constant $H_3=2I_3$ alone cannot recover
their marked cometric; (HS.8) proves that the full marked spectrum
does recover it. These are geometric and spectral assertions, not
a selection of a preferred dynamical metric or a field-theoretic mass.

\paragraph{Primary reference conventions.}
The nonnegative Laplacian and integer-valued dual-lattice convention
agree with A.~Gogolev, \emph{Spectra of flat tori}, Section~2
(27 November 2017), available at
\url{https://andreygogolev.wordpress.com/2017/11/27/spectra-of-flat-tori/}.
V.~Ivrii, \emph{Partial Differential Equations}, Appendix~4.B,
\url{https://www.math.utoronto.ca/courses/apm346h1/20181/PDE-textbook/Chapter4/S4.B.html},
instead places $2\pi$ in the definition of the dual lattice;
his dual is exactly $2\pi\mathcal N^{-\mathsf T}\mathbb Z^6$ here.
The horizontal-isometry definition and uniqueness of quotient metrics
are those in U.~Lang, \emph{Lecture Notes on Riemannian Geometry},
16 June 2020, Definition~4.1 and Proposition~4.4, printed pages~45
and~47,
\url{https://metaphor.ethz.ch/x/2020/fs/401-3532-08L/sc/DG2_16June2020.pdf}.
All specialized matrix identities and period/spectral asymptotics
used here have been proved above in the retained coordinates.

\subsection{The unrestricted and constrained scalar equations}

Equation~\eqref{q444:N15} gives the exact normal covector at $\rho=0$:
\begin{equation}\label{q444:normal-covector}
 d_{\rho}N_{15}(\delta\rho)
 =-2\operatorname{Re}(\delta\rho\,\chi^*z).
\end{equation}
It need not vanish. For $G=I_2$, $z=\chi=(1,0)^{\mathsf T}$,
$\lambda=3$, $\rho=t\in\mathbb R$,
\[
 N_{15}(t)=1-2t-3t^2.
\]
At $t=0$ this is a positive-definite quaternionic matrix:
$w=(0,1)^{\mathsf T}$ and the Schur excess is $3-1-1=1$.
Thus the nonzero normal covector occurs even in the interior of the
positive cone. This alone is a statement about the cubic; to determine
motion one must also retain the specified kinetic form.

\begin{theorem}[Normal force in the supplied scalar model]\label{q444:force}
Use the triality packet's fourteen-real-dimensional scalar variables and
potential
\begin{equation}\label{q444:potential}
 \begin{gathered}
 \Phi=\Phi^*,\quad\operatorname{tr}\Phi=0,\quad b\in\mathbb C^3,
 \quad s=\operatorname{tr}\Phi^2+2b^*b,\\
 V(\Phi,b)=\frac{m^2}{2}s+\frac{\Gamma}{4}s^2
                       -\kappa(\det\Phi-b^*\Phi b),\\
 L=\tfrac12\operatorname{tr}(\dot\Phi^2)+\dot b^*\dot b-V,
 \quad m^2\geq0,\quad\Gamma>0,\quad\kappa\in\mathbb R.
 \end{gathered}
\end{equation}
This is the explicitly specified spatially constant scalar mechanics with
gauge fields omitted, not a reduction claimed for arbitrary gauge fields.
For its solutions, with $\rho=u^*b$,
\begin{equation}\label{q444:rho-eom}
 \ddot b=-[(m^2+\Gamma s)I_3+\kappa\Phi]b,
 \qquad
 \ddot\rho=-(m^2+\Gamma s)\rho-\kappa u^*\Phi b.
\end{equation}
The phase-space constraint $\rho=\dot\rho=0$ is not invariant for
$\kappa\neq0$. At initial data
\[
 \Phi_0=\operatorname{diag}(1,-1,0),\quad
 b_0=(1,-1,0)^{\mathsf T},\quad \dot\Phi_0=0,\quad\dot b_0=0,
\]
one has $\ddot\rho(0)=-2\kappa$.
\end{theorem}
\begin{proof}
For a variation $\delta b$,
\[
 \delta_bV=2\operatorname{Re}
       \bigl(\delta b^*[(m^2+\Gamma s)I_3+\kappa\Phi]b\bigr).
\]
The kinetic variation has the same factor two, giving the first equation
in \eqref{q444:rho-eom}; multiply by $u^*$ for the second.
At the displayed initial data $u^*b_0=0$ but
$\Phi_0b_0=(1,1,0)^{\mathsf T}$, so the acceleration is $-2\kappa$.
The equations have polynomial right-hand sides and hence a unique local
solution for these initial data. Taylor's theorem gives
$\rho(t)=-\kappa t^2+o(t^2)$, which leaves the constraint.
For comparison with the marked cubic, if
$X^{-1}\Phi X^{-*}=\left(\begin{smallmatrix}G&z\\z^*&\lambda\end{smallmatrix}\right)$,
then
\[
 \operatorname{tr}(C^{-1}G)+3\lambda=0,\qquad
 u^*\Phi b=3(z^*\chi+\lambda\rho).
\]
Thus at $\rho=0$ the normal acceleration is $-3\kappa z^*\chi$.
This factor three is fixed by the marking and the kinetic form; it is not
removed by rewriting the cubic.
\end{proof}

\begin{proposition}[Exact constrained alternative]\label{q444:constrained}
Let $\Pi_u=uu^*/3$, $Q_u=I_3-\Pi_u$. Restrict the Lagrangian
\eqref{q444:potential} to $b\in u^\perp$ and velocities in $u^\perp$.
Its Euler--Lagrange equations are
\begin{align}\label{q444:projected}
 \ddot b&=-(m^2+\Gamma s)b-\kappa Q_u\Phi b,\notag\\
 \ddot\Phi&=-(m^2+\Gamma s)\Phi
       +\kappa\bigl(\operatorname{adj}\Phi-bb^*\bigr)_0,
 \qquad T_0=T-\tfrac13(\operatorname{tr}T)I_3.
\end{align}
The reaction force added to the unrestricted $b$ equation is exactly
$\kappa\Pi_u\Phi b$. The constrained solution is also an unrestricted
solution exactly when this force vanishes throughout the solution
(or automatically when $\kappa=0$).
\end{proposition}
\begin{proof}
Allowed $b$ variations are projected by the orthogonal matrix $Q_u$, so
projecting \eqref{q444:rho-eom} gives the first equation. In particular
$u^*\ddot b=0$, and initial data tangent to the constraint remain there.
For the $\Phi$ equation, variations are trace-free Hermitian and
\[
 \delta_\Phi V=
 \operatorname{tr}\bigl([(m^2+\Gamma s)\Phi
              -\kappa\operatorname{adj}\Phi+\kappa bb^*]\delta\Phi\bigr).
\]
Projection onto trace-free matrices gives the second equation. The
unrestricted $\Phi$ equation is unchanged. Subtracting the two $b$
equations gives the stated reaction, proving the last equivalence.
\end{proof}

The unrestricted model retains all six real coordinates of $b$; the
constrained model retains four and supplies the displayed reaction force.
Neither is excluded by the failure of the other identification. The old
complex sector $b=\dot b=0$ \emph{is} invariant by \eqref{q444:rho-eom}.
By contrast, a nonzero fixed complex subspace of $\mathbb C^3$ invariant
under all possible Hermitian $\Phi$ must be the whole space: traceless
Hermitian matrices together with $I_3$ complex-linearly span every matrix
unit, and the matrix units applied to a nonzero vector span $\mathbb C^3$.
Thus a four-real-dimensional complex-linear transverse sector cannot retain all
these unrestricted couplings without the explicit constraint or an
additional restriction on $\Phi$.

\subsection{Preserved interacting sectors of the unrestricted equations}

The failure of the whole slice to be invariant does not exclude invariant
motions within it. We now construct such motions using the original
equations, not by adding the reaction in Proposition~\ref{q444:constrained}.
Write $a(t)=m^2+\Gamma s(t)$, and retain the trace-free Hermitian variable
$\Phi$, the standard metric on $b\in\mathbb C^3$, and $u^*u=3$.

\begin{lemma}[Global existence in the scalar mechanics]\label{q444:global}
Every initial position and velocity in \eqref{q444:potential} has a unique
solution for all real time. Its conserved energy is
\[
 E=\tfrac12\operatorname{tr}\dot\Phi^2+\dot b^*\dot b+V(\Phi,b).
\]
\end{lemma}
\begin{proof}
The unrestricted equations for both variables are
\begin{equation}\label{q444:full-eom}
 \ddot\Phi=-a\Phi+\kappa(\operatorname{adj}\Phi-bb^*)_0,
 \qquad \ddot b=-ab-\kappa\Phi b.
\end{equation}
They form a polynomial first-order system on positions and velocities,
so local existence and uniqueness hold. Put $t_\Phi=\operatorname{tr}\Phi^2$
and $n_b=b^*b$. Direct differentiation of the kinetic energy gives
\[
 -a\operatorname{tr}(\dot\Phi\Phi)
 +\kappa\operatorname{tr}(\dot\Phi\operatorname{adj}\Phi)
 -\kappa b^*\dot\Phi b
 -2a\operatorname{Re}(\dot b^*b)
 -2\kappa\operatorname{Re}(\dot b^*\Phi b).
\]
The trace projection drops out because $\operatorname{tr}\dot\Phi=0$.
Differentiating the displayed potential gives precisely the negative of
this expression, including the two terms from differentiating $b^*\Phi b$.
Thus $E$ is constant.

For completeness, a bound in the unchanged scalar coordinates suffices
to rule out escape. The three real eigenvalues of $\Phi$ give
$|\det\Phi|\leq(t_\Phi/3)^{3/2}$ by the arithmetic--geometric mean
inequality applied to their squares. Cauchy--Schwarz and the Frobenius
norm give $|b^*\Phi b|\leq n_b\sqrt{t_\Phi}$. Since
$s=t_\Phi+2n_b$, these are the explicit inequalities
\[
 |\det\Phi-b^*\Phi b|
 \leq (t_\Phi/3)^{3/2}+n_b\sqrt{t_\Phi}
 \leq \left(\frac1{3\sqrt3}+\frac12\right)s^{3/2}.
\]
This bound is not asserted to be sharp and does not change $V$.
Consequently
\[
 V\geq \frac{\Gamma}{4}s^2+\frac{m^2}{2}s
       -|\kappa|\left(\frac1{3\sqrt3}+\frac12\right)s^{3/2}
       \longrightarrow+\infty\quad(s\longrightarrow\infty).
\]
The energy identity bounds positions; the global lower bound for $V$
then bounds velocities. A polynomial vector field and its derivative are
bounded on this compact phase-space set. A finite endpoint therefore has
a limiting initial state to which local existence applies, extending the
solution. Neither a finite forward nor a finite backward endpoint occurs.
\end{proof}

\begin{theorem}[Fixed-subspace invariant sectors]\label{q444:invariant-W}
Fix a complex subspace $W\subset u^\perp$ of dimension $r=1$ or $2$, and
let $E_W$ be its orthogonal projection in the original metric. Define
\[
 \mathcal T_W=\{(\Phi,b):\Phi=\Phi^*,\ \operatorname{tr}\Phi=0,
                       \ [\Phi,E_W]=0,\ b\in W\}.
\]
With tangent initial velocities, its tangent bundle $T\mathcal T_W$ is
invariant under the full equations \eqref{q444:full-eom} for all real time.
For each fixed $W$, the configuration dimensions are respectively $6$
and $8$, and the phase-space dimensions are $12$ and $16$. Every such
solution satisfies
\[
 \rho=0,\qquad u^*\Phi b=0,\qquad \kappa(uu^*/3)\Phi b=0
\]
identically, without any constraint force.
\end{theorem}
\begin{proof}
For Hermitian $\Phi$, preservation of $W$ also preserves $W^\perp$:
$\langle\Phi x,w\rangle=\langle x,\Phi w\rangle=0$ for
$x\in W^\perp$, $w\in W$. This is equivalent to $[\Phi,E_W]=0$.
The trace-free adjugate identity is
\begin{equation}\label{q444:adj-tracefree}
 \operatorname{adj}\Phi=\Phi^2-\tfrac12(\operatorname{tr}\Phi^2)I_3.
\end{equation}
Indeed the characteristic polynomial is
$z^3-\tfrac12(\operatorname{tr}\Phi^2)z-\det\Phi$;
Cayley--Hamilton and the adjugate identity prove
\eqref{q444:adj-tracefree} when $\Phi$ is invertible, and polynomial
continuity proves it at singular $\Phi$ as well. Thus the adjugate
commutes with $E_W$. So does $bb^*$: it takes $W$ to $W$ and vanishes
on $W^\perp$. Scalar trace subtraction preserves this property.
Both accelerations in \eqref{q444:full-eom} consequently remain in the
specified linear spaces. The position derivatives do so when the
velocities satisfy $[\dot\Phi,E_W]=0$, $\dot b\in W$.
The restricted polynomial vector field has a local solution; its inclusion
is a solution of the full system. Uniqueness and Lemma~\ref{q444:global}
give the claimed global invariance.

The Hermitian blocks on $W$ and $W^\perp$ have real dimensions $r^2$
and $(3-r)^2$. The trace imposes one real equation, and $b$ contributes
$2r$ real coordinates. Hence
$\dim_{\mathbb R}\mathcal T_W=r^2+(3-r)^2-1+2r$.
Velocities double this dimension. Finally $b,\Phi b\in W\subset u^\perp$
proves each asserted zero contraction. The parameter $W$ is fixed, not
an additional moving coordinate: for $r=1$ the choices form
$\mathbb P(u^\perp)\cong\mathbb{CP}^1$, and for $r=2$ the only choice is
$W=u^\perp$.
\end{proof}

Here is an explicit interacting member that also retains a nonzero marked
cross-coordinate $z$. Keep
\[
 v=(1,-1,0)^{\mathsf T},\quad w_0=(1,1,-2)^{\mathsf T},\quad
 E=vv^*/2,\quad F=w_0w_0^*/6,
\]
and take
\begin{equation}\label{q444:nonzero-sector}
 W=\mathbb Cv,\quad \Phi_0=E-F+w_0u^*+uw_0^*,\quad b_0=v,
 \quad \dot\Phi_0=0,\quad\dot b_0=0.
\end{equation}
The vectors $u,v,w_0$ are orthogonal with squared lengths $3,2,6$.
Thus $\operatorname{tr}\Phi_0=0$, $\Phi_0v=v$,
$\operatorname{tr}\Phi_0^2=38$, $s_0=42$, and
\[
 b_0^*\Phi_0b_0=2,\qquad
 \ddot b_0=-(m^2+42\Gamma+\kappa)v.
\]
For example $\kappa>0$ gives nonstationary interacting motion. In the
original marked coordinates, not a newly chosen basis,
\[
 \chi_0=R^*b_0=(1,0)^{\mathsf T},\qquad
 z_0=P\Phi_0u/3=(0,3)^{\mathsf T},\qquad\lambda_0=0.
\]
Hence $z_0\ne0$ while $z_0^*\chi_0=0$. Preservation of the scalar slice
does not require $z=0$ for every solution. The theorem constructs exact
families; it does not assert they exhaust all invariant initial data.

\begin{proposition}[Tangency through third order is insufficient]\label{q444:fourth-jet}
For any unrestricted solution put $c=u^*\Phi b$ and
$h=u^*\dot\Phi\dot b$. The exact identities are
\begin{align}\label{q444:jets}
 \ddot c&=-2ac-\frac{\kappa s}{3}\rho+2h,\notag\\
 \rho^{(4)}&=\left(-\ddot a+a^2+\frac{\kappa^2s}{3}\right)\rho
             -2\dot a\dot\rho+3a\kappa c-2\kappa h.
\end{align}
There are initial data with $\rho^{(j)}(0)=0$ for $0\leq j\leq3$
but $\rho^{(4)}(0)=-12\kappa$.
\end{proposition}
\begin{proof}
Equation~\eqref{q444:adj-tracefree} gives, with
$t_\Phi=\operatorname{tr}\Phi^2$ and $n_b=b^*b$,
\[
 (\operatorname{adj}\Phi-bb^*)_0
   =\Phi^2-bb^*+\frac{n_b-t_\Phi}{3}I_3,
 \qquad
 (\operatorname{adj}\Phi-bb^*)_0b=\Phi^2b-\frac{s}{3}b.
\]
Differentiate $c$ twice and substitute \eqref{q444:full-eom}:
\[
 \ddot c=u^*(\ddot\Phi b+2\dot\Phi\dot b+\Phi\ddot b)
         =-2ac-\frac{\kappa s}{3}\rho+2h.
\]
The two terms $\kappa u^*\Phi^2b$ cancel exactly. Twice differentiating
$\ddot\rho=-a\rho-\kappa c$ and substituting this identity proves the
second formula. In particular, at $\rho=\dot\rho=0$, vanishing second
and third derivatives requires $c=\dot c=0$ when $\kappa\ne0$;
then the fourth derivative is $-2\kappa h$.

For the promised initial data use the same $u,v,w_0$ as above, but now set
\[
 \Phi_0=0,\quad b_0=v,\quad\dot b_0=w_0,\quad
 \dot\Phi_0=(uw_0^*+w_0u^*)/3.
\]
This matrix is trace-free Hermitian, $\dot\Phi_0v=0$ and
$\dot\Phi_0w_0=2u$. Consequently
$\rho_0=\dot\rho_0=c_0=\dot c_0=0$, whereas $h_0=6$.
Formula~\eqref{q444:jets} gives $\rho^{(4)}(0)=-12\kappa$, and Taylor's
theorem gives $\rho(t)=-(\kappa/2)t^4+o(t^4)$.
For $\kappa\ne0$ this solution leaves the slice despite the lower-order
tangencies. For $\kappa=0$, by contrast, the full locus
$\rho=\dot\rho=0$ is invariant because its equation is the homogeneous
linear equation $\ddot\rho=-a(t)\rho$ with zero initial data.
\end{proof}

\subsection{Periodic motions outside every fixed-subspace family}

The families in Theorem~\ref{q444:invariant-W} do not exhaust the motions
inside $\rho=0$. The following exact relative equilibria rotate the
instantaneous kernel of $\Phi$ while retaining all the original couplings.

\begin{theorem}[A moving-kernel family]\label{q444:moving-family}
Choose $\Gamma>0$, $\kappa\ne0$, and real $x,a$ satisfying
\[
 0<\frac{x}{\kappa}\leq\frac1{6\Gamma},\qquad
 a>\frac{\kappa x}{3}.
\]
Choose $B,d>0$ with
\begin{equation}\label{q444:moving-parameters}
 B^2+d^2=\frac{3ax}{\kappa}-x^2,
 \qquad m^2=a\left(1-\frac{6\Gamma x}{\kappa}\right).
\end{equation}
The right-hand sides give a positive sum and $m^2\geq0$.
Retain $u,v,w_0$ from \eqref{q444:nonzero-sector}, and put
\[
 J_0=\begin{pmatrix}0&-1&1\\1&0&-1\\-1&1&0\end{pmatrix},
 \quad \Pi=uu^*/3,\quad Q_u=I_3-\Pi,\quad
 \Omega=\sqrt{a/3}\,J_0,
\]
\begin{equation}\label{q444:rotation}
 O(t)=\Pi+\cos(\sqrt a\,t)Q_u
                 +\frac{\sin(\sqrt a\,t)}{\sqrt3}J_0.
\end{equation}
Define the original-coordinate initial data and their motion by
\begin{align}\label{q444:moving-data}
 b_0&=\frac B{\sqrt2}v,\qquad
 \Phi_0=-\frac x6w_0w_0^*+\frac x3uu^*
          +\frac d{\sqrt{18}}(w_0u^*+uw_0^*),\notag\\
 b(t)&=O(t)b_0,\qquad \Phi(t)=O(t)\Phi_0O(t)^*.
\end{align}
These are global periodic solutions of \eqref{q444:full-eom}. They satisfy
\[
 s=\frac{6ax}{\kappa},\qquad m^2+\Gamma s=a,\qquad
 \Phi(t)b(t)=0,\qquad\rho(t)=0.
\]
No fixed complex $W\subset u^\perp$ contains this trajectory in
$T\mathcal T_W$.
\end{theorem}
\begin{proof}
Direct multiplication gives $J_0^*=-J_0$, $J_0u=0$,
$J_0^2=-3Q_u$, and $\Pi J_0=J_0\Pi=0$.
Consequently $O(t)^*O(t)=I_3$, $O(t)u=u$, and
$\dot O=\Omega O=O\Omega$. Its determinant is one by continuity from
$O(0)=I_3$. This proves that the displayed rotation is in the original
Euclidean metric, with its frequency $\sqrt a$ unchanged.

For the following verification only, use the explicit orthogonal matrix
\[
 K=\left(\frac v{\sqrt2},\frac{w_0}{\sqrt6},\frac u{\sqrt3}\right),
 \quad K^*K=I_3,
\]
and retain its inverse $K^*$ and all its factors. It gives
\[
 K^*\Phi_0K=
 H=\begin{pmatrix}0&0&0\\0&-x&d\\0&d&x\end{pmatrix},
 \quad K^*b_0=Be_1,\quad
 K^*\Omega K=A=\begin{pmatrix}0&-\sqrt a&0\\\sqrt a&0&0\\0&0&0\end{pmatrix}.
\]
Thus $s=2(x^2+d^2+B^2)=6ax/\kappa$, as asserted.
In these exact coordinates $HBe_1=0$ and
\[
 (\operatorname{adj}H-B^2e_1e_1^*)_0
    =\operatorname{diag}(-s/3,s/6,s/6).
\]
The required accelerations are explicitly
\[
 [A,[A,H]]=
 \begin{pmatrix}-2ax&0&0\\0&2ax&-ad\\0&-ad&0\end{pmatrix}
 =-aH+\kappa\operatorname{diag}(-s/3,s/6,s/6),
 \qquad A^2Be_1=-aBe_1.
\]
Conjugation by $K$ therefore verifies both equations at $t=0$.
The full force commutes with orthogonal conjugation on $\Phi$ and the
same orthogonal action on $b$: the adjugate, trace subtraction, norms,
and product $\Phi b$ all transform in this way. Since $O$ commutes with
$\Omega$, conjugation by $O(t)$ verifies both equations at every time,
not merely to a finite Taylor order. This also proves the stated kernel
and normal-coordinate identities.

At times $0$ and $\pi/(2\sqrt a)$, the two values of $b$ are nonzero
multiples of $v$ and $w_0$. They span $u^\perp$ over $\mathbb C$.
Any fixed $W$ containing the whole trajectory would consequently equal
$u^\perp$. But
$u^*\Phi_0w_0=\sqrt{18}\,d\ne0$, so $\Phi_0$ does not preserve this
plane. This rules out every fixed-$W$ family in the earlier theorem.
\end{proof}

A particularly small exact instance is
\begin{equation}\label{q444:periodic-instance}
 \begin{gathered}
 \Gamma=m^2=1,\quad\kappa=6,\quad x=\tfrac12,\quad a=2,
 \quad B=d=\frac1{2\sqrt2},\\
 b_0=\frac14\begin{pmatrix}1\\-1\\0\end{pmatrix},\quad
 \Phi_0=\frac14\begin{pmatrix}1&1&1\\1&1&1\\1&1&-2\end{pmatrix}.
 \end{gathered}
\end{equation}
Here $s=1$ and the period of $b$ is $2\pi/\sqrt2$. Its original marked
coordinates at zero are $\chi_0=(1/4,0)^{\mathsf T}$,
$z_0=(0,1/4)^{\mathsf T}$, and $\lambda_0=1/6$.
Although $\det\Phi-b^*\Phi b=0$ along this solution, its cubic force
does not vanish: at zero it is
$6(\operatorname{adj}\Phi_0-b_0b_0^*)_0\ne0$.
It is the coupled equations, not removal of the cubic, that sustain the
rotating-kernel solution.

\begin{proposition}[Conserved charge and the precise gauge comparison]\label{q444:charge}
In the scalar mechanics the anti-Hermitian matrix
\begin{equation}\label{q444:moment}
 \mathcal J=[\Phi,\dot\Phi]+b\dot b^*-\dot b b^*
\end{equation}
is constant. Its trace-free part is the $SU(3)$ moment matrix, using the
pairing $\operatorname{tr}(\mathcal J\xi)$ with
$\xi\in\mathfrak{su}(3)$. On \eqref{q444:moving-data},
\[
 \mathcal J=-\frac s2\Omega,
 \qquad T=\frac{as}{2},\qquad
 E=m^2s+\frac{3\Gamma}{4}s^2.
\]
In particular \eqref{q444:periodic-instance} has
$\mathcal J=-\Omega/2\ne0$, $T=1$, $V=3/4$, and $E=7/4$.
This scalar solution is not a source-free homogeneous gauge solution
with all $\mathcal A_\mu=0$: it fails the $\mathcal A_0$ field equation.
\end{proposition}
\begin{proof}
Using \eqref{q444:full-eom}, polynomial commutation of $\Phi$ with its
adjugate, and commutation with the scalar trace correction, gives
\[
 \frac{d}{dt}[\Phi,\dot\Phi]=-\kappa[\Phi,bb^*],\qquad
 \frac{d}{dt}(b\dot b^*-\dot b b^*)=\kappa[\Phi,bb^*].
\]
Their sum vanishes. Under the infinitesimal simultaneous symmetry
$\delta\Phi=[\xi,\Phi]$, $\delta b=\xi b$, the kinetic momentum paired
with this variation is
\[
 \operatorname{tr}(\dot\Phi[\xi,\Phi])
       +2\operatorname{Re}(\dot b^*\xi b)
       =\operatorname{tr}(\mathcal J\xi).
\]
This proves the pairing, including its coefficient and sign.
In the explicit frame $K$, compute
\[
 [H,[A,H]]+Be_1(ABe_1)^*-(ABe_1)(Be_1)^*
       =-(x^2+d^2+B^2)A=-\frac s2 A.
\]
It commutes with $A$, so conjugation by $O(t)$ retains the same constant
matrix. Similarly
$\tfrac12\operatorname{tr}[A,H]^2+|ABe_1|^2=as/2$.
The determinant and $b^*\Phi b$ both vanish on the trajectory; substitution
of $a=m^2+\Gamma s$ gives the displayed energy.

For the gauge comparison retain the packet's covariant derivatives
$D_0\Phi=\dot\Phi+[\mathcal A_0,\Phi]$,
$D_0b=\dot b+\mathcal A_0b$. Varying $\mathcal A_0$ by
$\xi\in\mathfrak{su}(3)$ in the scalar kinetic density gives
$\operatorname{tr}(\mathcal J\xi)$ at $\mathcal A_0=0$.
With homogeneous fields and every $\mathcal A_i=0$, the electric field
vanishes and the gauge kinetic variation contributes zero. Thus its
Euler--Lagrange equation requires the trace-free part of $\mathcal J$ to
vanish. The exhibited matrix is already trace-free and nonzero, so this
particular zero-connection lift fails. This is the Gauss equation obtained
by variation before imposing temporal gauge; compare D.~Tong,
\emph{Lectures on Gauge Theory}, Section~2.2.1, equation~(2.27),
\url{https://www.damtp.cam.ac.uk/user/tong/gaugetheory/gt.pdf}.
Our matter contribution has been derived here in the retained
anti-Hermitian convention, not imported from a different convention.
\end{proof}

There is an exact moving-frame map, not an assertion that these theories
are unrelated. With $U(t)=O(t)^*$ the gauge transformation rule is
\[
 \Phi'=U\Phi U^{-1}=\Phi_0,\quad b'=Ub=b_0,\quad
 \mathcal A'_0=U\mathcal A_0U^{-1}-\dot U U^{-1}=\Omega.
\]
The covariant velocities remain
$D'_0\Phi_0=[\Omega,\Phi_0]$ and $D'_0b_0=\Omega b_0$.
The same nonzero charge and Gauss residual therefore remain; making the
ordinary coordinates static does not remove them. A gauge-field lift
would have to supply the electric contribution to Gauss' equation and
satisfy the remaining coupled equations. Neither its existence nor its
nonexistence follows from failure of the zero-connection lift, and no
quantum spectral-gap conclusion is inferred from this periodic motion.

\subsection{The complete scalar vacuum calculation with its kinetic metric}

This subsection validates the homogeneous scalar vacuum theorem in the
triality packet. It does not assert a new discovery of that phase diagram,
and does not include the packet's subsequent gauge-orbit claims.
Set $Q=\Theta(\Phi,b)$ with $\operatorname{tr}Q=0$, and retain
\[
 g(Y,Z)=\operatorname{Re}\operatorname{tr}(YZ),\qquad
 s(Q)=g(Q,Q)=\operatorname{tr}\Phi^2+2b^*b,\qquad
 n(Q)=N(Q)=\det\Phi-b^*\Phi b.
\]
The kinetic term is $g(\dot Q,\dot Q)/2$, not the unweighted sum of
squares of the independent entries. Explicitly,
$g(Y,Y)=\sum_iY_{ii}^2+2\sum_{i<j}|Y_{ij}|^2$ for Hermitian $Y$.

\begin{lemma}[Sharp bound and all equality cases]\label{q444:sharp}
For every trace-free quaternionic Hermitian $Q$,
\[
 54n(Q)^2\leq s(Q)^3.
\]
For $Q\ne0$, equality holds exactly at
$Q=x(3P-I_3)$, $x\in\mathbb R\setminus\{0\}$, where
$P=P^*=P^2$ has quaternionic rank one. There
$s=6x^2$ and $n=2x^3$.
\end{lemma}
\begin{proof}
Here are the spectral and determinant details in the stated conventions.
The real-valued function $v^*Qv$ on the unit sphere of $\mathbb H^3$
attains a maximum. Its Lagrange-multiplier equation gives
$\operatorname{Re}(w^*(Qv-v\lambda))=0$ for every $w$, with real
$\lambda=v^*Qv$. Taking $w=Qv-v\lambda$ proves $Qv=v\lambda$.
The quaternionic orthogonal complement of $v$ is invariant because
$v^*Qw=\lambda v^*w$. Iteration gives a unitary eigenbasis with real
eigenvalues $\lambda_1,\lambda_2,\lambda_3$.

For quaternionic matrices, the real trace is cyclic by
$\operatorname{Re}(pq)=\operatorname{Re}(qp)$, entry by entry. Moreover,
the original cubic \eqref{q444:moore} obeys
\begin{equation}\label{q444:newton-cubic}
 6n(Q)=(\operatorname{tr}Q)^3
       -3\operatorname{tr}(Q)\operatorname{tr}(Q^2)
       +2\operatorname{tr}(Q^3).
\end{equation}
To check it without changing the quaternionic order, in the ordered
entries of \eqref{q444:moore} the cubic trace is the sum of the diagonal
cubes, plus $3(d+f)|q_1|^2+3(a+f)|q_2|^2+3(a+d)|q_3|^2$,
plus $6\operatorname{Re}(q_1q_2q_3)$. Substitution gives each term of
six times \eqref{q444:moore}. Unitary conjugation therefore preserves
$s,n$, which are respectively $\sum_i\lambda_i^2$ and
$\lambda_1\lambda_2\lambda_3$ in the eigenbasis. Their sum is zero.
The exact polynomial identity
\[
 s^3-54n^2=2(\lambda_1-\lambda_2)^2
                (\lambda_1-\lambda_3)^2(\lambda_2-\lambda_3)^2
\]
follows by substituting $\lambda_3=-\lambda_1-\lambda_2$ and expanding.
It proves the bound. For nonzero $Q$, equality means a repeated
eigenvalue, hence spectrum $(2x,-x,-x)$ for some $x\ne0$.
Letting $P$ project onto the simple eigenline gives $Q=x(3P-I_3)$.
The converse follows from this same spectrum.
\end{proof}

\begin{theorem}[All scalar global minima]\label{q444:vacua}
For $\Gamma>0$, $m^2\geq0$, and $\kappa\ne0$, put
$m_{\mathrm c}^2=\kappa^2/(27\Gamma)$.
The global minima of $V=m^2s/2+\Gamma s^2/4-\kappa n$ are exactly:
\begin{itemize}
\item $Q=0$ if $m^2>m_{\mathrm c}^2$;
\item $Q=0$ and $Q=\frac{\kappa}{9\Gamma}(3P-I_3)$ for every
rank-one quaternionic Hermitian projector $P$, if $m^2=m_{\mathrm c}^2$;
\item $Q=x_*(3P-I_3)$ for every such $P$, if $0\leq m^2<m_{\mathrm c}^2$,
where
\[
 x_* =\frac{\kappa+\operatorname{sgn}(\kappa)
              \sqrt{\kappa^2-24\Gamma m^2}}{12\Gamma}.
\]
\end{itemize}
Each nonzero component is diffeomorphic to $\mathbb{HP}^2$.
For $\kappa=0$, the origin is the unique global minimum.
\end{theorem}
\begin{proof}
For $\kappa=0$, the positive quartic proves the assertion, including
$m^2=0$. Otherwise at each fixed original norm $r=\sqrt s$ the sharp
bound is attained with
$x=\operatorname{sgn}(\kappa)r/\sqrt6$ and $Q=x(3P-I_3)$.
For $r>0$, these are exactly its equality and sign cases. Thus the full
minimization is equivalent, with all fibres of the norm map accounted for,
to minimizing
\[
 f(r)=\frac{m^2}{2}r^2+\frac\Gamma4r^4
                   -\frac{|\kappa|}{3\sqrt6}r^3,\qquad r\geq0.
\]
Set $r_{\mathrm c}=2|\kappa|/(3\sqrt6\Gamma)$. Exact factorization gives
\[
 f(r)=\frac\Gamma4r^2(r-r_{\mathrm c})^2
                    +\frac{m^2-m_{\mathrm c}^2}{2}r^2.
\]
This proves the first two cases, including every zero. Below the transition,
$f(r_{\mathrm c})<0$, so a global minimum is nonzero. Put
$a_0=|\kappa|/\sqrt6$. The derivative is
$f'(r)=r(\Gamma r^2-a_0r+m^2)$, with roots
\[
 r_\pm=\frac{a_0\pm\sqrt{a_0^2-4\Gamma m^2}}{2\Gamma}.
\]
Below the transition its discriminant is strictly positive. If $m^2>0$,
$f'$ has signs $+,-,+$ across these two positive roots, so $r_+$ is the
only positive local minimum. If $m^2=0$, $r_-=0$, and the signs on
$(0,r_+)$ and $(r_+,\infty)$ are $-,+$. Coercivity and the negative
value exclude the origin in both cases. Converting $r_+$ to $x$ proves
the third case and the completeness of its matrix equality locus.

For the manifold statement, a right quaternionic line represented by
$v\ne0$ has orthogonal projector $P_v=v(v^*v)^{-1}v^*$.
This is unchanged under $v\mapsto vh$, $h\in\mathbb H^\times$,
and every rank-one Hermitian projector is the unique projection onto
its image line. The formula and its inverse line charts are smooth.
For a fixed nonzero $x$, the maps $P\mapsto x(3P-I_3)$ and
$Q\mapsto(Q/x+I_3)/3$ are mutually inverse. This proves the claimed
$\mathbb{HP}^2$ identification without any gauge quotient.
\end{proof}

\begin{theorem}[Full metric Hessian]\label{q444:vacuum-hessian}
At every nonzero minimum $Q_*=x(3P-I_3)$, the Hessian relative to $g$
has eigenvalues
\[
 \begin{array}{c|c}
 \text{eigenvalue}&\text{real multiplicity}\\\hline
 x(12\Gamma x-\kappa)&1\\
 0&8\\
 3\kappa x&5
 \end{array}
\]
Both nonzero eigenvalues are positive. The entire kernel is the tangent
space of the nonzero minimum manifold. At the origin the metric Hessian
is $m^2I$ in all fourteen directions.
\end{theorem}
\begin{proof}
The preceding spectral proof and cubic trace identity show that
quaternionic unitary conjugation preserves both $g$ and $V$.
Transport the Hessian isometrically to $Q_*=xD$, $D=\operatorname{diag}(2,-1,-1)$.
Every trace-free Hermitian variation has the unique, $g$-orthogonal decomposition
\[
 Y=yD+\begin{pmatrix}0&v^*\\v&0_2\end{pmatrix}
       +\begin{pmatrix}0&0\\0&B_0\end{pmatrix},
 \qquad y\in\mathbb R,\quad v\in\mathbb H^2,
 \quad B_0=B_0^*,\quad\operatorname{tr}B_0=0.
\]
Write $B_0=\left(\begin{smallmatrix}\beta&\omega\\\bar\omega&-\beta\end{smallmatrix}\right)$,
$\beta\in\mathbb R$, $\omega\in\mathbb H$. The respective dimensions are
$1,8,5$, and the retained norm is
\[
 g(Y,Y)=6y^2+2v^*v+\|B_0\|^2,
 \qquad\|B_0\|^2=2\beta^2+2|\omega|^2.
\]
For $h=x+\varepsilon y$, the full cubic, including every mixed direction,
is
\begin{equation}\label{q444:hessian-cubic}
 n(Q_*+\varepsilon Y)
 =2h^3+\varepsilon^2h(v^*v-\|B_0\|^2)
                                      +\varepsilon^3v^*B_0v.
\end{equation}
Indeed the diagonal and lower off-diagonal contributions are
$2h(h^2-\varepsilon^2\beta^2)-2h\varepsilon^2|\omega|^2$.
The upper entries contribute
$h\varepsilon^2(|v_1|^2+|v_2|^2)
+\varepsilon^3\beta(|v_1|^2-|v_2|^2)$, and the triangular term is
$2\varepsilon^3\operatorname{Re}(\bar v_1\omega v_2)$.
These sum to \eqref{q444:hessian-cubic}, preserving the quaternionic order.
Also
\[
 s(Q_*+\varepsilon Y)=6x^2+12\varepsilon xy
            +\varepsilon^2(6y^2+2v^*v+\|B_0\|^2).
\]
The second derivative of the cubic is therefore
$12xy^2+2xv^*v-2x\|B_0\|^2$. Differentiating the potential gives the
complete quadratic form before using stationarity:
\begin{align*}
 D^2V_{xD}[Y,Y]
 ={}&6(m^2+18\Gamma x^2-2\kappa x)y^2\\
 &+2(m^2+6\Gamma x^2-\kappa x)v^*v\\
 &+(m^2+6\Gamma x^2+2\kappa x)\|B_0\|^2.
\end{align*}
Every nonzero minimum satisfies $m^2-\kappa x+6\Gamma x^2=0$.
Thus this becomes
\[
 6x(12\Gamma x-\kappa)y^2+3\kappa x\|B_0\|^2.
\]
There are no mixed quadratic terms; polarization therefore gives no mixed
bilinear blocks. Comparison with $g(Y,Y)$ proves every listed eigenvalue
and multiplicity, with no rescaling of the kinetic metric.

At coexistence the positive values are
$\kappa^2/(27\Gamma)$ and $\kappa^2/(3\Gamma)$.
Below it, $12\Gamma x-\kappa=
\operatorname{sgn}(\kappa)\sqrt{\kappa^2-24\Gamma m^2}$ and $x$ has
the same sign as $\kappa$, proving positivity. In the line chart
\[
 P(z)=\frac1{1+z^*z}\begin{pmatrix}1&z^*\\z&zz^*\end{pmatrix},
 \qquad DP_0[\zeta]=\begin{pmatrix}0&\zeta^*\\\zeta&0_2\end{pmatrix},
\]
the derivative of $x(3P(z)-I_3)$ spans exactly the eight-dimensional
zero block. Isometric conjugation transports this to every $P$.
At the origin, the cubic and quartic contribute zero Hessian, giving
$m^2g$ directly. In particular, at $m^2=\kappa=0$ the origin is still the
unique minimum, although its Hessian vanishes in all directions.
\end{proof}

These are scalar Hessian values of a specified finite-dimensional model.
Their exact multiplicities do not assert a gauge-reduced physical spectrum
or an interacting quantum mass gap. The current calculation includes the
full tangent and normal blocks, not solely the radial restriction.

\subsection{Gauge kinetic sign and precise scope of this audit}

In the packet's gauge lift, $\mathcal A_\mu\in\mathfrak{su}(3)$ and
$D_\mu b=\partial_\mu b+\mathcal A_\mu b$ fix anti-Hermitian gauge
generators. Retain Minkowski signature $(+---)$, consistent with the
positive time-derivative scalar term in that lift. Write
$\|T\|^2=-\operatorname{tr}(T^2)$ for anti-Hermitian $T$, and
$E_i=\mathcal F_{0i}$, $B_i=\tfrac12\epsilon_{ijk}\mathcal F_{jk}$.
Raising the time-space index pair changes its sign and summing both
orders of each antisymmetric pair gives
\[
 \operatorname{tr}(\mathcal F_{\mu\nu}\mathcal F^{\mu\nu})
 =2\sum_i(\|E_i\|^2-\|B_i\|^2).
\]
Consequently the packet's coefficient $-1/(2g_c^2)$ gives negative electric
kinetic energy in these conventions. Keeping its anti-Hermitian fields
and scalar kinetic terms, the positive-energy gauge density is
\[
 \mathcal L_{\mathrm{gauge}}
 =+\frac1{2g_c^2}\operatorname{tr}
       (\mathcal F_{\mu\nu}\mathcal F^{\mu\nu})
 =\frac1{g_c^2}\sum_i(\|E_i\|^2-\|B_i\|^2).
\]
For an immediate check take $\mathcal F_{01}=\operatorname{diag}(i,-i,0)$
and all other independent components zero: the trace contraction is $4$,
so the two signs give $-2/g_c^2$ and $+2/g_c^2$, respectively. This sign
correction changes neither the scalar potential nor its derivative
calculated above. It does matter before interpreting the full gauge lift
as a positive-energy physical model.

The equal-diagonal thirteen-space in the triality packet remains a
different embedded slice: its intersection with
$\operatorname{Herm}_3(\mathbb C)$ has dimension seven, whereas the
common Lorentz slice contains all nine complex-Hermitian coordinates.
Their precise common ambient space is $\operatorname{Herm}_3(\mathbb H)$;
Theorem~\ref{q444:identification} identifies the \emph{Lorentz} constructions,
not the equal-diagonal one. No result here identifies these cone-boundary
spheres with the claimed integrable complex structure on $S^6$, assigns
new spinor periods and toroidal fillings, or proves a continuum interacting
Yang--Mills spectral gap. The exact next data are the transition maps and
lattice extension for the additional spinor, followed by a physical state
space and Hamiltonian comparison retaining the couplings above.

% Ready-to-input fragment; no canonical files are modified by this artifact.
% Uses the existing q444:potential and q444:moving-family data and standard
% theorem/proposition/lemma environments. All new calculations are proved here.
\subsection{The full homogeneous gauge equations and the literal periodic orbit}

Retain signature $(+---)$, $\mathcal A_\mu\in\mathfrak{su}(3)$,
$\Phi=\Phi^*$, $\operatorname{tr}\Phi=0$, and $b\in\mathbb C^3$.
Write $g^2>0$ for the gauge coupling in this subsection and keep
\[
 s=\operatorname{tr}\Phi^2+2b^*b,\qquad
 a=m^2+\Gamma s,\qquad
 V=\frac{m^2}{2}s+\frac{\Gamma}{4}s^2-\kappa(\det\Phi-b^*\Phi b).
\]
Thus $m^2\geq0$, $\Gamma>0$, and $\kappa\in\mathbb R$ are unchanged.
For a matrix $Z$ write
$Z_{\mathrm{tf}}=Z-\operatorname{tr}(Z)I_3/3$. The full density is
\begin{equation}\label{q444:full-gauge-density}
 \mathcal L=\frac1{2g^2}\operatorname{tr}
 (\mathcal F_{\mu\nu}\mathcal F^{\mu\nu})
 +\frac12\operatorname{tr}(D_\mu\Phi D^\mu\Phi)
 +(D_\mu b)^*D^\mu b-V,
\end{equation}
where $\mathcal F_{\mu\nu}=\partial_\mu\mathcal A_\nu-
\partial_\nu\mathcal A_\mu+[\mathcal A_\mu,\mathcal A_\nu]$,
$D_\mu\Phi=\partial_\mu\Phi+[\mathcal A_\mu,\Phi]$, and
$D_\mu b=\partial_\mu b+\mathcal A_\mu b$.
We consider spatially homogeneous fields on a finite-volume flat
$\mathbb T^3$, or locally on $\mathbb R^3$. A nonzero constant energy
density is not asserted to have finite total energy on $\mathbb R^3$.

\begin{proposition}[Every homogeneous field equation]\label{q444:gauge-eom}
Set
\[
 E_i=\dot{\mathcal A}_i+[\mathcal A_0,\mathcal A_i],\qquad
 F_{ij}=[\mathcal A_i,\mathcal A_j],\qquad
 X=\dot\Phi+[\mathcal A_0,\Phi],\qquad v=\dot b+\mathcal A_0b,
\]
and define the lower-index matter currents
\[
 J_0=[\Phi,X]+bv^*-vb^*,\qquad
 J_i=[\Phi,[\mathcal A_i,\Phi]]
                -bb^*\mathcal A_i-\mathcal A_i bb^*.
\]
The complete homogeneous Euler--Lagrange system is
\begin{align}
 D_0X-\sum_i[\mathcal A_i,[\mathcal A_i,\Phi]]
       +a\Phi-\kappa(\operatorname{adj}\Phi-bb^*)_{\mathrm{tf}}&=0,
       \label{q444:gauge-phi}\\
 D_0v-\sum_i\mathcal A_i^2b+(aI_3+\kappa\Phi)b&=0,
       \label{q444:gauge-b}\\
 \frac2{g^2}\sum_i[\mathcal A_i,E_i]-(J_0)_{\mathrm{tf}}&=0,
       \label{q444:full-gauss}\\
 D_0E_i-\sum_j[\mathcal A_j,[\mathcal A_j,\mathcal A_i]]
       -\frac{g^2}{2}(J_i)_{\mathrm{tf}}&=0
       \quad(i=1,2,3). \label{q444:gauge-spatial}
\end{align}
Here $D_0X=\dot X+[\mathcal A_0,X]$,
$D_0v=\dot v+\mathcal A_0v$, and similarly for $E_i$.
\end{proposition}
\begin{proof}
First retain the homogeneous density without deleting $\mathcal A_0$:
\[
 \mathcal L_{\mathrm{hom}}
 =-\frac1{g^2}\sum_i\operatorname{tr}E_i^2
   +\frac1{g^2}\sum_{i<j}\operatorname{tr}F_{ij}^2
   +\frac12\operatorname{tr}X^2+v^*v
   -\frac12\sum_i\operatorname{tr}[\mathcal A_i,\Phi]^2
   -\sum_i|\mathcal A_i b|^2-V.
\]
Both signs in the gauge terms follow by raising indices in $(+---)$
and counting the two orders of each antisymmetric pair.
For trace-free Hermitian $\delta\Phi$ and arbitrary $\delta b$,
\[
 \delta V=
 \operatorname{tr}\!\left[
 \bigl(a\Phi-\kappa(\operatorname{adj}\Phi-bb^*)_{\mathrm{tf}}\bigr)
 \delta\Phi\right]
 +2\operatorname{Re}\bigl(\delta b^*(aI_3+\kappa\Phi)b\bigr).
\]
Covariant integration by parts therefore gives
$D_\mu D^\mu\Phi+a\Phi-
\kappa(\operatorname{adj}\Phi-bb^*)_{\mathrm{tf}}=0$ and
$D_\mu D^\mu b+(aI_3+\kappa\Phi)b=0$. Homogeneity gives the first
two displayed equations, including the minus signs in their spatial sums.

For the gauge variation, antisymmetry gives
\[
 \delta\mathcal L_{\mathrm{gauge}}
 =\frac2{g^2}\operatorname{tr}
       ((D_\mu\delta\mathcal A_\nu)\mathcal F^{\mu\nu}).
\]
After integration by parts its coefficient is
$-2D_\mu\mathcal F^{\mu\nu}/g^2$.
The scalar coefficient of $\delta\mathcal A_\nu$ is
\[
 J^\nu=[\Phi,D^\nu\Phi]
           +b(D^\nu b)^*-(D^\nu b)b^*.
\]
Indeed cyclicity and anti-Hermiticity respectively give
\begin{align*}
 \operatorname{tr}(D^\nu\Phi[\delta\mathcal A_\nu,\Phi])
 &=\operatorname{tr}([\Phi,D^\nu\Phi]\delta\mathcal A_\nu),\\
 2\operatorname{Re}((D^\nu b)^*\delta\mathcal A_\nu b)
 &=\operatorname{tr}\bigl((b(D^\nu b)^*-(D^\nu b)b^*)
                        \delta\mathcal A_\nu\bigr).
\end{align*}
Variation over $\mathfrak{su}(3)$ consequently yields
\[
 D_\mu\mathcal F^{\mu\nu}=\frac{g^2}{2}(J^\nu)_{\mathrm{tf}}.
\]
For $\nu=0$, $\mathcal F^{i0}=E_i$ and $J^0=J_0$.
For $\nu=i$, $\mathcal F^{0i}=-E_i$,
$\mathcal F^{ji}=[\mathcal A_j,\mathcal A_i]$, and $J^i=-J_i$.
These substitutions give exactly \eqref{q444:full-gauss} and
\eqref{q444:gauge-spatial}. In particular the $\mathcal A_0$ equation
survives after choosing temporal gauge. This is the same variational
principle emphasized by D.~Tong, \emph{Lectures on Gauge Theory},
Section~2.2.1, equation~(2.27); the matter currents, signs, and coefficient
$g^2/2$ above have been derived in the present conventions.
\end{proof}

\begin{theorem}[No temporal-gauge lift retaining the literal orbit]
\label{q444:temporal-no-lift}
Keep exactly the moving-kernel scalar functions in
Theorem~\ref{q444:moving-family}, with $B,d>0$ and $x\neq0$.
There are no homogeneous spatial fields $\mathcal A_i(t)$ for which
these same scalar functions, together with $\mathcal A_0=0$, solve
the full system \eqref{q444:gauge-phi}--\eqref{q444:gauge-spatial}.
\end{theorem}
\begin{proof}
The scalar functions already solve the unrestricted scalar equations.
Subtracting them from \eqref{q444:gauge-phi}--\eqref{q444:gauge-b},
with $\mathcal A_0=0$, gives
\[
 \sum_i[\mathcal A_i,[\mathcal A_i,\Phi]]=0,\qquad
 \sum_i\mathcal A_i^2b=0.
\]
Taking the trace pairing with $\Phi$ and the Hermitian pairing with $b$
gives respectively
\[
 -\sum_i\operatorname{tr}[\mathcal A_i,\Phi]^2=0,\qquad
 -\sum_i|\mathcal A_i b|^2=0.
\]
Each $[\mathcal A_i,\Phi]$ is Hermitian, so both are sums of
nonnegative squared norms with the displayed common minus sign. Thus
$[\mathcal A_i,\Phi]=0$ and $\mathcal A_i b=0$ separately for every $i$.

To classify this stabilizer use the retained orthogonal matrix
$K=(v/\sqrt2,w_0/\sqrt6,u/\sqrt3)$ from the moving-family proof, and
put $r^2=x^2+d^2>0$. At time zero the data are
\[
 H=K^*\Phi_0K=\begin{pmatrix}0&0&0\\0&-x&d\\0&d&x\end{pmatrix},
 \qquad K^*b_0=Be_1,\qquad
 A=K^*\Omega K=
 \begin{pmatrix}0&-\sqrt a&0\\\sqrt a&0&0\\0&0&0\end{pmatrix}.
\]
An anti-Hermitian traceless matrix killing $Be_1$ has zero first row
and column and a lower $\mathfrak{su}(2)$ block. Commutation with
$\left(\begin{smallmatrix}-x&d\\d&x\end{smallmatrix}\right)$, whose
distinct eigenvalues are $\pm r$, forces that block to be its real
multiple times $i$. Therefore the stabilizer Lie algebra is
$\mathbb R\,i\Phi(t)$. The full stabilizer group is the circle
\[
 \left\{P_b+\cos\theta(I_3-P_b)
          +i\frac{\sin\theta}{r}\Phi(t):\theta\in\mathbb R\right\},
 \qquad P_b=\frac{b(t)b(t)^*}{B^2}.
\]
Indeed it fixes $b$, has eigenvalues $1,e^{i\theta},e^{-i\theta}$,
and exhausts the unitary commutant fixing the zero eigenspace pointwise.
Consequently
\[
 \mathcal A_i=i c_i(t)\Phi(t),\qquad c_i(t)\in\mathbb R.
\]
The original orbit has
\[
 [\Phi,\dot\Phi]=-r^2\Omega,\qquad
 b\dot b^*-\dot b b^*=-B^2\Omega,\qquad
 J_0=-(r^2+B^2)\Omega.
\]
For verification, $H^2=r^2\operatorname{diag}(0,1,1)$,
$HAH=0$, and hence $[H,[A,H]]=-r^2A$; the vector term
is $-B^2A$. Conjugation back by $O(t)K$, with the rotation commuting
with $\Omega$, gives the stated original-coordinate identities.
Now
\[
 F_{ij}=0,\qquad
 E_i=i\dot c_i\Phi+i c_i\dot\Phi,\qquad
 \sum_i[\mathcal A_i,E_i]
          =r^2\left(\sum_i c_i^2\right)\Omega.
\]
Thus the full Gauss residual, with precisely the convention in
\eqref{q444:full-gauss}, is
\[
 \boxed{\left(\frac{2r^2}{g^2}\sum_i c_i^2+r^2+B^2\right)\Omega.}
\]
Its coefficient is strictly positive and $\Omega\neq0$, proving the
obstruction even before the spatial gauge equations are used.

For completeness those equations force all $c_i$ to vanish as well:
here $J_i=F_{ij}=0$, so they require
$i(\ddot c_i\Phi+2\dot c_i\dot\Phi+c_i\ddot\Phi)=0$.
At a fixed time conjugate by the retained moving frame. The $(1,1)$
entries of $H$ and $[A,H]$ vanish, while that of $[A,[A,H]]$ is
$-2ax\neq0$. Hence $c_i=0$. The remaining Gauss residual is then
$(r^2+B^2)\Omega$, still nonzero.
\end{proof}

\begin{proposition}[Arbitrary temporal connection and exact Gauss inertia]
\label{q444:gauss-inertia}
For arbitrary homogeneous scalar and spatial-gauge positions and
velocities, define
\[
 J_{\mathrm{sc}}=
 [\Phi,\dot\Phi]+b\dot b^*-\dot b b^*,\qquad
 \mathscr I_{\Phi,b}Z=
 \bigl([\Phi,[\Phi,Z]]+bb^*Z+Zbb^*\bigr)_{\mathrm{tf}},
\]
\[
 \mathscr K_{\mathcal A}Z
       =-\sum_i[\mathcal A_i,[\mathcal A_i,Z]].
\]
Then Gauss' equation is exactly
\begin{equation}\label{q444:inertia-equation}
 \left(\mathscr I_{\Phi,b}+\frac2{g^2}\mathscr K_{\mathcal A}\right)
 \mathcal A_0
 =(J_{\mathrm{sc}})_{\mathrm{tf}}
       -\frac2{g^2}\sum_i[\mathcal A_i,\dot{\mathcal A}_i].
\end{equation}
The operator on the left is self-adjoint and positive semidefinite for
$\langle Z,W\rangle=-\operatorname{tr}(ZW)$ on $\mathfrak{su}(3)$.
Its kernel is exactly
\[
 \{Z:[Z,\Phi]=0,\ Zb=0,\ [Z,\mathcal A_i]=0\text{ for every }i\}.
\]
The right-hand side always annihilates that kernel. Thus Gauss' equation
has a pointwise solution for $\mathcal A_0$, unique modulo this kernel.
This pointwise statement asserts neither a smooth selection across
kernel-rank changes nor a solution of the remaining field equations.
\end{proposition}
\begin{proof}
Expansion gives
$J_0=J_{\mathrm{sc}}-\mathscr I_{\Phi,b}\mathcal A_0$
after trace-free projection, and
\[
 \sum_i[\mathcal A_i,E_i]
 =\sum_i[\mathcal A_i,\dot{\mathcal A}_i]
          +\mathscr K_{\mathcal A}\mathcal A_0.
\]
Substitution proves \eqref{q444:inertia-equation}. Cyclicity gives
\[
 \langle Z,\mathscr I_{\Phi,b}Z\rangle
       =\operatorname{tr}[Z,\Phi]^2+2|Zb|^2,\qquad
 \langle Z,\mathscr K_{\mathcal A}Z\rangle
       =\sum_i\|[\mathcal A_i,Z]\|_{\mathfrak{su}}^2,
\]
where $\|T\|_{\mathfrak{su}}^2=-\operatorname{tr}T^2$.
More explicitly, cyclicity gives the bilinear identities
\[
 \langle Z,\mathscr I_{\Phi,b}W\rangle
 =\operatorname{tr}([Z,\Phi][W,\Phi])
                  +2\operatorname{Re}((Zb)^*Wb),\qquad
 \langle Z,\mathscr K_{\mathcal A}W\rangle
 =\sum_i\langle[\mathcal A_i,Z],[\mathcal A_i,W]\rangle.
\]
These are symmetric in $Z,W$ and prove self-adjointness; the sums of
squares prove the stated kernel. If $Z$ is in the kernel, the matter moment
paired with $Z$ vanishes because $[Z,\Phi]=Zb=0$, and
$\operatorname{tr}(Z[\mathcal A_i,\dot{\mathcal A}_i])
=\operatorname{tr}([Z,\mathcal A_i]\dot{\mathcal A}_i)=0$.
Thus the right-hand side is orthogonal to the kernel. In a finite
dimensional real inner-product space a self-adjoint operator has image
equal to the orthogonal complement of its kernel: this follows by
diagonalizing its real symmetric matrix. It gives existence and the
claimed exact ambiguity.
\end{proof}

For the literal moving orbit, subtracting the original scalar equations
while retaining arbitrary $\mathcal A_0$ gives the exact additional
equations
\begin{align*}
 [\dot{\mathcal A}_0,\Phi]+2[\mathcal A_0,\dot\Phi]
 +[\mathcal A_0,[\mathcal A_0,\Phi]]
 -\sum_i[\mathcal A_i,[\mathcal A_i,\Phi]]&=0,\\
 \dot{\mathcal A}_0b+2\mathcal A_0\dot b+\mathcal A_0^2b
                   -\sum_i\mathcal A_i^2b&=0.
\end{align*}
Their norm consequences are
\[
 \sum_i\operatorname{tr}[\mathcal A_i,\Phi]^2
       =\operatorname{tr}X^2-\operatorname{tr}\dot\Phi^2,\qquad
 \sum_i|\mathcal A_i b|^2=|v|^2-|\dot b|^2.
\]
For example the trace pairing of the first equation uses
$\operatorname{tr}(\Phi[\dot{\mathcal A}_0,\Phi])=0$ and
$\operatorname{tr}(\Phi[\mathcal A_0,\dot\Phi])
=-\operatorname{tr}([\mathcal A_0,\Phi]\dot\Phi)$; expansion of $X^2$
then gives the first identity. The real pairing of the second equation
with $b$ gives the second by anti-Hermiticity. These right-hand sides
need not be zero; temporal gauge cannot be imposed while silently keeping
the same literal scalar functions.

Here is an exact counterexample to extending the temporal-gauge
stabilizer deduction to nonzero $\mathcal A_0$. For real constants
$\alpha,\beta_1,\beta_2,\beta_3$ satisfying
$\alpha^2-\sum_i\beta_i^2=1$, keep the same scalar orbit and set
\[
 \mathcal A_0=(\alpha-1)\Omega,\qquad
 \mathcal A_i=\beta_i\Omega.
\]
All curvatures vanish, while
$D_0\Phi=\alpha[\Omega,\Phi]$, $D_0b=\alpha\Omega b$,
$D_i\Phi=\beta_i[\Omega,\Phi]$, and $D_i b=\beta_i\Omega b$.
Their covariant wave operators therefore equal the original second
derivatives, proving both scalar equations exactly. Taking
$\alpha=5/3$, $\beta_1=4/3$, $\beta_2=\beta_3=0$ gives
$\mathcal A_1b\neq0$: the spatial field does not stabilize the scalars.
Nevertheless $J_0=\alpha J_{\mathrm{sc}}\neq0$, since $|\alpha|\ge1$,
and the zero electric field cannot satisfy Gauss' equation. These are
scalar-compatible flat connections, not full gauge solutions.

There is also no lift of the moving orbit with all spatial
$\mathcal A_i=0$, even when $\mathcal A_0$ is arbitrary.
Indeed $J_{\mathrm{sc}}=-\mathscr I_{\Phi,b}\Omega$, so
\eqref{q444:inertia-equation} gives
$\mathscr I_{\Phi,b}(\mathcal A_0+\Omega)=0$.
Its kernel description implies $D_0\Phi=D_0b=0$.
Equation~\eqref{q444:gauge-b} would then require
$(aI_3+\kappa\Phi)b=ab=0$, impossible because $a>0$ and $B>0$.
This argument, and Theorem~\ref{q444:temporal-no-lift}, do not exclude
the general case with both nonzero $\mathcal A_0$ and nonzero spatial
fields satisfying all coupled equations.

% Insertion-ready continuation; no global macros or document wrappers.
% Requires the existing q444:moving-family and q444:gauge-eom sections.
\subsection{Constant-matter coordinates and an all-time real-connection obstruction}

Retain exactly the parameters and literal scalar functions of
Theorem~\ref{q444:moving-family}. In particular $a>0$, $B,d>0$, and
$x\ne0$. Keep the same original-coordinate $O(t)$, $\Omega$, and
the exact orthogonal map and inverse
\[
 K=\left(v/\sqrt2,w_0/\sqrt6,u/\sqrt3\right),\qquad K^{-1}=K^*.
\]
The squared lengths of $v,w_0,u$ are $2,6,3$, their pairwise inner
products vanish, and $\det(v,w_0,u)=6$; hence $K\in SO(3)$.
Put
\[
 H=K^*\Phi_0K=\begin{pmatrix}0&0&0\\0&-x&d\\0&d&x\end{pmatrix},
 \qquad q=K^*b_0=Be_1,\qquad
 A=K^*\Omega K=\begin{pmatrix}0&-\sqrt a&0\\\sqrt a&0&0\\0&0&0\end{pmatrix}.
\]
Thus $Hq=0$. Write $r^2=x^2+d^2$ and $\mu=r^2+B^2$.
The retained force, including the cubic term, is
\[
 F_H=-aH+\kappa(\operatorname{adj}H-qq^*)_{\mathrm{tf}}
 =\begin{pmatrix}-2ax&0&0\\0&2ax&-ad\\0&-ad&0\end{pmatrix}
 =[A,[A,H]],\qquad A^2q=-aq.
\]
Indeed $s=2(r^2+B^2)=6ax/\kappa$,
$\operatorname{adj}H=\operatorname{diag}(-r^2,0,0)$, and
$(\operatorname{adj}H-qq^*)_{\mathrm{tf}}
=\operatorname{diag}(-s/3,s/6,s/6)$, which give these identities
by direct multiplication.

\begin{proposition}[An exact constant-matter reformulation]
\label{q444:static-lift-system}
Set
\[
 C=K^*O^*\mathcal A_0OK+A,\qquad
 Y_i=K^*O^*\mathcal A_iOK,\qquad
 \mathcal E_i=\dot Y_i+[C,Y_i],
\]
\[
 \mathscr I_{H,q}Z=([H,[H,Z]]+qq^*Z+Zqq^*)_{\mathrm{tf}},\qquad
 \mathscr K_YZ=-\sum_i[Y_i,[Y_i,Z]].
\]
The prescribed literal scalar functions have a homogeneous gauge lift
if and only if the following complete system holds:
\begin{align}
 [\dot C,H]+[C,[C,H]]-\sum_i[Y_i,[Y_i,H]]&=F_H,
 \label{q444:static-lift-phi}\\
 (\dot C+C^2-\sum_iY_i^2)q&=-aq,
 \label{q444:static-lift-b}\\
 (\mathscr I_{H,q}+2\mathscr K_Y/g^2)C
   &=-\frac2{g^2}\sum_i[Y_i,\dot Y_i],
 \label{q444:static-lift-gauss}\\
 \dot{\mathcal E}_i+[C,\mathcal E_i]
 -\sum_j[Y_j,[Y_j,Y_i]]+\frac{g^2}{2}\mathscr I_{H,q}Y_i&=0
 \quad(i=1,2,3).
 \label{q444:static-lift-spatial}
\end{align}
The inverse connection maps are exactly
\[
 \mathcal A_0=OKCK^*O^*-\Omega,\qquad
 \mathcal A_i=OKY_iK^*O^*.
\]
Neither the literal functions nor their time scale has been changed.
\end{proposition}
\begin{proof}
For a spatially constant $U(t)\in SU(3)$, the transformations are
$\Phi'=U\Phi U^*$, $b'=Ub$,
$\mathcal A_0'=U\mathcal A_0U^*-\dot UU^*$, and
$\mathcal A_i'=U\mathcal A_iU^*$.
The product rule gives
$D'_\nu(Ub)=UD_\nu b$ and
$D'_\nu(U\Phi U^*)=U(D_\nu\Phi)U^*$.
Their commutators give $\mathcal F'_{\nu\rho}
=U\mathcal F_{\nu\rho}U^*$. Traces, trace projection, and the potential
are invariant under this unitary action; adjugate covariance follows
from the inverse formula on invertible matrices and polynomial
continuity. Thus all the equations of Proposition~\ref{q444:gauge-eom},
including Gauss' equation, transform covariantly.

Taking $U=K^*O^*$ gives the constant matter fields $H,q$ and the displayed
connection maps, because $-\dot UU^*=K^*\Omega K=A$.
The scalar covariant derivatives are $[C,H]$ and $Cq$; their second
derivatives are $[\dot C,H]+[C,[C,H]]$ and $(\dot C+C^2)q$.
The trace-free currents are
$(J_0)_{\mathrm{tf}}=-\mathscr I_{H,q}C$ and
$(J_i)_{\mathrm{tf}}=-\mathscr I_{H,q}Y_i$.
Finally
$\sum_i[Y_i,\mathcal E_i]
=\sum_i[Y_i,\dot Y_i]+\mathscr K_YC$.
Substitution gives every displayed equation with its stated sign and
factor. The explicitly invertible maps prove both directions.
\end{proof}

\begin{theorem}[No real-connection lift with any spatial time dependence]
\label{q444:real-no-lift}
The literal moving-kernel scalar functions admit no homogeneous lift
with $\mathcal A_\nu(t)\in\mathfrak{so}(3)$ for all four indices, on
any nonempty open time interval. The temporal connection is arbitrary;
the spatial connections may be arbitrary twice differentiable functions,
including rotating, multifrequency, and nonperiodic functions. No
commutativity or spatial-stationarity assumption is imposed.
\end{theorem}
\begin{proof}
Since $K,O,\Omega$ are real, the exact frame maps preserve real
skew-symmetry in both directions. Define
\[
 L_1=\begin{pmatrix}0&0&0\\0&0&-1\\0&1&0\end{pmatrix},\quad
 L_2=\begin{pmatrix}0&0&1\\0&0&0\\-1&0&0\end{pmatrix},\quad
 L_3=\begin{pmatrix}0&-1&0\\1&0&0\\0&0&0\end{pmatrix},\quad
 L(z)=\sum_jz_jL_j\quad(z\in\mathbb R^3).
\]
Direct multiplication gives
\[
 L(z)w=z\times w,\quad
 L(z)L(w)=wz^{\mathsf T}-(z\cdot w)I_3,\quad
 [L(z),L(w)]=L(z\times w),\quad
 -\operatorname{tr}(L(z)L(w))=2z\cdot w.
\]
These formulas specify the commutator orientation and every norm factor.
Write
\[
 C=L(c),\quad Y_i=L(y_i),\quad h=\dot c,\quad
 P=\sum_i y_iy_i^{\mathsf T},\quad
 Q=cc^{\mathsf T}-P=(q_{jk})_{j,k=1}^3.
\]
Here the scalar entries $q_{jk}$ do not replace the earlier fixed matter
vector $q=Be_1$. Since $A=\sqrt a L_3$, the residual in
\eqref{q444:static-lift-b}, divided by $B$, is
\[
 R_b=\begin{pmatrix}a-q_{22}-q_{33}\\h_3+q_{12}\\-h_2+q_{13}\end{pmatrix}.
\]
The residual in \eqref{q444:static-lift-phi} is a symmetric matrix
$R_H$ whose six independent entries are
\begin{align*}
 (R_H)_{11}&=2ax-4dq_{23}+2xq_{22}-2xq_{33},\\
 (R_H)_{12}&=dh_2+3dq_{13}+xh_3-3xq_{12},\\
 (R_H)_{13}&=-dh_3+3dq_{12}+xh_2+3xq_{13},\\
 (R_H)_{22}&=-2ax-2dh_1+2dq_{23}+4xq_{11}+2xq_{33},\\
 (R_H)_{23}&=ad-4dq_{11}-dq_{22}-dq_{33}-2xh_1,\\
 (R_H)_{33}&=2dh_1+2dq_{23}-4xq_{11}-2xq_{22}.
\end{align*}
These entries follow by multiplying the explicitly given matrices in
\[
 R_H=[L(h),H]+\sum_{j,k}q_{jk}[L_j,[L_k,H]]
                  -a[L_3,[L_3,H]],\qquad
 R_b=L(h)e_1+\sum_{j,k}q_{jk}L_jL_ke_1+ae_1.
\]
Thus no spatial time-dependence ansatz has entered this calculation.

The equation $R_b=0$ gives $q_{22}+q_{33}=a$, $h_3=-q_{12}$,
$h_2=q_{13}$. The two equations $(R_H)_{12}=(R_H)_{13}=0$ then give
$dq_{13}-xq_{12}=0$ and $dq_{12}+xq_{13}=0$.
Since $x^2+d^2>0$, they force
$q_{12}=q_{13}=h_2=h_3=0$.
Next $(R_H)_{11}=0$ gives $xq_{22}=dq_{23}$.
After these substitutions, $(R_H)_{23}=(R_H)_{33}=0$ give
$xh_1=-2dq_{11}$ and $dh_1=2xq_{11}$.
They imply $2(x^2+d^2)q_{11}=0$, and hence $q_{11}=h_1=0$.
Conversely the following values make every residual entry zero:
\begin{equation}\label{q444:real-scalar-covariance}
 \dot c=0,\qquad Q=
 \begin{pmatrix}0&0&0\\0&d\tau/x&\tau\\0&\tau&a-d\tau/x\end{pmatrix},
 \qquad \tau=q_{23}(t)\in\mathbb R.
\end{equation}
This is the exact necessary-and-sufficient scalar equation within the
real algebra. In particular $Qe_1=0$ and $\operatorname{tr}Q=a>0$.

The covariance $P=\sum_i y_iy_i^{\mathsf T}$ is positive semidefinite.
Suppose $c_1\ne0$ at any time. For $v\in e_1^\perp$ define
$z=v-(c\cdot v)e_1/c_1$. Then $c\cdot z=0$ and $Qe_1=0$, so
\[
 v^{\mathsf T}Qv=z^{\mathsf T}Qz
 =(c\cdot z)^2-z^{\mathsf T}Pz
 =-\sum_i(y_i\cdot z)^2\leq0.
\]
The lower $2\times2$ block of $Q$ would therefore be negative
semidefinite, contradicting its positive trace $a$.
Equivalently, $P_{11}=c_1^2$ and the Schur complement of $P_{11}$ in
$P=cc^{\mathsf T}-Q$ is exactly minus that lower block of $Q$.
Thus $c_1=0$ at every time. Since
$Q_{11}=c_1^2-\sum_i y_{i1}^2=0$, every $y_{i1}=0$ as well.
Differentiation on the open time interval now gives
\[
 c,y_i,\dot y_i\in e_1^\perp,\qquad
 \sum_i y_i\times\dot y_i\in\mathbb R e_1.
\]

The exact inertia on this real algebra is
\[
 \mathscr I_{H,q}L_1=4r^2L_1,\qquad
 \mathscr I_{H,q}L_2=\mu L_2,\qquad
 \mathscr I_{H,q}L_3=\mu L_3.
\]
To prove these formulas, $H^2=r^2\operatorname{diag}(0,1,1)$ and
$HL_1H=-r^2L_1$ give
$[H,[H,L_1]]=4r^2L_1$, while the $qq^*$ terms vanish.
For $j=2,3$, multiplication gives $HL_jH=0$,
$H^2L_j+L_jH^2=r^2L_j$, and
$qq^*L_j+L_jqq^*=B^2L_j$.
Every resulting real-skew matrix is already traceless. This proves the
three identities and gives $\mathscr I_{H,q}C=\mu C$.

For $Z,W\in\mathfrak{su}(3)$ set
$\langle Z,W\rangle=-\operatorname{tr}(ZW)$ and
$\|Z\|^2=\operatorname{tr}(Z^*Z)=-\operatorname{tr}Z^2$.
Cyclicity gives
$\langle Z,[Y,W]\rangle=-\langle[Y,Z],W\rangle$, so
\[
 \langle C,\mathscr K_YC\rangle
 =\sum_i\|[Y_i,C]\|^2
 =2\sum_i|y_i\times c|^2.
\]
Pair the full Gauss equation \eqref{q444:static-lift-gauss} with $C$.
Its right side vanishes because $c\perp\sum_i y_i\times\dot y_i$.
The left side, with the exact factor $2$ from the trace inner product, is
\[
 2\mu|c|^2+\frac4{g^2}\sum_i|y_i\times c|^2=0.
\]
Since $\mu>0$ and $g^2>0$, this forces $c=0$.
Then $Q=-P$ has
$\operatorname{tr}Q=-\sum_i|y_i|^2\leq0$, contradicting
$\operatorname{tr}Q=a>0$.
Every full solution must obey the two scalar equations and Gauss'
equation used here. Their contradiction already excludes a full lift,
without imposing or deleting any spatial gauge equation.
\end{proof}

\begin{corollary}[Constant spatial fields with arbitrary complex color]
\label{q444:static-complex-no-lift}
There is no lift with all $Y_i$ constant in the constant-matter frame,
even when $C(t)$ is arbitrary and the fields take values in the full
complex algebra $\mathfrak{su}(3)$.
\end{corollary}
\begin{proof}
Cyclicity and anti-Hermiticity give
\[
 \langle Z,\mathscr I_{H,q}Z\rangle
 =\operatorname{tr}[Z,H]^2+2|Zq|^2,\qquad
 \langle Z,\mathscr K_YZ\rangle=\sum_i\|[Y_i,Z]\|^2.
\]
For completeness, trace projection pairs trivially with traceless $Z$;
cyclicity gives
$-\operatorname{tr}(Z[H,[H,Z]])=\operatorname{tr}[Z,H]^2$, and
$-\operatorname{tr}(Z(qq^*Z+Zqq^*))=-2q^*Z^2q=2|Zq|^2$.
The second formula follows from the same commutator trace identity used
in the theorem. The matrix $[Z,H]$ is Hermitian, so all displayed
terms are nonnegative.

When $\dot Y_i=0$, Gauss' equation is
$(\mathscr I_{H,q}+2\mathscr K_Y/g^2)C=0$.
Its pairing with $C$ therefore forces $Cq=0$, $[C,H]=0$, and
$[C,Y_i]=0$ for every $i$ at every time.
Since $q$ is constant, $Cq=0$ implies $\dot Cq=0$ and $C^2q=0$.
The vector equation becomes
$(-\sum_iY_i^2+aI_3)q=0$. Pairing with $q$ gives
$\sum_i|Y_iq|^2+aB^2=0$, which contradicts $a>0$, $B>0$.
\end{proof}

\begin{corollary}[Matter-stabilizer gauge saturation]
\label{q444:real-stabilizer-no-lift}
There is no lift for which a differentiable real $\theta(t)$ makes
\[
 C'=SCS^*-i\dot\theta H,\qquad Y_i'=SY_iS^*,\qquad
 S(t)=\exp(i\theta(t)H)
\]
all real skew-symmetric. Nor is there a lift when such a transformation
makes all $Y_i'$ constant, without a reality restriction.
\end{corollary}
\begin{proof}
$H$ is Hermitian and traceless, so $S\in SU(3)$.
It commutes with $H$, and $Hq=0$ gives $Sq=q$.
Differentiation gives $\dot SS^*=i\dot\theta H$.
Thus these are exactly the gauge transformations proved in
Proposition~\ref{q444:static-lift-system}, retaining the same constant
matter fields. They would transform any claimed lift into one excluded
by Theorem~\ref{q444:real-no-lift} or
Corollary~\ref{q444:static-complex-no-lift}, respectively.
\end{proof}

These are exact exclusions for the stated classes, not an assertion
that every $\mathfrak{su}(3)$ connection is real or belongs to their
matter-stabilizer gauge saturation. Arbitrary time-dependent complex
$\mathfrak{su}(3)$ lifts outside those classes remain open here.

\subsection{Constructive homogeneous gauge solutions and their observations}

\begin{theorem}[A real radial family of full gauge solutions]
\label{q444:radial-gauge}
Let $D=\operatorname{diag}(-2,1,1)$ and retain
\[
 \Phi(t)=x(t)D,\qquad b(t)=\beta(t)e_1,\qquad
 x(t),\beta(t)\in\mathbb R,\qquad \mathcal A_\mu=0.
\]
This is a full homogeneous gauge--scalar solution exactly when
\begin{align*}
 \ddot x&=-a x-\kappa x^2+\frac{\kappa}{3}\beta^2,\\
 \ddot\beta&=-(a-2\kappa x)\beta,\qquad
 a=m^2+\Gamma(6x^2+2\beta^2).
\end{align*}
For every real initial position and velocity these equations have a
global solution. They give finite-energy homogeneous solutions on a
finite-volume flat $\mathbb T^3$, and local solutions on $\mathbb R^3$,
without asserting finite total energy on $\mathbb R^3$.
\end{theorem}
\begin{proof}
Here $[\Phi,\dot\Phi]=0$ and
$b\dot b^*-\dot b b^*=0$ because $x,\beta$ are real.
Thus every $J_\mu$ vanishes identically and so do all curvatures.
The gauge equations hold exactly, not after discarding Gauss' equation.
Directly
\[
 (\operatorname{adj}(xD)-\beta^2e_1e_1^*)_{\mathrm{tf}}
          =\left(-x^2+\frac{\beta^2}{3}\right)D,\qquad
 (aI_3+\kappa xD)\beta e_1=(a-2\kappa x)\beta e_1.
\]
The two scalar equations are consequently the displayed real system.
Its conserved density is
\[
 \mathcal E=3\dot x^2+\dot\beta^2+
 V(xD,\beta e_1),\qquad
 V=\frac{m^2}{2}s+\frac{\Gamma}{4}s^2
               -2\kappa x(\beta^2-x^2),\quad s=6x^2+2\beta^2.
\]
The stated conservation follows by differentiating and substituting the
two real equations; the $x$ kinetic coefficient is $3$, not $1/2$.
Moreover $|x|\le\sqrt{s/6}$ and $\beta^2+x^2\le s/2$, so
$|2x(\beta^2-x^2)|\le s^{3/2}/\sqrt6$ and
\[
 V\ge\frac{\Gamma}{4}s^2-\frac{|\kappa|}{\sqrt6}s^{3/2}.
\]
This is coercive and globally bounded below. Conservation bounds the
positions and velocities. The polynomial first-order vector field is
locally Lipschitz and bounded with bounded derivative on their compact
range; hence no finite-time escape is possible, proving global existence.
Multiplying the finite density by the finite volume of $\mathbb T^3$
gives finite total energy there.
\end{proof}

\begin{proposition}[A full-gauge vacuum and its marked Lorentz observation]
\label{q444:vacuum-observation}
At $\Gamma=m^2=1$, $\kappa=6$, let
\[
 x_0=\frac{3+\sqrt3}{12},\qquad
 \Phi=x_0\operatorname{diag}(-2,1,1),\qquad b=3x_0e_1.
\]
Together with $\mathcal A_\mu=0$ this is an equilibrium full gauge
solution at a global minimum of the scalar potential.
For $Q=\Theta(\Phi,b)$ its spectrum is $(4x_0,-2x_0,-2x_0)$.
Its spectral-floor matrix $H=Q+2x_0I_3$ has rank one and is positive.
The complex observation and hidden coordinate are
\[
 G=3x_0\begin{pmatrix}1&-1\\-1&2\end{pmatrix},
 \qquad \rho=3x_0,\qquad \det G=9x_0^2=|\rho|^2.
\]
More generally its $SU(3)$ orbit, parameterized by unit vectors
$e=ge_1$, has
\[
 \Phi=x_0(I_3-3ee^*),\quad b=3x_0e,\quad
 G=3x_0P(I_3-ee^*)P^*,\quad
 \det G=9x_0^2|u^*e|^2=|\rho|^2.
\]
For fixed $P$, this one gauge orbit contains both null and timelike
observations. The fixed-$P$ determinant is therefore not a
gauge-invariant mass observable.
\end{proposition}
\begin{proof}
For a nonzero real equilibrium with $\beta=3x>0$, the radial equations
reduce to
$24\Gamma x^2-2\kappa x+m^2=0$.
The displayed $x_0$ is the larger root at $(1,1,6)$.
Its quaternionic matrix is
\[
 Q=\begin{pmatrix}-2x_0&0&0\\0&x_0&3x_0\mathbf j\\
                       0&-3x_0\mathbf j&x_0\end{pmatrix}.
\]
The lower block has eigenvalues $x_0\pm3x_0$, proving the spectrum.
For an independent global-minimum check, every trace-free quaternionic
Hermitian matrix satisfies $54N(Q)^2\le s(Q)^3$: for eigenvalues summing
to zero, the exact identity is
$s^3-54N^2=2\prod_{i<j}(\lambda_i-\lambda_j)^2$.
Thus, with $r=\sqrt{s}\ge0$,
$V\ge r^2/2+r^4/4-2r^3/\sqrt6$, with equality at eigenvalues
$(2y,-y,-y)$, $y\ge0$, $r=\sqrt6y$.
On this equality family
$V=3y^2+9y^4-12y^3$ and
$dV/dy=6y(1-6y+6y^2)$.
The larger positive root is
$y=(3+\sqrt3)/6=2x_0$; the derivative is positive, negative,
then positive on the three successive intervals, and this minimum has
negative value, since at $y=2/3$ the potential is $-4/9$.
It is therefore the global minimum, attained by the displayed $Q$.

The complex part of $H$ is
$3x_0\operatorname{diag}(0,1,1)$, and compression by the retained $P$
gives the displayed $G$. The $SU(3)$ action on pairs is
$(\Phi,b)\mapsto(g\Phi g^*,gb)$; it is exactly quaternionic congruence
under $\Theta$, since $B(gb)=gB(b)g^{\mathsf T}$ for $g\in SU(3)$.
Consequently the spectrum and spectral-floor shift are unchanged.
Writing $v=Pe$ and using the determinant lemma in dimension two gives
\[
 \det(C-vv^*)=\det C\,(1-v^*C^{-1}v)
 =3\left(1-e^*\left(I_3-\frac{uu^*}{3}\right)e\right)
 =|u^*e|^2.
\]
This proves the general observation formula. A unit $e\in u^\perp$
gives determinant zero, while $e=e_1$ gives a positive determinant.
These vectors are related by $SU(3)$: extend either to a unitary basis
and adjust the phase of one of the remaining columns to make determinant
one. Thus the examples are genuinely on one gauge orbit.

The exact covariance map retains the observer as part of the data:
under $H'=gHg^*$, $P'=Pg^*$ one has
$P'H'P'^*=PHP^*$. Complex-component projection therefore has the same
value. Simultaneously $u'=gu$, $b'=gb$ gives $u'^*b'=u^*b$.
This is covariance of the marked observer, not invariance of a
fixed observer under an active transformation of the state.
\end{proof}

\begin{proposition}[Gauge averaging and distinct minimum orbits]
\label{q444:vacuum-average}
Let $d\mu$ be Haar measure on $SU(3)$ with its retained positive total
mass $V_\mu$, not set to one. For every complex Hermitian $A$,
\[
 \det(PAP^*)=u^*\operatorname{adj}(A)u,\qquad
 \int_{SU(3)}\det(PgAg^*P^*)\,d\mu(g)
                    =V_\mu\operatorname{tr}\operatorname{adj}(A).
\]
For a quaternionic positive rank-one matrix $H=\Theta(A,b)$,
$\operatorname{adj}(A)=bb^*$. Hence its averaged observation is
$V_\mu\|b\|^2$.

Put $\chi=2x_0>0$ for the preceding minimum and, for
$0\le\theta\le\pi/4$, define
\[
 v_\theta=(\cos\theta,\ \sin\theta\,\mathbf j,\ 0)^{\mathsf T},
 \qquad Q_\theta=\chi(3v_\theta v_\theta^*-I_3).
\]
Every $Q_\theta$ is a scalar global minimum, and its complex-plus-triplet
coordinate satisfies
\[
 b_\theta=(0,0,-3\chi\cos\theta\sin\theta)^{\mathsf T},\qquad
 \|b_\theta\|^2=\frac{9\chi^2}{4}\sin^2(2\theta).
\]
Consequently this minimum family contains pairwise $SU(3)$-inequivalent
states with invariant $\|b\|^2$ ranging from $0$ to $9\chi^2/4$.
The coefficient agrees with the supplied packet's invariant formula;
this calculation validates that coefficient rather than claiming it
as a newly discovered one.
\end{proposition}
\begin{proof}
The signed $2$-by-$2$ minors of $P$ form precisely $u=(1,1,1)^{\mathsf T}$.
The exterior-square form of the Cauchy--Binet expansion therefore gives
$\det(PAP^*)=u^*\operatorname{adj}(A)u$; equivalently this follows by
expanding the six entries of $PAP^*$ and collecting the cofactors.
The identity $\operatorname{adj}(gAg^*)=g\operatorname{adj}(A)g^*$
holds for invertible $A$ and then for every $A$ by polynomial continuity.
The Haar average of these conjugates commutes with every element of
$SU(3)$. Commutation with diagonal determinant-one phase matrices kills
its off-diagonal entries, and commutation with determinant-one cyclic
permutations makes its diagonal entries equal. Its trace fixes the
average at $V_\mu\operatorname{tr}(\operatorname{adj}A)I_3/3$.
Since $u^*u=3$, the stated integral follows with all factors retained.

Write $r=b^*b$. The skew convention gives directly
$B(b)\overline{B(b)}=bb^*-rI_3$.
The complex part of $H^2$ is therefore $A^2-bb^*+rI_3$,
and $e_2(H)=e_2(A)-r$.
For positive rank-one $H$, its spectral identity is
$H^2=(\operatorname{tr}H)H$ and $e_2(H)=0$.
Taking complex parts and using $\operatorname{tr}H=\operatorname{tr}A$
gives
$A^2-(\operatorname{tr}A)A+e_2(A)I_3=bb^*$.
The left side is $\operatorname{adj}(A)$, proving the assertion.

Finally $v_\theta^*v_\theta=1$, so $Q_\theta$ has fixed spectrum
$(2\chi,-\chi,-\chi)$ and hence the same minimum value.
Its $(1,2)$ quaternionic entry is
$-3\chi\cos\theta\sin\theta\,\mathbf j$, with the other quaternionic
off-diagonal coefficients zero. The retained definition of $B(b)$
therefore gives the displayed $b_\theta$ with its minus sign.
Taking its norm proves the formula. It is strictly increasing on the
given interval; since $b\mapsto gb$ preserves its norm, distinct
parameters there cannot lie on one $SU(3)$ orbit.
The positive spectral-floor matrix is $3\chi v_\theta v_\theta^*$,
so the averaging identity applies to the same triplet coordinate.
\end{proof}

The classical $SU(3)$ cohomogeneity-one geometry of $\mathbb{HP}^2$
is discussed by A.~Gambioli, \emph{Eight-dimensional $SU(3)$-manifolds
of cohomogeneity one}, arXiv:math/0611796, Section~4, the
$\mathbb{HP}^2=M(\mathbb{CP}^2,S^5)$ example,
\url{https://arxiv.org/abs/math/0611796}.
The explicit family and its separating invariant above are proved
directly; no new classification of this classical orbit space is claimed.
The averaging formula supplies the exact relation of that invariant to
the marked Lorentz determinant. A continuum of classical potential
minima is not, by itself, a calculation of a quantum Hamiltonian spectrum.

% Ready-to-input proof fragment. All edits are confined to agent_work.
% Requires amsmath, amssymb and the standard theorem/proposition/lemma
% environments. Labels beginning q444:vk- are private to this fragment.
\subsection{Vacuum kinetics and exact current-free continuations}

Retain the real scalar state space
\[
 \mathcal E=\operatorname{Herm}_3(\mathbb C)_0\oplus\mathbb C^3,
 \qquad
 \langle(\phi,c),(\psi,d)\rangle
 =\operatorname{tr}(\phi\psi)+2\operatorname{Re}(c^*d).
\]
The subscript \(0\) means trace zero, not a quotient. The scalar kinetic
density is exactly one half of this metric:
\[
 \mathcal L_{\rm sc}=\tfrac12\operatorname{tr}
       (D_\mu\Phi D^\mu\Phi)+(D_\mu b)^*D^\mu b-V,\qquad
 V=\tfrac{m^2}{2}s+\tfrac{\Gamma}{4}s^2
                   -\kappa(\det\Phi-b^*\Phi b),\quad
 s=\operatorname{tr}\Phi^2+2b^*b .
\]
The spacetime metric is \((+---)\), \(m^2\geq0\), \(\Gamma>0\);
the nonzero vacuum results below use \(\kappa>0\).
The gauge action is
\((\Phi,b)\mapsto(g\Phi g^*,gb)\), \(g\in SU(3)\).
Its anti-Hermitian connection is \(\mathcal A_\mu\), and the retained
gauge density is
\(\operatorname{tr}(\mathcal F_{\mu\nu}\mathcal F^{\mu\nu})/(2g^2)\),
with \(g^2>0\). All constructed connections below are identically zero.
We make no claim about connections not considered here.

For \(Z_{\rm tf}=Z-\operatorname{tr}(Z)I_3/3\), variation gives, at
\(\mathcal A_\mu=0\),
\begin{align}
 \Box\Phi&=-a\Phi+\kappa(\operatorname{adj}\Phi-bb^*)_{\rm tf},
       &a&=m^2+\Gamma s,\label{q444:vk-pde}\\
 \Box b&=-(aI_3+\kappa\Phi)b,\qquad&
 J_\mu&=[\Phi,\partial_\mu\Phi]
       +b(\partial_\mu b)^*-(\partial_\mu b)b^* .\notag
\end{align}
Indeed
\[
 dV(\delta\Phi,\delta b)
 =\operatorname{tr}\!\left(
 [a\Phi-\kappa(\operatorname{adj}\Phi-bb^*)_{\rm tf}]\delta\Phi\right)
 +2\operatorname{Re}\!\left[
 ((aI_3+\kappa\Phi)b)^*\delta b\right].
\]
This proves both scalar equations with their original metric factors.
Variation in \(\mathcal A_\mu\) pairs the matter variation with
\(\operatorname{tr}(J^\mu\delta\mathcal A_\mu)\).
Consequently \(J_\mu=0\) and \(\mathcal F_{\mu\nu}=0\) satisfy every
gauge equation, including Gauss, before any restriction of variations.

We retain the quaternionic convention
\[
 \Theta(\Phi,b)=\Phi+B(b)\mathbf j,\qquad
 B(b)=\begin{pmatrix}
 0&\bar b_3&-\bar b_2\\-\bar b_3&0&\bar b_1\\
 \bar b_2&-\bar b_1&0
 \end{pmatrix},\qquad \mathbf j z=\bar z\mathbf j .
\]
Its real trace metric is exactly the displayed metric on \(\mathcal E\):
the diagonal sum of \(B(b)\overline{B(c)}\) has real part
\(-2\operatorname{Re}(b^*c)\). Its cubic is
\(N(\Theta(\Phi,b))=\det\Phi-b^*\Phi b\). To verify the latter identity,
write \(\Phi_{23}=z_1,\Phi_{31}=z_2,\Phi_{12}=z_3\) and use the
Hermitian Moore formula
\[
 N\begin{pmatrix}a&q_3&\bar q_2\\\bar q_3&d&q_1\\q_2&\bar q_1&f\end{pmatrix}
 =adf-a|q_1|^2-d|q_2|^2-f|q_3|^2
                 +2\operatorname{Re}(q_1q_2q_3).
\]
Here \(q_i=z_i+\bar b_i\mathbf j\). The complex part of the triple
product is
\(z_1z_2z_3-\bar b_1b_2z_3-z_1\bar b_2b_3-\bar b_1\bar z_2b_3\).
Its extra real terms and the three extra diagonal terms are precisely
\(-b^*\Phi b\).

\begin{lemma}[Retained nonzero minimum and its curve]\label{q444:vk-vacuum}
Suppose \(0\leq m^2\leq\kappa^2/(27\Gamma)\), \(\kappa>0\), and put
\[
 \chi=\frac{\kappa+\sqrt{\kappa^2-24\Gamma m^2}}{12\Gamma}>0 .
\]
For a unit \(v\in\mathbb H^3\), \(Q=\chi(3vv^*-I_3)\) is a global
minimum of \(V\) on the trace-free quaternionic Hermitian state space.
The origin is an additional minimum only at
\(m^2=\kappa^2/(27\Gamma)\).
For \(0\leq\theta\leq\pi/4\), let
\[
 v_\theta=(\cos\theta,\sin\theta\,\mathbf j,0)^{\mathsf T},
 \qquad Q_\theta=\chi(3v_\theta v_\theta^*-I_3).
\]
Its pair coordinates and potential value are
\begin{align}
 \Phi_\theta&=\chi\operatorname{diag}
     (3\cos^2\theta-1,\,3\sin^2\theta-1,\,-1),\notag\\
 b_\theta&=-3\chi\cos\theta\sin\theta\,e_3,\label{q444:vk-curve}\\
 V_0&=3m^2\chi^2+9\Gamma\chi^4-2\kappa\chi^3
      =\kappa\chi^3-9\Gamma\chi^4 .\notag
\end{align}
\end{lemma}
\begin{proof}
The quaternionic spectral theorem gives real eigenvalues
\(\lambda_1+\lambda_2+\lambda_3=0\) and
\(s=\sum\lambda_i^2\). On each sphere \(s>0\), extrema of the product
have at most two distinct eigenvalues: subtracting two Lagrange equations
for the product gives
\((\lambda_i-\lambda_j)(-\lambda_k-2\beta)=0\).
If all three were distinct they would all equal \(-2\beta\), a
contradiction. The two-eigenvalue trace-free possibilities are
\((2t,-t,-t)\) and their negatives and permutations, where
\(t=\sqrt{s/6}>0\). Thus \(N\leq2t^3\), sharply, and
\[
 V\geq F(t)=3m^2t^2+9\Gamma t^4-2\kappa t^3,\qquad
 F'(t)=6t(m^2-\kappa t+6\Gamma t^2).
\]
For \(m^2>0\) the larger positive root is \(\chi\), the only positive
local minimum; the smaller root is a local maximum. For \(m^2=0\),
\(F'(t)<0\) on \(0<t<\chi\) and \(>0\) on \(t>\chi\).
Using \(m^2=\kappa\chi-6\Gamma\chi^2\) yields
\(F(\chi)=\chi^3(\kappa-9\Gamma\chi)\), negative below the displayed
threshold and zero at it. Comparing with \(F(0)=0\) proves the global
claim and all equality cases used here. The spectrum of \(Q_\theta\)
is \((2\chi,-\chi,-\chi)\). Its \((1,2)\) entry is
\(-3\chi\cos\theta\sin\theta\,\mathbf j\), giving the stated sign of
\(b_\theta\). Direct substitution gives the potential value.
\end{proof}

\begin{proposition}[Horizontal metric and normal acceleration]
\label{q444:vk-horizontal}
Write \(q_\theta=(\Phi_\theta,b_\theta)\),
\(T=\partial_\theta q_\theta\), \(N=\partial_\theta^2q_\theta\).
Then
\[
 \langle T,T\rangle=18\chi^2,\qquad
 \langle T,N\rangle=0,\qquad
 \langle N,N\rangle=72\chi^2 .
\]
The curve is \(SU(3)\)-horizontal and its induced kinetic density is
\(9\chi^2\dot\theta^2\). Every homogeneous unconstrained solution
confined to this fixed-\(\chi\) curve has constant \(\theta\).
The vector \(N\) is normal to the entire fixed-\(\chi\) minimum
manifold. Its radial projection is \(-3q_\theta\), and
\(\|N+3q_\theta\|^2=18\chi^2\); hence there is a nonzero normal
component in addition to the overall radial one.
\end{proposition}
\begin{proof}
With \(D=\operatorname{diag}(-1,1,0)\), differentiation gives
\[
 T_\Phi=3\chi\sin(2\theta)D,\quad T_b=-3\chi\cos(2\theta)e_3,
 \quad
 N_\Phi=6\chi\cos(2\theta)D,\quad N_b=6\chi\sin(2\theta)e_3.
\]
Since \(\operatorname{tr}D^2=2\), the three metric identities follow
without changing the coefficient of the triplet metric.
All commutators of the diagonal matrices vanish, and
\(b_\theta T_b^*-T_b b_\theta^*=0\) because both coefficients are real.
For every \(\xi\in\mathfrak{su}(3)\),
\[
 \langle T,([\xi,\Phi_\theta],\xi b_\theta)\rangle
 =\operatorname{tr}\bigl(
   ([\Phi_\theta,T_\Phi]+b_\theta T_b^*-T_b b_\theta^*)\xi\bigr)=0.
\]
Thus this nonzero tangent is horizontal and cannot itself be a
nonzero vertical tangent in this positive metric.
At a vacuum the potential gradient vanishes, so the homogeneous
equation is \(T\ddot\theta+N\dot\theta^2=0\).
Pairing with \(T\) first gives \(\ddot\theta=0\), and the nonzero
norm of \(N\) then gives \(\dot\theta=0\).
The exact vector that would have to be supplied as an extra force
for a constrained constant-speed motion is \(N\dot\theta^2\).
Its force covector is
\[
 \operatorname{tr}(N_\Phi\dot\theta^2\,\delta\Phi)
       +2\operatorname{Re}((N_b\dot\theta^2)^*\delta b).
\]

For the assertion about the whole minimum manifold,
\(v_\theta^*v'_\theta=0\), \(v_\theta'^*v'_\theta=1\), and
\(v_\theta''=-v_\theta\). Therefore
\[
 Q_\theta''=6\chi(v_\theta'v_\theta'^*-v_\theta v_\theta^*).
\]
Every tangent to the rank-one projector manifold has the form
\(3\chi(zv_\theta^*+v_\theta z^*)\), \(v_\theta^*z=0\).
The real trace pairing with both projectors in \(Q_\theta''\)
is zero, by \(v_\theta^*v_\theta'=0\) and cyclicity of real trace.
This proves normality. Finally
\(\|q_\theta\|^2=6\chi^2\) and
\(\langle N,q_\theta\rangle=-18\chi^2\); subtracting the radial
projection gives the asserted remaining norm.
Thus the constant-speed curve is an intrinsic geodesic of the
minimum manifold but not an unconstrained homogeneous solution.
\end{proof}

\begin{proposition}[Exact three-real-coordinate full-gauge sector]
\label{q444:vk-sector}
Let \(h,k,\beta\) be real fields and set
\[
 \Phi=\operatorname{diag}(h+k,h-k,-2h),\qquad b=\beta e_3,
 \qquad\mathcal A_\mu=0.
\]
This ansatz satisfies all full field equations if and only if
\begin{align}
 \Box h&=-ah-\kappa h^2+\tfrac{\kappa}{3}(k^2+\beta^2),\notag\\
 \Box k&=-(a-2\kappa h)k,\qquad
 \Box\beta=-(a-2\kappa h)\beta,\label{q444:vk-sector-eom}\\
 a&=m^2+\Gamma(6h^2+2k^2+2\beta^2).\notag
\end{align}
In particular it is an unconstrained invariant sector for homogeneous
initial data with tangent velocities. Its homogeneous density is
\[
 \mathcal L_{\rm hom}=3\dot h^2+\dot k^2+\dot\beta^2-V(h,k,\beta),
 \quad
 V=\tfrac{m^2}{2}s+\tfrac{\Gamma}{4}s^2
                    -2\kappa h(k^2+\beta^2-h^2),
 \quad s=6h^2+2k^2+2\beta^2 .
\]
Every homogeneous initial condition in this sector has a global
solution for \(m^2\geq0,\Gamma>0,\kappa\in\mathbb R\).
\end{proposition}
\begin{proof}
Direct multiplication gives
\[
 \det\Phi-b^*\Phi b=2h(k^2+\beta^2-h^2),\qquad
 \operatorname{adj}\Phi-bb^*
 =\operatorname{diag}(-2h(h-k),-2h(h+k),h^2-k^2-\beta^2).
\]
Its trace is \(-3h^2-k^2-\beta^2\). Removing one third of this trace
and substituting in \eqref{q444:vk-pde} gives exactly
\eqref{q444:vk-sector-eom}, including the factor \(1/3\) in the
\(h\) equation. The last component of the triplet force is
\(-(a-2\kappa h)\beta\); all other entries remain zero.
The currents vanish for each spacetime index: the matrices and
their derivatives are diagonal, while \(b\) and its derivative
are real multiples of \(e_3\). The zero connection therefore
satisfies all gauge equations. This also proves consistency
under unrestricted variations, not merely stationary points
of a restricted Lagrangian. As a check, that Lagrangian has
Euler--Lagrange weights \(6,2,2\) on \(h,k,\beta\), respectively.

For homogeneous fields the smooth polynomial vector field is tangent
to the displayed real linear sector, so local uniqueness proves
invariance. Its conserved energy density is
\[
 E=3\dot h^2+\dot k^2+\dot\beta^2+V.
\]
If \(r^2=k^2+\beta^2\), then
\[
 |2h(r^2-h^2)|\leq2|h|(h^2+r^2)
 \leq \frac{s^{3/2}}{\sqrt6},\qquad
 V\geq\frac{\Gamma}{4}s^2-\frac{|\kappa|}{\sqrt6}s^{3/2}.
\]
Here \(|h|\leq\sqrt{s/6}\) and \(h^2+r^2\leq s/2\).
This lower bound is coercive, so a fixed energy bounds all positions
and velocities. The polynomial first-order system remains in a
compact set on every finite interval and extends globally.
This is an ODE existence proof, not a general nonlinear PDE
global-existence claim.
\end{proof}

On \(r>0\), introduce \(k=r\cos(2\theta)\),
\(\beta=-r\sin(2\theta)\). The original kinetic density and homogeneous
equations become
\begin{align}
 \mathcal L_{\rm kin}
   &=3\dot h^2+\dot r^2+4r^2\dot\theta^2,\notag\\
 \ddot h&=-ah-\kappa h^2+\kappa r^2/3,\qquad
 \ddot r=4r\dot\theta^2-(a-2\kappa h)r,\label{q444:vk-polar}\\
 (r^2\dot\theta)^{\boldsymbol\cdot}&=0,\qquad
 a=m^2+\Gamma(6h^2+2r^2).\notag
\end{align}
These follow by differentiating \(k,\beta\), resolving the two
equations into radial and angular components, and using
\(k\dot\beta-\beta\dot k=-2r^2\dot\theta\).
The Cartesian sector includes \(r=0\); the polar chart does not.
Even permitting a varying positive amplitude
\(Q=\chi(t)(3v_\theta v_\theta^*-I_3)\) gives \(r=3h>0\).
Then the \(h\) equation implies
 \(3\ddot h=-(a-2\kappa h)r\), whereas the \(r\) equation
has the extra term \(4r\dot\theta^2\). Since \(\ddot r=3\ddot h\),
one obtains \(4r\dot\theta^2=0\). Thus changing this single amplitude
does not by itself make angular homogeneous motion consistent.

\begin{theorem}[Exact gauge-invariant periodic continuation]
\label{q444:vk-periodic}
Suppose \(0<m^2<\kappa^2/(27\Gamma)\), \(\Gamma>0,\kappa>0\).
Define
\[
 d_0=\sqrt{\kappa^2-24\Gamma m^2},\quad
 h_0=\frac{\kappa+d_0}{24\Gamma}=\frac{\chi}{2},\quad
 h_0<h<\frac{\kappa}{6\Gamma},\quad
 \Delta(h)=1-\frac{6\Gamma h}{\kappa}.
\]
Put
\[
 a(h)=\frac{m^2+12\Gamma h^2}{\Delta(h)},\qquad
 R^2(h)=\frac{3a(h)h}{\kappa}+3h^2,\qquad
 \omega^2(h)=a(h)-2\kappa h
   =\frac{m^2-2\kappa h+24\Gamma h^2}{\Delta(h)}.
\]
Take the positive roots \(R,\omega\). Then the sector fields
\[
 h(t)=h,\quad k(t)=R\cos(\omega t),\quad
 \beta(t)=R\sin(\omega t),\quad \mathcal A_\mu=0
\]
are nonconstant, spatially homogeneous, exact full gauge--scalar
solutions. The gauge-invariant observable
\(\|b(t)\|^2=R^2\sin^2(\omega t)\) is nonconstant.
Their conserved energy density is exactly
\begin{equation}\label{q444:vk-periodic-energy}
 E(h)=R^2\omega^2+
 m^2(3h^2+R^2)+\Gamma(3h^2+R^2)^2-2\kappa h(R^2-h^2).
\end{equation}
Let
\[
 \Delta_0=1-\frac{6\Gamma h_0}{\kappa}
          =\frac{3\kappa-d_0}{4\kappa}>0,\qquad
 V_0=12m^2h_0^2+144\Gamma h_0^4-16\kappa h_0^3 .
\]
As \(h\downarrow h_0\),
\begin{align}
 R&\longrightarrow3h_0,\qquad \omega\longrightarrow0,\notag\\
 \omega^2&=\frac{2d_0}{\Delta_0}(h-h_0)
                            +O((h-h_0)^2),\notag\\
 E(h)-V_0&=\frac{18h_0^2d_0}{\Delta_0}(h-h_0)
                            +O((h-h_0)^2)
          =9h_0^2\omega^2+O(\omega^4).\label{q444:vk-excess}
\end{align}
These are finite-volume classical statements. On a flat torus with
retained volume \(\mathcal V_3>0\), the total energy and its excess are
\(\mathcal V_3 E(h)\) and \(\mathcal V_3(E(h)-V_0)\).
The amplitude tends to the nonzero value \(3h_0\): the family approaches
the minimum manifold while traversing a finite range of its invariant
\(\|b\|^2\), rather than becoming a small oscillation about one fixed
vacuum. Its period tends to infinity.
\end{theorem}
\begin{proof}
The denominator is positive on the stated interval; \(a>0\) and
\(R^2>0\). The numerator of \(\omega^2\) has roots
\((\kappa-d_0)/(24\Gamma)\) and \(h_0\); with its positive leading
coefficient it is positive for \(h>h_0\). Moreover
\(d_0<\kappa\), so \(h_0<\kappa/(12\Gamma)<\kappa/(6\Gamma)\);
the interval is nonempty. The identities
\[
 a=m^2+\Gamma(6h^2+2R^2),\quad
 -ah-\kappa h^2+\kappa R^2/3=0,\quad
 \omega^2=a-2\kappa h
\]
follow by substitution. They are precisely the constant-\(h\)
and the two oscillatory equations in \eqref{q444:vk-sector-eom}.
The previous proposition proves all gauge equations.
The triplet norm is preserved by the full \(SU(3)\) action, so its
nonconstant value rules out a purely gauge motion.

Substitution in the conserved energy gives
\eqref{q444:vk-periodic-energy}. At \(h=h_0\),
\(m^2=2\kappa h_0-24\Gamma h_0^2\), so
\(a=2\kappa h_0\), \(R^2=9h_0^2\) and \(\omega^2=0\).
The limiting field at any phase is a global minimum: the
\(2\)-by-\(2\) quaternionic block has eigenvalues \(h_0\pm3h_0\),
and the last eigenvalue is \(-2h_0\); hence the full spectrum is
\((4h_0,-2h_0,-2h_0)\).
Writing the frequency numerator as \(n(h)\),
\(n(h_0)=0\) and \(n'(h_0)=2d_0\). Taylor expansion proves the
frequency estimate. The curve \(h\mapsto(h,R(h),0)\) is smooth
near \(h_0\), and its endpoint is a critical point of \(V\).
Consequently its potential differs from \(V_0\) by
\(O((h-h_0)^2)\). Its kinetic density is
\[
 R^2\omega^2
 =\left(9h_0^2+O(h-h_0)\right)
   \left(\frac{2d_0}{\Delta_0}(h-h_0)+O((h-h_0)^2)\right).
 \]
This proves both energy expansions. Since the leading coefficient
of \(\omega^2\) is strictly positive, replacing \(h-h_0\) by its
smooth inverse as a function of \(\omega^2\) gives the last equality.
For every \(h>h_0\), \(E(h)>V_0\) because the kinetic term is
positive and the potential is bounded below by \(V_0\).
The torus assertion is integration of a constant spatial density.
Nonzero homogeneous energy density on all of \(\mathbb R^3\) is
not asserted to have finite total energy.
\end{proof}

\begin{proposition}[Linearized physical tangent and exact null waves]
\label{q444:vk-wave}
At each fixed vacuum \(q_{\theta_0}\), the perturbation
\[
 \delta\Phi=T_\Phi(\theta_0)\eta,\qquad
 \delta b=T_b(\theta_0)\eta,\qquad
 \delta\mathcal A_\mu=0
\]
solves the full linearized equations precisely when \(\Box\eta=0\).
Its quadratic density is
\[
 9\chi^2\bigl((\partial_t\eta)^2-|\nabla\eta|^2\bigr).
\]
For a homogeneous mode on a torus of volume \(\mathcal V_3\), its
momentum and quadratic Hamiltonian are
\[
 \pi_\eta=18\chi^2\mathcal V_3\dot\eta,\qquad
 H_\eta=\frac{\pi_\eta^2}{36\chi^2\mathcal V_3}.
\]
There is also an exact nonlinear inhomogeneous continuation at fixed
\(\chi\): for any smooth real \(f\),
\[
 \theta(t,x)=f(t-x_1),\qquad q(t,x)=q_{\theta(t,x)},\qquad
 \mathcal A_\mu=0
\]
solves the full nonlinear field equations wherever these fields are
defined. On the rectangular flat torus with sides \(L_1,L_2,L_3\),
choose \(f\) \(L_1\)-periodic. Its excess over the vacuum energy is
\[
 E_{\rm wave}-\mathcal V_3V_0
 =18\chi^2L_2L_3\int_0^{L_1}(f'(u))^2\,du,\qquad
 \mathcal V_3=L_1L_2L_3.
\]
\end{proposition}
\begin{proof}
Differentiate \(\nabla V(q_\theta)=0\) in \(\theta\). The metric
Hessian annihilates \(T\). The linearized currents are
\[
 \bigl([\Phi_{\theta_0},T_\Phi]+b_{\theta_0}T_b^*
                         -T_b b_{\theta_0}^*\bigr)\partial_\mu\eta=0.
\]
Thus the linearized scalar equation is \(T\Box\eta=0\), and \(T\ne0\).
The gauge equations have zero current and zero field strength.
The kinetic norm \(18\chi^2\) proves the quadratic density and its
Legendre transform with the stated volume.

For the nonlinear assertion, the chain rule in signature \((+---)\)
gives
\[
 \Box q_\theta=T(\theta)\Box\theta
                 +N(\theta)\partial_\mu\theta\,\partial^\mu\theta .
\]
Both coefficients vanish for \(\theta=f(t-x_1)\). Every value remains
a critical point of \(V\); the trigonometric formulas define it for
all real \(\theta\), or one can restrict the range of \(f\) to the
displayed minimum section. The exact currents vanish for all
spacetime derivatives as in the sector proposition. This verifies
every nonlinear equation. The scalar energy density is the sum of
one temporal and one spatial term, each \(9\chi^2(f')^2\), plus \(V_0\).
Periodicity makes its integral independent of time and gives the
formula. For
\(f(u)=\theta_0+\varepsilon\cos(2\pi u/L_1)\) with range inside the
section, the excess is
\(9\chi^2\mathcal V_3\varepsilon^2(2\pi/L_1)^2\).
\end{proof}

\paragraph{Gauge-invariant coordinate and endpoint scope.}
The invariant
\[
 I=\|b_\theta\|^2=\frac{9\chi^2}{4}\sin^2(2\theta),\qquad
 \frac{dI}{d\theta}=\frac{9\chi^2}{2}\sin(4\theta)
\]
has nonzero derivative on \(0<\theta<\pi/4\). Therefore the displayed
linearized mode is visible through the genuine gauge-invariant
linear observable
\(\delta I=(9\chi^2/2)\sin(4\theta_0)\eta\) in the interior.
With \(M=9\chi^2/4\), its interior coordinate expression is
\[
 9\chi^2(\partial\theta)^2
       =\frac{9\chi^2}{16I(M-I)}(\partial I)^2 .
\]
The denominator vanishes at the endpoints because \(I\) ceases to be
a regular linear coordinate there, not because the norm of \(T\)
diverges.

For completeness the stabilizers on this section can be checked
without importing an orbit classification. In the interior the
three diagonal entries of \(\Phi_\theta\) are distinct, so a
stabilizer must be diagonal. Fixing the nonzero \(b_3\) and its
determinant then gives
\(\operatorname{diag}(e^{it},e^{-it},1)\), of dimension one.
At \(\theta=0\), \(b=0\) and the eigenvalues are \(2\chi,-\chi,-\chi\),
giving \(S(U(1)\times U(2))\), of dimension four.
At \(\theta=\pi/4\), fixing \(b_3\ne0\) and the two-dimensional
first eigenspace gives \(\operatorname{diag}(U,1)\), \(U\in SU(2)\),
of dimension three. The gauge-orbit dimensions are consequently
\(7,4,5\), respectively. The fixed-\(\chi\) minimum manifold has
dimension eight: the unit sphere in \(\mathbb H^3\) has dimension
eleven and the free unit-quaternion action \(v\mapsto va\) has
dimension three, with \(vv^*\) precisely its quotient.
Thus the horizontal tangent dimensions within that manifold are
\(1,4,3\), respectively. In particular the interior has exactly one
horizontal vacuum direction. At the endpoints the invariant
linear function \(I\) vanishes to first order; no assertion of a
smooth one-dimensional linear quotient there is made.
This agrees with the classical cohomogeneity-one description in
A.~Gambioli, \emph{Eight-dimensional \(SU(3)\)-manifolds of
cohomogeneity one}, arXiv:math/0611796, Section~4, Examples:
\(\mathbb{HP}^2\) has singular orbit types \(\mathbb{CP}^2\) and
\(S^5\). The kinetic and current statements above are independently
proved here, not imported from that classification.

\paragraph{Fixed observation and classical scope.}
Retain
\[
 P=\begin{pmatrix}1&-1&0\\0&1&-1\end{pmatrix},\qquad
 u=(1,1,1)^{\mathsf T}.
\]
The positive spectral-floor matrix \(H_\theta=Q_\theta+\chi I_3\)
has complex part \(A_\theta=3\chi
\operatorname{diag}(\cos^2\theta,\sin^2\theta,0)\). Consequently
\[
 G_\theta=PA_\theta P^*
 =3\chi\begin{pmatrix}1&-\sin^2\theta\\
                  -\sin^2\theta&\sin^2\theta\end{pmatrix},\qquad
 \rho=u^*b_\theta=-3\chi\cos\theta\sin\theta,\qquad
 \det G_\theta=|\rho|^2=I.
\]
These are direct matrix products with the original \(P,u,\chi\).
They do not assert gauge invariance of a fixed-\(P\) observation away
from this chosen section: under a general gauge transformation its
complex part is \(gA_\theta g^*\), so the observation is
\(PgA_\theta g^*P^*\), not generally \(PA_\theta P^*\).

On a rectangular torus the linearized wave equation has classical
Fourier frequencies
\(\omega_{\mathbf n}^2=4\pi^2\sum_i n_i^2/L_i^2\);
the spatial zero mode has zero linearized frequency. The exact
homogeneous periodic family additionally has arbitrarily small
positive classical frequency and energy excess at fixed volume.
The null waves give exact nonlinear inhomogeneous solutions with
arbitrarily small amplitude. None of these statements quantizes the
model, identifies a Hilbert-space spectrum, proves a quantum mass
gap or its absence, or assigns physical units that the retained
parameters have not been supplied. No finite-total-energy assertion
on all of \(\mathbb R^3\) is inferred from a finite-volume solution.

\subsection{The gauge-vector spectrum and an invariant vector observable}

We now compute masses in the linearized classical field equations of
\eqref{q444:full-gauge-density}, not by assigning a physical meaning to
a determinant in isolation. Retain a positive nonzero minimum amplitude
$\chi$ from Theorem~\ref{q444:vacua}; thus $\kappa>0$,
$0\leq m^2\leq\kappa^2/(27\Gamma)$ and
\[
 \chi=\frac{\kappa+\sqrt{\kappa^2-24\Gamma m^2}}{12\Gamma}.
\]
The nonzero branch is meant at coexistence. Keep $0\leq\theta\leq\pi/4$,
and introduce the notation $q=\cos(2\theta)$, not a change in the field
or its kinetic metric. The same minimum representatives are
\[
 \Phi_\theta=\chi\operatorname{diag}\left(
       \frac{1+3q}{2},\frac{1-3q}{2},-1\right),\qquad
 b_\theta=-\frac{3\chi}{2}\sin(2\theta)e_3,\qquad
 \beta^2:=b_\theta^*b_\theta=\frac{9\chi^2}{4}(1-q^2).
\]
Every constant representative together with $\mathcal A_\mu=0$
solves the full classical field equations.

\begin{theorem}[All eight gauge-inertia eigenvalues]
\label{q444:gauge-spectrum}
Use the positive inner product
$\langle Z,W\rangle=-\operatorname{tr}(ZW)$ on $\mathfrak{su}(3)$.
At the displayed minimum, the exact spectrum of
$\mathscr I_{\Phi_\theta,b_\theta}$ is
\[
\begin{array}{c|c|c}
\text{generator space}&\text{inertia eigenvalue}&
                         \text{real multiplicity}\\ \hline
i\operatorname{diag}(1,-1,0)&0&1\\
i\operatorname{diag}(1,1,-2)&3\chi^2(1-q^2)&1\\
\langle E_{12}-E_{21},\,i(E_{12}+E_{21})\rangle
                         &9\chi^2q^2&2\\
\langle E_{13}-E_{31},\,i(E_{13}+E_{31})\rangle
                         &\frac92\chi^2(1+q)&2\\
\langle E_{23}-E_{32},\,i(E_{23}+E_{32})\rangle
                         &\frac92\chi^2(1-q)&2
\end{array}                                                    \tag{GS.1}
\]
Here $E_{ij}$ are the actual matrix units. The unscaled generators in
the row order with the six off-diagonal generators first have kinetic
Gram matrix $\operatorname{diag}(2,2,2,2,2,2,2,6)$.
The squared masses of transverse solutions of the full linearized
field equations are exactly $g^2/2$ times the eigenvalues in (GS.1).
In particular the diagonal nonzero branch has
\[
 \mu_{\mathrm d}^{\,2}
   =\frac32 g^2\chi^2(1-q^2)=\frac{2g^2}{3}\|b_\theta\|^2.
                                                               \tag{GS.2}
\]
There is one zero eigenvalue for $0<q<1$, four at $q=1$, and three
at $q=0$. These are the Lie-algebra dimensions of the stabilizers
$U(1)_{1,-1,0}$, $S(U(1)\times U(2))$, and $SU(2)$, respectively.
\end{theorem}
\begin{proof}
For a diagonal Hermitian $\Phi=\operatorname{diag}(d_1,d_2,d_3)$,
the operator $[\Phi,[\Phi,Z]]$ multiplies the $ij$ off-diagonal entry
by $(d_i-d_j)^2$. Here $bb^*=\operatorname{diag}(0,0,\beta^2)$.
Consequently $\{bb^*,Z\}$ contributes zero on the $12$ root space
and $\beta^2$ on each of the $13$ and $23$ root spaces; these entries
have zero trace, so no trace-free correction changes them. Retaining
the displayed differences gives
\[
 (d_1-d_2)^2=9\chi^2q^2,\quad
 (d_1-d_3)^2+\beta^2=\frac92\chi^2(1+q),\quad
 (d_2-d_3)^2+\beta^2=\frac92\chi^2(1-q).
\]
For $Z=i\operatorname{diag}(z_1,z_2,z_3)$, $\sum z_i=0$, the
diagonal contribution is
\[
 \mathscr I Z=2i\beta^2z_3
       \left(E_{33}-\frac13I_3\right).
\]
It annihilates $i\operatorname{diag}(1,-1,0)$ and multiplies
$i\operatorname{diag}(1,1,-2)$ by $4\beta^2/3
=3\chi^2(1-q^2)$. The eight displayed generators are linearly
independent and span $\mathfrak{su}(3)$; multiplication and trace give
their stated diagonal Gram matrix. This proves the entire generalized
Hessian spectrum without replacing that metric by an identity matrix.

For the dynamical assertion, perturb only the connection at first order:
$\mathcal A_\nu=\varepsilon Z a_\nu$, with both scalar perturbations zero.
At the constant scalar background the first-order scalar equations
reduce to
$[Z,\Phi_\theta]\partial_\nu a^\nu=0$ and
$Zb_\theta\,\partial_\nu a^\nu=0$.
The first-order current is
$(J^\nu)_{\rm tf}=-\varepsilon\mathscr I(Z)a^\nu$.
For an inertia eigenvector with eigenvalue $\lambda$, the full gauge
equation is therefore
\[
 (\Box+g^2\lambda/2)a^\nu-\partial^\nu(\partial_\mu a^\mu)=0.
                                                               \tag{GS.3}
\]
Choose $a_0=0$ and
$a_j=e_j\cos(k\cdot x-\omega t)$ with $k\cdot e=0$.
Then $\partial_\mu a^\mu=0$ and all scalar equations hold. The temporal
gauge equation is $\sum_j\partial_j\dot a_j=0$, so Gauss is satisfied
as well, not omitted. The remaining equations give
$\omega^2=|k|^2+g^2\lambda/2$. At nonzero frequency or spatial momentum,
the corresponding field strength is nonzero. Thus these are actual
transverse solutions, not pure-gauge variations.
For a specified rectangular flat spatial torus,
$k_j=2\pi n_j/L_j$, with every length $L_j$ retained.
Their positive kinetic sign follows from
$\langle Z,Z\rangle/g^2>0$ in the original gauge density.

For $0<q<1$, every displayed eigenvalue except the first is positive.
At $q=1$ the second and the $23$ pair additionally vanish; at $q=0$
the $12$ pair additionally vanishes. A stabilizer satisfies
$[Z,\Phi_\theta]=0$ and $Zb_\theta=0$ by the proved inertia kernel.
For $0<q<1$ distinct diagonal entries and nonzero $b_\theta$ give
$\operatorname{diag}(e^{i\alpha},e^{-i\alpha},1)$.
At $q=1$, $b_\theta=0$ and the last two entries of $\Phi_\theta$
coincide, giving the stated $U(2)$ subgroup. At $q=0$ the first two
entries coincide and $b_\theta\ne0$ fixes $e_3$, giving the upper-left
$SU(2)$. This proves the full group statements, including the endpoints.
\end{proof}

\begin{proposition}[An exact invariant current and its mass branch]
\label{q444:current-mass}
The full nonlinear gauge--scalar theory has the real, gauge-invariant
conserved vector current
\[
 j_\nu=-i\left(b^*D_\nu b-(D_\nu b)^*b\right),\qquad
 \partial_\nu j^\nu=0.                                      \tag{GS.4}
\]
At any displayed vacuum with $\beta^2>0$, the transverse perturbation
$\mathcal A_\nu=\varepsilon i\operatorname{diag}(1,1,-2)a_\nu$,
with zero scalar perturbations, has
\[
 j_\nu=-4\varepsilon\beta^2 a_\nu+O(\varepsilon^2).
                                                               \tag{GS.5}
\]
It therefore detects the branch (GS.2) in the exact linearized system.
For the positive rank-one spectral-floor matrix
$H_\theta=Q_\theta+\chi I_3=\Theta(A_\theta,b_\theta)$,
\[
 \boxed{\displaystyle
 \mu_{\mathrm d}^{\,2}
 =\frac{2g^2}{3V_\mu}
       \int_{SU(3)}\det(PgA_\theta g^*P^*)\,d\mu(g).}          \tag{GS.6}
\]
The Haar volume $V_\mu>0$ and the original marking $P$ are retained.
At $\beta=0$ the linear coefficient in (GS.5) vanishes; that endpoint
does not supply a nonzero linear interpolator through this current.
\end{proposition}
\begin{proof}
Under $b\mapsto Ub$ and $D_\nu b\mapsto U D_\nu b$, each scalar product
in (GS.4) is unchanged. Their difference is purely imaginary, so $j_\nu$
is real. It is the Noether current for the constant phase action
$b\mapsto e^{i\alpha}b$: allowing $\alpha$ to depend on spacetime
produces $j^\nu\partial_\nu\alpha$ in the variation of the kinetic term.
The potential is phase invariant. Directly, the covariant product rule
gives
\[
 \partial_\nu j^\nu
 =-i\left(b^*D_\nu D^\nu b-(D_\nu D^\nu b)^*b\right).
\]
The exact equation is $D_\nu D^\nu b=-(aI_3+\kappa\Phi)b$;
the matrix in parentheses is Hermitian. The two scalar products
are therefore equal, proving conservation in the full interacting theory.

At the specified constant background,
$b_\theta^*\mathcal A_\nu b_\theta
=-2i\varepsilon\beta^2a_\nu$. Subtracting its conjugate and multiplying
by $-i$ gives $-4\varepsilon\beta^2a_\nu$, with no rescaling of the
Cartan generator. Every off-diagonal generator, and
$i\operatorname{diag}(1,-1,0)$, has zero $33$ entry, so none contributes
to this linear current when scalar perturbations are zero.
For an arbitrary scalar perturbation $\delta b$, its additional linear
term is $2\operatorname{Im}(b_\theta^*\partial_\nu\delta b)$,
a pure gradient at the constant background; the transverse sector
used in (GS.3) sets it to zero. Thus the nonzero transverse mode is
seen by a gauge-invariant observable.
Finally Proposition~\ref{q444:vacuum-average} proves the integral in
(GS.6) equals $V_\mu\beta^2$. Substitution into (GS.2) proves the
exact relation with all coefficients and domains fixed.
\end{proof}

This follows the classical Higgs-phase use of gauge-invariant vector
operators discussed by D.~Tong, \emph{Lectures on Gauge Theory},
Section~2.7.3, especially the fundamental-scalar example on printed
pages~107--108,
\url{https://www.damtp.cam.ac.uk/user/tong/gaugetheory/gt.pdf}.
The eight-eigenvalue calculation, current coefficient, and marked
determinant identity here are derived for this specified adjoint-plus-
fundamental model; they are not transferred from that example.
Equation (GS.6) is a relation to a linearized gauge-vector mass in this
model, not an identification of all entries of $G$ with spacetime momentum,
and not a spectral theorem for pure interacting quantum Yang--Mills.

\begin{corollary}[Variation of the massive branches along the minimum family]
\label{q444:massive-branch-bound}
For $0<q<1$, the smallest strictly positive inertia eigenvalue is
\[
 \lambda_{\rm pos}(q)=
 \begin{cases}
 9\chi^2q^2,&0<q\leq1/2,\\
 \frac92\chi^2(1-q),&1/2\leq q<1 .
 \end{cases}                                                 \tag{GS.7}
\]
It tends to zero at both ends and has unique maximum $9\chi^2/4$
at $q=1/2$. At that point five eigenvalues coincide at $9\chi^2/4$,
two equal $27\chi^2/4$, and one is zero. This positive degeneracy
does not enlarge the stabilizer. The retained spectral traces are
\[
 \operatorname{tr}\mathscr I=\chi^2(36-15\delta),\qquad
 \operatorname{tr}(\mathscr I^2)
       =\chi^4(324-405\delta+171\delta^2),\qquad \delta=1-q^2.
                                                               \tag{GS.8}
\]
\end{corollary}
\begin{proof}
The $13$ eigenvalue is at least the $23$ eigenvalue. Subtracting the
$12$ eigenvalue from the diagonal and $23$ eigenvalues gives
$3\chi^2(1-4q^2)$ and
$\frac92\chi^2(1-q-2q^2)$, respectively; both are nonnegative
exactly for $0\leq q\leq1/2$ on the stated interval.
For $q\geq1/2$, the diagonal minus the $23$ eigenvalue is
$\frac32\chi^2(1-q)(2q-1)\geq0$.
These comparisons prove (GS.7). The first expression increases strictly
and the second decreases strictly, proving the maximum and endpoint
limits. Substitution of $q=1/2$ in (GS.1) gives the multiplicities,
while its kernel remains one-dimensional. Adding the eight eigenvalues
and their squares gives $\chi^2(21+15q^2)$ and
$\chi^4(90+63q^2+171q^4)$, respectively. Replacing the notation
$q^2$ by $1-\delta$ proves (GS.8).
\end{proof}

There is always a zero transverse branch in this linearized system.
The strictly positive branches alone also have no positive lower bound
uniform over the minimum family. These statements concern the proved
quadratic spectrum and its exact parameters. They do not determine
radiative selection of a minimum, a renormalized continuum theory,
or the interacting Hamiltonian spectrum.

\subsection{The retained system and the scope of the result}

The literal constant-matter data are
\[
 H=\begin{pmatrix}0&0&0\\0&-1/2&1/(2\sqrt2)\\
 0&1/(2\sqrt2)&1/2\end{pmatrix},\qquad
 q=\frac{e_1}{2\sqrt2},\qquad
 A=\sqrt2\begin{pmatrix}0&-1&0\\1&0&0\\0&0&0\end{pmatrix}.
\]
The original scalar couplings are $m^2=\Gamma=1$, $\kappa=6$;
$s=\operatorname{tr} H^2+2q^*q=1$ and $a=m^2+\Gamma s=2$.
Write $g^2>0$ for the gauge coupling and keep $\omega\ne0$ real.
Every unknown $C,Z_1,Z_2$ belongs to $\mathfrak{su}(3)$, the real vector space
of traceless anti-Hermitian complex $3\times3$ matrices.
For any matrix $M$, set $M_{\mathrm{tf}}=M-(\operatorname{tr} M)I_3/3$, and define
\[
 \mathcal I X=([H,[H,X]]+qq^*X+Xqq^*)_{\mathrm{tf}},\qquad
 \mathcal KX=-[Z_1,[Z_1,X]]-[Z_2,[Z_2,X]],
\]
\[
 F=[A,[A,H]]=-2H+\operatorname{diag}(-2,1,1).
\]
The last identity follows from $Hq=0$,
$H^2=(3/8)\operatorname{diag}(0,1,1)$ and
$6(\operatorname{adj}H-qq^*)_{\mathrm{tf}}=\operatorname{diag}(-2,1,1)$.
The five equations to be retained are exactly
\begin{align}
 [C,[C,H]]-[Z_1,[Z_1,H]]-[Z_2,[Z_2,H]]&=F,\label{q444:cx:eq:phi}\\
 (C^2-Z_1^2-Z_2^2)q&=-2q,\label{q444:cx:eq:b}\\
 (\mathcal I+2\mathcal K/g^2)C&=(4\omega/g^2)[Z_1,Z_2],\label{q444:cx:eq:gauss}\\
 -\omega^2Z_1-2\omega[C,Z_2]+[C,[C,Z_1]]
             +\mathcal KZ_1+(g^2/2)\mathcal IZ_1&=0,\label{q444:cx:eq:sp1}\\
 -\omega^2Z_2+2\omega[C,Z_1]+[C,[C,Z_2]]
             +\mathcal KZ_2+(g^2/2)\mathcal IZ_2&=0.\label{q444:cx:eq:sp2}
\end{align}

This section does not solve or exclude all triples in $\mathfrak{su}(3)^3$.
It proves necessary identities for every such triple and a no-solution
theorem for the precisely stated complex block class below. No remaining
general equation is assumed true, and no numerical candidate is a result.

\subsubsection*{Exact map to the original literal fields}
For clarity, keep
\[
 u=(1,1,1)^{\mathsf T},\quad v_0=(1,-1,0)^{\mathsf T},\quad
 w_0=(1,1,-2)^{\mathsf T},\quad
 K=\left(v_0/\sqrt2,w_0/\sqrt6,u/\sqrt3\right),
\]
\[
 \Omega=\sqrt{2/3}\begin{pmatrix}0&-1&1\\1&0&-1\\-1&1&0\end{pmatrix},
 \quad O(t)=e^{t\Omega},\quad
 \Phi_0=\frac14\begin{pmatrix}1&1&1\\1&1&1\\1&1&-2\end{pmatrix},
 \quad b_0=\frac14(1,-1,0)^{\mathsf T}.
\]
The dot products and determinant give $K\in SO(3)$, and multiplication
gives $K^*\Phi_0K=H$, $K^*b_0=q$, $K^*\Omega K=A$.
The prescribed scalar fields remain
$\Phi(t)=O(t)\Phi_0O(t)^*$, $b(t)=O(t)b_0$; their frequency is $\sqrt2$.
For a solution of \eqref{q444:cx:eq:phi}--\eqref{q444:cx:eq:sp2}, set
\[
 Y_1=\cos(\omega t)Z_1-\sin(\omega t)Z_2,\quad
 Y_2=\sin(\omega t)Z_1+\cos(\omega t)Z_2,\quad Y_3=0,
\]
\begin{equation}\label{q444:cx:eq:original-inverse}
 \mathcal A_0=OKCK^*O^*-\Omega,\qquad
 \mathcal A_i=OKY_iK^*O^*.
\end{equation}
The inverse is $C=K^*O^*\mathcal A_0OK+A$ and
$Y_i=K^*O^*\mathcal A_iOK$. These are gauge transformations by
$K^*O^*$, including the inhomogeneous temporal term.

Here is a direct verification of the complete finite reduction. In the
constant-matter frame put $E_i=\dot Y_i+[C,Y_i]$. The two scalar
equations are
$[C,[C,H]]-\sum_i[Y_i,[Y_i,H]]=F$ and
$(C^2-\sum_iY_i^2)q=-2q$.
The projected currents are $J_0=-\mathcal IC$ and $J_i=-\mathcal IY_i$.
Gauss is $(2/g^2)\sum_i[Y_i,E_i]+\mathcal IC=0$, and the spatial
equations are
\[
 \dot E_i+[C,E_i]-\sum_j[Y_j,[Y_j,Y_i]]+(g^2/2)\mathcal IY_i=0.
\]
These follow by homogeneous substitution in the covariant equations with
$D_\mu=\partial_\mu+\mathcal A_\mu$, signature $(+---)$ and gauge
density $\operatorname{tr}(\mathcal F_{\mu\nu}\mathcal F^{\mu\nu})/(2g^2)$.
Specifically the adjoint derivative is a commutator, and
$J_\mu=[\Phi,D_\mu\Phi]+b(D_\mu b)^*-(D_\mu b)b^*$; at constant
matter it is minus the displayed inertia. Thus no temporal constraint
has been removed.
The rotation preserves $\sum_iY_i^2$ and every quadratic commutator sum,
and
$\sum_i[Y_i,\dot Y_i]=-2\omega[Z_1,Z_2]$.
Applying $(\partial_t+\operatorname{ad}_C)^2$ to the rotating pair
gives the terms $-\omega^2Z_i$, the two displayed terms with coefficient
$2\omega$, and $[C,[C,Z_i]]$. The third spatial equation is identically
zero. This proves equivalence with all five equations above at every time.
The product rule proves covariance under \eqref{q444:cx:eq:original-inverse},
so this equivalence also holds in the original coordinates.

\subsection{Universal necessary identities in the unrestricted complex class}

Use $\langle X,Y\rangle=-\operatorname{tr}(XY)$ and $\|X\|^2=-\operatorname{tr} X^2$ on
$\mathfrak{su}(3)$. This is a real positive inner product because $X^*=-X$.
Define the nonnegative real quantities
\[
 c=\langle C,\mathcal IC\rangle,\quad
 z=\sum_{j=1}^2\langle Z_j,\mathcal IZ_j\rangle,\quad
 k=\sum_{j=1}^2\|[C,Z_j]\|^2,\quad
 d=\|[Z_1,Z_2]\|^2,\quad z_0=\|Z_1\|^2+\|Z_2\|^2,
\]
and the signed real number $b=\langle C,[Z_1,Z_2]\rangle$.

\begin{proposition}\label{q444:cx:prop:universal}
Every solution of \eqref{q444:cx:eq:phi}--\eqref{q444:cx:eq:sp2} satisfies
\begin{equation}\label{q444:cx:eq:universal}
 c-z=2,\qquad (g^2/2)c+k=2\omega b,\qquad
 \omega^2z_0=k+2d+g^2c+(g^2/2)z.
\end{equation}
In particular $c,k,d,z_0$ are strictly positive, $\omega b>0$, and
\begin{equation}\label{q444:cx:eq:necessary-inequality}
 8kd+2g^2kz\geq g^4c^2.
\end{equation}
\end{proposition}
\begin{proof}
Cyclicity and anti-Hermiticity give
\[
 \langle X,\mathcal IX\rangle=\operatorname{tr}[X,H]^2+2|Xq|^2,
 \qquad
 \langle X,\mathcal KX\rangle=\sum_j\|[Z_j,X]\|^2.
\]
Trace projection pairs trivially with traceless $X$. The first
commutator term follows by cyclically moving $H$ twice, and the vector
term is $-2q^*X^2q=2|Xq|^2$. The second identity follows by applying
$\langle X,[Z,Y]\rangle=-\langle[Z,X],Y\rangle$ twice.
Because $[X,H]$ is Hermitian, all terms are nonnegative.

Pair \eqref{q444:cx:eq:phi} with $H$ using ordinary trace and pair
\eqref{q444:cx:eq:b} with $2q^*$, taking the real part. Their left sides add
to $-c+z$. Their right sides add to
$\operatorname{tr}(HF)-4q^*q=-2\operatorname{tr} H^2-4q^*q=-2$, proving the first identity.
Pair \eqref{q444:cx:eq:gauss} with $C$ to obtain the second.
Pair \eqref{q444:cx:eq:sp1} with $Z_1$ and \eqref{q444:cx:eq:sp2} with $Z_2$.
The gyroscopic terms sum to $4\omega b$, the double $C$ terms sum to
$-k$, and the spatial double commutators sum to $2d$.
The result is
$-\omega^2z_0+4\omega b-k+2d+(g^2/2)z=0$.
Substitute the second identity to obtain the third.

The first identity gives $c\geq2$. The second gives $\omega b>0$.
Thus $b\ne0$, which forces $d>0$, $z_0>0$, and $k>0$, since
$b=\langle[C,Z_1],Z_2\rangle$.
Apply Cauchy--Schwarz to the ordered pairs
$([C,Z_1],[C,Z_2])$ and $(Z_2,-Z_1)$ to obtain
$4b^2\leq kz_0$.
The second identity gives
$\omega^2z_0\geq(k+g^2c/2)^2/k$.
Subtract this lower bound from the third identity and multiply by $4k$;
the result is exactly \eqref{q444:cx:eq:necessary-inequality}.
\end{proof}

These identities rule out any attempted solution with zero spatial
commutator or zero electric coupling to $C$. They do not rule out a
general noncommuting complex triple.

\subsection{A new genuinely complex block exclusion}

Let $P=q q^*/(q^*q)=\operatorname{diag}(1,0,0)$ and put
\[
 \mathfrak p=\{X\in\mathfrak{su}(3):PXP=(I-P)X(I-P)=0\},\qquad
 \mathfrak k=\{X\in\mathfrak{su}(3):[X,P]=0\}.
\]
Thus $\mathfrak p$ consists of arbitrary complex matrices exchanging
the $q$-line with its orthogonal complement; $\mathfrak k$ consists of
the block-diagonal matrices preserving that line. Neither definition
imposes reality of the entries.

\begin{theorem}[Cross--cross--block exclusion]\label{q444:cx:thm:block}
There is no solution of all five equations with
$C,Z_1\in\mathfrak p$ and $Z_2\in\mathfrak k$.
There is likewise no solution for which a constant real rotation of
the ordered pair $(Z_1,Z_2)$ puts it in that class while
$C\in\mathfrak p$.
The assertion allows all complex phases and every $g^2>0$, $\omega\ne0$.
\end{theorem}

\subsubsection*{Exact coordinates used only for the proof}
Set $x=1/2$, $d_0=1/(2\sqrt2)$, $r=\sqrt{x^2+d_0^2}=\sqrt6/4$, and
\[
 S=\frac1{\sqrt{2r(r+x)}}
       \begin{pmatrix}d_0&-(r+x)\\r+x&d_0\end{pmatrix},
 \qquad U=\operatorname{diag}(1,S).
\]
Multiplication gives $S^{\mathsf T}S=I_2$, $\det S=1$,
$U^*HU=\operatorname{diag}(0,r,-r)$, and $U^*q=q$.
All connections and $A$ are conjugated by $U^*$, with inverse
conjugation by $U$. Brackets, trace projection, the force, and the
inertia transform by these exact invertible maps; hence all five
equations are equivalent. The subspaces $\mathfrak p,\mathfrak k$
are preserved because $U$ commutes with $P$.
No field, period, kinetic coefficient, or coupling is divided out.

For $v\in\mathbb C^2$ define
\[
 X(v)=\begin{pmatrix}0&-v^*\\v&0_{2\times2}\end{pmatrix}.
\]
In either of the proof frames below, $H$ has zero first row and column,
and its lower block squares to $r^2I_2$. Since $q^*q=1/8$,
\begin{equation}\label{q444:cx:eq:cross-inertia}
 \mathcal IX(v)=(r^2+1/8)X(v)=\tfrac12X(v).
\end{equation}
Indeed $HX(v)H=0$, $H^2X(v)+X(v)H^2=r^2X(v)$, and the two
$qq^*$ terms sum to $X(v)/8$; all these matrices are traceless.

\begin{lemma}[Scalar classification of the block generator]\label{q444:cx:lem:scalar}
Under the hypotheses of Theorem~\ref{q444:cx:thm:block}, the two scalar
equations put the block generator into exactly one of the following
two forms by an explicitly invertible constant unitary map fixing $q$:
\begin{align*}
 \textnormal{(I)}\quad&H=\operatorname{diag}(0,r,-r),
      &&Z_2=i\operatorname{diag}(0,\sigma,-\sigma),\quad\sigma\in\mathbb R;\\
 \textnormal{(II)}\quad&H=\begin{pmatrix}0&0&0\\0&0&r\\0&r&0\end{pmatrix},
      &&Z_2=i\operatorname{diag}(e,0,-e),\quad e\ne0.
\end{align*}
In both cases write $C=X(v)$, $Z_1=X(w)$ and
$U_v=vv^*-ww^*=(u_{jk})$. In case (I) scalar compatibility is exactly
\[
 u_{11}=1-\delta>0,\quad u_{22}=1+\delta>0,
 \qquad\delta=\frac1{2r}=\sqrt{2/3}.
\]
In case (II) it is exactly
\[
 \operatorname{tr} U_v=2+e^2,\qquad \operatorname{Re}u_{12}=-\delta\ne0.
\]
\end{lemma}
\begin{proof}
Start in the displayed diagonal-$H$ frame. An arbitrary element of
$\mathfrak k$ is $Z_2=i\operatorname{diag}(\eta,B)$, where $B$ is
Hermitian, $\operatorname{tr} B=-\eta$, and $\eta$ is real. Write it without deleting
any coefficient as
\[
 B=-\frac\eta2I_2+
 \begin{pmatrix}z&x_1-i x_2\\x_1+i x_2&-z\end{pmatrix},
 \qquad x_1,x_2,z\in\mathbb R,
 \qquad D=\operatorname{diag}(r,-r).
\]
The lower block of the adjoint scalar equation is
\[
 -\{U_v,D\}+[B,[B,D]]=I_2-2D.
\]
Matrix multiplication gives
\[
 [B,[B,D]]=4r
 \begin{pmatrix}
 x_1^2+x_2^2&-z(x_1-i x_2)\\
 -z(x_1+i x_2)&-(x_1^2+x_2^2)
 \end{pmatrix}.
\]
Its off-diagonal equation is $z(x_1-i x_2)=0$.
The two diagonal equations give
$u_{11}-u_{22}=-1/r$ and
$\operatorname{tr} U_v=2+4(x_1^2+x_2^2)$.
The vector equation gives $\operatorname{tr} U_v=2+\eta^2$.
Thus $\eta^2=4(x_1^2+x_2^2)$.

If $x_1=x_2=0$, then $\eta=0$ and $B=\operatorname{diag}(z,-z)$.
This is case (I) with $\sigma=z$, and its diagonal scalar equations
give the two values stated. They are positive because $\delta^2=2/3<1$.

Otherwise $z=0$ and $e:=\eta\ne0$.
Let $\gamma=2(x_1-i x_2)/\eta$, so $|\gamma|=1$.
Choose either complex square root $\zeta$ of $-\bar\gamma$ and set
\[
 W=\frac\zeta{\sqrt2}\begin{pmatrix}\gamma&\gamma\\1&-1\end{pmatrix}.
\]
Then $|\zeta|=1$, $W^*W=I_2$, $\det W=1$,
$W^*BW=\operatorname{diag}(0,-e)$, and $W^*DW=r\left(
\begin{smallmatrix}0&1\\1&0\end{smallmatrix}\right)$.
These identities follow by multiplying the two columns, which are the
zero and $-e$ eigenvectors of $B$, with the stated phases.
Conjugation by $\operatorname{diag}(1,W)^*$, with inverse its adjoint,
therefore gives case (II), fixes $q$, and preserves both block spaces.

With $\sigma_x=\left(\begin{smallmatrix}0&1\\1&0\end{smallmatrix}\right)$,
in this frame the lower scalar equation is
$-\{U_v,r\sigma_x\}+e^2r\sigma_x=I_2-2r\sigma_x$.
Its diagonal entries give $\operatorname{Re}u_{12}=-1/(2r)$ and
its off-diagonal entries give $\operatorname{tr} U_v=2+e^2$.
The vector equation gives the same trace. Conversely direct substitution
sets every scalar residual entry to zero in each of the two cases.
\end{proof}

\subsubsection*{All gauge equations in the two frames}
In either case write
\[
 Z_2=i\operatorname{diag}(\eta,d_1,d_2),\quad
 n_j=\eta-d_j,\quad V_j=|v_j|^2,\quad W_j=|w_j|^2,
\]
\[
 V=V_1+V_2,\quad W_0=W_1+W_2,\quad N=V+W_0,
\]
\[
 p=v^*w=\alpha+i b,\quad b_j=\operatorname{Im}(\bar v_jw_j),
\]
\[
 A_j=\frac{g^2}{4}+n_j^2+W_0>0,\quad
 P_j=3b+2\omega n_j,\quad B_j=A_j-(\omega^2+N).
\]
Here $b$ is only the imaginary inner product in this section, not the
triplet or the signed pairing in Proposition~\ref{q444:cx:prop:universal}.
Direct block multiplication gives
\begin{equation}\label{q444:cx:eq:block-double}
 -[X(w),[X(w),X(v)]]
       =X\bigl(W_0v+w(\bar p-2p)\bigr),\qquad
 [X(v),Z_2]=X\bigl(i\operatorname{diag}(n_1,n_2)v\bigr).
\end{equation}
Using \eqref{q444:cx:eq:cross-inertia}, the lower first-column equations of
Gauss and the first spatial equation are exactly
\begin{equation}\label{q444:cx:eq:cross-gauge}
 A_jv_j=(\alpha+iP_j)w_j,\qquad
 B_jw_j=-(\alpha-iP_j)v_j\quad(j=1,2).
\end{equation}
There are no other nonzero blocks in these two equations.

The second spatial equation has only a block-diagonal residual.
Write $\mathcal IZ_2=i\operatorname{diag}(\ell_0,\ell_1,\ell_2)$;
it is diagonal in both frames. Its lower diagonal entries are
\begin{equation}\label{q444:cx:eq:spatial-diagonal}
 -\omega^2d_j+2n_j(V_j-W_j)+4\omega b_j+(g^2/2)\ell_j=0.
\end{equation}
Its upper diagonal entry is
\begin{equation}\label{q444:cx:eq:spatial-top}
 -\omega^2\eta-2\sum_jn_j(V_j-W_j)-4\omega b
                                    +(g^2/2)\ell_0=0.
\end{equation}
Its lower off-diagonal entry is
\begin{equation}\label{q444:cx:eq:spatial-off}
 i(n_1+n_2)(v_1\bar v_2-w_1\bar w_2)
                    +2\omega(-v_1\bar w_2+w_1\bar v_2)=0.
\end{equation}
For verification, the lower block of $[C,Z_1]$ is $-vw^*+wv^*$;
the lower block contributed by $[C,[C,Z_2]]+\mathcal KZ_2$ has
entry $i(n_j+n_k)(v_j\bar v_k-w_j\bar w_k)$.
Its upper entry is minus its lower trace. These multiplications give
\eqref{q444:cx:eq:spatial-diagonal}--\eqref{q444:cx:eq:spatial-off} with every sign and
coefficient. They retain the entire second spatial equation.

\begin{lemma}\label{q444:cx:lem:phase}
Every solution of these gauge equations has $\alpha=0$ and
\begin{equation}\label{q444:cx:eq:energy-block}
 \omega^2W_0=\frac{g^2}{4}N+\sum_jn_j^2(V_j+W_j),
\end{equation}
\begin{equation}\label{q444:cx:eq:uj}
 V_j-W_j=\frac{V_j+W_j}{\omega^2+N}
     \left(\omega^2+V-W_0-\frac{g^2}{2}-2n_j^2\right).
\end{equation}
\end{lemma}
\begin{proof}
Conjugate the first equation in \eqref{q444:cx:eq:cross-gauge} and multiply by
$w_j$. Summing gives
$p=\sum_j(\alpha-iP_j)W_j/A_j$.
Taking real parts gives $\alpha=\alpha\sum_jW_j/A_j$.
If $W_0=0$ the sum is zero. Otherwise $A_j>W_0$ gives
$\sum_jW_j/A_j<\sum_jW_j/W_0=1$. Hence $\alpha=0$ in all cases.
The two equations are now $A_jv_j=iP_jw_j$ and $B_jw_j=iP_jv_j$.
Their real pairings with $v_j$ and $w_j$ give
$A_jV_j=-P_jb_j$ and $B_jW_j=P_jb_j$.
Adding them gives
$A_j(V_j+W_j)=(\omega^2+N)W_j$.
Summing and inserting $A_j$ proves \eqref{q444:cx:eq:energy-block}.
Subtracting twice this last formula divided by $\omega^2+N$ from
$V_j+W_j$ proves \eqref{q444:cx:eq:uj}. The denominator is positive because
$\omega\ne0$ and $N\geq0$.
\end{proof}

\begin{proof}[Proof of Theorem~\ref{q444:cx:thm:block}]
Use Lemma~\ref{q444:cx:lem:scalar} and retain every gauge equation as just derived.

In case (I), $\eta=0$, $n_1=-\sigma$, $n_2=\sigma$,
$\mathcal IZ_2=0$, and $V_j-W_j=1\mp\delta>0$.
If $\sigma=0$, then $Z_2=0$ and Gauss paired with $C$ gives
$\langle C,\mathcal IC\rangle+(2/g^2)\langle C,\mathcal KC\rangle=0$.
Both terms are nonnegative and \eqref{q444:cx:eq:cross-inertia} forces $C=0$,
contradicting the two positive scalar differences. Thus $\sigma\ne0$.
Both $v_j$ are nonzero. Gauss then forces $w_j\ne0$ and $P_j\ne0$.
Since $\alpha=0$,
\[
 w_j=i t_jv_j,\qquad t_j=-A_j/P_j\in\mathbb R\setminus\{0\}.
\]
Here $A_1=A_2=g^2/4+\sigma^2+W_0$.
Because $n_1+n_2=0$, \eqref{q444:cx:eq:spatial-off} becomes
$2i\omega(t_1+t_2)v_1\bar v_2=0$.
Thus $t_1+t_2=0$, which implies $P_1+P_2=0$, or $6b=0$.
But \eqref{q444:cx:eq:spatial-top} gives
\[
 b=\frac{\sigma((1-\delta)-(1+\delta))}{2\omega}
   =-\frac{\sigma\delta}{\omega}\ne0.
\]
This contradiction excludes case (I).

In case (II), $\eta=e$, $d_1=0$, $d_2=-e$,
$n_1=e$, $n_2=2e$, and scalar compatibility gives
$V-W_0=2+e^2$ and
$\operatorname{Re}(v_1\bar v_2-w_1\bar w_2)=-\delta\ne0$.
In particular $V_j+W_j>0$ for both $j$.
The transformations $(Z_1,Z_2)\mapsto(-Z_1,-Z_2)$ and
$(C,\omega)\mapsto(-C,-\omega)$ each preserve all five equations,
as substitution verifies. They allow $e>0$ and $\omega>0$ without
changing the given matter fields or their original frequency.

Direct multiplication in this frame gives
\[
 \mathcal IZ_2=ie\operatorname{diag}(1/6,2/3,-5/6).
\]
The adjoint and trace-projected vector parts, respectively, are
\[
 ie\operatorname{diag}(0,2r^2,-2r^2),\qquad
 ie\operatorname{diag}(1/6,-1/12,-1/12).
\]
Their sum proves the formula with $r^2=3/8$.
The right side of \eqref{q444:cx:eq:energy-block} is positive, so $W_0>0$.
Since $N=2W_0+2+e^2>2W_0$ and $n_j^2\geq e^2$, that identity gives
\[
 \omega^2W_0\geq(g^2/4+e^2)N
             >(g^2/2+2e^2)W_0,
 \qquad \omega^2>g^2/2+2e^2.
\]
Formula \eqref{q444:cx:eq:uj} consequently implies
\[
 V_1-W_1=
 \frac{V_1+W_1}{\omega^2+N}
       (\omega^2+2-e^2-g^2/2)>0.
\]
The $j=1$ equation in \eqref{q444:cx:eq:spatial-diagonal} is exactly
\[
 2e(V_1-W_1)+4\omega b_1+eg^2/3=0,
\]
so $b_1<0$. Gauss gives $b_1=-P_1W_1/A_1$.
Since $V_1>0$, it also forces $W_1>0$ and hence $P_1>0$.
Now $P_2=P_1+2e\omega>0$; because $V_2+W_2>0$, Gauss forces
both $V_2,W_2>0$ as well.

Taking imaginary parts in the proof of Lemma~\ref{q444:cx:lem:phase} gives
\[
 b\left(1+3\sum_j\frac{W_j}{A_j}\right)
       =-2\omega\sum_j\frac{n_jW_j}{A_j}<0.
\]
Therefore $b<0$ and $0<P_j=3b+2\omega n_j<2\omega n_j$.
Put $\rho_j=A_j/P_j>0$. Since $A_j>n_j^2$, one has the strict lower bound
\begin{equation}\label{q444:cx:eq:lower-rho}
 \rho_1+\rho_2>
 \frac{n_1+n_2}{2\omega}=\frac{3e}{2\omega}.
\end{equation}
Gauss also gives $w_j=-i\rho_jv_j$.
Substitution into the full off-diagonal spatial equation
\eqref{q444:cx:eq:spatial-off}, followed by division by the nonzero number
$iv_1\bar v_2$, gives
\[
 3e(1-\rho_1\rho_2)-2\omega(\rho_1+\rho_2)=0.
\]
As $\rho_1\rho_2>0$, this proves the contradictory strict upper bound
\[
 \rho_1+\rho_2=
       \frac{3e}{2\omega}(1-\rho_1\rho_2)<\frac{3e}{2\omega}.
\]
This excludes case (II) with arbitrary complex phases of $v,w$.

Finally a constant $SO(2)$ rotation of the ordered pair $(Z_1,Z_2)$
preserves their sum of squares, every quadratic commutator sum, and
$[Z_1,Z_2]$. It commutes with the spatial rotation generator
$\left(\begin{smallmatrix}0&-1\\1&0\end{smallmatrix}\right)$,
so the two spatial equations transform by the same invertible rotation.
It leaves $C,H,q$ unchanged. Thus a solution that could be put in the
excluded class by such a rotation would produce a solution already
excluded, proving the final assertion.
\end{proof}

\subsection{The precise unresolved part}

Theorem~\ref{q444:cx:thm:block} is a complex restriction, not the real-algebra
obstruction restated. In case (II) one has
$q^*Z_2q=ie/8\ne0$. A real-skew matrix paired with the real vector $q$
has $q^*Zq=0$, and any unitary transformation fixing $q$ preserves this
pairing. Thus these nonzero block generators cannot be turned into
real-skew generators by a matter-stabilizer transformation.

Nevertheless the theorem does not put arbitrary $C,Z_1,Z_2$ into its
block class. In particular it assumes $C\in\mathfrak p$ and the stated
decomposition of the spatial pair; neither follows here from all five
unrestricted equations. The general finite system remains unresolved.
The universal identities in Proposition~\ref{q444:cx:prop:universal} continue to
apply outside this class. There is no asserted full-complex lift and no
asserted full-complex no-lift theorem in this section.

The companion \path{check_complex_rotation.py} verifies the constant-frame
matrix reductions with arbitrary complex block coordinates, the two branch
inertia formulas, and the ratio identity by exact arithmetic. It also checks
the universal pairing formulas on one exact full-algebra test triple.
The universal proofs, branch inequalities, symmetry arguments, and original
time-dependent gauge-map equivalence are proved in the text. The checker
does not solve the unrestricted finite system.

\subsection{Exact unrestricted Gauss and frequency elimination}
\label{q444:complex-elimination-section}

The word unrestricted in this subsection refers to all complex colour
coordinates in the specified constant-$C$, two-spatial-rotation ansatz,
not to all time-dependent homogeneous connections. Retain the literal
$H,q,F$, the positive coupling $g^2$, and all five equations
\eqref{q444:cx:eq:phi}--\eqref{q444:cx:eq:sp2}. No block-class restriction
from Theorem~\ref{q444:cx:thm:block} is imposed here.
For arbitrary $Z_1,Z_2\in\mathfrak{su}(3)$ put
\[
 \mathscr B=[Z_1,Z_2],\qquad L=g^2\mathcal I+2\mathcal K.
\]

\begin{lemma}[Positive Gauss inverse on every possible solution]
\label{q444:complex-gauss-inverse}
If $\mathscr B\ne0$, then $L$ is a positive-definite self-adjoint
operator on the original real eight-dimensional metric space
$(\mathfrak{su}(3),-\operatorname{tr})$. Every solution of the five
equations has $\mathscr B\ne0$. Thus its Gauss equation is exactly
\[
 C=\omega T,\qquad T=4L^{-1}\mathscr B.
\]
\end{lemma}
\begin{proof}
Cyclicity gives the full bilinear identities
\begin{align*}
 \langle X,\mathcal IY\rangle
 &=\operatorname{tr}([X,H][Y,H])
                    +2\operatorname{Re}((Xq)^*Yq),\\
 \langle X,\mathcal KY\rangle
 &=\sum_{j=1}^2\langle[Z_j,X],[Z_j,Y]\rangle.
\end{align*}
They are real and symmetric. Their values at $Y=X$ are sums of
nonnegative squared norms, since $[X,H]$ is Hermitian. Consequently
\[
 \ker L=\{X:[X,H]=0,\ Xq=0,\ [X,Z_1]=[X,Z_2]=0\}.
\]
For the converse direction in this equality, each displayed vanishing
condition makes the corresponding double-commutator or vector term
zero directly. For the forward direction, vanishing of their positive
weighted sum forces every one of the norms to vanish.

The eigenvalues of $H$ are the three distinct real numbers $0,r,-r$,
$r=\sqrt6/4$, and $q$ spans its zero eigenline. Conjugate by the exact
matrix $U$ already displayed in the block proof. A matrix commuting with
this diagonal $H$ has no off-diagonal entries. Anti-Hermiticity makes
its diagonal entries imaginary; annihilation of $q$ sets the first
entry to zero; trace zero makes the last two opposite. Therefore
\[
 \{X:[X,H]=0,\ Xq=0\}=\mathbb R\,iH.
\]
If a nonzero element of this line also commuted with both $Z_j$,
both $Z_j$ would commute with $H$. The same distinct-eigenvalue
argument would make them diagonal in its eigenbasis, hence mutually
commuting. This contradicts $\mathscr B\ne0$. Thus $\ker L=0$,
and positivity and finite dimension give the asserted inverse.

For any solution, Proposition~\ref{q444:cx:prop:universal} proves
$\|[Z_1,Z_2]\|^2>0$. Equivalently, if that bracket were zero,
Gauss paired with $C$ would give $c=0$, whereas the scalar contraction
gives $c-z=2$ and $z\geq0$. Finally multiplying Gauss by $g^2$
gives $LC=4\omega\mathscr B$, exactly the stated inverse relation.
No temporal equation was omitted.
\end{proof}

\begin{theorem}[Full residual zero set and a two-sign inverse]
\label{q444:complex-exact-zero-set}
On the explicit open domain
\[
 \mathcal D=\{(Z_1,Z_2,g^2):Z_j\in\mathfrak{su}(3),\ g^2>0,
       \ [Z_1,Z_2]\ne0,\ Tq\ne0\},
\]
retain the just-defined $T$, and set
\begin{equation}\label{q444:complex-frequency}
 \tau=\frac{1/4+\|Z_1q\|^2+\|Z_2q\|^2}{\|Tq\|^2}>0.
\end{equation}
Define the four full residuals, with their original codomains,
\begin{align}
 R_H&=\tau[T,[T,H]]-\sum_{j=1}^2[Z_j,[Z_j,H]]-F
             \in\operatorname{Herm}_3(\mathbb C)_0,\notag\\
 R_q&=(\tau T^2-Z_1^2-Z_2^2)q+2q\in\mathbb C^3,\notag\\
 R_1&=\tau\bigl(-Z_1-2[T,Z_2]+[T,[T,Z_1]]\bigr)
                 +\mathcal KZ_1+(g^2/2)\mathcal IZ_1
                  \in\mathfrak{su}(3),\notag\\
 R_2&=\tau\bigl(-Z_2+2[T,Z_1]+[T,[T,Z_2]]\bigr)
                 +\mathcal KZ_2+(g^2/2)\mathcal IZ_2
                  \in\mathfrak{su}(3).
 \label{q444:complex-full-residuals}
\end{align}
Let $\mathcal Z\subset\mathcal D$ be their common zero set. The map
\begin{equation}\label{q444:complex-two-sign-map}
 (Z_1,Z_2,g^2,\varepsilon)\longmapsto
 \left(C=\varepsilon\sqrt\tau\,T,\ Z_1,Z_2,
       \omega=\varepsilon\sqrt\tau,\ g^2\right),
 \qquad\varepsilon\in\{+1,-1\},
\end{equation}
is a bijection from $\mathcal Z\times\{+1,-1\}$ onto all solutions
of the five retained equations. Its inverse forgets $C,\omega$ and
retains $\varepsilon=\operatorname{sign}(\omega)$.
\end{theorem}
\begin{proof}
For a solution, the preceding lemma gives $C=\omega T$ and a
nonzero spatial bracket. Pair its full vector equation with $q^*$.
Anti-Hermiticity of the three matrices gives exactly
\[
 \omega^2\|Tq\|^2
       =2q^*q+\|Z_1q\|^2+\|Z_2q\|^2
       =\tfrac14+\|Z_1q\|^2+\|Z_2q\|^2>0.
\]
Thus $Tq\ne0$ and $\omega^2=\tau$. The constant $1/4$ is the
original $2q^*q$; no scalar or frequency has been scaled to one.
This contraction determines the frequency but does not replace
the vector equation. Substitution of $C=\omega T$ into all the
scalar and spatial equations gives every entry of
\eqref{q444:complex-full-residuals}. Hence the original solution
determines a point of $\mathcal Z$ and its unique sign.

Conversely, start at any point of $\mathcal Z$ and choose either
sign in \eqref{q444:complex-two-sign-map}. By definition
$LT=4\mathscr B$, so $LC=4\omega\mathscr B$ and Gauss holds.
The substitutions $C=\omega T$, $\omega^2=\tau$ make each of the
other four original residuals identical to its displayed $R$.
They therefore vanish as full matrices or vectors, not only after
contraction. The constructions plainly recover all original
coordinates and signs, proving the bijection.
\end{proof}

This is an exact seventeen-real-variable formulation (sixteen spatial
matrix coordinates and the retained positive $g^2$), not a proof that
its zero set is nonempty or empty. The inverse returns all original
fields through \eqref{q444:cx:eq:original-inverse}. In particular the
nonzero denominator and bracket domain must be retained in any
polynomial elimination. The scalar contraction alone is insufficient:
the accompanying exact test point has that contraction zero while
its full $R_q$ is nonzero. No such point is reported as a solution.

% NEW insertion-ready module. Insert after the already-integrated canonical
% q444:complex-exact-zero-set theorem. Do not insert standalone_context.tex.
\subsection{Complete scalar-coordinate identities and an unrestricted signature constraint}
\label{q444:complex-scalar-coordinate-section}

The following calculation retains all eight real coordinates of each
of $C,Z_1,Z_2$. It is not a cross--cross--block restriction.
Let $x=1/2$, $d_0=1/(2\sqrt2)$, $r=\sqrt6/4$, and
\[
S=\frac1{\sqrt{2r(r+x)}}
\begin{pmatrix}d_0&-(r+x)\\r+x&d_0\end{pmatrix},\qquad
U=\operatorname{diag}(1,S),\qquad\delta=\frac1{2r}=\sqrt{2/3}.
\]
Both columns of the numerator of $S$ have squared length $2r(r+x)$,
their inner product is zero, and its determinant is $2r(r+x)>0$.
Multiplication by the lower block of the literal $H$ shows that the
columns are eigenvectors for $r,-r$ in that order. Thus
\[
U\in SO(3),\qquad U^*HU=H_d=\operatorname{diag}(0,r,-r),\qquad U^*q=q.
\]
Conjugate every connection and force by $U^*$, with inverse conjugation
by $U$. Products, commutators, the adjoint operation, trace projection,
and the five equations all transform by these exact inverse maps.
In this frame $F_d=\operatorname{diag}(-2,1-2r,1+2r)$.

For each original connection $X$ define its eight real coordinates by
\begin{equation}
\widetilde X=U^*XU=
\begin{pmatrix}
ia_X&-\overline{v_{X1}}&-\overline{v_{X2}}\\
v_{X1}&ib_X&z_X\\
v_{X2}&-\overline{z_X}&-i(a_X+b_X)
\end{pmatrix},\qquad
a_X,b_X\in\mathbb R,\quad v_{X1},v_{X2},z_X\in\mathbb C,
\label{q444:complex-eight-coordinates}
\end{equation}
and put $h_X=a_X+2b_X$. This representation and $X=U\widetilde XU^*$
are mutually inverse. In particular no entry or phase has been removed.
The stabilizer coefficient has the original-coordinate formula
\begin{equation}
h_X=\frac{\langle iH,X\rangle}{r}
=\frac{-i\operatorname{tr}(HX)}r.
\label{q444:complex-stabilizer-coefficient}
\end{equation}
Indeed $\operatorname{tr}(H_d\widetilde X)=ir(a_X+2b_X)$.

For a scalar expression $f(X)$ define its signed sum
\[
\mathscr S(f)=f(C)-f(Z_1)-f(Z_2).
\]
This notation retains the temporal plus sign and both spatial minus
signs; it is not a positive covariance. Define
\[
D_a=\mathscr S(a^2),\qquad D_z=\mathscr S(|z|^2),\qquad
D_j=\mathscr S(|v_j|^2)\quad(j=1,2),
\]
\[
m_j=\mathscr S(hv_j),\quad m_z=\mathscr S(hz),\qquad
n_1=\mathscr S(zv_2)+\frac i2\mathscr S(av_1),\quad
n_2=\mathscr S(\bar z v_1)-\frac i2\mathscr S(av_2).
\]

\begin{proposition}[All scalar equations in unrestricted coordinates]
\label{q444:complex-full-scalar-equivalence}
For arbitrary $C,Z_1,Z_2\in\mathfrak{su}(3)$, the two scalar equations
\eqref{q444:cx:eq:phi}--\eqref{q444:cx:eq:b}
hold if and only if the following three real and five complex identities
hold:
\begin{equation}
D_a=4D_z,\qquad
D_1+2D_z=1-\delta,\qquad D_2+2D_z=1+\delta,
\label{q444:complex-scalar-diagonal-identities}
\end{equation}
\begin{equation}
m_1=m_2=m_z=n_1=n_2=0.
\label{q444:complex-scalar-cross-identities}
\end{equation}
\end{proposition}
\begin{proof}
For a diagonal Hermitian matrix with entries $\lambda_j$,
\[
[X,[X,H_d]]_{jk}
=\sum_\ell X_{j\ell}X_{\ell k}
(\lambda_j+\lambda_k-2\lambda_\ell).
\]
Use this identity with $(\lambda_1,\lambda_2,\lambda_3)=(0,r,-r)$
and the full matrix \eqref{q444:complex-eight-coordinates}.
Let $Q_H$ and $Q_q$ denote respectively the left sides minus the right
sides of the two scalar equations in this diagonal frame. Their
diagonal and matter-line entries are
\begin{align*}
(Q_H)_{11}&=2r(D_1-D_2)+2,\\
(Q_H)_{22}&=-2rD_1-4rD_z-(1-2r),\\
(Q_H)_{33}&=2rD_2+4rD_z-(1+2r),\\
(Q_q)_1&=q_1(2-D_a-D_1-D_2),\qquad q_1=1/(2\sqrt2).
\end{align*}
The first of these three matrix entries is minus the sum of the next
two, since the full residual is traceless. The latter two vanish
exactly when $D_1+2D_z=1-\delta$ and $D_2+2D_z=1+\delta$.
Their sum gives $D_1+D_2=2-4D_z$; the vector's first entry vanishes
exactly when $D_a=4D_z$.

For clarity the individual lower matter-column products are
\[
(\widetilde X^2)_{21}=i(a_X+b_X)v_{X1}+z_Xv_{X2},\qquad
(\widetilde X^2)_{31}=-ib_Xv_{X2}-\bar z_Xv_{X1},
\]
and the three independent lower/upper matrix products are
\begin{align*}
[\widetilde X,[\widetilde X,H_d]]_{21}
&=r\bigl(i(a_X-b_X)v_{X1}+3z_Xv_{X2}\bigr),\\
[\widetilde X,[\widetilde X,H_d]]_{31}
&=r\bigl(-i(2a_X+b_X)v_{X2}+3\bar z_Xv_{X1}\bigr),\\
[\widetilde X,[\widetilde X,H_d]]_{23}
&=-2ir h_Xz_X.
\end{align*}
Since $b_X=(h_X-a_X)/2$, the complete remaining residual entries are
\begin{align*}
(Q_q)_2&=q_1(n_1+im_1/2),&
(Q_H)_{21}&=r(3n_1-im_1/2),\\
(Q_q)_3&=q_1(-n_2-im_2/2),&
(Q_H)_{31}&=r(3n_2-im_2/2),\\
&& (Q_H)_{23}&=-2ir m_z.
\end{align*}
The remaining off-diagonal entries of $Q_H$ are the conjugates, since
it is Hermitian. Both $q_1$ and $r$ are nonzero. For either $j$, the
two relevant equations first give $n_j=-im_j/2$ and then
$3n_j-im_j/2=-2im_j=0$. Thus $m_j=n_j=0$; the last matrix entry gives
$m_z=0$. Conversely these five complex equalities make every displayed
off-diagonal and vector entry vanish. Together with the three real
relations already proved, they set the entire scalar residuals to
zero. This proves both directions without a rank or phase restriction.
\end{proof}

\begin{theorem}[Quantitative spacelike stabilizer constraint]
\label{q444:complex-stabilizer-spacelike}
Every scalar-compatible unrestricted triple satisfies the following
inequalities. Put
\[
H_s=h_{Z_1}^2+h_{Z_2}^2,\qquad
S_j=\sum_{\ell=1}^2\bigl(|v_{Z_\ell j}|^2+2|z_{Z_\ell}|^2\bigr),
\qquad c_1=1-\delta>0,\quad c_2=1+\delta>0.
\]
Then
\begin{equation}
h_C^2(S_j+c_j)\le H_s S_j\quad(j=1,2).
\label{q444:complex-stabilizer-margin}
\end{equation}
In particular exactly one of the following possibilities holds:
\[
h_C=h_{Z_1}=h_{Z_2}=0,\qquad\text{or}\qquad
h_C^2<h_{Z_1}^2+h_{Z_2}^2.
\]
Equivalently, the three original-coordinate numbers
$\langle iH,C\rangle,\langle iH,Z_1\rangle,\langle iH,Z_2\rangle$
are all zero or have strictly negative squared norm for the retained
signature $(+--)$.
Every zero of the unrestricted reduced residuals consequently satisfies
\begin{equation}
\tau h_T^2\le H_s
\min\left\{\frac{S_1}{S_1+1-\delta},
\frac{S_2}{S_2+1+\delta}\right\}.
\label{q444:complex-reduced-stabilizer-filter}
\end{equation}
All denominators in this formula are strictly positive.
\end{theorem}
\begin{proof}
For each $j=1,2$ and each $X\in\{C,Z_1,Z_2\}$ put
$V_{Xj}=(v_{Xj},\sqrt2 z_X)\in\mathbb C^2$.
The identities $m_j=m_z=0$ give
\[
h_C V_{Cj}=h_{Z_1}V_{Z_1j}+h_{Z_2}V_{Z_2j}.
\]
Their squared norms and Cauchy--Schwarz on the ordered pair in
$\mathbb C^2\oplus\mathbb C^2$ imply
\[
h_C^2\|V_{Cj}\|^2
\le(h_{Z_1}^2+h_{Z_2}^2)
(\|V_{Z_1j}\|^2+\|V_{Z_2j}\|^2)=H_sS_j.
\]
The diagonal scalar identities give
$\|V_{Cj}\|^2-S_j=D_j+2D_z=c_j$, proving
\eqref{q444:complex-stabilizer-margin}.
Since $0<\delta<1$, each $c_j$ is strictly positive. If $H_s=0$,
the inequality forces $h_C=0$. If $H_s>0$, division by $S_j+c_j>0$
gives $h_C^2\le H_sS_j/(S_j+c_j)<H_s$.
This covers $S_j=0$ as well. The original-coordinate version follows
by multiplication by $r^2>0$ in
\eqref{q444:complex-stabilizer-coefficient}.
Finally $C=\varepsilon\sqrt\tau\,T$ gives $h_C^2=\tau h_T^2$,
so both inequalities give the last formula, independently of the
frequency sign.
\end{proof}

The quantitative inequalities are attained by scalar-compatible
triples; equality in them does not certify the gauge equations.
For an exact witness in the original coordinates retain $A,H,q$ and set
\[
C=\frac53A+\frac43iH,\qquad Z_1=0,\qquad
Z_2=\frac43A+\frac53iH.
\]
The equality $(5/3)^2-(4/3)^2=1$ makes the signed quadratic expressions
equal to those of $(A,0,iH)$: both cross terms cancel in
$C^2-Z_2^2$ and in $[C,[C,H]]-[Z_2,[Z_2,H]]$.
Since $Hq=0$ and $[iH,H]=0$, these are exactly the original two scalar
equations. The $A$ matrix has only matter-cross entries in the diagonal
frame. Their squared moduli are $1-\delta,1+\delta$, either by the
displayed $U$ or by applying the scalar diagonal identities to
$(A,0,0)$. Also $h_A=0$ and $h_{iH}=2r$ by
\eqref{q444:complex-stabilizer-coefficient}. Thus
\[
h_C=8r/3,\qquad h_{Z_2}=10r/3,\qquad
S_j=\frac{16}{9}c_j,
\]
and both sides of \eqref{q444:complex-stabilizer-margin} are equal.
This witness is not a full gauge solution: $[Z_1,Z_2]=0$ contradicts
the necessary noncommutation proved in
Lemma~\ref{q444:complex-gauss-inverse}.

There is a second unrestricted test using both spatial equations, not
only their trace pairing. We give all its coefficients and its exact
nonnegative discriminant below.

\subsection{A two-by-two stress obstruction before evaluating the full residuals}
\label{q444:complex-stress-section}

In the coordinates \eqref{q444:complex-eight-coordinates}, direct
multiplication gives the full inertia matrix
\begin{equation}
U^*(\mathcal IX)U=
\begin{pmatrix}
ia_X/6&-\overline{v_{X1}}/2&-\overline{v_{X2}}/2\\
v_{X1}/2&-ia_X/12&3z_X/2\\
v_{X2}/2&-3\bar z_X/2&-ia_X/12
\end{pmatrix}.
\label{q444:complex-full-inertia-coordinates}
\end{equation}
For off-diagonal entries the double commutator multiplies by the square
of the eigenvalue difference; these are $r^2=3/8$ and $4r^2=3/2$.
The $qq^*$ terms add $1/8$ to the two matter-cross entries.
For the diagonal entries those terms are $i\operatorname{diag}(a_X/4,0,0)$,
whose trace-free part is $i\operatorname{diag}(a_X/6,-a_X/12,-a_X/12)$.
This proves every entry and also
\[
\langle X,\mathcal IX\rangle
=a_X^2/4+|v_{X1}|^2+|v_{X2}|^2+3|z_X|^2.
\]
The eigenvalues of $\mathcal I$ are $0,1/6,1/2,3/2$ with respective
real multiplicities $1,1,4,2$. Indeed the $h$ diagonal direction, the
$a$ diagonal direction with $b=-a/2$, the four cross-entry directions,
and the two lower off-diagonal directions are mutually orthogonal and
are multiplied by precisely these scalars in the displayed formula.

For an arbitrary point of $\mathcal D$ from
Theorem~\ref{q444:complex-exact-zero-set}, define the real
symmetric two-by-two matrices and real scalars
\[
G_{jk}=\langle Z_j,Z_k\rangle,\qquad
D_{jk}=\langle[T,Z_j],[T,Z_k]\rangle,\qquad
M_{jk}=\langle Z_j,\mathcal IZ_k\rangle,
\]
\[
d=\|[Z_1,Z_2]\|^2>0,\qquad
b_T=\langle T,[Z_1,Z_2]\rangle,
\qquad N=G+D-2b_TI_2,\qquad B=dI_2+(g^2/2)M.
\]
All these objects are functions of the unrestricted spatial pair and
$g^2$, with no frequency sign selected. $M$ is positive semidefinite:
for $x\in\mathbb R^2$, $x^TMx=\langle\sum_jx_jZ_j,
\mathcal I\sum_jx_jZ_j\rangle\ge0$. Hence $B$ is positive definite.

\begin{proposition}[Exact stress matrix and nonnegative discriminant]
\label{q444:complex-stress-obstruction}
Every zero of the four full reduced residuals satisfies
\begin{equation}
\tau N=B.
\label{q444:complex-stress-matrix-equation}
\end{equation}
In particular $N$ must be positive definite. On the open domain
$N_{11}>0$, $\det N>0$, define
\begin{equation}
\Delta=(B_{11}N_{22}+B_{22}N_{11}-2B_{12}N_{12})^2
-4\det B\det N.
\label{q444:complex-stress-discriminant}
\end{equation}
Then $\Delta\ge0$, and the stress equation is equivalent to the two
scalar conditions
\begin{equation}
\Delta=0,\qquad \tau=B_{11}/N_{11}.
\label{q444:complex-stress-filter}
\end{equation}
These are necessary tests on the unrestricted reduced zero set, not
replacements for $R_H,R_q,R_1,R_2$.
\end{proposition}
\begin{proof}
For arbitrary $C,Z_1,Z_2,\omega,g^2$, let $Q_1,Q_2$ be the left sides
of the two original spatial equations. Write
$\beta=\langle C,[Z_1,Z_2]\rangle$ and
$D^C_{jk}=\langle[C,Z_j],[C,Z_k]\rangle$. Cyclicity gives the exact
matrix of pairings
\begin{equation}
\langle Z_j,Q_k\rangle
=-\omega^2G_{jk}-D^C_{jk}+(g^2/2)M_{jk}
+(d+2\omega\beta)\delta_{jk}.
\label{q444:complex-spatial-pairing-matrix}
\end{equation}
For the double $C$ commutator the term is $-D^C_{jk}$ by invariance
of the trace inner product. The spatial double commutators have
diagonal pairing $d$ and off-diagonal pairing zero: expand
$\langle Z_j,\mathcal KZ_k\rangle
=\sum_\ell\langle[Z_\ell,Z_j],[Z_\ell,Z_k]\rangle$.
The gyroscopic diagonal pairing in either equation is $2\omega\beta$;
the off-diagonal pairings vanish because $\langle X,[C,X]\rangle=0$.
These facts prove every entry of \eqref{q444:complex-spatial-pairing-matrix},
including its symmetry.
On a reduced zero, $C=\omega T$, $\omega^2=\tau$,
$D^C=\tau D$, and $\omega\beta=\tau b_T$. Setting the pairing matrix
to zero gives $\tau N=B$. Since $B$ is positive definite and $\tau>0$,
$N$ must be positive definite.

Put $n=N_{11}>0$, $b=N_{12}$, $c=N_{22}$, and $h=nc-b^2>0$.
Expansion of the following sum of squares is an exact polynomial identity:
\begin{align}
n^2\Delta={}&
\bigl(B_{11}(nc-2b^2)+2bnB_{12}-n^2B_{22}\bigr)^2\notag\\
&\quad+4h(nB_{12}-bB_{11})^2.
\label{q444:complex-stress-squares}
\end{align}
Indeed the coefficients of $B_{11}^2,B_{22}^2,B_{12}^2,
B_{11}B_{22},B_{11}B_{12},B_{22}B_{12}$ on either side are respectively
$n^2c^2,n^4,4n^3c,-2n^3c+4n^2b^2,-4n^2bc,-4n^3b$.
Thus $\Delta\ge0$. Equality holds exactly when
$nB_{12}=bB_{11}$ and
$B_{11}(nc-2b^2)+2bnB_{12}-n^2B_{22}=0$.
Substitution of the first equality in the second gives
$n(cB_{11}-nB_{22})=0$. Since $n>0$, the two equalities are exactly
$B=(B_{11}/n)N$. Conversely this proportionality makes the sum of
squares zero. Finally $B_{11}>0$, so this proportionality factor is
positive, and $\tau N=B$ holds exactly when $\tau=B_{11}/n$.
This proves the claimed equivalence and every denominator condition.
\end{proof}

The scalar-coordinate equivalence, the quantitative stabilizer
inequality, and the stress discriminant apply to the full complex
finite system. They neither force a cross--cross--block ansatz nor
establish a zero or emptiness of the unrestricted reduced zero set.
Their exact inverse coordinate maps and the reconstruction
\eqref{q444:complex-two-sign-map} preserve all original fields,
couplings, signs, and both frequency branches.

% Insertion-ready after q444:complex-stress-obstruction.
% This module proves a new bounded finite-fibre elimination theorem.
% It re-proves, but does not claim novelty for, the signed norm identity
% already present in coordinate_constraints/orthogonality_normal_form.tex.
\subsection{An exact bounded finite coupling fibre for every unrestricted spatial pair}
\label{q444:finite-coupling-section}

Retain the literal matrices $H,q,F$, all eight real coordinates of each
$Z_j\in\mathfrak{su}(3)$, and the original positive coupling $g^2$.
The inner product remains $\langle X,Y\rangle=-\operatorname{tr}(XY)$.
The letter $t$ below is a polynomial indeterminate whose evaluated value is
the original $g^2$; no physical constant or amplitude is set to one.
This subsection constructs the complete fibre over a fixed ordered pair
$(Z_1,Z_2)$ of the four-residual zero set
\eqref{q444:complex-full-residuals}. The ordered pair is not constrained
to a real subalgebra, a block class, or a prescribed phase relation.

\begin{proposition}[An explicit upper bound for the original coupling]
\label{q444:finite-coupling-bound}
For an arbitrary fixed ordered pair with $[Z_1,Z_2]\ne0$, define
\[
 z_0=\|Z_1\|^2+\|Z_2\|^2,\qquad
 z=\sum_{j=1}^2\langle Z_j,\mathcal IZ_j\rangle,\qquad
 d=\|[Z_1,Z_2]\|^2,
\]
\[
 c_* =2+z,\qquad k_*=2z_0(z_0+4),\qquad
 t_*=
 \frac{k_*z+\sqrt{k_*^2z^2+8k_*d c_*^2}}{c_*^2}.
\]
All these numbers are real, with $z\ge0$ and
$z_0,d,c_*,k_*,t_*>0$.
Every full solution over this ordered pair satisfies
\begin{equation}
 \|C\|^2\le z_0+4,\qquad
 0<g^2\le t_*,\qquad
 0<\omega^2\le
 \frac{k_*+2d+t_*(c_*+z/2)}{z_0}.
 \label{q444:finite-coupling-bounds}
\end{equation}
The same bounds apply to both signs of the original frequency.
\end{proposition}
\begin{proof}
We include the signed norm calculation so that the amplitude estimate has
no unproved input from a proposed coordinate restriction. In the full
coordinates \eqref{q444:complex-eight-coordinates}, put $h_X=a_X+2b_X$.
Multiplication of the diagonal and off-diagonal entries gives
\[
 \|X\|^2=\tfrac32a_X^2+\tfrac12h_X^2
       +2(|v_{X1}|^2+|v_{X2}|^2+|z_X|^2).
\]
The scalar identities
\eqref{q444:complex-scalar-diagonal-identities} give
$D_a=4D_z$ and $D_1+D_2=2-4D_z$. Consequently
\begin{align*}
 \|C\|^2-z_0
 &=\tfrac32D_a+\tfrac12(h_C^2-h_{Z_1}^2-h_{Z_2}^2)
                    +2(D_1+D_2+D_z)\\
 &=4+\tfrac12(h_C^2-h_{Z_1}^2-h_{Z_2}^2).
\end{align*}
For completeness, the sign of the last correction follows directly from
the full scalar cross identities. Write
$V_X=(v_{X1},\sqrt2 z_X)$,
$S_1=\|V_{Z_1}\|^2+\|V_{Z_2}\|^2$, and
$H_s=h_{Z_1}^2+h_{Z_2}^2$. Those identities give
$h_C V_C=h_{Z_1}V_{Z_1}+h_{Z_2}V_{Z_2}$, while the diagonal identities
give $\|V_C\|^2=S_1+1-\sqrt{2/3}$.
Cauchy--Schwarz therefore proves
\[
 h_C^2(S_1+1-\sqrt{2/3})\le H_s S_1.
\]
Since $1-\sqrt{2/3}>0$, this implies $h_C^2\le H_s$, also when
$S_1=0$ or $H_s=0$. This proves $\|C\|^2\le z_0+4$.

For arbitrary anti-Hermitian matrices $X,Y$, diagonalize $X$ unitarily,
writing its eigenvalues as $i\lambda_a$ with real $\lambda_a$.
In that same basis,
\[
 \|[X,Y]\|^2
 =\sum_{a,b}(\lambda_a-\lambda_b)^2|Y_{ab}|^2
 \le2\Bigl(\sum_a\lambda_a^2\Bigr)
               \Bigl(\sum_{a,b}|Y_{ab}|^2\Bigr)
 =2\|X\|^2\|Y\|^2.
\]
Here $(\lambda_a-\lambda_b)^2\le2(\lambda_a^2+\lambda_b^2)
\le2\sum_a\lambda_a^2$ proves the inequality term by term.
Thus the retained $k=\sum_j\|[C,Z_j]\|^2$ satisfies $k\le k_*$.
The universal identities \eqref{q444:cx:eq:universal} give
$c=c_*$ and
\[
 (g^2)^2c_*^2\le8kd+2g^2kz
                    \le8k_*d+2g^2k_*z.
\]
The quadratic on moving the right side to the left has positive leading
coefficient and roots
$\bigl(k_*z\pm\sqrt{k_*^2z^2+8k_*d c_*^2}\bigr)/c_*^2$.
Its negative root is strictly negative, so its nonpositive region at
$g^2>0$ is exactly $g^2\le t_*$. Finally the universal spatial trace
identity is
\[
 \omega^2z_0=k+2d+g^2(c_*+z/2).
\]
Substitution of the two upper bounds proves the last assertion.
Nonzero bracket implies $z_0,d>0$; positivity of $\mathcal I$ and the
definitions prove all remaining strict and weak sign assertions.
\end{proof}

We now eliminate the temporal stabilizer component exactly before
clearing any frequency denominator. This preserves the original
stabilizer vector $J=iH$, whose squared norm is $3/4$.
Set
\[
 W=\{X\in\mathfrak{su}(3):\langle J,X\rangle=0\},\qquad
 \Pi X=X-\frac{\langle J,X\rangle}{3/4}J,
 \qquad \mathscr B=[Z_1,Z_2].
\]
For the fixed spatial pair define
\[
 \eta=\langle J,\mathcal KJ\rangle,\qquad
 w=\Pi\mathcal KJ,\qquad \beta=\langle J,\mathscr B\rangle,
\]
and operators and a vector on the original metric subspace $W$ by
\[
 \mathcal I_7=\mathcal I|_W,\qquad
 \mathcal S_7=\Pi\mathcal K|_W-\frac{w\otimes w}{\eta},\qquad
 f_7=\Pi\mathscr B-\frac{\beta}{\eta}w.
\]
Here $(w\otimes w)x=w\langle w,x\rangle$; no choice of basis
alters any of these maps.

\begin{lemma}[Full stabilizer elimination and polynomial Gauss numerator]
\label{q444:finite-coupling-schur}
If $\mathscr B\ne0$, then $\eta>0$, $\mathcal S_7\succeq0$, and
$\mathcal I_7\succ0$. Define the polynomial endomorphism and its
determinant and adjugate by
\[
 \mathscr M(t)=t\mathcal I_7+2\mathcal S_7,\qquad
 D_7(t)=\det_W\mathscr M(t),\qquad
 V_7(t)=4\operatorname{adj}_W(\mathscr M(t))f_7,
\]
\begin{equation}
 N_7(t)=V_7(t)+
       \frac{2\beta D_7(t)-\langle w,V_7(t)\rangle}{\eta}J.
 \label{q444:finite-coupling-N}
\end{equation}
Then $D_7(t)>0$ for $t>0$, $D_7$ has degree exactly seven and leading
coefficient $3/128$, $\deg V_7\le6$, and $\deg N_7\le7$.
For every original $g^2>0$,
\begin{equation}
 T=4(g^2\mathcal I+2\mathcal K)^{-1}\mathscr B
   =\frac{N_7(g^2)}{D_7(g^2)},\qquad
 Tq=\frac{V_7(g^2)q}{D_7(g^2)}.
 \label{q444:finite-coupling-T}
\end{equation}
In particular all temporal coordinates, including the entire stabilizer
coefficient, are reconstructed by the displayed formula.
\end{lemma}
\begin{proof}
The positive form of $\mathcal K$ gives
$\eta=\sum_j\|[Z_j,J]\|^2\ge0$. If it were zero, both $Z_j$ would
commute with $H$. The three distinct eigenvalues $0,\sqrt6/4,-\sqrt6/4$
of $H$ would make both matrices diagonal in its eigenbasis; hence
$[Z_1,Z_2]=0$, a contradiction. Thus $\eta>0$.
The kernel calculation in Lemma~\ref{q444:complex-gauss-inverse} gives
$\ker\mathcal I=\mathbb RJ$. Self-adjointness makes $W$ invariant, and
positivity then makes $\mathcal I_7$ positive definite.

For $x\in W$ and $s\in\mathbb R$, use
$\langle J,\mathcal Kx\rangle=\langle w,x\rangle$ to complete the square:
\[
 \langle x+sJ,\mathcal K(x+sJ)\rangle
 =\langle x,\mathcal S_7x\rangle
   +\eta\left(s+\frac{\langle w,x\rangle}{\eta}\right)^2.
\]
Taking $s=-\langle w,x\rangle/\eta$ and using positivity of
$\mathcal K$ proves $\mathcal S_7\succeq0$.
Thus $\mathscr M(t)$ is positive definite for every $t>0$.

Write the full Gauss unknown as $T=x+sJ$ with $x\in W$, with no
restriction on $s$. Pair
$(t\mathcal I+2\mathcal K)T=4\mathscr B$ with $J$ and then apply $\Pi$.
These operations give respectively
\[
 s=\frac{2\beta-\langle w,x\rangle}{\eta},\qquad
 (t\mathcal I_7+2\mathcal S_7)x=4f_7.
\]
Conversely these two equations recover the $J$ component and the $W$
projection of Gauss; the orthogonal decomposition
$\mathfrak{su}(3)=W\oplus\mathbb RJ$ therefore recovers the entire
equation. Inverting the positive operator on $W$ yields
$x=V_7(t)/D_7(t)$, and the formula for $s$ yields
\eqref{q444:finite-coupling-N}. Since $Jq=0$, it also yields the second
formula in \eqref{q444:finite-coupling-T}.

Determinants and adjugates may be calculated in any real basis of $W$;
similarity preserves their resulting scalar and endomorphism. The
inertia computation \eqref{q444:complex-full-inertia-coordinates} gives
the seven eigenvalues $1/6,1/2,1/2,1/2,1/2,3/2,3/2$ on $W$.
Thus $\det_W\mathcal I_7=(1/6)(1/2)^4(3/2)^2=3/128$.
The determinant consequently has degree seven with this leading
coefficient; its positivity for $t>0$ follows from positive definiteness.
The adjugate of a seven-by-seven matrix with affine entries has entries
of degree at most six, and \eqref{q444:finite-coupling-N} proves the
remaining degree bound.
\end{proof}

Put
\[
 \alpha=\frac14+\|Z_1q\|^2+\|Z_2q\|^2>0,\qquad
 Q_7(t)=\|V_7(t)q\|^2.
\]
The real polynomial $Q_7$ is nonnegative at every real $t$ and has degree
at most twelve. Its identically-zero case has the following finite test,
which does not require computing a polynomial determinant.

\begin{proposition}[Three exact vector tests for an empty Gauss domain]
\label{q444:finite-coupling-domain-test}
For every fixed noncommuting spatial pair, $Q_7$ is identically zero if
and only if the following three vector equalities all hold:
\begin{equation}
 f_7q=0,\qquad (\mathcal S_7f_7)q=0,\qquad
                 (\mathcal S_7^2f_7)q=0.
 \label{q444:finite-coupling-three-tests}
\end{equation}
When they hold, no positive original coupling over that spatial pair
belongs to the original four-residual domain.
\end{proposition}
\begin{proof}
Put $\mathcal N=\{X\in W:Xq=0\}$. In the complete coordinates
\eqref{q444:complex-eight-coordinates}, membership of $W$ sets $h_X=0$;
the first column condition $Xq=0$ sets $a_X=v_{X1}=v_{X2}=0$.
The only remaining entry is the arbitrary complex $z_X$. Thus
$\dim_{\mathbb R}\mathcal N=2$, and the full inertia formula gives
$\mathcal I_7|_{\mathcal N}=(3/2)\operatorname{id}_{\mathcal N}$.
Self-adjointness also makes $\mathcal N^\perp\cap W$ invariant.
Consequently $\mathcal I_7^{-1}y\in\mathcal N$ if and only if
$y\in\mathcal N$, and on this subspace the inverse is multiplication
by $2/3$.

Suppose first that $Q_7$ is identically zero. For every $t>0$ the vector
$x(t)=4(t\mathcal I_7+2\mathcal S_7)^{-1}f_7$ then belongs to
$\mathcal N$. Set $B_7=\mathcal I_7^{-1}$.
Factoring the pencil and using the exact finite inverse identity
$(I+A)^{-1}=I-A+A^2-A^3(I+A)^{-1}$ gives
\[
 x(t)=\frac4tB_7f_7-\frac8{t^2}B_7\mathcal S_7B_7f_7
       +\frac{16}{t^3}B_7\mathcal S_7B_7\mathcal S_7B_7f_7
       +O(t^{-4}).
\]
Here the remainder has bounded norm after multiplication by $t^4$:
for $A=2B_7\mathcal S_7/t$ and $t>4\|B_7\mathcal S_7\|$, the inverse
has norm at most two, since
$\|(I+A)y\|\ge(1-\|A\|)\|y\|\ge\|y\|/2$.
The finite-dimensional subspace $\mathcal N$ is closed. Multiplying the
expansion by $t$ and taking the limit shows $B_7f_7\in\mathcal N$, hence
$f_7\in\mathcal N$. Subtracting the first term, multiplying by $t^2$
and taking the limit shows $B_7\mathcal S_7B_7f_7\in\mathcal N$,
hence $\mathcal S_7f_7\in\mathcal N$. Subtracting the next term and
multiplying by $t^3$ similarly gives
$\mathcal S_7^2f_7\in\mathcal N$. In the latter two steps every occurrence
of $B_7$ on the already established vectors in $\mathcal N$ is exactly
multiplication by $2/3$. These are precisely the three displayed tests.

Conversely suppose the three tests hold. If $f_7=0$, then $x(t)=0$ for
all $t>0$ and the conclusion follows. If $f_7\ne0$ and
$\mathcal S_7f_7$ is dependent on $f_7$, their one-dimensional span is
invariant under $\mathcal S_7$. Otherwise those two vectors span all of
$\mathcal N$, and the third test makes that two-dimensional span
invariant as well. In either case let $V$ be the resulting invariant
span. It is contained in $\mathcal N$, so it is invariant under
$\mathcal I_7$ and therefore under the positive pencil
$t\mathcal I_7+2\mathcal S_7$.
Its restriction to $V$ is injective and, in finite dimension, onto $V$;
the full inverse thus carries $f_7\in V$ back into $V$. Hence
$x(t)q=0$ for all $t>0$, and the polynomial $Q_7$ vanishes on the
entire positive half-line. A polynomial with infinitely many roots is
identically zero. Formula \eqref{q444:finite-coupling-T} then proves
the asserted exclusion from the original domain for every $g^2>0$.
\end{proof}

Define four polynomial residuals with the original
matrix and vector codomains:
\begin{align}
 P_H(t)&=\alpha[N_7,[N_7,H]]
     -Q_7\left(\sum_j[Z_j,[Z_j,H]]+F\right),\notag\\
 P_q(t)&=\alpha N_7^2q+Q_7(2I_3-Z_1^2-Z_2^2)q,\notag\\
 P_1(t)&=\alpha\left(-D_7^2Z_1-2D_7[N_7,Z_2]
                       +[N_7,[N_7,Z_1]]\right)
                  +Q_7\left(\mathcal KZ_1+\frac t2\mathcal IZ_1\right),\notag\\
 P_2(t)&=\alpha\left(-D_7^2Z_2+2D_7[N_7,Z_1]
                       +[N_7,[N_7,Z_2]]\right)
                  +Q_7\left(\mathcal KZ_2+\frac t2\mathcal IZ_2\right).
 \label{q444:finite-coupling-polynomial-residuals}
\end{align}
All polynomials on the right are evaluated at the same $t$.
Their codomains are respectively
$\operatorname{Herm}_3(\mathbb C)_0$, $\mathbb C^3$,
$\mathfrak{su}(3)$, and $\mathfrak{su}(3)$ for real $t$.
Define the single real polynomial
\begin{equation}
 p_Z(t)=\operatorname{tr}(P_H(t)^2)+P_q(t)^*P_q(t)
                  -\operatorname{tr}(P_1(t)^2)-\operatorname{tr}(P_2(t)^2).
 \label{q444:finite-coupling-one-polynomial}
\end{equation}

% Independent strengthening of finite_coupling_fibres.tex.
% Uses only the definitions stated explicitly in that insertion.
% Canonical and checkpoint files were not edited by the auditor.
\begin{proposition}[The exact degree and positive leading coefficient]
\label{q444:finite-coupling-exact-degree}
Retain an arbitrary fixed ordered pair
$Z_1,Z_2\in\mathfrak{su}(3)$ with $\mathscr B=[Z_1,Z_2]\ne0$,
and all original definitions above: $J,\eta,\beta,\mathcal I_7$,
$D_7,V_7,N_7,Q_7$, $\alpha$, and $P_H,P_q,P_1,P_2,p_Z$.
Put
\[
 \delta_7=\det_W\mathcal I_7=\frac3{128},\qquad
 a_J=\frac{2\beta}{\eta},\qquad A_JX=[J,X],
\]
and define an endomorphism of the original real product metric space
$\mathfrak{su}(3)\oplus\mathfrak{su}(3)$ by
\[
 \mathcal L_{a_J}(X_1,X_2)
    =(X_1+a_J A_JX_2,\ X_2-a_J A_JX_1).
\]
Here $a_J$ is an explicitly defined scalar belonging to the fixed spatial
pair; it is not a new coupling or a rescaling of the original fields.
The coefficient of $t^{28}$ in the original polynomial $p_Z$ is
\begin{equation}
 [t^{28}]p_Z
 =\alpha^2\delta_7^4
   \bigl\|\mathcal L_{a_J}^2(Z_1,Z_2)\bigr\|^2.
 \label{q444:audit-exact-leading-coefficient}
\end{equation}
If $Q_7$ is not identically zero, this coefficient is strictly positive.
Consequently $\deg p_Z=28$ on every such spatial pair.
\end{proposition}
\begin{proof}

First write $N_7=V_7+s_7J$. The previously proved bounds are
$\deg V_7\le6$, $\deg s_7\le7$, and $\deg Q_7\le12$.
Since $[J,H]=0$ and $Jq=0$, the exact expressions are
\[
 [N_7,[N_7,H]]=[V_7,[V_7,H]]+s_7[J,[V_7,H]],\qquad
 N_7^2q=V_7^2q+s_7JV_7q.
\]
Hence $\deg P_H,\deg P_q\le13$. The full spatial formulas give
$\deg P_1,\deg P_2\le14$, including the $tQ_7\mathcal I/2$ terms.
Taking their displayed squared norms yields $\deg p_Z\le28$.
The established degree statements give
\[
 D_7(t)=\delta_7t^7+\text{terms of degree at most six},\qquad
 V_7(t)=\text{terms of degree at most six}.
\]
The defining full stabilizer numerator is
\[
 N_7(t)=V_7(t)+
    \frac{2\beta D_7(t)-\langle w,V_7(t)\rangle}{\eta}J.
\]
Its coefficient of $t^7$ is therefore exactly $a_J\delta_7J$,
also when $\beta=0$ and that coefficient is zero. Since
$\deg Q_7\le12$, its contributions
$Q_7(\mathcal KZ_j+t\mathcal IZ_j/2)$ have degree at most thirteen.
It follows directly from the full spatial numerator formulas that
\begin{align*}
 [t^{14}]P_1
 &=\alpha\delta_7^2
       \bigl(-Z_1-2a_J A_JZ_2+a_J^2A_J^2Z_1\bigr),\\
 [t^{14}]P_2
 &=\alpha\delta_7^2
       \bigl(-Z_2+2a_J A_JZ_1+a_J^2A_J^2Z_2\bigr).
\end{align*}
Direct composition, with both signs retained, gives
\[
 \mathcal L_{a_J}^2(Z_1,Z_2)
 =(Z_1+2a_J A_JZ_2-a_J^2A_J^2Z_1,
   Z_2-2a_J A_JZ_1-a_J^2A_J^2Z_2).
\]
The pair of leading spatial coefficients is thus
$-\alpha\delta_7^2\mathcal L_{a_J}^2(Z_1,Z_2)$.
The scalar numerator degrees satisfy $\deg P_H,\deg P_q\le13$,
so their squared norms contribute no term of degree twenty-eight.
The two spatial squared norms therefore give exactly
\eqref{q444:audit-exact-leading-coefficient}.

It remains to prove strict positivity under the stated domain-polynomial
hypothesis. Invariance of $\langle X,Y\rangle=-\operatorname{tr}(XY)$
gives
\[
 \langle X,A_JY\rangle=-\langle A_JX,Y\rangle.
\]
For arbitrary $X=(X_1,X_2)$ and $Y=(Y_1,Y_2)$, this implies
\begin{align*}
 \langle X,\mathcal L_{a_J}Y\rangle
 &=\langle X_1,Y_1\rangle+\langle X_2,Y_2\rangle
   +a_J\langle X_1,A_JY_2\rangle
   -a_J\langle X_2,A_JY_1\rangle\\
 &=\langle X_1+a_J A_JX_2,Y_1\rangle
   +\langle X_2-a_J A_JX_1,Y_2\rangle
 =\langle\mathcal L_{a_J}X,Y\rangle.
\end{align*}
Thus $\mathcal L_{a_J}$ is self-adjoint on this positive metric space.
If the coefficient in \eqref{q444:audit-exact-leading-coefficient}
were zero, then $\mathcal L_{a_J}^2Z=0$ for $Z=(Z_1,Z_2)$.
Consequently
\[
 \|\mathcal L_{a_J}Z\|^2
 =\langle Z,\mathcal L_{a_J}^2Z\rangle=0,
\]
so $\mathcal L_{a_J}Z=0$. If $a_J=0$, this gives $Z_1=Z_2=0$,
contradicting $\mathscr B\ne0$. Hence $a_J\ne0$, and the two component
equations are
\[
 a_J[J,Z_1]=Z_2,\qquad a_J[J,Z_2]=-Z_1.
\]
The full spatial operator now satisfies the exact identity
\begin{align*}
 a_J\mathcal KJ
 &=-a_J[Z_1,[Z_1,J]]-a_J[Z_2,[Z_2,J]]\\
 &=a_J[Z_1,[J,Z_1]]+a_J[Z_2,[J,Z_2]]\\
 &=[Z_1,Z_2]+[Z_2,-Z_1]=2\mathscr B.
\end{align*}
Apply the original orthogonal projection $\Pi$ to this equality and
use $w=\Pi\mathcal KJ$ and $a_J/2=\beta/\eta$. It yields
\[
 f_7=\Pi\mathscr B-\frac{\beta}{\eta}w
    =\frac{a_J}2w-\frac{a_J}2w=0.
\]
The exact definition $V_7(t)=4\operatorname{adj}_W(\mathscr M(t))f_7$
then makes $V_7$ identically zero, and hence makes
$Q_7(t)=\|V_7(t)q\|^2$ identically zero. This contradicts the
hypothesis. The coefficient is nonnegative by its squared-norm formula,
so it must be strictly positive. Since the global degree bound is
twenty-eight, the degree is exactly twenty-eight.
\end{proof}

\begin{theorem}[Complete finite coupling fibre and its exact inverse]
\label{q444:finite-coupling-fibre}
For every ordered pair $Z_1,Z_2\in\mathfrak{su}(3)$ the complete fibre of
the four-residual zero set is as follows.
If $[Z_1,Z_2]=0$, it is empty. Otherwise form the explicitly defined
polynomials above. If $Q_7$ is identically zero, the fibre is empty.
If $Q_7$ is not identically zero, then $p_Z$ is not identically zero,
$\deg p_Z=28$, and the complete set of admissible couplings is exactly
\begin{equation}
 \left\{t\in(0,t_*]:Q_7(t)>0,\ p_Z(t)=0\right\}.
 \label{q444:finite-coupling-root-set}
\end{equation}
This set has at most fourteen elements. Equivalently, the projection
$\mathcal Z\longrightarrow\mathfrak{su}(3)^2$ that forgets $g^2$ has
fibres of cardinality at most fourteen, including empty fibres.

For each root in \eqref{q444:finite-coupling-root-set} and each
$\varepsilon\in\{+1,-1\}$, the complete original finite-system solution is
\begin{equation}
 g^2=t,\qquad
 \omega=\varepsilon\frac{\sqrt\alpha\,D_7(t)}{\sqrt{Q_7(t)}},\qquad
 C=\varepsilon\frac{\sqrt\alpha\,N_7(t)}{\sqrt{Q_7(t)}},\qquad
 (Z_1,Z_2)=(Z_1,Z_2).
 \label{q444:finite-coupling-reconstruction}
\end{equation}
Thus every fibre of the five-equation solution set over the fixed ordered
spatial pair has at most twenty-eight elements, counting both frequency
signs. No root with $Q_7(t)=0$ is included.
\end{theorem}
\begin{proof}
A zero spatial bracket is impossible by
Lemma~\ref{q444:complex-gauss-inverse}. For a nonzero bracket,
\eqref{q444:finite-coupling-T} gives
$\|Tq\|^2=Q_7(t)/D_7(t)^2$ for every $t>0$.
Thus the exact original domain condition $Tq\ne0$ is precisely
$Q_7(t)>0$. If $Q_7$ is identically zero, the domain has no point over
this pair. On the remaining domain the unchanged frequency is
\[
 \tau=\frac{\alpha D_7(t)^2}{Q_7(t)}.
\]
Substitution in all four full residuals gives, entry by entry,
\begin{equation}
 (R_H,R_q,R_1,R_2)
   =\frac1{Q_7(t)}(P_H(t),P_q(t),P_1(t),P_2(t)).
 \label{q444:finite-coupling-exact-residual-map}
\end{equation}
Every denominator here is strictly positive on the asserted domain.
For real $t$, each summand in \eqref{q444:finite-coupling-one-polynomial}
is a squared norm in its stated real inner-product space. The Hermitian
and anti-Hermitian assertions follow directly from the commutators in
\eqref{q444:finite-coupling-polynomial-residuals}; they do not require a
residual to vanish. Therefore $p_Z(t)\ge0$ for every real $t$, and its
vanishing is equivalent to vanishing of every full polynomial residual.
This proves the exact correspondence of roots and full residual zeros,
before the bound on the number of roots is established.

To prove the degree bounds, write $N_7=V_7+s_7J$, with
$\deg V_7\le6$ and $\deg s_7\le7$. Since $[J,H]=0$ and $Jq=0$,
\[
 [N_7,[N_7,H]]=[V_7,[V_7,H]]+s_7[J,[V_7,H]],
 \qquad
 N_7^2q=V_7^2q+s_7JV_7q.
\]
Their degrees are at most thirteen. Hence $\deg P_H,\deg P_q\le13$.
In the two spatial residuals the terms $D_7^2Z_j$,
$D_7[N_7,Z_k]$, and $[N_7,[N_7,Z_j]]$ have degree at most fourteen;
the terms involving $Q_7\mathcal K$ and $tQ_7\mathcal I$ have degree
at most twelve and thirteen. Thus $\deg P_1,\deg P_2\le14$ and
$\deg p_Z\le28$.

Suppose now that $Q_7$ is not identically zero but $p_Z$ is identically
zero. A nonzero univariate polynomial has only finitely many real
roots, so there exists $t>t_*$ with $Q_7(t)\ne0$. Nonnegativity of
$Q_7$ gives $Q_7(t)>0$. The established exact residual correspondence
then makes this $t$ a full zero; the two-sign inverse of
Theorem~\ref{q444:complex-exact-zero-set} reconstructs a solution of all
five original equations. This contradicts
Proposition~\ref{q444:finite-coupling-bound}. Consequently $p_Z$ is a
nonzero polynomial. Every real root of a real polynomial that is
nonnegative on all of $\mathbb R$ has even multiplicity: at a root
$t_0$, factor it as $(t-t_0)^m u(t)$ with $u(t_0)\ne0$; continuity
makes $u$ have a constant nonzero sign near $t_0$, and an odd $m$ would
give opposite signs on the two sides. Thus a nonzero polynomial of
degree at most twenty-eight of this type has at most fourteen distinct
real roots. Imposing the explicit interval and denominator conditions
can only decrease that number.

The stronger equality $\deg p_Z=28$ follows from the independent
leading-coefficient proof in Proposition~\ref{q444:finite-coupling-exact-degree}.

Finally $D_7(t)>0$ identifies the sign of the displayed $\omega$ with
$\varepsilon$ and gives $C=\omega T$ exactly. The frequency square is
the unchanged $\tau$. Formula
\eqref{q444:finite-coupling-exact-residual-map} proves that all remaining
equations vanish, while the Schur inverse proves the full Gauss equation.
Conversely every original solution determines its fixed pair, its root
$t=g^2$, and its unique sign, and these formulas recover $C,\omega$.
This proves the stated bijection and both cardinality bounds.
\end{proof}

The original spacetime reconstruction is retained literally: with the
already specified $O(t_{\rm phys})$, $K$, and $\Omega$, put
\[
 \begin{aligned}
 Y_1(t_{\rm phys})&=\cos(\omega t_{\rm phys})Z_1
                         -\sin(\omega t_{\rm phys})Z_2,\\
 Y_2(t_{\rm phys})&=\sin(\omega t_{\rm phys})Z_1
                         +\cos(\omega t_{\rm phys})Z_2,\qquad Y_3=0.
 \end{aligned}
\]
\[
 \mathcal A_0=OKCK^*O^*-\Omega,\qquad
 \mathcal A_i=OKY_iK^*O^*.
\]
The prescribed scalar fields remain
$\Phi(t_{\rm phys})=O\Phi_0O^*$ and $b(t_{\rm phys})=Ob_0$, with their
original frequency $\sqrt2$. The polynomial indeterminate $t=g^2$ is
distinct from the physical time $t_{\rm phys}$.

This theorem solves the coupling fibre exactly as a bounded finite root
set while keeping every unrestricted spatial coordinate. It does not
establish that any ordered spatial pair has a nonempty fibre, or that
all such fibres are empty. Existence or emptiness of the complete
sixteen-coordinate spatial projection remains unresolved by this result.

% Insert after the full finite coupling fibre theorem and reconstruction.
\subsection{A proper algebraic exceptional set for the full spatial projection}
\label{q444:generic-projection-section}

We now retain all sixteen spatial coordinates simultaneously. The new
statement is about the entire spatial projection, not only a selected
block class or one value of the coupling. The five original equations and
their prescribed rotating ansatz remain unchanged.

Use the exact fixed matrix $U$ of
\eqref{q444:complex-eight-coordinates}. For each $j=1,2$, write
$\widetilde Z_j=U^*Z_jU$ using the eight real coordinates
$(a_j,h_j,\operatorname{Re}v_{j1},\operatorname{Im}v_{j1},
\operatorname{Re}v_{j2},\operatorname{Im}v_{j2},
\operatorname{Re}z_j,\operatorname{Im}z_j)$, with $b_j=(h_j-a_j)/2$.
This is an invertible linear coordinate map on the original spatial
carrier. Denote the resulting sixteen-coordinate vector by $\xi$.
No coordinate is held at zero in the theorem below.

\begin{lemma}[Polynomial diagonal residuals with retained denominators]
\label{q444:generic-polynomial-diagonals}
Let $\eta=\langle J,\mathcal KJ\rangle$ be the original nonnegative
polynomial in $\xi$. On $\eta>0$ define
\[
 F_\xi(t)=\eta^{16}(U^*P_H(t)U)_{11},\qquad
 G_\xi(t)=\eta^{16}(U^*P_H(t)U)_{22}.
\]
Both extend as polynomials in $(\xi,t)$ with real algebraic
coefficients, and their degree in $t$ is at most twelve.
On the stated domain the inverse formulas are exactly
$(U^*P_HU)_{11}=F_\xi/\eta^{16}$ and
$(U^*P_HU)_{22}=G_\xi/\eta^{16}$.
\end{lemma}
\begin{proof}
All real coordinate components of $\mathcal K,w,\eta,\beta,\Pi\mathscr B$
are polynomials in $\xi$ with fixed real algebraic coefficients; the
projection denominator $\langle J,J\rangle=3/4$ is a fixed nonzero
constant. Choose once and for all a fixed basis of $W$ with real
algebraic coordinates in the displayed diagonal-frame carrier.
The entries of $\mathscr M=t\mathcal I_7+2\mathcal S_7$ have
denominator dividing $\eta$, and those of $f_7$ have the same property.
The determinant expansion shows that $\eta^7D_7$ is polynomial.
Each adjugate entry is a six-by-six minor, so $\eta^7V_7$ is
polynomial as well. The full formula
\eqref{q444:finite-coupling-N} therefore makes $\eta^8N_7$
polynomial. Complex conjugation here conjugates fixed coefficients and
the displayed complex coordinate combinations, while every component of
$\xi$ and $t$ is real. Thus $\eta^{14}Q_7$ is polynomial.
The full formula for $P_H$ now proves polynomial extension after the
displayed multiplication by $\eta^{16}$. No value at $\eta=0$
is treated as an inverse Gauss solution.

To prove the sharper degree bound, work through the exact conjugation
by $U^*$, so $H_d=\operatorname{diag}(0,r,-r)$ and
$J_d=iH_d$ are both diagonal. Write $N_7=V_7+s_7J$ as before.
Since $[J,H]=0$,
\[
 [N_7,[N_7,H]]=[V_7,[V_7,H]]+s_7[J,[V_7,H]].
\]
For every matrix $Y$, $([J_d,Y])_{jj}=0$ follows by subtracting
$(J_d)_{jj}Y_{jj}$ from itself. Consequently the second term contributes
zero to each diagonal entry after conjugation. The first term has
degree at most twelve since $\deg_t V_7\le6$. The other term in
$P_H$ is $Q_7$ times a matrix independent of $t$, and
$\deg_t Q_7\le12$. Multiplication by $\eta^{16}$ does not
change the degree in $t$. This proves every assertion.
\end{proof}

Define a fixed-size polynomial matrix $\mathscr S(\xi)$ by the
linear map
\begin{equation}
 \mathbb R[t]_{<12}\oplus\mathbb R[t]_{<12}
 \longrightarrow\mathbb R[t]_{<24},\qquad
 (A,B)\longmapsto AF_\xi+BG_\xi,
 \qquad \mathscr R(\xi)=\det\mathscr S(\xi).
 \label{q444:generic-sylvester-map}
\end{equation}
Use the increasing-power bases $1,t,\ldots$ and put the $A$ summand
before the $B$ summand. This specifies the determinant sign and its
twenty-four columns, including when either polynomial loses degree.
This is the classical Sylvester resultant map; see the Stacks Project,
\href{https://stacks.math.columbia.edu/tag/00UA}{Tag 00UA}.
The argument below proves the properties needed here directly, without
replacing either polynomial by a monic one.

\begin{theorem}[Generic exclusion in all sixteen spatial coordinates]
\label{q444:generic-projection-obstruction}
The explicitly defined polynomial $\mathscr R$ is not identically
zero. The spatial projection of every full solution of
\eqref{q444:cx:eq:phi}--\eqref{q444:cx:eq:sp2} is contained in
\[
 \{\xi\in\mathbb R^{16}:\mathscr R(\xi)=0\}.
\]
This exceptional set is closed, has empty interior, and has
sixteen-dimensional Lebesgue measure zero. Thus the complement is an
open dense set of unrestricted spatial pairs, each admitting no positive
coupling and no frequency sign solving the five retained equations.
This theorem does not decide whether any point of the exceptional set
admits a solution.
\end{theorem}
\begin{proof}
Any full solution has $[Z_1,Z_2]\ne0$, hence $\eta>0$ by
Lemma~\ref{q444:finite-coupling-schur}. It also has $Q_7(t)>0$ at
the original $t=g^2$, and the exact residual identity
\eqref{q444:finite-coupling-exact-residual-map} makes the entire
$P_H(t)$ zero. Therefore $F_\xi(t)=G_\xi(t)=0$. Evaluation at
that real $t$ is a nonzero linear functional on
$\mathbb R[t]_{<24}$ (it sends $1$ to $1$), and it annihilates
the image of \eqref{q444:generic-sylvester-map}. That image cannot
be all of the codomain. The determinant is zero. This proof includes
every degree-drop case and does not use a converse resultant criterion
outside the original denominator domain.

Here is an exact witness that the determinant polynomial is nonzero.
Write
\[
 M(a,h,v_1,v_2,z)=
 \begin{pmatrix}
 ia&-\bar v_1&-\bar v_2\\
 v_1&i(h-a)/2&z\\
 v_2&-\bar z&-i(a+h)/2
 \end{pmatrix}
\]
and set, in the displayed diagonal coordinates,
\[
 \widetilde Z_1=M(1,2,1+i,-1+2i,2-i),\qquad
 \widetilde Z_2=M(-2,1,2+3i,1-i,-1+2i).
\]
The original pair is $Z_j=U\widetilde Z_jU^*$, not a replacement
of the prescribed $H,q,F$. Every one of its sixteen displayed real
spatial coordinates is nonzero. Substitution in the original definitions
of $D_7,V_7,N_7,Q_7,P_H$ gives two diagonal polynomials of degree
twelve in $\mathbb Q(\sqrt6)[t]$. Their coprimality can be verified
by the following small certificate, with all coefficient lists in
descending powers. Reduce their coefficients by the explicit map
$\sqrt6\mapsto39$ to $\mathbb F_{101}$, where $39^2=6$.
Each rational coefficient denominator is nonzero modulo $101$;
the exact coefficient file and replay test check them individually.
The two reduced polynomials $f,g$ and the Bezout coefficients $a,b$ are
\begin{align*}
 f &: (39,85,79,93,67,20,38,98,26,56,8,25,28),\\
 g &: (1,87,96,7,76,88,75,93,53,73,64,41,16),\\
 a &: (4,60,31,59,74,31,34,18,59,97,95,96),\\
 b &: (46,85,47,5,14,36,38,21,35,48,34,53).
\end{align*}
Multiplication gives $af+bg=1$ in $\mathbb F_{101}[t]$.
Both $f$ and $g$ have degree exactly twelve. If a pair
$(a_0,b_0)$ of degrees less than twelve were in the kernel of their
Sylvester map, coprimality would imply that $f$ divides $b_0$ and
$g$ divides $a_0$. The degree bounds force both to be zero. Thus that
finite-field map, between spaces of equal dimension twenty-four, has
nonzero determinant. Reduction commutes with the determinant polynomial.
It follows that the Sylvester determinant of the two original diagonal
polynomials is nonzero in $\mathbb Q(\sqrt6)$ as well. This uses
the localized coefficient ring with all its denominators inverted,
not a homomorphism from the whole field $\mathbb Q$ to a finite field.

For explicit verification at the witness, multiplication gives
\[
 [\widetilde Z_1,\widetilde Z_2]_{21}=9-\tfrac72 i\ne0,\qquad
 \eta=\tfrac34(7+15)+3(5+5)=\tfrac{93}2>0.
\]
Multiplying both diagonal polynomials by the same nonzero scalar
$\eta^{16}$ multiplies each of the twenty-four columns by
$\eta^{16}$, and hence multiplies their determinant by
$\eta^{384}\ne0$. Consequently $\mathscr R$ is nonzero at
this exact witness and cannot be the zero polynomial. The finite-field
calculation is solely a certificate of this real algebraic statement;
it does not change the physical carrier or any coupling.

As a separate certificate, write $j_0=\langle J,J\rangle=3/4$.
In $W\oplus\mathbb RJ$, without changing the length of $J$, the
full Gauss endomorphism has block matrix
\[
 L=\begin{pmatrix}
 t\mathcal I_7+2\Pi\mathcal K|_W&2w\\
 2\langle w,\cdot\rangle/j_0&2\eta/j_0
 \end{pmatrix}.
\]
Subtracting $j_0w/\eta$ times its last row from the upper block rows
preserves the determinant and makes the upper-right block zero; its
upper-left block becomes $\mathscr M(t)$. Thus
$\det L=(2\eta/j_0)D_7$ exactly.
The independent replay reconstructs both original diagonal polynomials
from that matrix modulo $101$. It obtains Sylvester determinant $40$;
after the retained $\eta^{16}$ multiplication it obtains
$40\eta^{384}=60\ne0$ in $\mathbb F_{101}$. These residues refer
to the exact basis order in \eqref{q444:generic-sylvester-map}.

For completeness, a nonzero real polynomial on $\mathbb R^n$ has a
measure-zero zero set by induction on $n$. For $n=1$ it has finitely
many roots. For $n>1$, write it as a polynomial in the last variable
whose coefficients are polynomials in the first $n-1$ variables.
Choose one nonzero coefficient. The set where all coefficients vanish
is contained in the measure-zero zero set of that coefficient, by
induction. At every other base point the last-variable polynomial has
only finitely many roots. Applying Fubini on each bounded box shows
that the zero set in that box has measure zero; countably many boxes
cover $\mathbb R^n$. Its complement is open because the polynomial
is continuous. Every nonempty open ball has positive measure, so no
ball is contained in the zero set, and the complement is dense.
This proves the assertions at $n=16$.
\end{proof}

The exact large-coefficient identity over $\mathbb Q(\sqrt6)$ and
the small finite-field identity are retained separately in
\path{PAIR_BEZOUT_CERTIFICATE.json} and
\path{PAIR_MODULAR_CERTIFICATE.json}. The replay
\path{check_pair_bezout.py} constructs the original full pair and every
Gauss and scalar numerator, then checks the identities; the integer-only
certificate replay checks the displayed finite-field products separately.
Neither a finite failed search nor two failed test couplings proves this
theorem: the proof uses an explicit nonzero polynomial vanishing on
every possible spatial projection, with its nonzeroness certified at
the literal witness.


% Shared canonical-ready fragment derived from two_cross_exceptional_exclusion.tex.
% Input immediately after q444:generic-projection-section and before the
% existing Sources, independent checks, and reproducibility subsection.
% No preamble, document environment, or global macro definitions.
% Stage the companion proof/checker/receipt at the relative evidence path below.
\subsection{An eight-dimensional exceptional spatial class with no full solution}
\label{q444:twocross:section}

\subsubsection{Literal system and the explicitly selected class}
\label{q444:twocross:literal-system}

All vector spaces in this note retain their real structures. In particular,
$\mathfrak{su}(3)$ is the eight-real-dimensional space of traceless anti-Hermitian
complex $3\times3$ matrices, with positive inner product
$\langle X,Y\rangle=-\operatorname{tr}(XY)$. Retain the original matrices
\[
H=\begin{pmatrix}0&0&0\\0&-1/2&1/(2\sqrt2)\\
0&1/(2\sqrt2)&1/2\end{pmatrix},\quad
q=\frac{e_1}{2\sqrt2},\quad
A=\sqrt2\begin{pmatrix}0&-1&0\\1&0&0\\0&0&0\end{pmatrix},
\]
and the original scalar constants $m^2=\Gamma=1$, $\kappa=6$.
For $Y_{\rm tf}=Y-\operatorname{tr}(Y)I_3/3$, define
\[
\mathcal IX=([H,[H,X]]+qq^*X+Xqq^*)_{\rm tf},\qquad
\mathcal KX=-[Z_1,[Z_1,X]]-[Z_2,[Z_2,X]],
\]
\[
F=[A,[A,H]]=-2H+\operatorname{diag}(-2,1,1).
\]
The last identity is verified by the two displayed matrix products; also
$H^2=(3/8)\operatorname{diag}(0,1,1)$ and $q^*q=1/8$.
The indeterminate/value $t=g^2$ is the original positive gauge coupling,
not physical time. The original frequency satisfies $\omega\in\mathbb R
\setminus\{0\}$. The five retained equations are
\begin{align}
[C,[C,H]]-[Z_1,[Z_1,H]]-[Z_2,[Z_2,H]]&=F,\label{q444:twocross:eq:H}\\
(C^2-Z_1^2-Z_2^2)q&=-2q,\label{q444:twocross:eq:q}\\
(t\mathcal I+2\mathcal K)C&=4\omega[Z_1,Z_2],\label{q444:twocross:eq:G}\\
-\omega^2Z_1-2\omega[C,Z_2]+[C,[C,Z_1]]
+\mathcal KZ_1+(t/2)\mathcal IZ_1&=0,\label{q444:twocross:eq:S1}\\
-\omega^2Z_2+2\omega[C,Z_1]+[C,[C,Z_2]]
+\mathcal KZ_2+(t/2)\mathcal IZ_2&=0.\label{q444:twocross:eq:S2}
\end{align}
Multiplying the original Gauss equation by its retained positive $t$
gives \eqref{q444:twocross:eq:G}; division by that same $t$ is its exact inverse.

Put $P=qq^*/(q^*q)=\operatorname{diag}(1,0,0)$ and define
\[
\mathfrak p=\{X\in\mathfrak{su}(3):PXP=(I_3-P)X(I_3-P)=0\},\qquad
\mathfrak k=\{X\in\mathfrak{su}(3):[X,P]=0\}.
\]
Thus $\mathfrak p$ exchanges the matter line and its orthogonal complement, while
$\mathfrak k$ preserves them. Every $X\in\mathfrak p$ has the unique form
\[
X(v)=\begin{pmatrix}0&-v^*\\v&0_{2\times2}\end{pmatrix},
\qquad v\in\mathbb C^2,
\]
whose inverse takes the lower first column. This is a real linear
isomorphism $\mathbb C^2\to\mathfrak p$, so $\mathfrak p\times\mathfrak p$ has real dimension eight.
Block decomposition gives the orthogonal direct sum $\mathfrak{su}(3)=\mathfrak k\oplus\mathfrak p$,
with both summands of real dimension four.

\begin{theorem}[Two-cross spatial exclusion]\label{q444:twocross:thm:exclusion}
There is no solution of \eqref{q444:twocross:eq:H}--\eqref{q444:twocross:eq:S2} with both
$Z_1,Z_2\in\mathfrak p$. No hypothesis on $C$ is imposed in this statement.
Every complex coordinate of each of the two cross vectors is allowed,
and both signs of $\omega\ne0$ and every original $t>0$ are retained.
\end{theorem}

This is an explicitly restricted spatial class, not a claim that arbitrary
spatial pairs can be put into this class. The proof below also shows that
the entire class lies in the previously defined real algebraic exceptional
zero set. Consequently the result excludes an eight-dimensional linear
subspace inside that zero set, but does not decide existence in its
remaining part.

\subsubsection{Gauss and the complete scalar equations force the temporal form}
\label{q444:twocross:temporal-form}

Block multiplication gives
$[\mathfrak p,\mathfrak p]\subset\mathfrak k$, $[\mathfrak k,\mathfrak p]\subset\mathfrak p$.
Because $H$ and $qq^*$ are block diagonal, $\mathcal I$ preserves
both summands. Direct multiplication using $H^2$ and $q^*q$ gives
\begin{equation}
\mathcal IX(v)=\left(\frac38+\frac18\right)X(v)=\frac12X(v).
\label{q444:twocross:eq:Ip}
\end{equation}
In this computation $HX(v)H=0$, the two $H^2$ products sum to
$(3/8)X(v)$, the two $qq^*$ products sum to $X(v)/8$, and every
term is traceless. If $Z_1,Z_2\in\mathfrak p$, their double commutators
therefore make $\mathcal K$ preserve both summands as well.
Cyclicity gives, for every $Y\in\mathfrak{su}(3)$,
\[
\langle Y,\mathcal KY\rangle=\sum_{j=1}^2\|[Z_j,Y]\|^2\ge0.
\]
Project \eqref{q444:twocross:eq:G} onto $\mathfrak p$. Its right side lies in $\mathfrak k$, hence
\[
\left(\frac t2\operatorname{id}_{\mathfrak p}+2\mathcal K|_{\mathfrak p}\right)C_{\mathfrak p}=0.
\]
Pairing with $C_{\mathfrak p}$ gives a sum whose first term is
$(t/2)\|C_{\mathfrak p}\|^2$ and whose other terms are nonnegative.
Since $t>0$, $C_{\mathfrak p}=0$. Thus Gauss forces $C\in\mathfrak k$; this was
not an additional restriction on the temporal unknown.

Here are all frame changes, together with their inverses. Put
\[
x=\frac12,\quad d_0=\frac1{2\sqrt2},\quad r=\frac{\sqrt6}{4},
\quad \delta=\frac1{2r}=\sqrt{\frac23},
\]
\[
S=\frac1{\sqrt{2r(r+x)}}
\begin{pmatrix}d_0&-(r+x)\\r+x&d_0\end{pmatrix},\qquad
U=\operatorname{diag}(1,S).
\]
The relation $r^2=x^2+d_0^2$ shows that the numerator columns have
squared length $2r(r+x)$, are orthogonal, and have determinant
$2r(r+x)>0$. Multiplication by the lower block of $H$ gives their
eigenvalues $r,-r$ in that order. Hence
\[
U\in SO(3),\qquad U^*HU=\operatorname{diag}(0,r,-r),\qquad U^*q=q.
\]
For every matrix $Y$ the change is $Y\mapsto U^*YU$, with inverse
$Y'\mapsto UY'U^*$. It is applied to $H,F,C,Z_1,Z_2$ simultaneously.
It commutes with products, brackets, adjoints and trace projection,
and fixes $q$. Thus it intertwines $\mathcal I,\mathcal K$ and each
of the five full residuals exactly. It preserves $\mathfrak p,\mathfrak k$ because
it commutes with $P$. No physical constant is rescaled.

In this first frame write the full temporal block matrix as
\[
C_d=\begin{pmatrix}
ia&0&0\\0&i(h-a)/2&z\\0&-\bar z&-i(a+h)/2
\end{pmatrix},\qquad a,h\in\mathbb R,\quad z\in\mathbb C.
\]
Write the spatial matrices as $X(v),X(w)$ and put
$M=vv^*+ww^*$ in this frame only for the following calculation.
The vector scalar residual is exactly
\[
(C_d^2-X(v)^2-X(w)^2)q+2q
=(-a^2+\operatorname{tr}M+2)q.
\]
The complete nonzero block of the adjoint scalar residual, including
its upper entry, is
\[
\begin{pmatrix}
2-2r(M_{11}-M_{22})&0&0\\
0&-4r|z|^2+2rM_{11}-1+2r&-2irhz\\
0&2irh\bar z&4r|z|^2-2rM_{22}-1-2r
\end{pmatrix}.
\]
For example, if the lower Hermitian block $B=-i(C_d)_{\rm lower}$
and $D=\operatorname{diag}(r,-r)$ are used, its lower block is
$-[B,[B,D]]+\{M,D\}-(I_2-2D)$; multiplying these two-by-two
matrices gives every entry just displayed.
The two scalar equations are therefore equivalent to
\begin{equation}
a^2=2+\operatorname{tr}M,\qquad
2|z|^2-M_{11}=1-\delta,\qquad
2|z|^2-M_{22}=1+\delta,\qquad hz=0.
\label{q444:twocross:eq:first-scalar}
\end{equation}
The upper entry is minus the sum of the two lower diagonal entries,
so no scalar entry is left unchecked. Adding the middle two equations
and comparing with the first gives $a^2=4|z|^2$. Since
$M_{22}\ge0$ and $1+\delta>0$, $z\ne0$. Consequently
\begin{equation}
h=0,\qquad a\ne0,\qquad |z|=|a|/2.
\label{q444:twocross:eq:temporal-forced}
\end{equation}

Set $e=a$, retaining either sign, and $\gamma=-2iz/e$, so
$|\gamma|=1$. Choose either square root $\zeta$ of $-\bar\gamma$ and put
\[
W=\frac{\zeta}{\sqrt2}
\begin{pmatrix}\gamma&\gamma\\1&-1\end{pmatrix},\qquad
V=\operatorname{diag}(1,W).
\]
Then $|\zeta|=1$, $W^*W=I_2$ and
$\det W=-\zeta^2\gamma=\bar\gamma\gamma=1$.
The lower Hermitian block of $C_d$ is
\[
\frac e2\begin{pmatrix}-1&\gamma\\\bar\gamma&-1\end{pmatrix}.
\]
Its products with $(\gamma,1)^{\mathsf T}$ and
$(\gamma,-1)^{\mathsf T}$ have eigenvalues $0,-e$, respectively.
Those same columns give
$W^*\operatorname{diag}(r,-r)W=r\left(\begin{smallmatrix}0&1\\1&0\end{smallmatrix}\right)$.
Thus the second exact change $Y_d\mapsto V^*Y_dV$ gives
\begin{equation}
C'=ie\operatorname{diag}(1,0,-1),\quad
H'=\begin{pmatrix}0&0&0\\0&0&r\\0&r&0\end{pmatrix},\quad q'=q,
\quad F'=-2H'+\operatorname{diag}(-2,1,1).
\label{q444:twocross:eq:second-frame}
\end{equation}
The inverse for every original matrix is, explicitly,
\begin{equation}
Y=UVY'V^*U^*.
\label{q444:twocross:eq:original-inverse}
\end{equation}
For spatial vectors the forward change is the corresponding lower
unitary action $v\mapsto W^*v$ and its inverse is $v'\mapsto Wv'$.
The unitary parameters are retained; the assertion is not that $H$ can
be changed while the other data are left fixed. As with $U$, this
simultaneous conjugation intertwines all five equations. In the remainder
of the proof primes are omitted only to ease typography; all matrices
are the exact objects in \eqref{q444:twocross:eq:second-frame} and have inverse
\eqref{q444:twocross:eq:original-inverse}.

\subsubsection{Every gauge entry in the classified frame}
\label{q444:twocross:full-gauge-entries}

Write $Z_1=X(v)$, $Z_2=X(w)$ with arbitrary $v,w\in\mathbb C^2$,
and define, afresh in the second frame,
\[
V_0=v^*v,\quad W_0=w^*w,\quad p_0=v^*w=\alpha_0+ib,
\quad M=vv^*+ww^*,\quad N=i(vw^*-wv^*).
\]
Both $M$ and $N$ are Hermitian; $M\succeq0$. Put
\[
u=M_{11},\quad v_0=M_{22},\quad m=M_{12},\quad
d_M=\det M=|v_1w_2-v_2w_1|^2\ge0.
\]
The determinant identity follows by expanding the four products in
$M_{11}M_{22}-M_{12}M_{21}$ and collecting the squared modulus.
The symbol $v_0$ here is a real diagonal entry, not the vector $v$.
All eight real components of the two vectors remain independent.

Multiplication gives the full scalar residuals in this frame:
\[
Q_H=\begin{pmatrix}
2-4r\operatorname{Re}m&0&0\\
0&2r\operatorname{Re}m-1&r(\operatorname{tr}M-e^2+2)\\
0&r(\operatorname{tr}M-e^2+2)&2r\operatorname{Re}m-1
\end{pmatrix},
\quad Q_q=(\operatorname{tr}M-e^2+2)q.
\]
Consequently scalar compatibility is exactly
\begin{equation}
e^2=2+\operatorname{tr}M,\qquad \operatorname{Re}m=\delta>0.
\label{q444:twocross:eq:second-scalar}
\end{equation}
In particular $m\ne0$. No phase of $m,v,w$ has been set to a chosen value.

Set
\[
Q=\operatorname{diag}(1,2),\qquad
E=\operatorname{diag}(2/3,-5/6).
\]
Direct multiplication, keeping $r^2=3/8$ and $q^*q=1/8$, gives
\begin{equation}
\mathcal IC=ie\operatorname{diag}(1/6,2/3,-5/6),\qquad
\mathcal IX(v)=\tfrac12X(v).
\label{q444:twocross:eq:inertia-classified}
\end{equation}
For the first formula the two contributions, respectively from $H$
and from the trace-free $q$ term, are
\[
 ie\operatorname{diag}(0,3/4,-3/4),\qquad
 ie\operatorname{diag}(1/6,-1/12,-1/12).
\]
The second formula is \eqref{q444:twocross:eq:Ip}, or follows from
$(H')^2=(3/8)\operatorname{diag}(0,1,1)$ directly.
The entire spatial bracket and the entire $\mathcal KC$ are
\[
[Z_1,Z_2]=\begin{pmatrix}-2ib&0\\0&iN\end{pmatrix},\qquad
\mathcal KC=\begin{pmatrix}2ie\operatorname{tr}(QM)&0\\
0&-ie(QM+MQ)\end{pmatrix}.
\]
For example $[C,X(v)]=X(-ieQv)$, and the two double cross products
give the displayed anticommutator. Thus the full Gauss equation is
equivalent to the upper entry and lower Hermitian block
\begin{equation}
\frac{et}{6}+4e\operatorname{tr}(QM)+8\omega b=0,\qquad
etE-2e(QM+MQ)-4\omega N=0.
\label{q444:twocross:eq:gauss-block}
\end{equation}
The upper entry is the negative lower trace, since
$\operatorname{tr}N=2b$. No off-block entries occur. The entries needed
below are therefore exact consequences of the full Gauss equation:
\begin{equation}
b=-\frac e{2\omega}\left(u+2v_0+\frac t{24}\right),\qquad
N_{12}=-\frac{3e}{2\omega}m,\qquad
N_{22}=-\frac e\omega\left(2v_0+\frac{5t}{24}\right).
\label{q444:twocross:eq:gauss-needed}
\end{equation}
All divisions use the original $\omega\ne0$.

Put
\[
D=\left(\frac t4-\omega^2\right)I_2-e^2Q^2.
\]
The complete spatial residuals are $X(s_1),X(s_2)$, respectively, where
\begin{align}
s_1&=(D+W_0I_2)v+
 \left(-\alpha_0I_2+i(2\omega eQ-3bI_2)\right)w,\label{q444:twocross:eq:s1}\\
s_2&=(D+V_0I_2)w+
 \left(-\alpha_0I_2-i(2\omega eQ-3bI_2)\right)v.\label{q444:twocross:eq:s2}
\end{align}
To verify every term, $[C,X(v)]=X(-ieQv)$ gives the two gyroscopic
terms and the term $-e^2Q^2$. The inertia gives $tI_2/4$.
For the remaining double commutators, block multiplication gives
\[
\mathcal KX(v)=X\bigl(W_0v+(\bar p_0-2p_0)w\bigr),\qquad
\mathcal KX(w)=X\bigl(V_0w+(p_0-2\bar p_0)v\bigr).
\]
Since $\bar p_0-2p_0=-\alpha_0-3ib$ and
$p_0-2\bar p_0=-\alpha_0+3ib$, these are exactly
\eqref{q444:twocross:eq:s1}--\eqref{q444:twocross:eq:s2}. Their upper rows are the negative
adjoints of their lower first columns, and all other entries are zero.
Thus $s_1=s_2=0$ is equivalent to the two full spatial equations.

We now use a consequence of those equations; no converse is needed.
Multiplying \eqref{q444:twocross:eq:s1} by $v^*$, \eqref{q444:twocross:eq:s2} by $w^*$ and adding
gives the exact matrix identity
\begin{equation}
s_1v^*+s_2w^*
=DM+d_MI_2+(4bI_2-2\omega eQ)N.
\label{q444:twocross:eq:outer-product}
\end{equation}
Indeed $\operatorname{tr}(M)M-M^2=d_MI_2$ follows by multiplying
the two-by-two matrix $M$. Its off-diagonal entries are zero, and its
diagonal entries are both $M_{11}M_{22}-M_{12}M_{21}$.
The cubic terms on the left of \eqref{q444:twocross:eq:outer-product} are
\[
W_0vv^*+V_0ww^*-\alpha_0(vw^*+wv^*)+3ib(vw^*-wv^*).
\]
Those in $\operatorname{tr}(M)M-M^2$ have $-ib$ instead of $3ib$;
their difference is $4bN$. The gyroscopic terms sum to
$-2\omega eQN$. These expansions prove \eqref{q444:twocross:eq:outer-product}
entry by entry, without a reality or orthogonality restriction on $v,w$.

\begin{proof}[Proof of Theorem~\ref{q444:twocross:thm:exclusion}]
Gauss and the scalar equations have already forced the exact temporal
form and \eqref{q444:twocross:eq:second-scalar}. In particular $e\ne0$, $m\ne0$.
For a full solution, \eqref{q444:twocross:eq:outer-product} is zero. Set
\[
\chi=\frac{3e}{2\omega},\qquad B_2=4b-4\omega e,
\quad D_2=\frac t4-\omega^2-4e^2.
\]
Its $(2,1)$ entry is
$D_2\bar m+B_2N_{21}=\bar m(D_2-\chi B_2)=0$ by Gauss.
Since $m\ne0$, it follows that $D_2=\chi B_2$.
Its $(2,2)$ entry is therefore
\[
0=d_M+B_2(\chi v_0+N_{22}).
\]
Using exactly \eqref{q444:twocross:eq:gauss-needed} gives
\[
\chi v_0+N_{22}=-\frac e{2\omega}\left(v_0+\frac{5t}{12}\right),
\qquad
B_2=-\frac{2e}{\omega}
 \left(u+2v_0+\frac t{24}+2\omega^2\right).
\]
Consequently every proposed full solution in this class would satisfy
\begin{equation}
0=d_M+\frac{e^2}{\omega^2}
\left(u+2v_0+\frac t{24}+2\omega^2\right)
\left(v_0+\frac{5t}{12}\right).
\label{q444:twocross:eq:positive-contradiction}
\end{equation}
But $u,v_0\ge0$, $d_M\ge0$, $t>0$, $\omega^2>0$ and $e^2>0$.
Both parentheses are strictly positive. The second term on the right
is strictly positive for either sign of $e$ and either sign of $\omega$,
contradicting \eqref{q444:twocross:eq:positive-contradiction}. This proves the theorem
without selecting a frequency or deleting any cross-vector phase.
\end{proof}

\begin{corollary}[A further necessary condition on every full solution]
Every solution of the unrestricted five equations satisfies
\[
\sum_{j=1}^2\left(
\|PZ_jP\|_{\rm HS}^2+
\|(I_3-P)Z_j(I_3-P)\|_{\rm HS}^2\right)>0.
\]
\end{corollary}
\begin{proof}
Each summand is a squared Hilbert--Schmidt norm and is nonnegative.
The sum vanishes if and only if both spatial matrices lie in $\mathfrak p$,
which is excluded by Theorem~\ref{q444:twocross:thm:exclusion}.
\end{proof}

The scalar-compatible part of the class is nonempty even with a nonzero
spatial bracket. This is not a full solution: in the frame
\eqref{q444:twocross:eq:second-frame}, choose the exact values
\[
v=\sqrt\delta(1,1)^{\mathsf T},\qquad w=(i,0)^{\mathsf T},
\qquad e=\sqrt{3+2\delta}.
\]
Then $\operatorname{tr}M=1+2\delta$, $M_{12}=\delta$, and
$[X(v),X(w)]_{11}=-2i\sqrt\delta\ne0$.
Thus both full scalar equations hold by their displayed residuals.
To embed this example into the literal frame, use the fixed
$W=i\left(\begin{smallmatrix}1&1\\1&-1\end{smallmatrix}\right)/\sqrt2$
and \eqref{q444:twocross:eq:original-inverse}. The choice of the component $i$ is a
specified example, not a normalization of the general class. The theorem
proves that no retained $t>0,\omega\ne0$ can complete this example to all
five equations.

\subsubsection{Exact inclusion of the selected class in the exceptional locus}
\label{q444:twocross:exceptional-inclusion-section}

We give the determinant map explicitly to specify the comparison with the
full sixteen-coordinate projection, including its degree conventions.
In this section the symbols $H,q,F$ again denote their literal original
matrices in \eqref{q444:twocross:eq:H}--\eqref{q444:twocross:eq:S2},
identical to the earlier retained data of
\eqref{q444:cx:eq:phi}--\eqref{q444:cx:eq:sp2}. For any such matrix write
$\widetilde Y=U^*YU$ when using its first, $H$-diagonal-frame value.
The inverse remains $Y=U\widetilde YU^*$. Put
\[
J=iH,\quad j_0=\langle J,J\rangle=\frac34,\quad
\mathscr W=J^\perp,\quad
\Pi X=X-\frac{\langle J,X\rangle}{j_0}J,\quad \mathscr B=[Z_1,Z_2],
\]
\[
\eta=\langle J,\mathcal KJ\rangle,\quad w_J=\Pi\mathcal KJ,
\quad\beta=\langle J,\mathscr B\rangle,\quad
\mathcal I_7=\mathcal I|_{\mathscr W},
\]
\[
\mathcal S_7=\Pi\mathcal K|_{\mathscr W}
-\frac{w_J\otimes w_J}{\eta},\qquad
f_7=\Pi\mathscr B-\frac\beta\eta w_J.
\]
Here $(w_J\otimes w_J)X=w_J\langle w_J,X\rangle$ and the original length of
$J$ is retained. On $\eta>0$ define
\[
\mathscr M_7(t)=t\mathcal I_7+2\mathcal S_7,\qquad
D_7=\det_{\mathscr W}\mathscr M_7,\qquad
V_7=4\operatorname{adj}_{\mathscr W}(\mathscr M_7)f_7,
\]
\[
N_7=V_7+\frac{2\beta D_7-\langle w_J,V_7\rangle}{\eta}J,
\qquad Q_7=\|V_7q\|^2,
\quad\alpha=\frac14+\|Z_1q\|^2+\|Z_2q\|^2,
\]
\[
P_H(t)=\alpha[N_7,[N_7,H]]
-Q_7\left(\sum_{j=1}^2[Z_j,[Z_j,H]]+F\right).
\]
Let $F_\xi=\eta^{16}(\widetilde P_H)_{11}$ and
$G_\xi=\eta^{16}(\widetilde P_H)_{22}$ denote the exact polynomials used for
the original sixteen-coordinate exceptional set, with inverse division
by $\eta^{16}$ on $\eta>0$. The matrix $P_H$ in its definition is
computed with the literal original $H,q,F$; $\widetilde P_H=U^*P_HU$
is then its exact conjugate, not a second change of already-conjugated data.
They extend polynomially in the sixteen real spatial coordinates $\xi$.
For clarity, their polynomial extension follows from the definitions:
$\eta^7D_7$, $\eta^7V_7$, $\eta^8N_7$ and $\eta^{14}Q_7$ are polynomial,
by determinant and minor expansions. Hence $\eta^{16}P_H$ is polynomial.
All other denominators are fixed nonzero constants. Their two selected
diagonal entries are real because $P_H$ is Hermitian for real $t$.
The universal degree bounds are $\deg V_7\le6$ and
$N_7=V_7+s_7J$ with $\deg s_7\le7$. Since
$[\widetilde J,\widetilde H]=0$ and every diagonal entry of
$[\widetilde J,Y]$ is zero, the diagonal entries of
$[\widetilde N_7,[\widetilde N_7,\widetilde H]]$ have degree at most
twelve; also $\deg Q_7\le12$. Thus $F_\xi,G_\xi$ have degree at most
twelve in $t$, including every specialization where their degrees drop.

The exceptional polynomial is the determinant $\mathscr R(\xi)$ of the
fixed map
\begin{equation}
\mathbb R[t]_{<12}\oplus\mathbb R[t]_{<12}\longrightarrow
\mathbb R[t]_{<24},\qquad (A_0,B_0)\longmapsto A_0F_\xi+B_0G_\xi,
\label{q444:twocross:eq:fixed-map}
\end{equation}
in increasing-power bases, with the first summand before the second.
Theorem~\ref{q444:generic-projection-obstruction} proves this polynomial
is nonzero on the ambient sixteen-dimensional carrier and that all full
solutions project to its zero set. The following inclusion is proved directly, independently
of whether any point of that zero set admits a solution.

\begin{proposition}\label{q444:twocross:prop:exceptional-inclusion}
The exact linear embedding $\mathfrak p\times\mathfrak p\hookrightarrow\mathfrak{su}(3)^2$ has image
inside $\{\mathscr R=0\}$. On every nonzero pair in this class the two
diagonal Gauss-eliminated polynomials have the common factor $D_4(t)^2$,
where
\[
D_4(t)=\det_{\mathfrak p}\left(\frac t2\operatorname{id}_{\mathfrak p}+2\mathcal K|_{\mathfrak p}\right)
\]
has degree exactly four and leading coefficient $1/16$.
\end{proposition}
\begin{proof}
The first-frame $\widetilde J=i\operatorname{diag}(0,r,-r)$ satisfies
$[\widetilde J,X(v)]=X(i\operatorname{diag}(r,-r)v)$.
The change by $U$ preserves $\mathfrak p$ and the inner product.
Thus, for the selected class,
\[
\eta=\sum_j\|[Z_j,J]\|^2
=r^2(\|Z_1\|^2+\|Z_2\|^2).
\]
It is positive except at the zero pair. This also shows why this class is
compatible with the stabilizer decomposition: $-\operatorname{ad}_J^2$
is exactly $r^2\operatorname{id}$ on $\mathfrak p$.

Put $\mathfrak k_0=\mathfrak k\cap J^\perp$. The orthogonal decomposition is
$\mathscr W=\mathfrak k_0\oplus\mathfrak p$, with dimensions three and four.
The already-proved block preservation of $\mathcal I,\mathcal K$,
and $J\in\mathfrak k$, give $w_J,f_7\in\mathfrak k_0$.
Consequently $\mathcal S_7$ preserves both summands and equals
$\mathcal K$ on $\mathfrak p$. Thus
\[
\mathscr M_7=\begin{pmatrix}\mathscr M_3&0\\0&\mathscr M_4\end{pmatrix},
\qquad \mathscr M_4=\frac t2\operatorname{id}_{\mathfrak p}+2\mathcal K|_{\mathfrak p},
\quad D_7=D_3D_4.
\]
Its four-dimensional determinant has degree four and leading coefficient
$(1/2)^4=1/16$. Moreover $\mathscr M_4$ is positive for $t>0$.
For completeness, $\mathcal I$ has kernel $\mathbb RJ$: its quadratic
form is $\operatorname{tr}[X,H]^2+2\|Xq\|^2$; a zero commutes with the
three-distinct-eigenvalue $H$ and annihilates $q$, so has diagonal
entries $0,is,-is$. The same formula proves positivity on $\mathscr W$.
Completing the square in $\langle X+sJ,\mathcal K(X+sJ)\rangle$ gives
\[
\langle X,\mathcal S_7X\rangle
+\eta\left(s+\frac{\langle w_J,X\rangle}{\eta}\right)^2\ge0,
\]
so choosing the displayed square to be zero proves
$\mathcal S_7\succeq0$. Thus $\mathscr M_7$ and its three-dimensional
block are positive for every $t>0$.

Define $V_3=4\operatorname{adj}(\mathscr M_3)f_7\in\mathfrak k_0$.
The block inverse for $t>0$ gives
\[
V_7=D_4V_3,\qquad
N_7=D_4N_3,\qquad
N_3=V_3+\frac{2\beta D_3-\langle w_J,V_3\rangle}\eta J.
\]
These are polynomial identities in $t$: the two sides agree for every
positive $t$, while a nonzero polynomial has only finitely many roots,
as follows by repeated division by its linear root factors.
Therefore $Q_7=D_4^2Q_3$, where $Q_3=\|V_3q\|^2$, and
\begin{equation}
P_H=D_4^2P_{H,3},\qquad
P_{H,3}=\alpha[N_3,[N_3,H]]
-Q_3\left(\sum_j[Z_j,[Z_j,H]]+F\right).
\label{q444:twocross:eq:factor}
\end{equation}
The three-dimensional minors give $\deg V_3\le2$, $\deg D_3\le3$.
Writing $N_3=V_3+s_3J$, one has $[J,H]=0$. In the diagonal frame,
$[J,Y]_{jj}=0$ for every matrix $Y$, by subtracting the identical
diagonal products. Hence the selected diagonal entries of
$[N_3,[N_3,H]]=[V_3,[V_3,H]]+s_3[J,[V_3,H]]$ have degree at most four.
Also $\deg Q_3\le4$. Thus at each fixed pair with $\eta>0$,
\[
F_\xi=D_4^2F_3,\qquad G_\xi=D_4^2G_3,\qquad
\deg F_3,\deg G_3\le4,
\]
where $F_3,G_3$ retain the scalar multiplier $\eta^{16}$.
If they are not both zero, $(A_0,B_0)=(G_3,-F_3)$ is a nonzero
kernel element of \eqref{q444:twocross:eq:fixed-map}, with both degrees less than
twelve. If both are zero, $(1,0)$ is such a kernel element instead.
The determinant therefore vanishes for every nonzero pair in the class,
including every degree-drop case. At the zero pair it vanishes by
continuity of the original polynomial $\mathscr R$ along this linear
subspace. This proves the entire asserted inclusion without using a
converse resultant test or interpreting $\eta=0$ as an inverse solution.
\end{proof}

\subsubsection{An exact certificate retaining nonzero residuals}
\label{q444:twocross:quantitative-section}

The exclusion has a quantitative extension. It also gives a polynomial
identity before the Gauss or spatial residuals are set to zero. Use the
entire second-frame variables above, with every entry of $v,w$ retained.
More precisely, return here to the matrices and inverse in
\eqref{q444:twocross:eq:second-frame}--\eqref{q444:twocross:eq:original-inverse},
not the literal-frame $H$ used in the preceding determinant comparison.
All of $M,N,m,u,v_0,b,s_1,s_2$ in this subsection are computed in
that same second frame as in \eqref{q444:twocross:eq:s1}--\eqref{q444:twocross:eq:s2}.
Define the full spatial residual matrix and the Hermitian Gauss components
\[
 \mathcal A=s_1v^*+s_2w^*,\qquad
 \gamma_0=et/6+4e(u+2v_0)+8\omega b,
 \qquad \mathcal G=etE-2e(QM+MQ)-4\omega N.
\]
Here $\mathcal A$ is not the original generator $A$: its calligraphic
symbol denotes exactly this two-by-two residual matrix. The full Gauss
residual is $i\operatorname{diag}(\gamma_0,\mathcal G)$, with the
second block of size two. Set $B_2=4b-4\omega e$ and
\[
 \mathcal P=d_M+\frac{e^2}{\omega^2}
       (u+2v_0+t/24+2\omega^2)(v_0+5t/12).
\]
This expression is defined for the original nonzero $\omega$; the
identity below has no such denominator after multiplication.

\begin{proposition}[Polynomial residual certificate and strict error bound]
\label{q444:twocross:quantitative}
For every $v,w\in\mathbb C^2$ and real $e,t,\omega$ in the displayed
second-frame matrices, the following polynomial identity holds:
\begin{align}
 &4\omega^2(\bar m\mathcal A_{22}-v_0\mathcal A_{21})
   +\omega B_2(\bar m\mathcal G_{22}-v_0\mathcal G_{21})\notag\\
 &\hspace{8mm}+e(v_0+5t/12)\bar m\gamma_0\notag\\
 &=4\bar m\bigl[\omega^2d_M+
      e^2(u+2v_0+t/24+2\omega^2)(v_0+5t/12)\bigr].
 \label{q444:twocross:polynomial-certificate}
\end{align}
No scalar, Gauss or spatial equation is assumed in this identity.

Now restrict to the exactly specified set of original tuples with
$Z_1,Z_2\in\mathfrak p$, $t>0$, $\omega\ne0$, for which both
scalar equations and Gauss hold. Let $R_1,R_2\in\mathfrak{su}(3)$
be their two full spatial residuals. Then
\begin{equation}
 \bigl(\|R_1\|_{\mathrm{HS}}^2+
             \|R_2\|_{\mathrm{HS}}^2\bigr)^{1/2}
 \geq
 \frac{\sqrt2\,\mathcal P}
 {\sqrt{\operatorname{tr}M\,(1+v_0^2/|m|^2)}}
 >\frac{5\sqrt2\,e^2t}
 {6\sqrt{\operatorname{tr}M\,(1+v_0^2/|m|^2)}}.
 \label{q444:twocross:residual-lower-bound}
\end{equation}
The full matrix norms are preserved by the two exact unitary frame maps.
In particular, on any such set with $\operatorname{tr}M\leq L$ and
$t\geq t_0>0$, where $L>0$, the strictly positive uniform bound is
\begin{equation}
 \bigl(\|R_1\|_{\mathrm{HS}}^2+
             \|R_2\|_{\mathrm{HS}}^2\bigr)^{1/2}
 >\frac{5\sqrt2\,t_0}
 {3\sqrt{L(1+L^2/\delta^2)}},\qquad \delta=1/(2r).
 \label{q444:twocross:uniform-bound}
\end{equation}
The complete scalar--Gauss solution set is parameterized in
Theorems~\ref{sgd:parametrization} and~\ref{sgd:literal-addon:quotient},
and Corollary~\ref{sgd:every-coupling} proves it is nonempty at every
original $t=g^2>0$. No sequence in its displayed bounded portion can
make both spatial residuals tend to zero.
\end{proposition}

\begin{proof}
The complete residual formulas above give
\[
 \mathcal A_{21}=D_2\bar m+B_2N_{21},\qquad
 \mathcal A_{22}=D_2v_0+d_M+B_2N_{22},\quad
 D_2=t/4-\omega^2-4e^2,
\]
and, without any equation imposed,
\[
 \mathcal G_{21}=-6e\bar m-4\omega N_{21},\qquad
 \mathcal G_{22}=-5et/6-8ev_0-4\omega N_{22}.
\]
Subtracting the two indicated spatial entries cancels $D_2$ exactly:
\[
 \bar m\mathcal A_{22}-v_0\mathcal A_{21}
 =\bar m d_M+B_2(\bar mN_{22}-v_0N_{21}).
\]
Multiplication of the indicated Gauss entries gives
\[
 \bar m\mathcal G_{22}-v_0\mathcal G_{21}
 =-e(2v_0+5t/6)\bar m
      -4\omega(\bar mN_{22}-v_0N_{21}).
\]
These cancel the two $N$ entries on the left of
\eqref{q444:twocross:polynomial-certificate}. The remaining factor is
\[
 \gamma_0-2\omega B_2
 =4e(u+2v_0+t/24+2\omega^2).
\]
It yields exactly its stated right side. Each displayed step is a
polynomial identity, including at $m=0$ or $\omega=0$; no inverse or
solution is assigned at those values.

On the stated scalar--Gauss zero set the preceding proof forces the
second-frame form, $e^2=2+\operatorname{tr}M$ and
$\operatorname{Re}m=\delta>0$. Hence $m\ne0$ and
$\operatorname{tr}M>0$. Setting $\gamma_0=\mathcal G=0$ in the
polynomial identity and dividing by its retained nonzero factors gives
\[
 \mathcal A_{22}-\frac{v_0}{\bar m}\mathcal A_{21}
 =\mathcal P>0.
\]
Cauchy--Schwarz on these two entries implies
\[
 \mathcal P\leq
  \sqrt{1+v_0^2/|m|^2}\,\|\mathcal A\|_{\mathrm{HS}}.
\]
Let $S=(s_1\ s_2)$ and $V=(v\ w)$ be the two column matrices.
Then $\mathcal A=SV^*$. Applying Cauchy--Schwarz to each entry and
summing proves $\|SV^*\|_{\mathrm{HS}}^2
\leq\|S\|_{\mathrm{HS}}^2\|V\|_{\mathrm{HS}}^2$.
The right factors are $\|s_1\|^2+\|s_2\|^2$ and
$\operatorname{tr}M$, respectively. Directly from the two off-diagonal
blocks of $X(s)$, $\|X(s)\|_{\mathrm{HS}}^2=2\|s\|^2$.
Thus $R_j=X(s_j)$ proves the first inequality in
\eqref{q444:twocross:residual-lower-bound}.

The exact expansion
\begin{align*}
 &(u+2v_0+t/24+2\omega^2)(v_0+5t/12)\\
 &=5t\omega^2/6
   +v_0(u+2v_0+t/24+2\omega^2)
   +(5t/12)(u+2v_0+t/24)
\end{align*}
has nonnegative remainder terms, the last strictly positive for $t>0$.
Together with $d_M\geq0$ and $e^2>0$ it gives
$\mathcal P>5e^2t/6$, proving the strict second inequality.
Finally $e^2=2+\operatorname{tr}M\geq2$,
$0\leq v_0\leq\operatorname{tr}M\leq L$, and $|m|\geq\delta$
give \eqref{q444:twocross:uniform-bound}. Since unitary conjugation
preserves the Hilbert--Schmidt norm, the same inequalities hold for
the complete spatial residuals in the original literal matter frame.
\end{proof}


\subsubsection{Scope, provenance, and exact replay}
\label{q444:twocross:scope-replay}

Theorem~\ref{q444:twocross:thm:exclusion} and Proposition~\ref{q444:twocross:prop:exceptional-inclusion}
together prove an empty admissible fibre at every point of the explicit
eight-dimensional two-cross linear subspace inside the full exceptional set.
The result does not prove global emptiness or produce a full solution.
It does not assert a result for arbitrary time-dependent gauge connections,
and it makes no quantum mass-gap assertion. The earlier cross--cross--block exclusion in
Theorem~\ref{q444:cx:thm:block} instead assumes $C,Z_1\in\mathfrak p$, $Z_2\in\mathfrak k$;
here only $Z_1,Z_2\in\mathfrak p$ were assumed and $C\in\mathfrak k$ was proved.

The retained literal system is
\eqref{q444:cx:eq:phi}--\eqref{q444:cx:eq:sp2}; its complete scalar
coordinates are given in \eqref{q444:complex-eight-coordinates} and
Proposition~\ref{q444:complex-full-scalar-equivalence}.
The ambient projection statement and determinant convention are
Theorem~\ref{q444:generic-projection-obstruction} and
\eqref{q444:generic-sylvester-map}. The complete source proof, exact
checker, and replay receipt for this additional exclusion are retained
under \path{supporting_materials/workbench/evidence/two_cross_exceptional_20260905/}.
All matrix identities needed for the theorem are derived above; no
external classification theorem, numerical root search, or unproved
literature identity is used.

The companion \path{check_two_cross_exceptional.py} checks 26 exact groups:
the original force and constants, both frame maps (the second over the
entire coordinate ring $\gamma\bar\gamma=1$), every scalar residual before
and after temporal classification, both inertia maps, the full bracket and
Gauss residual, both complete spatial residuals, the nonnegative determinant
identity, the full outer-product identity and the final positivity formula.
Every real component of $v,w$ is an independent symbol. It uses no floating
point acceptance criterion and writes no files. One serial run passed all
26 groups in $0.416$ seconds of measured checker time.
The determinant block-factor proof and all positivity/domain arguments
are analytic proofs in the preceding sections, not claims inferred from
sampled determinant values.

The additional denominator-free residual identity and the full spatial
error bound are proved in Proposition~\ref{q444:twocross:quantitative}.
The companion \path{check_residual_certificate.py} verifies its polynomial
identity and exact positive remainder. The Cauchy--Schwarz estimates and
their scalar--Gauss domain are proved in the text, not inferred from tests.

% Shared complete scalar/Gauss proof, with literal-domain quotient addendum.
% Original sources and exact transformation receipt are in
% supporting_materials/workbench/evidence/scalar_gauss_domain_20260905/.
% No preamble, global macros, or hidden mathematical input.
\subsection{The complete scalar--Gauss domain at every positive coupling}
\label{q444:sgd:section}

\subsubsection{Literal data, residuals, and coordinate maps}

All complex vector spaces below retain their underlying real structures.
The adjoint is denoted by $*$, and $[X,Y]=XY-YX$. Keep exactly
\[
 H_0=\begin{pmatrix}0&0&0\\0&-1/2&1/(2\sqrt2)\\
 0&1/(2\sqrt2)&1/2\end{pmatrix},\qquad
 q=\frac1{2\sqrt2}\begin{pmatrix}1\\0\\0\end{pmatrix},\qquad
 A_0=\sqrt2\begin{pmatrix}0&-1&0\\1&0&0\\0&0&0\end{pmatrix}.
\]
The original constants are $m^2=\Gamma=1$ and $\kappa=6$.
For $Y_{\rm tf}=Y-(\operatorname{tr} Y)I_3/3$, define
\[
 \mathcal I_HY=([H,[H,Y]]+qq^*Y+Yqq^*)_{\rm tf},\qquad
 \mathcal K_{Z_1,Z_2}Y=-\sum_{j=1}^2[Z_j,[Z_j,Y]].
\]
Direct multiplication gives
\[
 F_0=[A_0,[A_0,H_0]]=-2H_0+\operatorname{diag}(-2,1,1),\quad
 H_0^2=\frac38\operatorname{diag}(0,1,1),\quad q^*q=\frac18.
\]
The scalar--Gauss domain in question consists of
\[
 (t,\omega,C_0,Z_{1,0},Z_{2,0}),\quad t=g^2>0,\quad
 \omega\in\mathbb R\setminus\{0\},\quad C_0,Z_{j,0}\in\mathfrak{su}(3),
\]
where $P=qq^*/(q^*q)=\operatorname{diag}(1,0,0)$ and
\[
 PZ_{j,0}P=(I_3-P)Z_{j,0}(I_3-P)=0,
\]
satisfying exactly the following three full residual equations:
\begin{align}
 [C_0,[C_0,H_0]]-\sum_j[Z_{j,0},[Z_{j,0},H_0]]-F_0&=0,
 \label{sgd:literal-H}\\
 (C_0^2-Z_{1,0}^2-Z_{2,0}^2)q+2q&=0,
 \label{sgd:literal-q}\\
 (t\mathcal I_{H_0}+2\mathcal K_{Z_{1,0},Z_{2,0}})C_0
             -4\omega[Z_{1,0},Z_{2,0}]&=0.
 \label{sgd:literal-G}
\end{align}
Equation~\eqref{sgd:literal-G} is the original Gauss equation multiplied
by its retained positive $t$; division by the same $t$ is its inverse.
Neither spatial equation is imposed anywhere in the construction.

The exact maps to and from the classified frame are as follows. Retain
\[
 \begin{gathered}
 r=\frac{\sqrt6}{4},\quad \delta=\frac1{2r}=\sqrt{\frac23},\quad
 x=\frac12,\quad d_0=\frac1{2\sqrt2},\\
 S=\frac1{\sqrt{2r(r+x)}}
       \begin{pmatrix}d_0&-(r+x)\\r+x&d_0\end{pmatrix},\quad
 U=\operatorname{diag}(1,S).
 \end{gathered}
\]
The identity $r^2=x^2+d_0^2$ gives $S^*S=I_2$, $\det S=1$;
multiplication of $H_0$ by its two displayed lower columns gives
\[
 U^*H_0U=\operatorname{diag}(0,r,-r),\qquad U^*q=q.
\]
Choose any $\zeta\in\mathbb C$ with $|\zeta|=1$, and set
\[
 \gamma=-\zeta^{-2},\qquad
 W=\frac\zeta{\sqrt2}\begin{pmatrix}\gamma&\gamma\\1&-1\end{pmatrix}
  =\frac1{\sqrt2}\begin{pmatrix}-\zeta^{-1}&-\zeta^{-1}\\
                         \zeta&-\zeta\end{pmatrix},\qquad
 V=\operatorname{diag}(1,W),\quad B=UV.
\]
Here $|\gamma|=1$ and $\zeta^2=-\overline\gamma$, retaining exactly
the original phase relation. Every unit $\gamma$ occurs, with its two
possible square roots $\zeta$. The displayed entries give
\[
 W^*W=I_2,\quad \det W=1,\quad
 W^*\operatorname{diag}(r,-r)W=\begin{pmatrix}0&r\\r&0\end{pmatrix}.
\]
Thus, simultaneously for all data,
\begin{equation}
 Y'=B^*Y_0B,\qquad Y_0=BY'B^*,\qquad
 H'=\begin{pmatrix}0&0&0\\0&0&r\\0&r&0\end{pmatrix},\quad
 q'=q,\quad F'=-2H'+\operatorname{diag}(-2,1,1).
 \label{sgd:frame}
\end{equation}
These are inverse maps, with no constant or unknown rescaled. In fact
$B^*(XY)B=(B^*XB)(B^*YB)$, the same equality holds for brackets,
trace and adjoints, and $Bq=q$. Therefore $\mathcal I$ and $\mathcal K$
and every full residual in \eqref{sgd:literal-H}--\eqref{sgd:literal-G}
are conjugated exactly (the vector residual is multiplied by $B^*$).
This proves that zero residuals in the classified frame give zero residuals
in the literal frame, with the displayed inverse recovering every datum.
The block-diagonal $B$ commutes with $P$, preserving the two-cross class.

\subsubsection{Every scalar and Gauss entry in the classified frame}

For $v,w\in\mathbb C^2$ write
\[
 X(v)=\begin{pmatrix}0&-v^*\\v&0_{2\times2}\end{pmatrix},\qquad
 Z_1=X(v),\quad Z_2=X(w),\quad
 C=ie\operatorname{diag}(1,0,-1),\quad e\in\mathbb R\setminus\{0\}.
\]
The inverse of $v\mapsto X(v)$ is the lower first column; every
$X(v)$ is traceless anti-Hermitian. In this section all unprimed matrices
are these classified-frame matrices, with literal inverse~\eqref{sgd:frame}.
Define
\[
 \begin{gathered}
 M=vv^*+ww^*=\begin{pmatrix}u&m\\\bar m&v_0\end{pmatrix},\quad
 N=i(vw^*-wv^*),\\
 b=\operatorname{Im}(v^*w),\quad
 Q=\operatorname{diag}(1,2),\quad E=\operatorname{diag}(2/3,-5/6).
 \end{gathered}
\]
The subscript on the real scalar $v_0$ distinguishes it from the vector $v$.

Here are the complete residuals, rather than only contractions of them:
\begin{align}
 Q_H&=\begin{pmatrix}
 2-4r\operatorname{Re}m&0&0\\
 0&2r\operatorname{Re}m-1&r(\operatorname{tr} M-e^2+2)\\
 0&r(\operatorname{tr} M-e^2+2)&2r\operatorname{Re}m-1
 \end{pmatrix},\label{sgd:scalar-H}\\
 Q_q&=(\operatorname{tr} M-e^2+2)q,\label{sgd:scalar-q}\\
 Q_G&=i\operatorname{diag}(\gamma_0,\mathcal G),\label{sgd:Gauss-full}\\
 \gamma_0&=et/6+4e(u+2v_0)+8\omega b,\qquad
 \mathcal G=etE-2e(QM+MQ)-4\omega N.\label{sgd:Gauss-entries}
\end{align}
In \eqref{sgd:Gauss-full} the second diagonal block has size two.
To derive every displayed entry, multiplication gives
\[
 X(v)^2=\operatorname{diag}(-v^*v,-vv^*),\quad
 [C,X(v)]=X(-ieQv),\quad
 [X(v),X(w)]=\operatorname{diag}(-2ib,iN).
\]
The first two identities and $H'_{23}=H'_{32}=r$ give
\eqref{sgd:scalar-H} and \eqref{sgd:scalar-q}; all their omitted entries
are zero. The two contributions to $\mathcal I_{H'}C$ are
\[
 ie\operatorname{diag}(0,3/4,-3/4),\qquad
 ie\operatorname{diag}(1/6,-1/12,-1/12),
\]
so $\mathcal I_{H'}C=ie\operatorname{diag}(1/6,2/3,-5/6)$. Also
\[
 \mathcal KC=\operatorname{diag}\bigl(2ie\operatorname{tr}(QM),-ie(QM+MQ)\bigr).
\]
Substitution of these full blocks gives \eqref{sgd:Gauss-full}.
Finally $\operatorname{tr} N=2b$, since $\operatorname{tr}(vw^*)=w^*v=\overline{v^*w}$.
Thus $\operatorname{tr}\mathcal G=-\gamma_0$; the upper Gauss entry vanishes
whenever the full lower block does, and this implication is an identity,
not an omitted equation. The two scalar equations are equivalent to
\begin{equation}
 e^2=2+u+v_0,\qquad \operatorname{Re}m=\delta.
 \label{sgd:scalar-conditions}
\end{equation}

\subsubsection{The exact realization of the two Hermitian matrices}

\begin{lemma}\label{sgd:realization-necessity}
For every pair $v,w$, let $D=(v\ w)$ and
$J=\left(\begin{smallmatrix}0&i\\-i&0\end{smallmatrix}\right)$.
Then $M=DD^*$, $N=DJD^*$ and
\[
 M+N=(v-iw)(v-iw)^*,\qquad
 M-N=(v+iw)(v+iw)^*.
\]
If $M$ is positive definite, then
\begin{equation}
 \operatorname{tr}(M^{-1}N)=0,\qquad \det N=-\det M.
 \label{sgd:trace-det}
\end{equation}
\end{lemma}
\begin{proof}
Multiplication of the displayed two-by-two blocks gives all three
factorizations. If $M>0$, $D$ is invertible because
$z^*Mz=\|D^*z\|^2$ has zero kernel. Consequently
$M^{-1}N=D^{-*}JD^*$ is similar to $J$, whose trace is zero
and determinant is $-1$. Taking the trace and determinant proves
\eqref{sgd:trace-det}, with every factor retained.
\end{proof}

\begin{lemma}\label{sgd:rank-one-impossible}
Every classified scalar--Gauss solution has $M>0$.
\end{lemma}
\begin{proof}
The scalar equation gives $m=\delta+iy\ne0$, so $u,v_0>0$
and $M\ne0$. It is already positive semidefinite as $DD^*$.
If $M$ had rank one, the range of $N=DJD^*$ would be contained
in its one-dimensional range. A Hermitian matrix with that range
and its orthogonal complement in its kernel is a real scalar multiple
of $M$, so $N=\nu M$ for some real $\nu$.
Set $k=e/\omega\ne0$. The full Gauss equation reads
\begin{equation}
 N=kT,\qquad
 T=\frac t4 E-\frac12(QM+MQ)
  =\begin{pmatrix}t/6-u&-3m/2\\-3\bar m/2&-5t/24-2v_0\end{pmatrix}.
 \label{sgd:T}
\end{equation}
Its off-diagonal entry gives $\nu=-3k/2$ because $m\ne0$.
The $(1,1)$ entry then gives $t/6-u=-3u/2$, or $t/6=-u/2$,
which contradicts $t>0$ and $u>0$. Thus $M$ has rank two and is
positive definite.
\end{proof}

\begin{theorem}[Complete classified-frame parameterization]
\label{sgd:parametrization}
Define the open real parameter set
\[
 \mathcal D=\{(u,v_0,y)\in\mathbb R^3:
      u>0,\quad d=uv_0-\delta^2-y^2>0,\quad c=4v_0-5u>0\}.
\]
For $(u,v_0,y)\in\mathcal D$, retain
\begin{align}
 m&=\delta+iy,&
 M&=\begin{pmatrix}u&m\\\bar m&v_0\end{pmatrix},&
 t&=\frac{72d}{c},\label{sgd:parameters1}\\
 T&=\frac t4E-\frac12(QM+MQ),&
 k&=\sigma\sqrt{\frac d{-\det T}},&
 e&=\epsilon\sqrt{2+u+v_0},\label{sgd:parameters2}\\
 \omega&=e/k,&N&=kT,& \epsilon,\sigma&\in\{1,-1\}.
 \label{sgd:parameters3}
\end{align}
The denominator $-\det T$ is strictly positive. The two matrices
$A_+=M+N$ and $A_-=M-N$ are positive semidefinite of rank one,
and their $(1,1)$ entries $h_+,h_-$ are strictly positive. Define
\begin{equation}
 a_\pm=\begin{pmatrix}\sqrt{h_\pm}\\(A_\pm)_{21}/\sqrt{h_\pm}\end{pmatrix},
 \quad p_\pm\in\mathbb C,\quad |p_\pm|=1,\quad
 v=\frac{p_+a_++p_-a_-}{2},\quad
 w=\frac{i(p_+a_+-p_-a_-)}{2}.
 \label{sgd:vectors}
\end{equation}
These formulas give exactly all classified-frame solutions of
\eqref{sgd:scalar-H}--\eqref{sgd:Gauss-entries}. In fact they are a
bijection from $\mathcal D\times\{1,-1\}^2\times U(1)^2$ to
that solution set with coordinates $(t,\omega,e,v,w)$.
\end{theorem}
\begin{proof}
First consider any solution. Lemma~\ref{sgd:rank-one-impossible}
gives $M>0$ and $d>0$. Inverting its two-by-two matrix gives
\[
 M^{-1}=\frac1d\begin{pmatrix}v_0&-m\\-\bar m&u\end{pmatrix}.
\]
The trace condition in \eqref{sgd:trace-det}, divided by the retained
nonzero $k$, now gives exactly
\begin{align}
 0=d\operatorname{tr}(M^{-1}T)
 &=v_0(t/6-u)+u(-5t/24-2v_0)+3|m|^2\notag\\
 &=\frac{t(4v_0-5u)}{24}-3d.\label{sgd:trace-condition}
\end{align}
Hence $c>0$ and $t=72d/c$. The determinant condition gives
$k^2\det T=-d$, so $\det T<0$ and $k$ is precisely one of
the two choices in \eqref{sgd:parameters2}. The scalar conditions
give precisely the two signs of $e$ stated there. The original
$\omega=e/k$ is unchanged, with sign $\epsilon\sigma$.

Conversely, start with parameters in $\mathcal D$. Positivity of $M$
follows from the exact factorization
\[
 M=LL^*,\qquad
 L=\begin{pmatrix}\sqrt u&0\\\bar m/\sqrt u&\sqrt d/\sqrt u\end{pmatrix},
 \qquad \det L=\sqrt d>0.
\]
The specified positive $t$ makes \eqref{sgd:trace-condition} zero.
Put $R=L^{-1}TL^{-*}$. This is Hermitian and
$\operatorname{tr} R=\operatorname{tr}(M^{-1}T)=0$, so its full form is
\[
 R=\begin{pmatrix}a&z\\\bar z&-a\end{pmatrix},\qquad a\in\mathbb R.
\]
The nonzero $m$ makes $T_{12}=-3m/2\ne0$, hence $R\ne0$.
Thus $\det R=-a^2-|z|^2<0$ and $\det T=d\det R<0$.
This proves that both signs of $k$ are well defined. Let $J_0=kR$.
It is Hermitian, has trace zero and determinant $-1$, and multiplying
the displayed matrix shows $J_0^2=I_2$.
For either sign and every column $z_0$, one has
\[
 z_0^*(I_2\pm J_0)z_0=\frac12\|(I_2\pm J_0)z_0\|^2\ge0.
\]
Also $\det(I_2\pm J_0)=0$ and $\operatorname{tr}(I_2\pm J_0)=2$,
so these matrices have rank one. Therefore
\[
 A_\pm=M\pm N=L(I_2\pm J_0)L^*
\]
are positive semidefinite of rank one, as claimed.

Their first diagonal entries are nonnegative. Suppose $h_s=0$ for
a sign $s\in\{1,-1\}$, where $A_s=M+sN$. Its determinant is
zero, so $|(A_s)_{12}|^2=0$. But
\[
 (A_s)_{12}=m(1-3sk/2),\quad
 h_s=u+sk(t/6-u).
\]
Since $m\ne0$, this forces $sk=2/3$ and then
$h_s=u/3+t/9>0$, a contradiction. Thus $h_\pm>0$.
Using $\det A_\pm=0$ gives
$(A_\pm)_{22}=|(A_\pm)_{21}|^2/h_\pm$; consequently
$A_\pm=a_\pm a_\pm^*$ with exactly the vectors in
\eqref{sgd:vectors}.

For those vectors, $v-iw=p_+a_+$ and $v+iw=p_-a_-$.
Their two outer products are $A_+,A_-$ because the phases have unit
modulus. Adding and subtracting the identities in
Lemma~\ref{sgd:realization-necessity} proves
\[
 vv^*+ww^*=M,\qquad i(vw^*-wv^*)=N.
\]
Equation~\eqref{sgd:scalar-conditions} holds by construction, so both
full scalar residuals vanish. Substitution of $N=kT$ and
$\omega k=e$ gives, with all coefficients retained,
\[
 \mathcal G=etE-2e(QM+MQ)
        -4e\left(\frac t4E-\frac12(QM+MQ)\right)=0.
\]
The exact identity $\gamma_0=-\operatorname{tr}\mathcal G$ then checks the
remaining upper Gauss entry. Thus every entry of all three original
residuals vanishes after applying \eqref{sgd:frame}.

For bijectivity, $u,v_0,y$ are recovered from $M$, $\epsilon$ is
the sign of $e$, and $\sigma$ the sign of $e/\omega$. The two
phase factors are recovered without any zero denominator by
\[
 p_+=\frac{(v-iw)_1}{\sqrt{h_+}},\qquad
 p_-=\frac{(v+iw)_1}{\sqrt{h_-}}.
\]
Indeed a column $c$ with $cc^*=a_\pm a_\pm^*$ and $h_\pm>0$
has $|c_1|=\sqrt{h_\pm}$, and the $(2,1)$ entry forces its second
component to be the same unit multiple of $(a_\pm)_2$.
These formulas prove both surjectivity and the exact inverse, without
setting any spatial phase to a selected value.
\end{proof}

\subsubsection{The complete literal scalar--Gauss domain and its phase quotient}
\label{sgd:literal-addon:section}

The classified-frame bijection of Theorem~\ref{sgd:parametrization}
has an exact extension to the entire literal two-cross domain.
We give the forced temporal classification, the full parameter map,
and the complete fibres of that map. All spaces retain their underlying
real structures, $[X,Y]=XY-YX$, and $*$ denotes the adjoint.

\paragraph{The literal domain and the retained constants}

Retain exactly
\[
 H_0=\begin{pmatrix}0&0&0\\0&-1/2&1/(2\sqrt2)\\
 0&1/(2\sqrt2)&1/2\end{pmatrix},\qquad
 q=\frac1{2\sqrt2}\begin{pmatrix}1\\0\\0\end{pmatrix},\qquad
 A_0=\sqrt2\begin{pmatrix}0&-1&0\\1&0&0\\0&0&0\end{pmatrix},
\]
\[
 m^2=\Gamma=1,\qquad \kappa=6,\qquad
 F_0=[A_0,[A_0,H_0]]=-2H_0+\operatorname{diag}(-2,1,1).
\]
Here $\mathfrak{su}(3)$ consists of the traceless anti-Hermitian complex
$3\times3$ matrices, with positive real inner product
$\langle X,Y\rangle=-\operatorname{tr}(XY)$.
For $Y_{\rm tf}=Y-(\operatorname{tr}Y)I_3/3$, define
\[
 \mathcal I_{H_0}Y
 =([H_0,[H_0,Y]]+qq^*Y+Yqq^*)_{\rm tf},\qquad
 \mathcal K_{Z_1,Z_2}Y=-\sum_{j=1}^2[Z_j,[Z_j,Y]].
\]
Multiplication of the retained matrices gives
$H_0q=0$, $H_0^2=(3/8)\operatorname{diag}(0,1,1)$ and $q^*q=1/8$.
Consequently $\operatorname{tr}H_0^2+2q^*q=1$, so the original
scalar coefficient $m^2+\Gamma(\operatorname{tr}H_0^2+2q^*q)$ is $2$.
Also $\operatorname{adj}H_0=\operatorname{diag}(-3/8,0,0)$ gives
$6(\operatorname{adj}H_0-qq^*)_{\rm tf}
=\operatorname{diag}(-2,1,1)$, agreeing with the displayed force.

Put
\[
 P=\frac{qq^*}{q^*q}=\operatorname{diag}(1,0,0),\qquad
 \mathfrak p=\{X\in\mathfrak{su}(3):
 PXP=(I_3-P)X(I_3-P)=0\}.
\]
Let $\mathcal S_{\rm sg}$ be exactly the set of tuples
\[
 (t,\omega,C_0,Z_{1,0},Z_{2,0}),\qquad
 t=g^2>0,\quad \omega\in\mathbb R\setminus\{0\},\quad
 C_0\in\mathfrak{su}(3),\quad Z_{j,0}\in\mathfrak p,
\]
for which the following three full residuals vanish:
\begin{align}
 [C_0,[C_0,H_0]]
 -\sum_j[Z_{j,0},[Z_{j,0},H_0]]-F_0&=0,
 \label{sgd:literal-addon:H}\\
 (C_0^2-Z_{1,0}^2-Z_{2,0}^2)q+2q&=0,
 \label{sgd:literal-addon:q}\\
 (t\mathcal I_{H_0}+2\mathcal K_{Z_{1,0},Z_{2,0}})C_0
 -4\omega[Z_{1,0},Z_{2,0}]&=0.
 \label{sgd:literal-addon:G}
\end{align}
Their codomains are, respectively,
$\operatorname{Herm}_3(\mathbb C)_0$, $\mathbb C^3$, and
$\mathfrak{su}(3)$. The original Gauss residual is
\[
 (\mathcal I_{H_0}+2\mathcal K_{Z_{1,0},Z_{2,0}}/g^2)C_0
 -(4\omega/g^2)[Z_{1,0},Z_{2,0}].
\]
Multiplication by the same retained positive $t=g^2$ gives
\eqref{sgd:literal-addon:G}, and division by $t$ is its inverse.
No spatial residual is required to vanish in the definition of
$\mathcal S_{\rm sg}$.

\paragraph{Temporal classification and all inverse frame maps}

For $v\in\mathbb C^2$, define
\[
 X(v)=\begin{pmatrix}0&-v^*\\v&0_2\end{pmatrix},\qquad
 \mathfrak k=\{X\in\mathfrak{su}(3):[X,P]=0\}.
\]
Every element of $\mathfrak p$ has exactly this form, whose inverse
is the lower first column. Define the fixed constants and matrix
\[
 r=\frac{\sqrt6}{4},\qquad \delta=\frac1{2r}=\sqrt{\frac23},
 \qquad x=\frac12,\qquad d_0=\frac1{2\sqrt2},
\]
\[
 S=\frac1{\sqrt{2r(r+x)}}
 \begin{pmatrix}d_0&-(r+x)\\r+x&d_0\end{pmatrix},
 \qquad U=\operatorname{diag}(1,S).
\]
The identity $r^2=x^2+d_0^2$ shows that the numerator columns
are orthogonal, have squared length $2r(r+x)$, and have determinant
$2r(r+x)>0$. Multiplication by the lower block of $H_0$ gives
their respective eigenvalues $r,-r$. Thus
\[
 U^*U=I_3,\qquad \det U=1,\qquad
 U^*H_0U=\operatorname{diag}(0,r,-r),\qquad Uq=q.
\]

\begin{lemma}[The temporal form is forced by scalars and Gauss]
\label{sgd:literal-addon:temporal}
Every tuple in $\mathcal S_{\rm sg}$ has $C_0\in\mathfrak k$.
In the fixed first frame $C_d=U^*C_0U$, there are unique
$e\in\mathbb R\setminus\{0\}$ and $z\in\mathbb C$ with
$|z|=|e|/2$ such that
\[
 C_d=\begin{pmatrix}
 ie&0&0\\0&-ie/2&z\\0&-\bar z&-ie/2
 \end{pmatrix}.
\]
Set $\gamma=-2iz/e$. Then $|\gamma|=1$. For either square root
$\zeta^2=-\bar\gamma$, define
\[
 W(\zeta)=\frac{\zeta}{\sqrt2}
 \begin{pmatrix}\gamma&\gamma\\1&-1\end{pmatrix}
 =\frac1{\sqrt2}
 \begin{pmatrix}-\zeta^{-1}&-\zeta^{-1}\\\zeta&-\zeta\end{pmatrix},
 \qquad B(\zeta)=U\operatorname{diag}(1,W(\zeta)).
\]
The simultaneous change and its inverse are
\begin{equation}
 Y'=B(\zeta)^*Y_0B(\zeta),\qquad
 Y_0=B(\zeta)Y'B(\zeta)^*.
 \label{sgd:literal-addon:frame}
\end{equation}
They give
\[
 C'=ie\operatorname{diag}(1,0,-1),\qquad
 H'=\begin{pmatrix}0&0&0\\0&0&r\\0&r&0\end{pmatrix},\qquad
 q'=q,\qquad F'=-2H'+\operatorname{diag}(-2,1,1).
\]
These maps preserve the spatial class and intertwine all three full
residuals in \eqref{sgd:literal-addon:H}--\eqref{sgd:literal-addon:G}.
\end{lemma}

\begin{proof}
Block multiplication gives the orthogonal direct sum
$\mathfrak{su}(3)=\mathfrak k\oplus\mathfrak p$, together with
$[\mathfrak p,\mathfrak p]\subset\mathfrak k$ and
$[\mathfrak k,\mathfrak p]\subset\mathfrak p$.
Since $H_0,qq^*$ are block diagonal, $\mathcal I_{H_0}$ preserves
both summands. More explicitly, $H_0X(v)H_0=0$; the two $H_0^2$
products sum to $(3/8)X(v)$, and the two $qq^*$ products sum
to $X(v)/8$. Each term is traceless. Thus
\[
 \mathcal I_{H_0}|_{\mathfrak p}
 =\tfrac12\operatorname{id}_{\mathfrak p}.
\]
The two bracket inclusions make $\mathcal K=\mathcal K_{Z_{1,0},Z_{2,0}}$
preserve both summands as well. Cyclicity gives, for every
$Y\in\mathfrak{su}(3)$,
\[
 \langle Y,\mathcal KY\rangle
 =-\sum_j\operatorname{tr}([Z_{j,0},Y]^2)
 =\sum_j\|[Z_{j,0},Y]\|^2\ge0.
\]
The $\mathfrak p$ projection of Gauss is therefore
\[
 (\tfrac t2\operatorname{id}_{\mathfrak p}
       +2\mathcal K|_{\mathfrak p})(C_0)_{\mathfrak p}=0.
\]
Pairing with $(C_0)_{\mathfrak p}$ gives nonnegative summands,
including $(t/2)\|(C_0)_{\mathfrak p}\|^2$. Since $t>0$,
$(C_0)_{\mathfrak p}=0$, proving $C_0\in\mathfrak k$.

The arbitrary first-frame temporal block consequently has the form
\[
 C_d=\begin{pmatrix}
 ia&0&0\\0&i(h-a)/2&z\\0&-\bar z&-i(a+h)/2
 \end{pmatrix},\qquad a,h\in\mathbb R,\quad z\in\mathbb C.
\]
Write the spatial matrices in this first frame as
$X(\widetilde v),X(\widetilde w)$, and define
$\widetilde M=\widetilde v\widetilde v^*
                +\widetilde w\widetilde w^*$.
The whole vector residual is
\[
 (-a^2+\operatorname{tr}\widetilde M+2)q,
\]
and the whole adjoint residual is
\[
 \begin{pmatrix}
 2-2r(\widetilde M_{11}-\widetilde M_{22})&0&0\\
 0&-4r|z|^2+2r\widetilde M_{11}-1+2r&-2irhz\\
 0&2irh\bar z&4r|z|^2-2r\widetilde M_{22}-1-2r
 \end{pmatrix}.
\]
For completeness, if $D=\operatorname{diag}(r,-r)$, the lower
spatial contribution is $\widetilde MD+D\widetilde M$, whose
diagonal entries are $2r\widetilde M_{11},-2r\widetilde M_{22}$
and whose off-diagonal entries vanish. The temporal double
commutator has diagonal entries $-4r|z|^2,4r|z|^2$ and
off-diagonal entries $-2irhz,2irh\bar z$.
Subtracting the lower force $I_2-2D$ gives the displayed lower block.
The upper spatial contribution and force give
$2-2r(\widetilde M_{11}-\widetilde M_{22})$.
Finally $X(v)^2=\operatorname{diag}(-v^*v,-vv^*)$ gives the
displayed vector residual, with all its other components zero.

Thus both scalar equations are equivalent to
\[
 a^2=2+\operatorname{tr}\widetilde M,\qquad
 2|z|^2-\widetilde M_{11}=1-\delta,\qquad
 2|z|^2-\widetilde M_{22}=1+\delta,\qquad hz=0.
\]
The upper adjoint entry is minus the lower trace.
Adding the middle equations gives
$4|z|^2=2+\operatorname{tr}\widetilde M=a^2$.
Since $\widetilde M_{22}\ge0$ and $1+\delta>0$, $z\ne0$.
Hence $h=0$ and $a\ne0$. This proves the claimed temporal block
with $e=a$, where its entries recover $e,z$ uniquely.

Now $|\gamma|=1$ and each root $\zeta^2=-\bar\gamma$ has
$|\zeta|=1$, with $\gamma=-\zeta^{-2}$. The displayed entries give
$W^*W=I_2$ and $\det W=-\zeta^2\gamma=1$.
The lower Hermitian block of $-iC_d$ is
\[
 \frac e2\begin{pmatrix}-1&\gamma\\\bar\gamma&-1\end{pmatrix}.
\]
Its products with $(\gamma,1)^{\mathsf T}$ and
$(\gamma,-1)^{\mathsf T}$ have eigenvalues $0,-e$,
respectively. The same columns give
$W^*DW=\left(\begin{smallmatrix}0&r\\r&0\end{smallmatrix}\right)$.
This proves the stated classified matrices.

The block-diagonal unitary $B(\zeta)$ fixes $q$ and commutes with
$P$. Its conjugation therefore preserves $\mathfrak p$ and
$\mathfrak k$. Unitarity gives
$B^*(XY)B=(B^*XB)(B^*YB)$; adjoints, commutators and trace
projection satisfy the corresponding equalities.
It follows that $\mathcal I,\mathcal K$ and every full matrix
residual are intertwined by conjugation. The vector residual is
multiplied by $B^*$. Equation~\eqref{sgd:literal-addon:frame} gives
the exact inverse on every datum, with no rescaling.
\end{proof}

\paragraph{The full parameter map and the quotient bijection}

We restate the classified parameters so that the literal map is explicit.
Put
\[
 Q=\operatorname{diag}(1,2),\qquad E=\operatorname{diag}(2/3,-5/6),
\]
\[
 \mathcal D=\{(u,v_0,y)\in\mathbb R^3:
 u>0,\ d=uv_0-\delta^2-y^2>0,\ c=4v_0-5u>0\}.
\]
For $(u,v_0,y)\in\mathcal D$, independent
$\epsilon,\sigma\in\{1,-1\}$, and $p_+,p_-\in U(1)$, define
\begin{align}
 m&=\delta+iy,&
 M&=\begin{pmatrix}u&m\\\bar m&v_0\end{pmatrix},&
 t&=\frac{72d}{c},\notag\\
 T&=\frac t4E-\frac12(QM+MQ),&
 k&=\sigma\sqrt{\frac d{-\det T}},&
 e&=\epsilon\sqrt{2+u+v_0},\notag\\
 \omega&=e/k,&N&=kT,&A_\pm&=M\pm N.
 \label{sgd:literal-addon:parameters}
\end{align}
Theorem~\ref{sgd:parametrization} proves that $\det T<0$,
that $A_\pm$ are positive semidefinite of rank one, and that
$h_\pm=(A_\pm)_{11}>0$. Thus the following vectors are defined
on the whole displayed domain:
\begin{equation}
 a_\pm=\begin{pmatrix}
 \sqrt{h_\pm}\\(A_\pm)_{21}/\sqrt{h_\pm}
 \end{pmatrix},\qquad
 v=\frac{p_+a_++p_-a_-}{2},\qquad
 w=\frac{i(p_+a_+-p_-a_-)}{2}.
 \label{sgd:literal-addon:vectors}
\end{equation}
In particular $vv^*+ww^*=M$ and $i(vw^*-wv^*)=N$.
The cited theorem proves that these parameters are in bijection
with all solutions of the two full scalar equations and full Gauss
in the frame $C'=ie\operatorname{diag}(1,0,-1)$,
$H'_{23}=H'_{32}=r$, $q'=q$, and
$F'=-2H'+\operatorname{diag}(-2,1,1)$.
It also proves the exact inverse phases
\begin{equation}
 p_+=\frac{(v-iw)_1}{\sqrt{h_+}},\qquad
 p_-=\frac{(v+iw)_1}{\sqrt{h_-}}.
 \label{sgd:literal-addon:phase-inverse}
\end{equation}

Define the full raw parameter space
\[
 \mathcal P_{\rm lit}
 =\mathcal D\times\{1,-1\}^2\times U(1)^2\times U(1)
\]
with coordinates $(u,v_0,y,\epsilon,\sigma,p_+,p_-,\zeta)$.
For the last unit phase set $\gamma=-\zeta^{-2}$ and use the
matrix $B(\zeta)$ of Lemma~\ref{sgd:literal-addon:temporal}.
Define the map $\mathfrak F:\mathcal P_{\rm lit}\to\mathcal S_{\rm sg}$
by \eqref{sgd:literal-addon:parameters}--\eqref{sgd:literal-addon:vectors}
and the following entire literal tuple:
\begin{equation}
 \begin{gathered}
 \mathfrak F(u,v_0,y,\epsilon,\sigma,p_+,p_-,\zeta)
 =(t,\omega,C_0,Z_{1,0},Z_{2,0}),\\
 \begin{aligned}
 C_0&=B(\zeta)\bigl(ie\operatorname{diag}(1,0,-1)\bigr)B(\zeta)^*,\\
 Z_{1,0}&=B(\zeta)X(v)B(\zeta)^*,\\
 Z_{2,0}&=B(\zeta)X(w)B(\zeta)^*.
 \end{aligned}
 \end{gathered}
 \label{sgd:literal-addon:forward}
\end{equation}
The codomain assertion follows from the classified theorem and the
exact residual intertwining in Lemma~\ref{sgd:literal-addon:temporal}.

\begin{theorem}[Exact quotient parameterization of the entire literal domain]
\label{sgd:literal-addon:quotient}
The map $\mathfrak F$ is onto $\mathcal S_{\rm sg}$.
The map
\begin{equation}
 \iota(u,v_0,y,\epsilon,\sigma,p_+,p_-,\zeta)
 =(u,v_0,y,\epsilon,\sigma,-p_+,-p_-,-\zeta)
 \label{sgd:literal-addon:involution}
\end{equation}
is a free involution of $\mathcal P_{\rm lit}$.
Each fibre of $\mathfrak F$ is exactly one two-element
$\iota$-orbit. Consequently \eqref{sgd:literal-addon:forward}
induces a bijection of sets
\begin{equation}
 \overline{\mathfrak F}:
 \mathcal P_{\rm lit}/\langle\iota\rangle
 \ \longrightarrow\ \mathcal S_{\rm sg}.
 \label{sgd:literal-addon:bijection}
\end{equation}
The inverse of this bijection is recovered from the literal matrices
by the formulas in the proof, with both square-root representatives
retained.
\end{theorem}

\begin{proof}
We first construct a preimage of every literal tuple, recording all
inverse coordinates. Given
$(t,\omega,C_0,Z_{1,0},Z_{2,0})\in\mathcal S_{\rm sg}$, compute
\[
 C_d=U^*C_0U,\qquad
 e=-i(C_d)_{11},\qquad z=(C_d)_{23},\qquad
 \gamma=-2iz/e.
\]
Lemma~\ref{sgd:literal-addon:temporal} proves $e\in\mathbb R\setminus\{0\}$
and $|\gamma|=1$. Choose either square root
$\zeta^2=-\bar\gamma$. Define
\[
 Z'_j=B(\zeta)^*Z_{j,0}B(\zeta),\qquad
 v=\begin{pmatrix}(Z'_1)_{21}\\(Z'_1)_{31}\end{pmatrix},\qquad
 w=\begin{pmatrix}(Z'_2)_{21}\\(Z'_2)_{31}\end{pmatrix}.
\]
The lemma proves that these are classified solutions with the original
$t,\omega,e$. Compute
\[
 M=vv^*+ww^*,\qquad N=i(vw^*-wv^*),\qquad
 u=M_{11},\quad v_0=M_{22},\quad y=\operatorname{Im}M_{12},
\]
\[
 \epsilon=\operatorname{sign}e,\qquad
 \sigma=\operatorname{sign}(e/\omega).
\]
Theorem~\ref{sgd:parametrization} proves
$(u,v_0,y)\in\mathcal D$, $M_{12}=\delta+iy$, and all equations
in \eqref{sgd:literal-addon:parameters} with this unchanged
$t,\omega,e$. It proves $h_\pm=(M\pm N)_{11}>0$, so
\eqref{sgd:literal-addon:phase-inverse} recovers unique unit
$p_+,p_-$. The vector formulas recover the displayed $v,w$
and the inverse frame map then recovers all three literal matrices.
Thus \eqref{sgd:literal-addon:forward} sends the recovered raw
parameter to the original tuple. This proves surjectivity.

Next we prove that \eqref{sgd:literal-addon:involution} preserves
each fibre. It preserves the parameter domain, and applying it
twice is the identity. It has no fixed point because
$\zeta=-\zeta$ would imply $\zeta=0$, contrary to $|\zeta|=1$.
Replacing $\zeta$ by $-\zeta$ preserves $\gamma=-\zeta^{-2}$ and
gives $W(-\zeta)=-W(\zeta)$. Put
$D_0=\operatorname{diag}(1,-I_2)$. Then
 \begin{align*}
 B(-\zeta)&=B(\zeta)D_0,\\
 D_0\bigl(ie\operatorname{diag}(1,0,-1)\bigr)D_0
 &=ie\operatorname{diag}(1,0,-1),\\
 D_0X(v)D_0&=X(-v).
 \end{align*}
Replacing $(p_+,p_-)$ by $(-p_+,-p_-)$ replaces $(v,w)$ by
$(-v,-w)$ in \eqref{sgd:literal-addon:vectors}, while leaving
$M,N,t,e,k,\omega$ fixed. The preceding matrix equalities
therefore give $\mathfrak F\circ\iota=\mathfrak F$ on the entire
raw parameter domain. This proves that $\overline{\mathfrak F}$
is well defined.

To determine every fibre, suppose two raw parameters have the same
literal image. The common temporal matrix $C_d$ recovers the same
$e,z,\gamma$ by the inverse formulas above. Their frame roots
therefore satisfy either $\zeta_2=\zeta_1$ or
$\zeta_2=-\zeta_1$, since they square to the same nonzero complex
number. If $\zeta_2=\zeta_1$, the classified temporal and spatial
matrices obtained by that common inverse frame map coincide.
The classified bijection of Theorem~\ref{sgd:parametrization}
then forces equality of all other raw parameters.

If $\zeta_2=-\zeta_1$, the identity
$B(-\zeta)=B(\zeta)D_0$ implies
\[
 Z'_{j,2}=D_0Z'_{j,1}D_0,\qquad
 v_2=-v_1,\qquad w_2=-w_1.
\]
Hence $M_2=M_1$ and $N_2=N_1$, so the inverse coordinates
$u,v_0,y$ are equal. The original $e,\omega$ are equal, which
also gives equal $\epsilon,\sigma$. Consequently $h_\pm$ agree,
and \eqref{sgd:literal-addon:phase-inverse} gives
\[
 p_{+,2}=-p_{+,1},\qquad p_{-,2}=-p_{-,1}.
\]
Thus the two raw parameters differ by exactly $\iota$.
No further case is possible. Each fibre is therefore precisely
one two-element orbit. This proves injectivity of
$\overline{\mathfrak F}$ as well as its already-proved surjectivity.

Finally the inverse formulas starting from a literal tuple give
the same orbit for either root choice, by the second fibre case.
Starting from a raw parameter, applying $\mathfrak F$ and then those
inverse formulas gives its orbit, by the two fibre cases.
Starting from a literal tuple, applying the inverse formulas and
then $\overline{\mathfrak F}$ gives that literal tuple by the first
paragraph. These prove both compositions are identities on their
stated domains, completing the quotient bijection.
\end{proof}

The construction retains
$\operatorname{sign}e=\epsilon$ and
$\operatorname{sign}\omega=\epsilon\sigma$.
The involution changes neither sign. The two independent vector phases
and the frame phase remain coordinates, subject only to the proved
simultaneous-sign identification.

For explicit scope, the additional full spatial residuals would be
\begin{align*}
 R_1={}&-\omega^2Z_{1,0}-2\omega[C_0,Z_{2,0}]
 +[C_0,[C_0,Z_{1,0}]]
 +\mathcal K_{Z_{1,0},Z_{2,0}}Z_{1,0}
 +(t/2)\mathcal I_{H_0}Z_{1,0},\\
 R_2={}&-\omega^2Z_{2,0}+2\omega[C_0,Z_{1,0}]
 +[C_0,[C_0,Z_{2,0}]]
 +\mathcal K_{Z_{1,0},Z_{2,0}}Z_{2,0}
 +(t/2)\mathcal I_{H_0}Z_{2,0}.
\end{align*}
Neither $R_1=0$ nor $R_2=0$ is part of
$\mathcal S_{\rm sg}$ or used in either proof above.
The bijection therefore classifies exactly the two scalar equations
and Gauss; it does not assert a solution of these two additional
spatial equations.


\subsubsection{A short exact family with every equation checked}

\begin{proposition}\label{sgd:explicit-point}
Choose arbitrary $p_+,p_-\in U(1)$ and independent
$\epsilon,\sigma\in\{1,-1\}$. Put
\[
 M=\delta\begin{pmatrix}1&1\\1&2\end{pmatrix},\quad
 t=24\delta,\quad
 e=\epsilon\sqrt{2+3\delta},\quad
 k=\frac{2\sigma}{3\sqrt{13}},\quad
 \omega=\frac{3\sigma\sqrt{13}}2e,
\]
\[
 N=\frac{\sigma\delta}{\sqrt{13}}
        \begin{pmatrix}2&-1\\-1&-6\end{pmatrix},\qquad
 a_\pm=\sqrt\delta\begin{pmatrix}
 \sqrt{1\pm2\sigma/\sqrt{13}}\\[1mm]
 (1\mp\sigma/\sqrt{13})/\sqrt{1\pm2\sigma/\sqrt{13}}
 \end{pmatrix}.
\]
Define $v,w$ by \eqref{sgd:vectors}. For any retained unit frame
phase $\zeta$, these give original tuples by
\[
 C_0=B\bigl(ie\operatorname{diag}(1,0,-1)\bigr)B^*,\qquad
 Z_{1,0}=BX(v)B^*,\qquad Z_{2,0}=BX(w)B^*.
\]
Every such tuple satisfies \eqref{sgd:literal-H}--\eqref{sgd:literal-G}
at the positive original coupling $g^2=24\delta$.
\end{proposition}
\begin{proof}
The first diagonal factors are positive because $2/\sqrt{13}<1$.
For a sign $s$, the exact identity
\[
 (1-s\sigma/\sqrt{13})^2
  =(1+2s\sigma/\sqrt{13})(2-6s\sigma/\sqrt{13})
\]
follows by multiplying out and using $\sigma^2=1$. Hence
$a_\pm a_\pm^*=M\pm N$. The construction therefore has exactly
the displayed $M,N$ for every pair of phases. In particular
\[
 u=\delta,\quad v_0=2\delta,\quad m=\delta,\quad
 \operatorname{tr} M=3\delta,\quad d=\delta^2,\quad
 b=\tfrac12\operatorname{tr} N=-\frac{2\sigma\delta}{\sqrt{13}}.
\]
The scalar identities are $e^2=2+3\delta$ and
$2r\delta=1$; substitution in \eqref{sgd:scalar-H} and
\eqref{sgd:scalar-q} makes every entry zero.
For the full Gauss block, direct multiplication gives
\[
 QM+MQ=\delta\begin{pmatrix}2&3\\3&8\end{pmatrix},\qquad
 tE-2(QM+MQ)
   =\delta\begin{pmatrix}12&-6\\-6&-36\end{pmatrix}.
\]
The same matrix is $(4\omega/e)N$, since
$\omega/e=3\sigma\sqrt{13}/2$. Thus $\mathcal G=0$ entry
by entry, and the upper entry is explicitly
\[
 \gamma_0=4e\delta+20e\delta
       +8\left(\frac{3\sigma\sqrt{13}}2e\right)
             \left(-\frac{2\sigma\delta}{\sqrt{13}}\right)=0.
\]
Finally \eqref{sgd:frame} proves all three literal equations exactly.
As an additional check on the construction, its determinant data are
\[
 T=\frac{3\delta}{2}\begin{pmatrix}2&-1\\-1&-6\end{pmatrix},\quad
 \det T=-\frac{117}{4}\delta^2,\quad k^2=\frac4{117},
\]
agreeing with the general parameterization. Both signs of the original
frequency are present because $\operatorname{sign}\omega=\epsilon\sigma$.
\end{proof}

\begin{corollary}[Every positive original coupling occurs]
\label{sgd:every-coupling}
For every specified original $\tau=g^2>0$, the literal two-cross
scalar--Gauss domain contains tuples at $t=\tau$, with either original
frequency sign and both temporal signs. It contains continuous phase
families at each such coupling.
\end{corollary}
\begin{proof}
For $\rho>1/\sqrt2$ choose
\[
 M_\rho=\delta\begin{pmatrix}\rho&1\\1&2\rho\end{pmatrix}.
\]
Its parameters obey $u=\delta\rho>0$,
$d=\delta^2(2\rho^2-1)>0$, $c=3\delta\rho>0$ and $m=\delta$.
Thus Theorem~\ref{sgd:parametrization} gives the complete solution
family at
\[
 t=\frac{72\delta^2(2\rho^2-1)}{3\delta\rho}
   =\frac{24\delta(2\rho^2-1)}\rho.
\]
For a given $\tau>0$, set
\[
 \rho=\frac{\tau+\sqrt{\tau^2+4608\delta^2}}{96\delta}.
\]
Because $\sqrt{4608}\delta=48\sqrt2\delta$, this is strictly
greater than $1/\sqrt2$. Squaring the displayed defining equality
gives $48\delta\rho^2-\tau\rho-24\delta=0$, which is precisely
$t=\tau$. No limiting passage or change of coupling is used.
The independent choices of $\epsilon,\sigma,p_+,p_-$ supply the
stated signs and phases; the inverse in Theorem~\ref{sgd:parametrization}
ensures different phase pairs yield different classified tuples.
\end{proof}

\subsubsection{Exact scope and reproducibility}

The results concern the simultaneous zero set of both scalar equations
and Gauss. The construction never imposes either spatial equation and
does not claim a full rotating-ansatz solution. It proves that the
previously specified scalar--Gauss domain of the two-cross quantitative
residual certificate is nonempty, in fact at every original positive
coupling. This is an algebraic statement about the retained classical
residuals and makes no quantum assertion.

The companion \texttt{check\_scalar\_gauss\_domain.py} is a finite exact
symbolic replay. It checks the full scalar and Gauss matrix formulas with
eight independent real spatial coordinates, the determinant/trace
identities, the explicit factors and both signs, and the original
unitary frame maps. The rank, positivity, parameter domains, completeness
and inverse-map proofs are given in this document; they are not inferred
from sampled values or from a numerical residual tolerance.

The exact full literal parameter map and its inverse are proved in
Theorem~\ref{sgd:literal-addon:quotient}; its only identification changes
the signs of the frame phase and both vector phases together. Combining
Corollary~\ref{sgd:every-coupling} with
Proposition~\ref{q444:twocross:quantitative} makes the stated spatial-error
bound non-vacuous at every positive original coupling. This does not alter
the full-solution exclusion in Theorem~\ref{q444:twocross:thm:exclusion}.

The complete source versions, the independently checked literal-domain
addendum and the exact symbolic replay are retained under
\path{supporting_materials/workbench/evidence/scalar_gauss_domain_20260905/}.
The retained equations are those derived from the parent model in
\eqref{q444:cx:eq:phi}--\eqref{q444:cx:eq:sp2}; this proof neither changes
that model nor replaces its two spatial equations by the Gauss equation.

% Complete smooth literal-domain refinement; exact staged source and review
% are retained in evidence/s6_global_smooth_reaudit_20260905/.
\subsection{Smooth literal inverse charts and all positive-coupling fibres}
\label{sgd:smooth-addon:section}

The exact algebraic bijections of
Theorems~\ref{sgd:parametrization} and
\ref{sgd:literal-addon:quotient} admit a smooth refinement.
We prove the parameter domain's convexity, give the phase quotient
its explicit smooth coordinates, and construct smooth inverse maps
from the literal matrix entries. We then give a global product
description in which the first continuous coordinate is the unchanged
original coupling $t=g^2$. The topology of every fixed-coupling fibre
will follow from explicit inverse formulas.

\subsubsection{Retained objects and the exact parameter domain}

Keep the literal matrices and constants
\[
 H_0=\begin{pmatrix}0&0&0\\0&-1/2&1/(2\sqrt2)\\
 0&1/(2\sqrt2)&1/2\end{pmatrix},\quad
 q=\frac1{2\sqrt2}\begin{pmatrix}1\\0\\0\end{pmatrix},\quad
 A_0=\sqrt2\begin{pmatrix}0&-1&0\\1&0&0\\0&0&0\end{pmatrix},
\]
\[
 m^2=\Gamma=1,\qquad \kappa=6,\qquad
 F_0=[A_0,[A_0,H_0]]=-2H_0+\operatorname{diag}(-2,1,1),
\]
and set
\[
 r=\frac{\sqrt6}{4},\quad\delta=\frac1{2r}=\sqrt{\frac23},
 \quad x=\frac12,\quad d_0=\frac1{2\sqrt2},\quad
 U=\operatorname{diag}\left(
 1,\frac1{\sqrt{2r(r+x)}}
 \begin{pmatrix}d_0&-(r+x)\\r+x&d_0\end{pmatrix}\right).
\]
Thus $U^*U=I_3$, $\det U=1$, $Uq=q$, and
$U^*H_0U=\operatorname{diag}(0,r,-r)$, as proved in
Lemma~\ref{sgd:literal-addon:temporal}.
Write
\[
 P=\operatorname{diag}(1,0,0),\qquad
 X(v)=\begin{pmatrix}0&-v^*\\v&0_2\end{pmatrix},\qquad
 \mathfrak p=X(\mathbb C^2)\subset\mathfrak{su}(3).
\]
The original real ambient manifold for this selected spatial class is
\[
 \mathcal E=(0,\infty)\times(\mathbb R\setminus\{0\})
             \times\mathfrak{su}(3)\times\mathfrak p\times\mathfrak p.
\]
It has real dimension $1+1+8+4+4=18$.
For $Y_{\rm tf}=Y-(\operatorname{tr}Y)I_3/3$, put
\[
 \mathcal I_{H_0}Y=([H_0,[H_0,Y]]+qq^*Y+Yqq^*)_{\rm tf},
 \qquad \mathcal K_{Z_1,Z_2}Y=-\sum_{j=1}^2[Z_j,[Z_j,Y]].
\]
The subset $\mathcal S_{\rm sg}\subset\mathcal E$ is exactly the
simultaneous zero set
\begin{align*}
 [C_0,[C_0,H_0]]
     -\sum_j[Z_{j,0},[Z_{j,0},H_0]]-F_0&=0,\\
 (C_0^2-Z_{1,0}^2-Z_{2,0}^2)q+2q&=0,\\
 (t\mathcal I_{H_0}+2\mathcal K_{Z_{1,0},Z_{2,0}})C_0
                   -4\omega[Z_{1,0},Z_{2,0}]&=0.
\end{align*}
Neither additional spatial residual is included in this definition.
The topology of $\mathcal S_{\rm sg}$ initially means its subspace
topology in $\mathcal E$; its embedded smooth structure will be proved.

For real coordinates $(u,v_0,y)$ define
\[
 h=\delta^2+y^2,\qquad m=\delta+iy,\qquad
 M=\begin{pmatrix}u&m\\\bar m&v_0\end{pmatrix},\qquad
 d=uv_0-h,\qquad c=4v_0-5u,
\]
\[
 \mathcal D=\{(u,v_0,y)\in\mathbb R^3:u>0,\ d>0,\ c>0\}.
\]
Here $h$ is a scalar, distinct from the two factor entries
$h_\pm$ used below.

\begin{lemma}[Convexity of the entire continuous parameter domain]
\label{sgd:smooth-addon:convex}
The set $\mathcal D$ is a nonempty open convex subset of
$\mathbb R^3$. In particular it is connected and has real dimension three.
\end{lemma}
\begin{proof}
The inequalities are strict inequalities for continuous functions,
so $\mathcal D$ is open. For $u>0$ and $(a,b)\in\mathbb C^2$,
the exact quadratic-form identity is
\[
 \begin{pmatrix}\bar a&\bar b\end{pmatrix}
 M\begin{pmatrix}a\\b\end{pmatrix}
 =u\left|a+\frac m u b\right|^2+\frac d u|b|^2.
\]
It proves $M>0$ when $u>0,d>0$. Conversely, $M>0$ gives
$u>0$ by the first coordinate vector, and evaluation at
$(a,b)=(-m/u,1)$ gives $d/u>0$.
Thus the conditions $u>0,d>0$ are exactly positive definiteness
of the Hermitian matrix in the affine slice $\operatorname{Re}M_{12}=\delta$.

If $M_0,M_1$ are positive definite and $0\le\lambda\le1$,
then for every nonzero $z\in\mathbb C^2$,
\[
 z^*((1-\lambda)M_0+\lambda M_1)z
 =(1-\lambda)z^*M_0z+\lambda z^*M_1z>0.
\]
The affine slice is preserved by this combination, and the additional
linear inequality $4v_0-5u>0$ is preserved as well.
The coordinate map $(u,v_0,y)\mapsto M$ is affine and injective,
so these facts prove convexity in the original real coordinates.
Finally $(u,v_0,y)=(\delta,2\delta,0)$ has
$u=\delta>0$, $d=\delta^2>0$, and $c=3\delta>0$.
It supplies a point and, by openness, a three-dimensional neighbourhood.
Straight segments inside a convex set join any two points,
proving connectedness.
\end{proof}

\subsubsection{The phase quotient with every original phase retained}

Write $\mathbb T=U(1)=\{z\in\mathbb C:|z|=1\}$.
On the original three phases, in the order $(\zeta,p_+,p_-)$,
define
\begin{equation}
 \Psi:\mathbb T^3\longrightarrow\mathbb T^3,\qquad
 \Psi(\zeta,p_+,p_-)
       =\left(\zeta^2,\frac{p_+}{\zeta},\frac{p_-}{\zeta}\right)
       =(\eta,\alpha,\beta).
 \label{sgd:smooth-addon:phase-map}
\end{equation}

\begin{lemma}[Exact smooth phase quotient]
\label{sgd:smooth-addon:phase-quotient}
The map $\Psi$ is a surjective smooth group homomorphism with kernel
$\{(1,1,1),(-1,-1,-1)\}$.
It induces a diffeomorphism
\[
 \mathbb T^3/\{(1,1,1),(-1,-1,-1)\}\ \longrightarrow\ \mathbb T^3.
\]
The quotient smooth charts have explicit local inverses retaining
both original square-root representatives.
\end{lemma}
\begin{proof}
Products, inverses on $\mathbb T$, and the square map are smooth;
the displayed formulas commute with coordinatewise multiplication.
An element is in the kernel precisely when $\zeta^2=1$ and
$p_+=p_-=\zeta$, giving exactly the two listed elements.
Given any $(\eta,\alpha,\beta)\in\mathbb T^3$, choose either root
$\zeta^2=\eta$ and set $p_+=\zeta\alpha$, $p_-=\zeta\beta$.
This gives all preimages, so $\Psi$ is onto and each fibre is
exactly one simultaneous-sign pair.

Here are the smooth local charts, rather than only a set quotient.
Fix $\eta_0\in\mathbb T$ and one root $\zeta_0^2=\eta_0$.
On the open arc $\eta=\eta_0e^{i\theta}$ with $-\pi<\theta<\pi$,
the two lifts are exactly
\[
 \zeta_\pm=\pm\zeta_0e^{i\theta/2},\qquad
 p_{+,\pm}=\zeta_\pm\alpha,\qquad
 p_{-,\pm}=\zeta_\pm\beta.
\]
Each formula is smooth. The two root arcs are disjoint and open
in the preimage, and each maps diffeomorphically to the displayed
arc times $\mathbb T^2$, using \eqref{sgd:smooth-addon:phase-map}
as its inverse. Their simultaneous-sign interchange identifies
exactly these two charts.
The quotient topology agrees with these charts: the quotient
projection is open because the inverse image of the image of any
open set is that set united with its simultaneous-sign translate.
On overlapping arcs the quotient coordinates are the same
$(\eta,\alpha,\beta)$, so their transitions are the ordinary smooth
torus transitions. These charts prove the asserted diffeomorphism.
They do not select a global square root or discard either representative.
\end{proof}

\subsubsection{Global formulas for literal matrices and their smooth inverses}

We record all the algebraic parameters used by the smooth map.
For $(u,v_0,y)\in\mathcal D$, set
\[
 Q=\operatorname{diag}(1,2),\qquad E=\operatorname{diag}(2/3,-5/6),
 \qquad t=\frac{72d}{c},\qquad
 T=\frac t4E-\frac12(QM+MQ).
\]
For independent signs $\epsilon,\sigma\in\{1,-1\}$ put
\[
 e=\epsilon\sqrt{2+u+v_0},\qquad
 k=\sigma\sqrt{\frac d{-\det T}},\qquad
 \omega=e/k,\qquad N=kT,\qquad A_\pm=M\pm N.
\]
Theorem~\ref{sgd:parametrization} proves on this entire domain that
$-\det T>0$, that $A_\pm$ are positive semidefinite of rank one,
and that $h_\pm=(A_\pm)_{11}>0$. Consequently
\[
 a_\pm=\begin{pmatrix}\sqrt{h_\pm}\\
                 (A_\pm)_{21}/\sqrt{h_\pm}\end{pmatrix}
\]
are defined everywhere in $\mathcal D$ for both signs.

For $(\eta,\alpha,\beta)\in\mathbb T^3$ put
\[
v^\flat=\frac{\alpha a_++\beta a_-}{2},\qquad
w^\flat=\frac{i(\alpha a_+-\beta a_-)}{2}.
\]
\[
L(\eta)=\frac1{\sqrt2}\begin{pmatrix}-1&-1\\\eta&-\eta\end{pmatrix},
 \qquad B^\flat(\eta)=U\operatorname{diag}(1,L(\eta)).
\]
The superscript $\flat$ names these exact intermediate coordinates;
it does not denote a metric dual or a change of magnitude.
Multiplication gives
\[
 L(\eta)^*L(\eta)=I_2,\qquad
 \det L(\eta)=\eta,\qquad \det B^\flat(\eta)=\eta.
\]
In particular these auxiliary matrices are unitary with their
displayed determinant; they are not asserted to have determinant one.
The precise relation to the original determinant-one frame is,
for either root $\zeta^2=\eta$,
\[
 \gamma=-\zeta^{-2},\qquad
 W(\zeta)=\frac1{\sqrt2}
 \begin{pmatrix}-\zeta^{-1}&-\zeta^{-1}\\\zeta&-\zeta\end{pmatrix},
 \qquad L(\eta)=\zeta W(\zeta).
\]
If $p_+=\zeta\alpha$, $p_-=\zeta\beta$, the original vectors satisfy
$v=\zeta v^\flat$ and $w=\zeta w^\flat$.
Thus every phase and every determinant factor is retained.

Let
\[
 \mathcal X=\{1,-1\}^2\times\mathcal D\times\mathbb T^3
\]
have coordinates $(\epsilon,\sigma,u,v_0,y,\eta,\alpha,\beta)$.
Define $\mathcal F^\flat:\mathcal X\to\mathcal E$ by
\begin{equation}
 \begin{gathered}
 \mathcal F^\flat(\epsilon,\sigma,u,v_0,y,\eta,\alpha,\beta)
       =(t,\omega,C_0,Z_{1,0},Z_{2,0}),\\
 C_0=B^\flat(\eta)\bigl(ie\operatorname{diag}(1,0,-1)\bigr)
                                    B^\flat(\eta)^*,\\
 \begin{aligned}
 Z_{1,0}&=B^\flat(\eta)X(v^\flat)B^\flat(\eta)^*,\\
 Z_{2,0}&=B^\flat(\eta)X(w^\flat)B^\flat(\eta)^*.
 \end{aligned}
 \end{gathered}
 \label{sgd:smooth-addon:literal-forward}
\end{equation}

\begin{theorem}[Four smooth connected components in literal coordinates]
\label{sgd:smooth-addon:literal-diffeomorphism}
The image of $\mathcal F^\flat$ is exactly $\mathcal S_{\rm sg}$,
and $\mathcal F^\flat$ is a diffeomorphism onto an embedded
six-real-dimensional submanifold of $\mathcal E$.
The submanifold has exactly four connected components, indexed by
\[
 \epsilon=\operatorname{sign}e,\qquad
 \sigma=\operatorname{sign}(e/\omega).
\]
Each component is diffeomorphic to $\mathcal D\times\mathbb T^3$.
The smooth structure and inverse charts refer to the literal
matrix subset; no quotient of literal configurations by a physical
gauge action is being taken.
\end{theorem}
\begin{proof}
First verify exactly which literal matrices the displayed formulas give.
Let $D_\zeta=\operatorname{diag}(1,\zeta I_2)$ and let
$B(\zeta)=U\operatorname{diag}(1,W(\zeta))$ be the original frame.
Then
\begin{align*}
 B^\flat(\eta)&=B(\zeta)D_\zeta,\\
 D_\zeta\bigl(ie\operatorname{diag}(1,0,-1)\bigr)D_\zeta^*
       &=ie\operatorname{diag}(1,0,-1).
\end{align*}
\[
 D_\zeta X(v^\flat)D_\zeta^*=X(\zeta v^\flat)=X(v),
 \qquad D_\zeta X(w^\flat)D_\zeta^*=X(w).
\]
Thus \eqref{sgd:smooth-addon:literal-forward} gives exactly the same
literal tuple as the original map for each of its two phase lifts.
Lemma~\ref{sgd:smooth-addon:phase-quotient} and the already-proved
literal quotient bijection, Theorem~\ref{sgd:literal-addon:quotient},
therefore show that its image is $\mathcal S_{\rm sg}$ and that
$\mathcal F^\flat$ is bijective.

Here is its inverse directly in literal matrix entries.
Given a tuple in $\mathcal S_{\rm sg}$, compute
\begin{equation}
 C_d=U^*C_0U,\qquad
 e=-i(C_d)_{11},\qquad z=(C_d)_{23},\qquad
 \eta=-\frac{2i\bar z}{e}.
 \label{sgd:smooth-addon:literal-temporal-inverse}
\end{equation}
The temporal classification proves that $e$ is real and nonzero,
$|z|=|e|/2$, and hence $|\eta|=1$.
The sign is checked by multiplying the lower temporal block in
\eqref{sgd:smooth-addon:literal-forward}:
\[
 (C_d)_{\rm lower}
 =-\frac{ie}{2}\begin{pmatrix}1&\bar\eta\\\eta&1\end{pmatrix}.
\]
In particular $z=-ie\bar\eta/2$, so
$-2i\bar z/e=\eta$ exactly.

Define the two matrices and their entire lower first columns by
\[
 Z_j^\flat=B^\flat(\eta)^*Z_{j,0}B^\flat(\eta),\qquad
 v^\flat=\begin{pmatrix}(Z_1^\flat)_{21}\\(Z_1^\flat)_{31}\end{pmatrix},
 \qquad
 w^\flat=\begin{pmatrix}(Z_2^\flat)_{21}\\(Z_2^\flat)_{31}\end{pmatrix}.
\]
Then recover
\[
 M=v^\flat(v^\flat)^*+w^\flat(w^\flat)^*,\qquad
 N=i\bigl(v^\flat(w^\flat)^*-w^\flat(v^\flat)^*\bigr),
\]
\begin{equation}
 \begin{gathered}
 u=M_{11},\quad v_0=M_{22},\quad y=\operatorname{Im}M_{12},
 \qquad \epsilon=\operatorname{sign}e,\quad
 \sigma=\operatorname{sign}(e/\omega),\\
 h_\pm=(M\pm N)_{11},\qquad
 \alpha=\frac{(v^\flat-iw^\flat)_1}{\sqrt{h_+}},\qquad
 \beta=\frac{(v^\flat+iw^\flat)_1}{\sqrt{h_-}}.
 \end{gathered}
 \label{sgd:smooth-addon:literal-spatial-inverse}
\end{equation}
Indeed multiplication by the common unit $\zeta$ preserves both
$M$ and $N$. The original phase inverse therefore gives
$\alpha=p_+/\zeta$, $\beta=p_-/\zeta$ in these formulas, with
$h_\pm>0$, $|\alpha|=|\beta|=1$.
The algebraic inverse proves that all recovered coordinates lie in
their exact domains. The displayed formulas recover the original
parameter from \eqref{sgd:smooth-addon:literal-forward}, and the
forward map recovers the original literal tuple from them.

All forward functions are real smooth on each sign sector:
$c>0$, $d>0$, $-\det T>0$, $2+u+v_0>0$, and $h_\pm>0$
are precisely the denominators and square-root arguments appearing
above, while $k\ne0$. Matrix products and conjugates are real smooth.
In the inverse, $e\ne0$ and $h_\pm>0$ ensure smoothness of every
non-discrete formula. The signs are locally constant.
The inverse formulas in particular give a continuous inverse for
the subspace topology on $\mathcal S_{\rm sg}$.

We check embeddedness, rather than just transferring an abstract
smooth structure through a bijection.
At any one image point choose open torus arcs around
$\eta,\alpha,\beta$, with their real angle coordinates.
The inverse formulas extend as real smooth functions to an open
neighbourhood in $\mathcal E$ as follows. Keep $e\ne0$ and the
computed $h_\pm>0$ by taking a sufficiently small neighbourhood.
The same matrix formulas define the complex quantities
$\eta,\alpha,\beta$ there, even when the tuple is off the zero set.
For a unit base point $\eta_0$, on the complex neighbourhood
$\operatorname{Re}(\eta/\eta_0)>0$, its local angle is the smooth
real function
\[
 \arctan\!\left(
 \frac{\operatorname{Im}(\eta/\eta_0)}
      {\operatorname{Re}(\eta/\eta_0)}\right).
\]
Use the corresponding expressions about $\alpha_0,\beta_0$.
On the zero set these are exactly its three local torus coordinates.
Together with $u,v_0,y$, they give a real smooth six-coordinate
map from an ambient neighbourhood whose composition with
$\mathcal F^\flat$ on the parameter chart is the identity.
The derivative of that composition is the identity, so the
derivative of $\mathcal F^\flat$ has rank six at every point.

For clarity, the real inverse function theorem
\cite[\S1, pp.~4--5]{sgd:Milnor1965} now applies with
its full-rank hypothesis verified: choose six ambient linear
coordinate projections whose composite with this rank-six
parametrization has invertible derivative. Locally the image is a
smooth graph over those six coordinates.
The continuous inverse already proved prevents other parameter
patches from contributing additional image points arbitrarily near
this point outside its chart: the inverse image of the chart is
open in the subspace topology. Thus the graph is a neighbourhood
of the point in the whole subset $\mathcal S_{\rm sg}$.
This proves it is an embedded six-dimensional submanifold and that
the displayed inverse is smooth for its embedded structure.

It remains to count components exactly. By
Lemma~\ref{sgd:smooth-addon:convex}, $\mathcal D$ is path connected.
Each circle is path connected: for chosen real arguments of two
unit endpoints, interpolating those arguments gives an explicit
path on the circle. Products of these paths prove $\mathbb T^3$
is path connected. Hence each of the four sign sectors in
$\mathcal X$ is nonempty and connected.
On the literal domain the two sign functions in
\eqref{sgd:smooth-addon:literal-spatial-inverse} are continuous
with values in $\{1,-1\}$ because $e,\omega$ never vanish.
They cannot change on a connected subset, so distinct sectors
cannot belong to the same component. This proves there are
exactly four connected components.

Finally the only quotient used to construct the coordinates was
the two-element kernel of the parameter map.
Its two representatives give identical literal matrices, as the
first displayed equalities proved. Distinct literal tuples remain
distinct under the bijection. No further gauge equivalence or
spatial-pair equivalence has been imposed.
\end{proof}

\subsubsection{The original coupling is a surjective submersion}

On $\mathcal D$ the original coupling function is
\[
 t(u,v_0,y)=\frac{72(uv_0-\delta^2-y^2)}{4v_0-5u}
          =\frac{72d}{c}.
\]

\begin{theorem}[Exact differential and global coupling coordinates]
\label{sgd:smooth-addon:coupling}
The function $t:\mathcal D\to(0,\infty)$ is a smooth surjective
submersion. Its exact partial derivatives are
\begin{equation}
 \begin{gathered}
 \frac{\partial t}{\partial u}
 =\frac{72(4v_0^2-5h)}{c^2},\qquad
 \frac{\partial t}{\partial v_0}
 =\frac{72(4h-5u^2)}{c^2},\\
 \frac{\partial t}{\partial y}=-\frac{144y}{c},
 \qquad h=\delta^2+y^2.
 \end{gathered}
 \label{sgd:smooth-addon:derivatives}
\end{equation}
In fact $\partial t/\partial u>0$ throughout $\mathcal D$.

Define the fixed open set
\[
 \mathcal B=\{(v_0,y)\in\mathbb R^2:
        v_0>\tfrac{\sqrt5}{2}\sqrt{\delta^2+y^2}\}.
\]
The map
\begin{equation}
 \Theta:\mathcal D\longrightarrow(0,\infty)\times\mathcal B,
 \qquad \Theta(u,v_0,y)=(t(u,v_0,y),v_0,y)
 \label{sgd:smooth-addon:trivialization}
\end{equation}
is a diffeomorphism, with exact inverse
\begin{equation}
 u=\frac{4\tau v_0+72(\delta^2+y^2)}{72v_0+5\tau},
 \qquad (\tau,v_0,y)\in(0,\infty)\times\mathcal B.
 \label{sgd:smooth-addon:coupling-inverse}
\end{equation}
Moreover $\mathcal B$ is diffeomorphic to $\mathbb R^2$ by the
explicit inverse pair
\begin{equation}
 \begin{aligned}
 (y,s)&\longmapsto
 \left(\tfrac{\sqrt5}{2}\sqrt{\delta^2+y^2}+e^s,\ y\right),\\
 (v_0,y)&\longmapsto
 \left(y,\log\!\left(v_0-\tfrac{\sqrt5}{2}\sqrt{\delta^2+y^2}\right)\right).
 \end{aligned}
 \label{sgd:smooth-addon:B-chart}
\end{equation}
Here $\tau=t=g^2$ is the unchanged positive original coupling.
\end{theorem}

\begin{proof}
Since $c>0$, the quotient defining $t$ is smooth.
The derivatives $d_u=v_0$, $d_{v_0}=u$, $d_y=-2y$,
$c_u=-5$, $c_{v_0}=4$, $c_y=0$ give
\[
 \frac{\partial t}{\partial u}
 =\frac{72(v_0c+5d)}{c^2}
 =\frac{72(4v_0^2-5h)}{c^2},
 \qquad
 \frac{\partial t}{\partial v_0}
 =\frac{72(uc-4d)}{c^2}
 =\frac{72(4h-5u^2)}{c^2},
\]
and the last derivative in \eqref{sgd:smooth-addon:derivatives}
is $72(-2y)/c$.
On $\mathcal D$, $v_0=(d+h)/u>0$, $c>0$, and $d>0$.
Hence $v_0c+5d>0$, proving $\partial t/\partial u>0$.
The differential therefore has rank one everywhere, proving the
submersion assertion independently of any fibre-topology argument.
In particular its three partial derivatives cannot simultaneously
vanish. The proposed algebraic simultaneous-zero test also agrees:
the first two zero equations would give
$4v_0^2=5h$ and $4h=5u^2$, so positivity of $u,v_0$ would imply
$v_0=5u/4$, and then both $c=0$ and $d=0$, excluded from $\mathcal D$.

We now prove the exact global inverse and surjectivity.
For a point of $\mathcal D$, the identity
\[
 4v_0^2-5h=v_0c+5d>0
\]
and $v_0>0$ imply $(v_0,y)\in\mathcal B$.
Writing $\tau=72d/c>0$, multiply by the retained $c$ to obtain
\[
 \tau(4v_0-5u)=72(uv_0-h),
 \qquad
 (72v_0+5\tau)u=4\tau v_0+72h.
\]
Its denominator $72v_0+5\tau$ is strictly positive, giving
\eqref{sgd:smooth-addon:coupling-inverse}.

Conversely fix any $\tau>0$ and $(v_0,y)\in\mathcal B$, and use
\eqref{sgd:smooth-addon:coupling-inverse}. Here $h>0$, $v_0>0$,
and the numerator and denominator of $u$ are positive.
Exact multiplication, with no sign or factor removed, gives
\begin{align}
 c=4v_0-5u
 &=\frac{72(4v_0^2-5h)}{72v_0+5\tau}>0,\notag\\
 d=uv_0-h
 &=\frac{\tau(4v_0^2-5h)}{72v_0+5\tau}>0.
 \label{sgd:smooth-addon:inverse-domain}
\end{align}
Thus the recovered point is in $\mathcal D$, and
$72d/c=\tau$ exactly.
These calculations prove both compositions of
\eqref{sgd:smooth-addon:trivialization} and
\eqref{sgd:smooth-addon:coupling-inverse} are identities.
Both maps are smooth on their stated open domains.

The set $\mathcal B$ is nonempty; for example
$(v_0,y)=(\sqrt5\delta/2+1,0)$ belongs to it.
Thus every $\tau>0$ has a preimage, proving surjectivity.
Finally $\delta>0$ makes $\sqrt{\delta^2+y^2}$ smooth for all real
$y$, and the defining inequality of $\mathcal B$ makes the
logarithm argument in \eqref{sgd:smooth-addon:B-chart} strictly
positive. Exponential and logarithm are inverse smooth functions
on their stated real domains. Substitution into
\eqref{sgd:smooth-addon:B-chart} proves both compositions are
identities, establishing the diffeomorphism with $\mathbb R^2$.
\end{proof}

\subsubsection{Exact scope and topology of the positive-coupling level sets}

\begin{corollary}[All positive-coupling levels and their literal inverse charts]
\label{sgd:smooth-addon:levels}
In the literal embedded manifold $\mathcal S_{\rm sg}$, the
original coupling projection
\[
 \pi:\mathcal S_{\rm sg}\longrightarrow(0,\infty),\qquad
 (t,\omega,C_0,Z_{1,0},Z_{2,0})\longmapsto t
\]
is a smooth surjective submersion. There is an explicit diffeomorphism
\begin{equation}
 \mathcal S_{\rm sg}
 \ \cong\
 \{1,-1\}^2\times(0,\infty)\times\mathbb R^2\times\mathbb T^3
 \label{sgd:smooth-addon:global-product}
\end{equation}
under which $\pi$ is projection onto the displayed unchanged
$(0,\infty)$ coordinate.
For every specified original $\tau>0$, the complete fixed-coupling
level $\pi^{-1}(\tau)$ is an embedded five-real-dimensional
submanifold, with exactly four connected components.
Each component is diffeomorphic to $\mathbb R^2\times\mathbb T^3$.
These statements concern precisely the scalar--Gauss domain with
both original signs and all three phase coordinates retained.
\end{corollary}
\begin{proof}
Compose the literal diffeomorphism of
Theorem~\ref{sgd:smooth-addon:literal-diffeomorphism} with
the diffeomorphisms of
Theorem~\ref{sgd:smooth-addon:coupling}. More explicitly, given
\[
 (\epsilon,\sigma,\tau,y,s,\eta,\alpha,\beta)
 \in\{1,-1\}^2\times(0,\infty)\times\mathbb R^2\times\mathbb T^3,
\]
first put
\[
 v_0=\frac{\sqrt5}{2}\sqrt{\delta^2+y^2}+e^s,\qquad
 u=\frac{4\tau v_0+72(\delta^2+y^2)}{72v_0+5\tau}.
\]
Then apply \eqref{sgd:smooth-addon:literal-forward} with these
$u,v_0,y$ and the retained signs and phases.
Equation~\eqref{sgd:smooth-addon:inverse-domain} proves that its
coupling is exactly $\tau$ and every denominator is positive.
The inverse begins with the literal matrix formulas
\eqref{sgd:smooth-addon:literal-temporal-inverse} and
\eqref{sgd:smooth-addon:literal-spatial-inverse}, keeps
$\tau=t$ from the original tuple, and recovers
\[
 s=\log\!\left(v_0-\frac{\sqrt5}{2}\sqrt{\delta^2+y^2}\right).
\]
The preceding theorems prove that these maps are smooth, mutually
inverse, and defined on every point of their stated domains.
This proves \eqref{sgd:smooth-addon:global-product}.

In these exact product coordinates $\pi$ is the projection onto
$\tau$. Its derivative is onto the one-dimensional tangent space
of $(0,\infty)$, and every $\tau>0$ occurs.
The fixed-coordinate slice is a smooth embedded codimension-one
submanifold. Transport through the diffeomorphism makes
$\pi^{-1}(\tau)$ an embedded five-dimensional submanifold of
$\mathcal S_{\rm sg}$, and hence of the original ambient manifold
$\mathcal E$. Its inclusion into the fixed-coupling ambient
hyperplane has the same smooth structure, since that coordinate is
identically $\tau$.

The formula identifies the entire fibre, without a missing
parameter portion, with
$\{1,-1\}^2\times\mathbb R^2\times\mathbb T^3$.
The real plane is path connected by line segments; the three
circles are path connected by the explicit angle paths used
earlier. Their product is therefore connected.
The two continuous nonvanishing sign invariants prevent any
connection between distinct sign sectors. This proves both
the exact count of four components and their stated diffeomorphism
type. The fibre topology follows from these explicit inverse maps,
not merely from the fact that $\pi$ is a submersion.
\end{proof}

These manifolds are subsets of the space of ordered literal tuples.
The only phase identification was equality of two parameter
representatives producing the same tuple. No spatial equation has
been imposed, and no physical gauge quotient, frequency-sign
identification, or exchange of $Z_{1,0}$ and $Z_{2,0}$ has been made.
The construction and all level-set statements use the open domain
$t>0$; they assign no boundary fibre at $t=0$.


\subsection{Sources, independent checks, and reproducibility}

The evidence is under \path{supporting_materials/workbench/evidence/},
in subdirectory \path{quaternionic_4441_audit_20260904/}.
The supplied files are preserved unchanged in its \path{source/} directory.
The source claims used here are the complex-plus-triplet and Lorentz-slice
sections of \emph{Triality 4-4-4-1 Dynamic Completion}, and the coordinates,
symmetry, positive fibres, and Jordan-completion sections of
\emph{Quaternionic 4441 Continuation}, both dated 4 September 2026.
Their source TeX and original verification results remain distinguishable
from this audit and from the independent checker \path{verify_exact.py}
in that evidence directory.
It implements quaternionic matrix multiplication independently, expands
the full fifteen-real-variable determinant identity, checks the marked
skew transport and charge conjugation, and verifies the positive-cone
counterexample, retained Jordan coefficient, group composition, actual
scalar acceleration, and electric sign. These checks supplement the proofs
above. The separate \path{verify_invariant_sectors.py} expands the
trace-free adjugate and the full complex tangency identity, and checks the
retained nonzero-sector coordinates and fourth-order escape coefficient.
The invariant-family theorem is proved by tangency and uniqueness, not
inferred from finite tests. These checks do not substitute for checking
other packet claims such as the gauge-orbit part of the vacuum theorem
or the finite-lattice quantization theorem. The separate moving-sector and
vacuum verifiers test the exact periodic equations, conserved charge,
fourteen-coordinate Hessian cubic, and metric eigenvalues. The scalar
global-minimum theorem is proved above, independently of those gauge and
quantum claims.
The portable \path{verify_gauge_lift.py} checks all fourteen potential
derivatives symbolically and all forty-six real homogeneous variational
components at a retained exact two-jet, as well as the orbit, radial,
and observation identities. This finite variational test supplements,
and does not replace, the full variational proof above.
The portable \path{exact_period_metric_checker.py} verifies the unchanged
periods, inverse real period map, full affine fibres, monodromy, rank-loss
constants, and a nonreal rational positive-definite example.
The smoothness and operator inequalities are proved above rather than
inferred from finite tests.
The further portable checks
\path{exact_hermitian_metric_spectrum_checker.py},
\path{verify_vacuum_kinetics.py}, \path{verify_mass_spectrum.py}, and
\path{verify_arbitrary_gauge_lift.py}
test the retained metric/cometric factors, integral quotient and spectral
identities, all current-free continuation equations, gauge-generator
Gram matrix and eigenvalues, and the complete real-connection residual
elimination. Their exact assertion counts and replay logs are in
\path{CLASSICAL_SPECTRUM_EXACT_REPLAYS.json}. Full proofs of existence,
spectral completeness, limits and exclusions are given above; finite
symbolic tests do not replace those arguments.
The further portable \path{check_complex_rotation.py} verifies the
constant-frame blocks, retained original-coordinate map, weighted
Gauss inverse, and both sign reconstructions at an exact unrestricted
test pair. \path{check_complex_constraints.py} expands the full
unrestricted scalar entries, inertia, spatial stress and polynomial
discriminant. Neither checker certifies a zero of the remaining system.
The strict inequalities, equivalences, domains and exclusion proofs are
the analytic arguments above.

The 5 September finite-fibre and projection evidence is in
\path{supporting_materials/workbench/evidence/finite_coupling_20260905/}.
The two original symbolic checkers retain the full eight-dimensional
Gauss equation and all residual codomains. The distinct
\path{check_projection_independent.py} uses only standard-library
rational-pair and finite-field arithmetic, independently reconstructs
the full eight-dimensional Gauss matrix, and reproduces the nonzero
Sylvester determinant. Its 101 evaluations verify degree-twelve
polynomial identities, not a sampled real coupling exclusion.
Complete definitions and proofs are in
Theorems~\ref{q444:finite-coupling-fibre} and
\ref{q444:generic-projection-obstruction}; no existence inside the
exceptional algebraic set, or general complex time-dependent lift,
is inferred. The original scripts and certificates are retained in
the evidence directory's \path{checkers/} subdirectories.

The separately dated scalar--Gauss continuation and independent review are
retained in
\path{supporting_materials/workbench/evidence/scalar_gauss_domain_20260905/}.
The maintained proof is \path{scalar_gauss_domain_shared.tex}; it includes
the full temporal-classification and literal-phase-quotient addendum.
The portable \path{check_scalar_gauss_domain.py} replay passes all 25
exact groups, including every scalar and Gauss entry, both sign choices,
the continuous phases, inverse frames and the prescribed-coupling formula.
Positivity, rank exclusion, completeness and the two-element fibres are
proved analytically in Theorems~\ref{sgd:parametrization}
and~\ref{sgd:literal-addon:quotient}. This continuation supersedes the
earlier absence of a scalar--Gauss nonemptiness assertion; it does not
alter the historical proof snapshot: the two spatial residuals still
cannot vanish simultaneously.
The unrestricted exceptional-set existence problem remains open, and no
quantum spectral conclusion follows from these classical residuals.

The subsequent shared source \path{scalar_gauss_smooth_shared.tex}
proves the smooth phase quotient, embedded literal inverse charts, and
the global trivialization over the unchanged positive coupling axis
(Theorems~\ref{sgd:smooth-addon:literal-diffeomorphism}
and~\ref{sgd:smooth-addon:coupling},
Corollary~\ref{sgd:smooth-addon:levels}). These analytic proofs establish
the component and fibre topology beyond the earlier set bijections;
the earlier exact replay receipt retains its original scope.
